%%%%%%%%%%%%%%%%%%%%%%%%%%%%%%%%%%%%%%%%% Preamble %%%%%%%%%%%%%%%%%%%%%%%%%%%%%%%%%%%%%%%%%

\documentclass[11pt]{report}

\usepackage{titling} % makes following commands available
\title{Catalytic properties at the nanoscale probed by surface X-ray diffraction and coherent diffraction imaging}
\author{David Simonne}
\date{\today}

%%%%%%%%%%%%%%%%%%%%%%%%%%%%%%%%%%%%%%%%%% Colours %%%%%%%%%%%%%%%%%%%%%%%%%%%%%%%%%%%%%%%%%%

\usepackage{xcolor}

\definecolor{ArgonColor}{HTML}{008FD5}
\definecolor{AmmoniaColor}{HTML}{DD65B0}
\definecolor{OxygenColor}{HTML}{FC4F30}
\definecolor{NitrogenColor}{HTML}{810F7C}
\definecolor{NitrousOxideColor}{HTML}{8B8B8B}
\definecolor{NitrogenOxideColor}{HTML}{6D904F}
\definecolor{WaterColor}{HTML}{E5AE38}

\definecolor{Ostwald}{HTML}{800080}
\definecolor{Haber}{HTML}{158466}

\definecolor{LightBlue}{RGB}{30,75,180}
\definecolor{LUBlue}{RGB}{0,0,128}
\definecolor{LUGrey}{RGB}{77,76,68}
\definecolor{Prune}{RGB}{99,0,60}
\definecolor{Important}{HTML}{980000}

\definecolor{DarkFern}{HTML}{407428}
\definecolor{DarkCharcoal}{HTML}{4D4944}
\definecolor{DarkBlue}{HTML}{23373B}
\definecolor{DarkOrange}{HTML}{EA7500}
\definecolor{LightOrange}{HTML}{F9D6B3}
\definecolor{Blue}{HTML}{13406a}

\colorlet{lightblue}{blue!60}
\colorlet{lightorange}{orange!60}
\colorlet{lightgreen}{green!60}
\colorlet{lightred}{red!60}
\colorlet{lightgray}{gray!60}
\colorlet{lightbrown}{brown!60}
\colorlet{lightpink}{pink!60}
\colorlet{lightviolet}{violet!60}

\colorlet{shadecolor}{gray!40}

%%%%%%%%%%%%%%%%%%%%%%%% Page layout %%%%%%%%%%%%%%%%%%%%%%%%

\usepackage[
    a4paper,
    width=155mm,
    headheight=110pt,
    top=25mm,
    bottom=25mm
]{geometry}

\usepackage{etoolbox} % Apply different styles
\usepackage{fancyhdr}
\pagestyle{fancy}
\renewcommand{\chaptermark}[1]{\markboth{#1}{}}
\renewcommand{\sectionmark}[1]{\markright{\thesection\ #1}}

% Define fancy style for thesis
\fancypagestyle{main}{
%    \fancyhead[L]{\chaptername \, \thechapter:  \small{\textsf{\leftmark}}}
    \fancyhead[L]{\chaptername \, \thechapter}
    \fancyhead[R]{\small{\textsf{\rightmark}}}
    \fancyfoot[R]{\small{\thepage}}
    \fancyfoot[L]{\small{\theauthor}}
    \fancyfoot[C]{}
    \renewcommand{\headrulewidth}{0.4pt}%
    \renewcommand{\footrulewidth}{0.4pt}%
}
\appto\mainmatter{\pagestyle{main}}

% Define plain style for table of contents and acknowledgements
\fancypagestyle{plain}{
    \fancyhf{}
    \fancyfoot[C]{\thepage}
    \renewcommand{\headrulewidth}{0pt}%
    \renewcommand{\footrulewidth}{0pt}%
}
\appto\frontmatter{\pagestyle{plain}}

\usepackage{sectsty} % To colour sections, raise an error on underbar
\chapterfont{\color{Prune}}
\sectionfont{\color{LUGrey}}
\subsectionfont{\color{LUGrey}}
\subsubsectionfont{\color{LUGrey}}

%%%%%%%%%%%%%%%%%%%%%%%%%% Citation %%%%%%%%%%%%%%%%%%%%%%%%%%

\usepackage[utf8]{inputenc} % Use accents in source code, load before biblatex
\usepackage[T1]{fontenc} % Show accents in output

\usepackage[style=authoryear, isbn=false, date=year, sorting=nyt, maxcitenames=1, % mincitenames=1,
    eprint=false, giveninits=true, url=false]{biblatex}
\addbibresource{/home/david/Documents/PhD/PhdThesis/references.bib}
\AtEveryBibitem{\clearfield{pages}}

\usepackage[nottoc,numbib]{tocbibind}
\usepackage{tocloft}
% \setlength{\cftbeforetoctitleskip}{0em}
% \setlength{\cftaftertoctitleskip}{0em}

%%%%%%%%%%%%%%%%%%%%%%% Other packages %%%%%%%%%%%%%%%%%%%%%%%
\usepackage{lipsum} % Latin blabla
\usepackage{lmodern} % to avoid errors in font size
\usepackage[british]{babel}
\usepackage{setspace} % Space between lines
\usepackage{mathtools}
\usepackage{tikz} % Drawing
\usetikzlibrary{shapes.geometric, arrows, positioning, calc}
\usepackage{colortbl} % Colour the lines in the tables
% \usepackage{minted} % Format code
\usepackage[version=4]{mhchem}
\usepackage{animate} % for animations in 3D
\usepackage{pdfpages}
\pdfminorversion=6 % to avoid errors when including pdf
\usepackage{epigraph} % for quote

\usepackage{textcomp} % To avoid warnings in gensymb
\usepackage{gensymb} % for degree symbol

\usepackage{siunitx} % Units
\DeclareSIUnit\angstrom{\text{Å}}
\DeclareSIUnit\barn{b}
\DeclareSIUnit\photons{\text{photons}}
\DeclareSIUnit\bar{bar}
\usepackage{wasysym} % for diameter symbol

\usepackage{graphicx} % Insert images
\graphicspath{{/home/david/Documents/PhD/Figures/}}

\usepackage[hypcap=true,font={small,it},justification=centering]{caption}
\captionsetup{belowskip=2pt,aboveskip=2pt}

%\usepackage{floatrow}
\usepackage{sidecap}
\sidecaptionvpos{figure}{c}
\usepackage{scrextend} % Forcer la 4eme  de couverture en page pair
\usepackage{mdframed} % Entourer un paragraphe
\usepackage{multirow} % Mettre un texte sur plusieurs rangées
\usepackage{multicol} % Mettre un texte sur plusieurs colonnes
\usepackage{array} % Create tables as arrays
\usepackage{csquotes}
\usepackage{booktabs} % for \midrule and such
\usepackage{lscape}
\usepackage{amsmath} % for maths
\usepackage{amssymb}

\usepackage[page,toc,titletoc,title]{appendix}

\usepackage[
    colorlinks = true,
    linkcolor = LightBlue,
    urlcolor  = LightBlue,
    citecolor = LightBlue,
    anchorcolor = LightBlue,
]{hyperref} % load last to avoid errors

%%%%%%%%%%%%%%%%%%%%%% Commands %%%%%%%%%%%%%%%%%%%%%

\newcommand{\argon}{\ce{Ar}}
%{\textcolor{ArgonColor}

\newcommand{\ammonia}{\ce{NH_3}}
%{\textcolor{AmmoniaColor}

\newcommand{\nitrousoxide}{\ce{N_2O}}
%{\textcolor{NitrousOxideColor}

\newcommand{\water}{\ce{H_2O}}
%{\textcolor{WaterColor}

\newcommand{\nitrogen}{\ce{N_2}}
%{\textcolor{NitrogenColor}

\newcommand{\nitricacid}{\ce{HNO_3}}
%{\textcolor{DarkOrange}

\newcommand{\nitricoxide}{\ce{NO}}
%{\textcolor{NitrogenOxideColor}

\newcommand{\nitrogendioxide}{\ce{NO_2}}

\newcommand{\carbondioxide}{\ce{CO_2}}

\newcommand{\dioxygen}{\ce{O_2}}
%{\textcolor{OxygenColor}

\newcommand{\ptthreeofour}{\ce{Pt_3O_4}}
%{\textcolor{Important}

\newcommand{\ammoniumnitrate}{\ce{NH_4NO_3}}
%{\textcolor{Blue}

\newcommand{\urea}{\ce{CO(NH_2)_2}}

\newcommand{\yes}{{\textcolor{green}{0}}}

\newcommand{\no}{{\textcolor{red}{x}}}
\newcommand{\ra}[1]{\renewcommand{\arraystretch}{#1}}

%%%%%%%%%%%%%%%%%%%%%% Document %%%%%%%%%%%%%%%%%%%%%%

\begin{document}

    \begin{titlepage}

\newgeometry{left=6cm,bottom=1cm, top=1cm, right=1cm}

\tikz[remember picture,overlay] \node[opacity=1,inner sep=0pt] at (-13mm,-135mm){\includegraphics{logo/Bandeau_UPaS.pdf}};

% fonte sans empattement pour la page de titre
\fontfamily{fvs}\fontseries{m}\selectfont

%*****************************************************
%******** NUMÉRO D'ORDRE DE LA THÈSE À COMPLÉTER *****
%******** POUR LE SECOND DÉPOT                   *****
%*****************************************************

\color{white}

\begin{picture}(0,0)

\put(-110,-743){\rotatebox{90}{NNT: 2020UPASA000}}
\end{picture}

\vspace{-12mm}
\flushright \includegraphics[height=2.1cm]{logo/CEA.png} \hspace{7mm}
\includegraphics[height=2.1cm]{logo/SOLEIL.png}

%*****************************************************
%******************** TITRE **************************
%*****************************************************

\flushright
\vspace{10mm}
\color{Prune}
\fontfamily{cmss}\fontseries{m}\fontsize{22}{26}\selectfont
Propriétés catalytiques à l'échelle nanométrique sondées par diffraction des rayons X de surface et imagerie de  diffraction cohérente.

\normalsize
\color{black}
\Large{\textit{Catalytic properties at the nanoscale probed by surface x-ray diffraction and coherent diffraction imaging}}
%*****************************************************

\vspace{1.5cm}
\normalsize
\textbf{Thèse de doctorat de l'Université Paris-Saclay}

\vspace{15mm}

École doctorale n$^{\circ}$ 000, dénomination et sigle\\
\small Spécialité de doctorat: Physique\\
\footnotesize Unité de recherche: voir annexe\\
\footnotesize Référent: : voir annexe\\
\vspace{15mm}

\textbf{Thèse présentée et soutenue à .....,\\ le .... 2023, par}\\
\bigskip
\Large {\color{Prune} \textbf{David SIMONNE}}

%************************************
\vspace{\fill} % ALIGNER LE TABLEAU EN BAS DE PAGE
%************************************

\bigskip
\flushleft
\scriptsize
\begin{tabular}{|p{7cm}l}
\arrayrulecolor{Prune}
{\footnotesize \textbf{Composition du jury}}\\
& \\
\textbf{Prénom Nom} &   Président/e\\ 
Titre, Affiliation & \\
\textbf{Prénom Nom} &  Rapportrice \\ 
Titre, Affiliation   &   \\ 
\textbf{Prénom Nom} &  Rapporteur \\ 
Titre, Affiliation  &   \\ 
\textbf{Prénom Nom} &  Examinatrice \\ 
Titre, Affiliation   &   \\ 
\textbf{Prénom Nom} &  Examinateur \\ 
Titre, Affiliation   &   \\ 
\textbf{Prénom Nom} &  Examinateur \\ 
Titre, Affiliationt   &   \\ 

\end{tabular} 

\medskip
\begin{tabular}{|p{7cm}l}
\arrayrulecolor{Prune}
{\footnotesize \textbf{Direction de la thèse}}\\
& \\
\textbf{Alessandro Coati} & Directeur\\
Dr., Synchrotron SOLEIL & \\
\textbf{Andrea Resta} & Codirecteur\\
Dr., Synchrotron SOLEIL & \\
\textbf{Marie-Ingrid Richard} & Coencadrante\\
Dr., CEA Grenoble & \\

\end{tabular} 

\end{titlepage}
    \newpage
    %%%%%%%%%%%%%%%%%%%%%%%%%%%%%%%%%%%%%%%%%%%%%%%%%%%%%%%%%%%%%%%
\thispagestyle{empty}

\lhead{}
\rhead{}
\rfoot{}
\cfoot{}
\lfoot{}

\noindent 
\includegraphics[height=2.4cm]{/home/david/Documents/PhD/Figures/Logos/logo_usp_PIF.png}
\vspace{0.5cm}

\small

\begin{mdframed}[linecolor=Prune,linewidth=1]

\textbf{Titre:} Propriétés catalytiques à l’échelle nanométrique sondées par diffraction des rayons X de surface et imagerie de diffraction cohérente

\noindent \textbf{Mots clés:} Diffraction de rayons X, Catalyse hétérogène, Surface, Structure cristalline, Déformation, Oxydation de l'ammoniac

\vspace{-.5cm}
\begin{multicols}{2}
\noindent \textbf{Résumé:}

Le principal objectif de ce travail est d'étudier des catalyseurs hétérogènes \textit{in situ} et \textit{operando} pendant l'oxydation de l'ammoniac en se rapprochant des valeurs de température et pression industrielles.
Actuellement, ce processus catalytique et les changements structurels associés sont mal compris.
Nous proposons d'utiliser différents échantillons en platine - nanoparticules et monocristaux - afin de réduire l'écart entre les études scientifiques sur échantillons modèles et les catalyseurs utilisés en industrie.
L'activité catalytique des différents échantillons est mesurée pour lier structure et sélectivité durant la réaction, qui peut être focalisée vers la production d'azote (\ce{N_2}), d'oxyde nitrique (\ce{NO}), ou de protoxyde d'azote (\ce{N_2O}).
Le développement d'une catalyse hétérogène avec une sélectivité ciblant les \qty{100}{\percent} est un défi constant, ainsi que la compréhension de la durabilité, du vieillissement, et de la désactivation du catalyseur.
Trois techniques ont été principalement utilisées, l'imagerie par diffraction cohérente de Bragg, la diffraction des rayons X de surface et la spectroscopie photoélectronique des rayons X, combinées à des mesures par spectrométrie de masse.
Ces techniques sont compatibles avec des conditions de pression ambiantes, permettant de réduire l'écart de pression entre expérience et industrie en catalyse hétérogène.
Mesurer la structure de nanoparticules à l'échelle nanométrique permet de révéler les effets de volume, de tension et de compression de surface et d'interface, ainsi que l'existence de différents types de défauts.
En complément des études d'imagerie par diffraction cohérente de rayons X en condition de Bragg sur des nanoparticules individuelles, l'étude d'un ensemble de nanoparticules est effectuée \textit{via} la diffraction des rayons X à incidence rasante.
La diffraction cohérente de rayons X en condition de Bragg étant une technique récente, une organisation typique de la réduction et de l'analyse des données est proposée.
L'importance d'étudier plusieurs particules a été mise en perspective par l'étude comparative de deux particules durant l'oxydation de l'ammoniac à \qty{300}{\degreeCelsius} et \qty{400}{\degreeCelsius}, en fonction du rapport entre \ce{O_2} et \ce{NH3}.
Différents comportements ont été révélés sur les deux nanoparticules, qui présentent une taille, une forme, des facettes et un état de déformation initial différents, alors qu'aucun important changement n'a été mesuré en dessous de \qty{600}{\degreeCelsius} par diffraction de surface.
Une particule a montré une diminution/augmentation réversible de la déformation homogène lorsque l'ammoniac a été introduit/retiré du réacteur.
L'apparition d'un défaut à \qty{400}{\degreeCelsius} est liée à une augmentation non réversible et importante de la déformation homogène lors de l'oxydation de l'ammoniac pour la seconde particule, qui continue d'augmenter en fonction du rapport entre ammoniac et oxygène.
Cette évolution structurelle est clairement visible en 3D, avec des volumes de densité électronique de Bragg manquants.
De plus, deux type de surfaces présente sur les nanoparticules (\{111\} et \{100\}) ont été étudiées à l'aide de monocristaux par diffraction des rayons X en surface, et spectroscopie photoélectronique par rayons X.
De ce fait, la structure de surface ainsi que la présence d'espèces adsorbées peuvent être reliées à l'activité catalytique mesurée, permettant une meilleure compréhension du mécanisme de réaction.
Différentes reconstructions de surfaces ont été mesurées sur le Pt(100) pendant la réaction, ce qui n'est pas le cas pour le Pt(111).
La présence d'oxygène est liée à une importante rugosité de surface, alors que la présence d'ammoniac à une baisse de celle-ci.
De plus, différentes oxydes, ainsi que des structures transitoires ont été identifiés sur les deux monocristaux.
La mesure de différents niveaux par spectroscopie photoélectronique à rayons X a permis de lier une sélectivité accrue du Pt(100) vers NO par rapport au Pt(111) avec une plus importante présence d'oxygène adsorbée sur la surface.
\end{multicols}

\end{mdframed}

\newpage
\thispagestyle{empty}

\lhead{}
\rhead{}
\rfoot{}
\cfoot{}
\lfoot{}

\noindent
\includegraphics[height=2.4cm]{/home/david/Documents/PhD/Figures/Logos/logo_usp_PIF.png}
\vspace{0.5cm}

\small

\begin{mdframed}[linecolor=Prune,linewidth=1]

\textbf{Title:} Catalytic properties at the nanoscale probed by surface X-ray diffraction and coherent diffraction imaging

\noindent \textbf{Keywords:} X-ray diffraction, Heterogeneous catalysis, Surface, Crystal structure, Strain, Ammonia oxidation

\vspace{-.5cm}
\begin{multicols}{2}
\noindent \textbf{Abstract:}
The main objective of this work is to study heterogeneous catalysts \textit{in situ} and \textit{operando} during the oxidation of ammonia by approaching industrial temperature and pressure values.
Currently, this catalytic process and the associated structural changes are poorly understood, and we propose to use different samples in platinum, nanoparticles and single crystals in order to reduce the gap between scientific studies on model samples and catalysts used in industry.
The catalytic activity of the different samples is measured to link structure and selectivity during the reaction, which can be focused towards the production of nitrogen (\ce{N2}), nitric oxide (\ce{NO}), or nitrous oxide (\ce{N2O}).
Developing heterogeneous catalysis with selectivity targeting \qty{100}{\percent} is an ongoing challenge, as is understanding the durability, ageing, and deactivation of the catalyst itself.
Three techniques were mainly used, Bragg coherent diffraction imaging, surface X-ray diffraction and X-ray photoelectron spectroscopy, combined with mass spectrometry measurements.
These techniques are compatible with ambient pressure, making it possible to reduce the difference in pressure between surface science studies and industrial conditions for heterogeneous catalysis.
Measuring the structure of nanoparticles at the nanoscale makes it possible to reveal the effects of volume, surface and interface tension and compression, as well as the existence of different types of defects.
In addition to imaging studies by coherent X-ray diffraction in Bragg conditions on individual nanoparticles, the study of a set of nanoparticles is carried out \textit{via} grazing incidence X-ray diffraction.
Coherent X-ray diffraction in Bragg conditions being a recent technique, a typical organisation of data reduction and analysis is proposed.
The importance of studying several particles was put into perspective by the comparative study of two particles during the oxidation of ammonia at \qty{300}{\degreeCelsius} and \qty{400}{\degreeCelsius} as a function of the pressure ratio between \ce{O2} and \ce{NH3}.
Different behaviours were revealed on the two nanoparticles, which present a different size, shape, facet coverage and initial deformation state, while no significant changes were measured below \qty{600}{\degreeCelsius} by diffraction of surface.
One particle showed a reversible decrease/increase in homogeneous strain when ammonia was introduced/removed from the reactor.
The appearance of a defect at \qty{400}{\degreeCelsius} is linked to a non-reversible and significant increase in the homogeneous deformation during the oxidation of ammonia for a second particle, which continues to increase as a function of the ratio between ammonia and oxygen.
This structural evolution is clearly visible in 3D, with missing Bragg electron density volumes.
In addition, two types of surfaces present on the nanoparticles (\{111\} and \{100\}) were also studied using single crystals by surface X-ray diffraction and X-ray photoelectron spectroscopy.
Therefore, the surface structure as well as the presence of adsorbed species can be linked to the measured catalytic activity, allowing a better understanding of the reaction mechanism.
Different surface reconstructions were measured on Pt(100) during the reaction, which is not the case for Pt(111).
The presence of oxygen is linked to significant surface roughness, while the presence of ammonia is linked to a reduction in roughness.
Furthermore, different oxides, as well as transient structures, were identified on the two single crystals.
The measurement of different levels by X-ray photoelectron spectroscopy made it possible to link an increased selectivity of Pt(100) towards NO compared to Pt(111) to a greater presence of oxygen adsorbed on the surface.
\end{multicols}

\end{mdframed}

\normalsize

\vspace{\fill}

\newpage\thispagestyle{empty}\null\newpage

    % Optional Chapters
        \chapter*{Dedication}
        I would like to dedicate this thesis to Brittany who has pushed and supported me every day since we've met in a crystallography course.
        \chapter*{Declaration}
        I hereby declare that this thesis is my original work and has not been presented and will not be presented to any other university or institution for a similar or any other degree or professional qualification.
I have acknowledged all sources of information that have been used in the thesis, including any assistance that has been received during the research and writing process.
Any contribution by individuals or sources other than my own has been cited and acknowledged.
I understand that any act of plagiarism or academic dishonesty is a violation of the academic integrity policy of the University Paris-Saclay, and I am aware of the potential consequences.

\begin{figure}[!htb]
    \centering
    \includegraphics[width=0.495\textwidth]{/home/david/Documents/PhD/Figures/signature.png}
\end{figure}

        \chapter*{Acknowledgements}
        I extend my deepest gratitude to Andrea for his unwavering support throughout the duration of this thesis.
You have contaminated me with your love of Linux and open source technology, and set a strong example in terms of commuting to work by bicycle, which I failed to follow in rainy and cold days.
Thank you very much for all the feedback that you gave me about this thesis, down to the details in the figure placement, fontsize or linewidth, and sentence length.
Pursuing your interest in the study of the ammonia oxidation has allowed me to explore a large subject.
This work has given me a sense of purpose, as it feels like my contribution is ever so slightly adding to the global endeavor in combating global warming and climate change.

I am immensely grateful for the shared beamtimes, the camaraderie during strenuous long shifts in summer or winter, skipping breakfast or dinner, compiling data while monitoring the sample heater.

Thank you Alessandro for being engaged in everything you do, at work and outside, for taking care of all the tedious work so that I could have the best possible experience, for inviting me to have meals and sometimes bringing me home after work, and for providing the SixS beamline with good italian coffee.

Thank you Marie-Ingrid for being so driven by your love of science that everybody around wants to participate, for arranging my stay in Grenoble and bringing all of the padawans to conferences, for always being in a good mood and ready to discuss new challenging experiments.

Thank you to my three advisors for being good mentors, allowing me to teach, to participate in schools and numerous beamtimes, for having my best interest in mind at all times and for trusting me.
I have learned an incredible amount from each of you and particularly appreciated the beamtime at Diamond, during which the four of us collected many electrons, but also shared meaningful conversations around a pint of beer and delicious stew.

I am very grateful to the research staff of SixS, to Alina and Yves for helping me when I was lost in the reciprocal space, to Benjamin without which none of the experiments could be possible, and to Michèle for her numerous advices.
Thank you also Frédéric for spending a lot of time in helping us regarding BCDI and SXRD code development.
Thank you to all the young researchers at SOLEIL for creating a group so that Saclay did not feel too remote and secluded.

I of course have a specific thought for the Grenoble legions, half of whom I don't know half as well as I should like, and half of which I know from half my life spent in beamtime.
Maxime, for introducing me to the world of political satire videos during ammonia oxidation cycles at SixS; Corentin, for the delightful culinary explorations in Grenoble, and for lending me the ID01 bicycle; Clément, for your constant positivity and our conversations about the promising future of PhD students; Nikita and Mattia, for the shared moments over a refreshing beer and the intense MSSBB games during Hercules; Ewen, for representing Bretagne alongside me; Edoardo, for the initiative of LAN parties, and the unforgettable Quake games; and Noor, for being an exceptional office (and coffee) mate in Grenoble.

I would also like to thank Steven for his good advice and for kind conversations, as well as Tobias and Joël for facilitating my stay in Grenoble for a few months during this thesis.
Thank you Vincent for letting me participate in the engaging ESRF tutorials, and Jérôme for starting the large work of bringing our BCDI community together in code development.
My heartfelt appreciation to Sarah and Stéphane for their warm hospitality in Marseille and their engaging discussions.

I feel very grateful towards Virginie and Christian who have followed me during the thesis by participating in the CSI, even available on weekends so that I would not miss registrations deadlines.

Thank you Andrea, Davide, Elisa and all the others from the university of Torino for making me want to pursue an academic path.
Thank you Davide again for hosting me in Grenoble after the ESRF user meeting because my train got cancelled, thank you Alessia and Giorgio for welcoming me.

Thank you to all the MaMaSELF members without which I would never have kept studiying physics and material science, and thus missed on studying in many countries, discovering new cultures, and mindsets.
Thank you Michael for bringing me to the world of neutron diffraction at Munich, and for letting me participate in my first beamtime at the ILL with Ulrike, a master and PhD student that had to learn a lot fast, but sadly not fast enough to handle the beam polarisation well.

Thank you to all the professors from Rennes that transmitted me their love of physics, especially Pr. Guérin that permitted me to study one year in Japan, thank you to Pr. Iwai for accepting me as part of his group in Sendaï.

Finally, to my family, whose unwavering support has bolstered my every endeavor, to Lucas, whose shared passion for science has been a source of motivation, and to my parents, for their trust and encouragement in my global pursuits.
Thank you to all my friends that have seen me being projected in numerous states during those three years, and without which I would probably not have kept the same level of sanity.

I extend my thanks to the esteemed members of the jury for their invaluable feedback and constructive critique, shaping the course of this work.

Furthermore, I acknowledge the invaluable support and resources provided by synchrotron SOLEIL, the CEA, and the ERC Carine, without which this research endeavor would not have been feasible.

Lastly, I express my gratitude to all those who have, in various capacities, contributed to the completion of this thesis.

\newpage
        \chapter*{Résumé étendu en français}
        % chapter 1 / intro
Dans le premier chapitre de cette thèse, l'oxydation de l'ammoniac (\ce{NH_3}) ainsi que son importance est présentée.
La réaction peut être décrite par trois équations qui, en fonction du rapport stoechiométrique entre \ce{NH_3} et \ce{O_2}, ont des produits différents.

\begin{align}
    \label{eq:AmmoniaOxidationNitrogenFr}
    4 \ammonia + 3 \ce{O_2} & \rightarrow 6 \water + 2 \nitrogen \\
    \label{eq:AmmoniaOxidationNitrousOxideFr}
    4 \ammonia + 4 \ce{O_2} & \rightarrow 6 \water + 2 \nitrousoxide \\
    \label{eq:AmmoniaOxidationNitricOxideFr}
    4 \ammonia + 5 \ce{O_2} & \rightarrow 6 \water + 4 \nitricoxide
\end{align}

Cette réaction et ses produits associés ont exercé une profonde influence au cours des 20$^{ème}$ et 21$^{ème}$ siècles, jouant un rôle central par exemple dans l'industrie des engrais (procédé d'Ostwald, \textit{via} la synthèse acide nitrique) et les changements démographiques associés.
Malgré son importance industrielle cruciale, les diverses applications industrielles liées au procédé d'Ostwald ont contribué de manière significative à la fois au changement climatique et à la pollution toujours croissante de nos écosystèmes.

Bien qu’il s’agisse d’une réaction catalytique majeure, les mécanismes d’action exacts de l'oxydation de l'ammoniac ne sont pas encore entièrement compris.
Notamment, l'origine d'importants changements morphologiques durant l'utilisation des catalyseurs en industrie (fig. \ref{fig:GauzesFr}), liés à une \textit{activation} de l'échantillon permettant une meilleure sélectivité vers \ce{NO}, est encore inconnue.

\begin{figure}[!htb]
    \centering
    \includegraphics[width=\textwidth]{/home/david/Documents/PhD/Figures/Sample/EtchedGauzeAndReconstructedGauze.pdf}
    \caption{
    Images obtenues par microscope électronique à balayage, gazes reconstruits (Pt-Rh) montrant l'existence de motifs en forme de chou-fleur après utilisation en industrie (droite).
    Adapté depuis Bergene et al. \parencite*{Bergene1996}.
    La barre horizontale mesure \qty{0.1}{\mm}.
    }
    \label{fig:GauzesFr}
\end{figure}

Obtenir une compréhension complète de ces mécanismes est essentiel pour contrôler la sélectivité des réactions.
De plus, une compréhension plus approfondie du mécanisme de réaction est prometteuse pour le développement de nouveaux catalyseurs capables de s’affranchir de la dépendance aux métaux précieux coûteux, ou de fonctionner à plus basse température.

Cette thèse vise à étudier l'oxydation catalytique de l'ammoniac à pression ambiante en couplant plusieurs techniques de diffraction et de spectroscopie.
Pour ce faire, trois techniques principales ont été utilisées, l'imagerie par diffraction cohérente de Bragg, la diffraction des rayons X de surface et la spectroscopie photoélectronique des rayons X, combinées à des mesures par spectrométrie de masse (tab. \ref{tab:TechniquesFr}).
Ces techniques sont compatibles avec des conditions de pression proches de l'ambiante, permettant de réduire l'écart de pression en catalyse hétérogène.
Le réacteur utilisé pour étudier la réaction catalytique hétérogène sur la ligne SixS (SOLEIL) est déjà compatible avec les environnements hautement oxydants \parencite{VanRijn2010, Resta2020a}, et permet des expériences de diffraction des rayons X en surface à haute pression et d'imagerie par diffraction cohérente de Bragg.

Pour combler le fossé matériel, des nanoparticules et monocristaux de platine ont été utilisés.
Tout au long de ce travail, de faibles rapports de pression partielle \ce{O_2}/\ce{NH_3} et des basses températures ont été liés à la sélectivité vers \ce{N_2}, alors que des rapports \ce{O_2}/\ce{NH_3} et températures élevés ont été liés à une sélectivité accrue vers \ce{NO} et \ce{N_2O}.
Ceci est cohérent avec la sélectivité rapportée pour les catalyseurs industriels, qui soutient l'utilisation de nanoparticules et de monocristaux comme catalyseurs modèles pour améliorer la compréhension du système.

Les éléments optiques nécessaires à l'utilisation d'un faisceau cohérent focalisé ont été implémentés et caractérisés.
L'extension du panel de techniques disponibles contribuera à combler le fossé entre les matériaux et pressions utilisés en industrie par rapport à ceux utilisés en sciences de surfaces.
Dans le futur, les utilisateurs du synchrotron auront la possibilité d'étudier les réactions catalytiques dans des conditions et environnements équivalents, mais avec des techniques différentes.

L'expérience de spectroscopie photoélectronique à rayons X a été réalisée sur la ligne de lumière B07 (Diamond), également compatible avec les hautes pressions \parencite{Held2020}.
L'écart entre le matériau et la pression est partiellement comblé en fonctionnant à des températures supérieures à l'amorçage du catalyseur, à une pression presque industrielle.

\begin{table}[!htb]
\centering
\resizebox{\textwidth}{!}{%
    \begin{tabular}{l|l|l|l}
        Technique    & Diffraction des rayons    & Diffraction cohérente & Spectroscopie photoélectronique \\
                     & x en surface              & des rayons x en       & par rayons x              \\
                     &                           & condition de Bragg    &                           \\
        \midrule
        Échantillons & Monocristaux de           & Particules de         & Monocristaux de           \\
                     & Pt(111) et Pt(100),       & platine isolées       & Pt(111) et Pt(100),       \\
                     & particules de platine     &                       &                           \\
        \midrule
        Information  & Structure de surface,      & Forme, champs de              & Présence d'espèces de      \\
                     & rugosité, relaxation \&    & déformation et de déplacement & surfaces, quantité \& état \\
                     & phases crystallographiques & de l'objet unique             & d'oxydation                \\
        \midrule
        Ligne de lumière & SixS (SOLEIL)          & SixS (SOLEIL)                 & B-07 (Diamond Light Source) \\
    \end{tabular}
    }
    \caption{
        Techniques utilisées dans le cadre de cette thèse employant les rayons X.
    }
    \label{tab:TechniquesFr}
\end{table}

% chapter 2
Le deuxième chapitre fourni un aperçu concis de la catalyse hétérogène et des principes fondamentaux régissant l'interaction entre les rayons X et la matière.
Cette section met l'accent sur les origines, les avantages et les contraintes de chaque technique utilisée dans l'étude, en soulignant comment les rayons X peuvent observer indirectement les réactions catalytiques en détectant des marqueurs uniques laissés sur les matériaux.
Le concept de sites actifs est introduit ainsi que le lien entre adsorption et déformation de surface, crucial pour l'étude de la réaction catalytique avec les techniques de diffraction des rayons X, intrinsèquement sensible à la déformation du réseau dans les cristaux.

La ligne de lumière SixS du synchrotron SOLEIL, principal outil pour la plupart des expériences, est présentée, mettant l'accent sur les dernières avancées en matière de techniques expérimentales et de matériel spécifique à la ligne de lumière.

De plus, la thèse présente les différents logiciels utilisés pour la réduction et l'analyse des données développés au cours de ce travail, en se concentrant sur l'analyse complète pour l'imagerie par diffraction cohérente de Bragg, une technique qui n'a pas encore atteint son plein potentiel grâce au développement de synchrotrons de $4^{ème}$-génération et de clusters de calcul puissants.
Une description du logiciel \textit{Gwaihir} \parencite{Simonne2022}, qui vise à faciliter la réduction et l'analyse des données dans le langage de programmation \textit{Python}, tout en encourageant une adoption plus large de la technique à travers une interface graphique, est fournie.

% chapter 3
Les troisième et quatrième chapitres présentent les résultats de la thèse durant l'étude de différents échantillons pendant l'oxydation de l'ammoniac.
L'étude de la réaction est réalisée \textit{operando}, en explorant l'espace de paramètres multidimensionnel défini par la température de fonctionnement, la pression totale et le rapport \ce{O2} / \ce{NH3}.
Pour explorer la corrélation entre la structure de surface et l'activité catalytique, des données de spectrométrie de masse sont collectées simultanément avec toutes les mesures.

Dans le chapitre 3, des nanoparticules de platine sont utilisées pour explorer la relation entre morphologie/structure et sélectivité à pression ambiante et à haute température, au moyen d'une vaste gamme de rapports de pression partielle \ce{O2}/\ce{NH3}, favorisant soit la production de \ce{N2} ou \ce{NO}.
Ces nanoparticules présentes une distribution de taille entre \qty{350}{\nm} et \qty{800}{\nm}.

\begin{figure}[!htb]
    \centering
    \includegraphics[width=\textwidth]{/home/david/Documents/PhD/PhDScripts/SixS_2021_03_SXRD_NH3/figures/facets/facet_signal_evolution_together.pdf}
    \caption{
    Évolution de l'intensité diffusée prise le long d'une zone carrée perpendiculaire à la direction [111], à trois positions différentes dans le plan ($\vec{q}_x, \vec{q}_z$) pour sonder l'évolution de l'intensité des tiges de troncature cristalline sur une assemblée de particules en diffraction de surface.
    }
    \label{fig:FacetSignalFr}
\end{figure}

La possibilité que différents phénomènes se produisent dans les mêmes conditions de réaction a été révélée.
D'importantes différences ont en effet été mesurées entre les informations moyennes obtenues par diffraction de surface sur un échantillon exhibant des milliers de particules, et les informations obtenues par diffraction cohérente sur des particules uniques.

Les nanoparticules de platine épitaxiées sur saphir ont d'abord été mesurées à \qty{300}{\degreeCelsius}, \qty{500}{\degreeCelsius} et \qty{600}{\degreeCelsius} par diffraction de surface lors de l'oxydation de l'ammoniac, en commençant par l'introduction de l'ammoniac dans le réacteur, et suivie d'une augmentation progressive de la pression d'oxygène (\ce{O_2}/\ce{NH_3} = \{0, 0,5, 1, 2, 8\}).
La forme moyenne des nanoparticules de Pt est principalement constituée de facettes de types \{111\}, \{110\}, \{100\} et \{113\}.

Un re-facetage des particules est notamment révélé à \qty{600}{\degreeCelsius} lors des conditions de réaction.
Une fois que \ce{O_2}/\ce{NH_3} = 1, et en augmentant les rapports oxygène/ammoniac, le taux de couverture des facettes \{110\} et \{113\} diminue, remplacée par des facettes de type \{111\}, et \{100\} qui sont favorisées (fig. \ref{fig:FacetSignalFr}).

Une particule a ensuite été imagée par imagerie cohérente en condition de Bragg sous atmosphère inerte, à des températures comprises entre \qty{25}{\degreeCelsius} et \qty{600}{\degreeCelsius}.
Une dislocation présente à l'interface depuis la température ambiante jusqu'à \qty{125}{\degreeCelsius} a été éliminée avec succès par recuit au-dessus de \qty{800}{\degreeCelsius}.
Il a été démontré que cette dislocation a un impact sur la relaxation thermique de la particule.
De plus, les facettes \{111\} proches de l'interface évoluent entre des facettes de type \{211\}, \{221\} et \{110\} en fonction de la température de l'échantillon et de la déformation interfaciale.
L'introduction d'ammoniac dans le réacteur à \qty{600}{\degreeCelsius} a entraîné une inversion de la déformation des facettes, probablement à l'origine de la création d'un réseau de dislocations interfaciales, mettant en évidence le lien entre facette et déformation interfaciale dans les nanoparticules.

L'importance de la forme des particules, de la taille des facettes, de leur taux de couverture ainsi que de l'état de déformation initial a été mise en perspective par l'étude comparative de deux particules durant l'oxydation de l'ammoniac à \qty{300}{\degreeCelsius} et \qty{400}{\degreeCelsius}, en fonction du rapport entre \ce{O_2} et \ce{NH3}.
Différents comportements ont été révélés sur les deux nanoparticules, qui présentent une taille, une forme, une couverture de facettes et un état de déformation initial différents, alors qu'aucun important changement n'a été mésuré en dessous de \qty{600}{\degreeCelsius} par diffraction de surface.
La surface de la plus grande nanoparticule (particule \textit{C}, \qty{800}{\nm} de largeur) est constituée de facettes de type \{111\}, \{110\} et \{100\}.

\begin{figure}[!htb]
    \centering
    \includegraphics[width=0.32\textwidth]{/home/david/Documents/PhD/PhDScripts/SixS_2021_06_BCDI_NH3/reconstructions/vti_good_files/D6/300/300_ammonia_100.png}
    \includegraphics[width=0.32\textwidth]{/home/david/Documents/PhD/PhDScripts/SixS_2021_06_BCDI_NH3/reconstructions/vti_good_files/D6/300/300_ammonia_010.png}
    \includegraphics[width=0.32\textwidth]{/home/david/Documents/PhD/PhDScripts/SixS_2021_06_BCDI_NH3/reconstructions/vti_good_files/D6/300/300_ammonia_001.png}
    \caption{
        Surface de la particule C reconstruite à \qty{300}{\degreeCelsius} sous atmosphère d'argon (inerte).
        La surface est colorée par les valeurs de la déformation hétérogène hors plan pour chaque voxel de surface, la limite de chaque facette est délimitée par d'épaisses lignes blanches.
        Mesures effectuées en diffraction cohérente des rayons x.
    }
    \label{fig:D6FacetsFr}
\end{figure}

La particule \textit{C} a montré une diminution/augmentation réversible de la déformation homogène lorsque l'ammoniac a été introduit/retiré du réacteur, et n'a pas pu être reconstruite sans la présence d'ammoniac.
Une augmentation non réversible de la déformation hétérogène a été mesurée à \qty{300}{\degreeCelsius} lors de la première exposition de l'échantillon à des conditions de réaction (\ce{O_2}/\ce{NH_3} = 0,5), aucune augmentation de ce type n'a été reproduite aux mêmes conditions à \qty{400}{\degreeCelsius}.

L'augmentation de déformation homogène liée à la présence/absence d'ammoniac n'a pas été reproduite sur la deuxième nanoparticule (particule \textit{B}, \qty{300}{\nm} de large).
Contrairement à la particule précedente (\textit{C}), la particule \textit{B} à des facettes de type \{113\} également présentes à sa surface.

\begin{figure}[!htb]
    \centering
    \includegraphics[width=0.32\textwidth]{/home/david/Documents/PhD/PhDScripts/SixS_2021_06_BCDI_NH3/reconstructions/vti_good_files/B7/RT/RT_100_B7.png}
    \includegraphics[width=0.32\textwidth]{/home/david/Documents/PhD/PhDScripts/SixS_2021_06_BCDI_NH3/reconstructions/vti_good_files/B7/RT/RT_010_B7.png}
    \includegraphics[width=0.32\textwidth]{/home/david/Documents/PhD/PhDScripts/SixS_2021_06_BCDI_NH3/reconstructions/vti_good_files/B7/RT/RT_001_B7.png}
    \caption{
        Surface de la particule B reconstruite à \qty{25}{\degreeCelsius} sous atmosphère d'argon (inerte).
        La surface est colorée par les valeurs de la déformation hétérogène hors plan pour chaque voxel de surface, la limite de chaque facette est délimitée par d'épaisses lignes blanches.
    }
    \label{fig:B7FacetsFr}
\end{figure}

Une augmentation non réversible de la déformation hétérogène a également été mesurée à \qty{300}{\degreeCelsius}, mais induite par la présence d'ammoniac, sans oxygène.
La possibilité d'avoir un échantillon oxydé peut expliquer cette évolution sous atmosphère réductrice, mais n'est pas observée pour la particule \textit{C}.
L'apparition d'un défaut à \qty{400}{\degreeCelsius} est liée à une augmentation non réversible et importante de la déformation homogène lors de l'oxydation de l'ammoniac, qui continue d'augmenter en fonction du rapport entre ammoniac et oxygène.
Cette évolution structurelle est clairement visible en 3D, avec des volumes de densité électronique de Bragg manquants (ann. \ref{app:AppendixBCDIOperando}).
La présence de défauts pourrait jouer un rôle important dans le champ de déformation du catalyseur et donc dans ses propriétés catalytiques.
Il est possible que cette transformation correspond à l'origine de la transformation observée dans les catalyseurs industriels.
Une activation globale du catalyseur à \qty{300}{\degreeCelsius} et \qty{400}{\degreeCelsius} a été rapportée tandis que le rapport \ce{O_2}/\ce{NH_3} était maintenu égal à 2, avec une production accrue de \ce{NO}, et dans une moindre mesure \ce{N2O}, au détriment de \ce{N2}.

En conclusion, il est primordial de sonder différentes nanoparticules avant de tirer une conclusion sur leur comportement global lors d’une réaction catalytique hétérogène.

Pour mieux comprendre le rôle de chaque facette dans le comportement des nanoparticules de Pt, des expériences de diffraction de surface et de spectroscopie photoélectronique à rayons X ont été réalisées à \qty{450}{\degreeCelsius} sur des monocristaux de Pt(111) et Pt(100).
Les échantillons ont d'abord été exposés à une atmosphère riche en oxygène (\qty{80}{\milli\bar}) pour oxyder les surfaces de Pt (pression totale toujours maintenue à \qty{500}{\milli\bar} grâce à l'utilisation d'argon).
Deux rapports oxygène/ammoniac différents sont utilisés après l'oxydation de surface, en introduisant d'abord \qty{10}{\milli\bar} d'ammoniac (\ce{O_2}/\ce{NH_3} = \num{8}), puis en réduisant la pression d'oxygène à \qty{5}{\milli\bar} (\ce{O_2}/\ce{NH_3} = \num{0.5}).
L'expérience de spectroscopie photoélectronique à rayons X a été réalisée avec les mêmes rapports ammoniac/oxygène, mais à des pressions partielles plus faibles (environ \qty{10}{\percent}), et sans garder la pression totale constante.

La présence d'oxydes de platine est surveillée afin de comprendre son importance dans le mécanisme réactionnel, ainsi que dans d'éventuelles reconstructions du catalyseurs.
Les différents résultats sont résumés dans le tableau \ref{tab:RecapSXRDFr}.

\begin{table}[!htb]
\centering
\resizebox{\textwidth}{!}{%
\begin{tabular}{@{}llll@{}}
    \toprule
    \ce{O2} (\unit{\milli\bar}) & \ce{NH3} (\unit{\milli\bar}) & Pt(111) & Pt(100) \\
    \midrule
    0 & 0 & Pas de couche d'oxyde / pas de reconstruction & Pas de couche d'oxyde / pas de reconstruction\\
    80 & 0 & Deux superstructures hexagonales tournés & Couche épaisse épitaxiée de \\
     & & Pt(111)-($6\times6$)-R\ang{\pm8.8} (mono-couche) & \ce{Pt3O4}, Pt(100)-($2\times2$) \\
     & & Après \qty{9}{\hour}\qty{30}{\minute}: superstructure & et signaux décalés en H ou K \\
     & & Pt(111)-($8\times8$) (multicouche) & (mono-couche) \\
    \midrule
    80 & 10 & Pas de couche d'oxyde / pas de reconstruction & Pt(100)-($10\times10$) reconstruction\\
     & & Signaux faibles au niveau O 1s & (multicouche) \\
     & & & et signaux décalés en H ou K \\
     & & & (mono-couche) \\
     & & & Sélectivité vers \ce{NO} supérieure \\
     & & & Signal important au niveau O 1s en \\
     & & & comparaison avec Pt(111) \\
    10 & 10 & Pas de couche d'oxyde / pas de reconstruction & Structure Pt(100)-Hex (mono-couche) \\
     & & Aucun signal au niveau O 1s & Aucun signal au niveau O 1s \\
     & & \ce{N_a} et \ce{NH_{3,a}} au niveau N 1s & Seulement \ce{N_a} au niveau N 1s \\
    \midrule
    0 & 10 & Pas de couche d'oxyde / pas de reconstruction & Suppression progressive de la structure Pt(001)-Hex \\
    0 & 0 & Pas de couche d'oxyde / pas de reconstruction & Pas de couche d'oxyde / pas de reconstruction\\
    \midrule
    5 & 0 & Retour aux mêmes structures qu'à \qty{80}{\milli\bar} & Structures transitoires, différentes des signaux \\
       & & de \ce{O2}, mais avec une cinétique diminuée & observée à \qty{80}{\milli\bar} de \ce{O2} \\
    \bottomrule
\end{tabular}%
}
\caption{Résumé des structures de surface identifiées par diffraction de surface et des changements pertinents dans les niveaux mesurés par spectroscopie photoélectronique des rayons X, combinées à des mesures par spectrométrie de masse.}
\label{tab:RecapSXRDFr}
\end{table}

L'oxydation des deux surfaces est d'abord discutée.
Une couche épaisse de \ce{Pt_3O_4} a été identifiée sur Pt(100), dans un arrangement Pt(100)-($2\times2$) et d'épaisseur moyenne égale à \qty{16}{\angstrom} (\textit{i.e.} 3 cellules unitaires).
Des signaux décalés dans l'espace réciproque par rapport aux pics \ce{Pt_3O_4} sont également mesurés.
Les signaux liés au \ce{Pt_3O_4} et décalés ont été mesurés \qty{1}{\hour} après l'introduction de l'oxygène, mais n'ont pas pu être détectés sous une atmosphère réduite en oxygène (\qty{5}{\milli\bar}), même après plusieurs heures.
Des signaux transitoires sont en effet mesurés, soulignant l’importance de la pression partielle d’oxygène dans l’oxydation de surface.

Le \ce{Pt3O4} n'a pas été observé sur Pt(111).
Néanmoins, une superstructure commensurable Pt(111)-($8\times8$) a été clairement identifiée après \qty{9}{\hour}\qty{30}{\minute} de temps écoulé sous une atmosphère à forte teneur en oxygène, et après \qty {23}{\hour}\qty{30}{\minute} sous une atmosphère réduite en oxygène.
À partir de la distribution d’intensité des signaux hors plan associés, il a été déterminé que cette structure avait quelques couches d’épaisseur.
Des signaux supplémentaires dans le plan sont également détectés dès que de l'oxygène est introduit dans le réacteur, à basse et haute pression d'oxygène, liés à une structure Pt(111)-($6\times6$)-R\ang{\pm 8.8}.
L'intensité de ces signaux diminue une fois que la structure Pt(111)-($8\times8$) est présente, ce qui soutient un lien précurseur entre les deux structures.
% De plus, lors de la description des signaux avec la matrice Pt(111)-p$\begin{pmatrix} 1.08 & -0.21 \\ -0.21 & 1.08 \end{pmatrix}$, ce qui donne un angle spatial réel de \ang{ 137.4}, un signal de second ordre est partagé.
Les signaux mesurés dans le niveau O 1s dans une atmosphère réduite en oxygène et attribués aux espèces d'oxygène de surface sont plus élevés pour le Pt (100) que pour le Pt (111), montrant que la surface du Pt (100) est plus facilement oxydée que le Pt (111).
L'oxydation des surface Pt(111) et du Pt(100) est aussi liée à une rugosité de surface accrue et à une déformation compressive hors plan par rapport à l'atmosphère inerte pour le Pt(111).

Différents comportements ont été mesurés sur les deux surfaces lors de conditions de réaction.
Sous un rapport \ce{O_2}/\ce{NH_3} élevé, qui favorise la production de \ce{NO}, les oxydes de surface sont directement éliminés du Pt(111), mais reconstruits avec un arrangement (10x10) sur Pt(100) .
Les signaux des espèces oxygénées au niveau de O 1s sont faibles pour le Pt(111), difficiles à dissocier du fond, alors que les pics sont clairement détectés pour le Pt(100).
La différence de présence d’oxygène en surface est déjà observée lors de la précédente oxydation de surface.
La reconstruction de la surface du Pt(100) est liée à la présence persistante d'espèces oxygénées de surface lors de l'oxydation de l'ammoniac, différant ainsi de la surface du Pt(111).
De plus, la sélectivité envers \ce{NO} est augmentée pour le Pt(100), également liée à la présence plus importante d'oxygène en surface, qui est prédite comme cruciale dans la production de \ce{NO} lors du mécanisme réactionnel \parencite{NovellLeruth2005, Offermans2006, Offermans2007, Imbihl2007, NovellLeruth2008}.
Une rugosité de surface importante est également observée dans ces conditions pour les deux surfaces.

Le Pt(111) et Pt(100) montrent tout deux une sélectivité similaire vers \ce{N_2} en abaissant le rapport \ce{O_2}/\ce{NH_3} à \num{0,5}.
L'ampleur de la déformation hors plan (relaxation de surface) et la rugosité de la surface sont déjà réduites dans ces conditions pour le Pt(111), mais restent à des valeurs similaires pour Pt (100).
Sur les deux surfaces, les espèces d'oxygène sont absentes du niveau O 1 et l'azote atomique adsorbé est mesuré.
L'ammoniac adsorbé est également observé sur le Pt(111).

La différence de déformation hors plan (relaxation de surface) lors de la diminution du rapport \ce{O2}/\ce{NH3} de \num{8} à \num{0,5} est du même ordre de grandeur, mais de nature différente, environ \qty{0,06}{\percent} de déformation en tension / compression sur le Pt(100) / Pt(111).
Plus important encore, il est clair que la contrainte est contenue dans les couches supérieures des catalyseurs au platine.
Si l'on considère un voxel de surface Pt(100) de \qty{10}{\nm} d'épaisseur, la déformation moyenne correspondante du voxel serait égale à \qty{0,002}{\percent} par rapport au réseau global.
Par conséquent, le changement de déformation entre les conditions de réaction devient difficile à résoudre, ce qui peut expliquer pourquoi aucune différence n'est observée sur les particules \textit{C} lors de l'oxydation de l'ammoniac.
Un autre facteur à prendre en compte est l'augmentation de la rugosité de surface liée à la forte pression d'oxygène, qui a pour effet de diminuer l'intensité des photons diffusés loin du pic de Bragg, et va ainsi réduire la résolution expérimentale.
Des études supplémentaires à haute résolution pourraient contribuer à comprendre l’effet des adsorbats sur la relaxation de surface.

La thèse démontre l’importance de combiner les mesures de particules individuelles et des assemblages de particules.
Par imagerie de diffraction cohérente de Bragg, nous avons démontré qu'en fonction de la morphologie (forme, taille, type de facette et couverture, \textit{etc.}) et de l'état de déformation initial des particules de Pt, différentes évolutions structurelles sont observées (en termes de déformation, de morphologie et de défauts) lors de l'oxydation de l'ammoniac.
Une meilleure représentation des différents comportements suivis par les particules de platine au cours de la réaction, en mesurant plusieurs particules individuelles, semble indispensable pour une compréhension globale de la relation structure-activité.
Il est à noter que la déformation à l’interface particule/support semble exercer une influence prononcée sur le comportement des particules.
Les nanoparticules servent de plate-forme pour explorer simultanément l’évolution structurelle de diverses facettes cristallographiques, alors que les monocristaux permettent de mieux isoler les comportements des facettes uniques sur les nanoparticules.

En conclusion, l'utilisation de trois techniques de rayons x différentes a permis de mieux appréhender les mécanismes à l'oeuvre lors de l'oxydation de l'ammoniac.

%************************************
\vspace{\fill} % ALIGNER EN BAS DE PAGE
%************************************

\newpage\thispagestyle{empty}\null\newpage

    \frontmatter
    {\hypersetup{linkcolor=black}
        \tableofcontents
        % \listoffigures
        % \listoftables
    }
    \newpage

    \mainmatter
    \chapter{Introduction}
    The use of catalysts has several advantages such as faster, selective, and more energy-efficient chemical reactions, directed towards producing higher amounts of the desired product while reducing unwanted byproducts.
Over the years, scientists have developed specialised catalysts for various real-world applications, today \qty{90}{\percent} of chemical processes involve catalysts in at least one of their steps \parencite{Weiner1998, DeVries2012}.
Notable advancements in catalysis have led to the production of biodegradable plastics, novel pharmaceuticals, and eco-friendly fuels and fertilisers \parencite{Fechete2012}.

\begin{figure}[!htb]
    \centering
    \includegraphics[width=\textwidth]{/home/david/Documents/PhD/Figures/ammonia/ParisNO2English.png}
    \caption{
        $NO_2$ levels in Paris near the main traffic roads remained on average twice superior to the annual limit of \qty{40}{\ug \per \m^3} \parencite{AirParis} between 2012 and 2018 despite a global decrease of since 2012.
    }
    \label{fig:NO2Paris}
\end{figure}

Several new challenges have emerged in the field of catalysis related to improving efficiency, reducing environmental impact, and developing sustainable processes.
First, environmental challenges concern minimising and/or managing by-products, reducing contamination in effluents/wastewaters, and using sustainable sources of raw materials \parencite{Ludwig2017, Lange2021} and energy supplies.
Secondly, economical challenges imply using cheaper, readily available raw materials, increased productivity, and decreased lag-time between discovery to commercialisation \parencite{Keisuke2019, Gunay2021}.
For example, recent studies suggest that alternative, more economical catalysts, such as non-noble metals \parencite{Zhong2021} and other derived metal-based compounds, need to be tested as possible substitutes for the most frequently used noble metals, which are very efficient but expensive.

Catalysis also has a role to play to combat pollution and create cleaner energy with for example the development of efficient water-splitting technologies \parencite{Ahmad2015}, and enhancing the use of biomass and other energy vectors such as ammonia \parencite{Fang2022}.
Finally, challenges also arise in automotive exhaust where catalysts participate in the reduction of the emissions of toxic gases and nanoparticles \parencite{WHOAirPollution, Heck2001, Gandhi2003}.
Some of the major air pollutants such as nitrogen oxides, (\ce{NO_x}), and particulate matter (PM) are emitted by road traffic (\qty{65}{\percent} of \ce{NO_x}, \qty{\approx 35}{\percent} of PM), mainly by diesel vehicles, and directly inhaled by nearby major city inhabitants.
To set a striking example, in Paris in 2018, 700 000 inhabitants were exposed to \nitrogendioxide concentrations exceeding the regulations (fig. \ref{fig:NO2Paris}), 60 000 inhabitants for PM$_{10}$, and all Parisians were concerned by exceeding the World Health Organisation (WHO) recommendations for PM$_{2.5}$ \parencite{AirParis}.
Air quality is the main environmental concern of Ile-de-France residents (\qty{65}{\percent} of total mentions) ahead of climate change (\qty{63}{\percent}) and food (\qty{38}{\percent} - \cite{AirParis}).

\section{The oxidation of Ammonia}

\subsection{The Haber-Bosch and Ostwald processes}

\epigraph{Today, about \qty{50}{\percent} of the world population relies on nitrogen-based fertilisers to produce the food necessary to their alimentation.}%{\textit{Nature Catalysis editorial, \cite*{NatureEditorial2018}.}}

The story of ammonia begins in the early $20^{th}$ century with the discovery in 1902 of the Ostwald process that permitted the synthesis of fertilisers from ammonia.
Wilhelm Ostwald later received the Nobel price in 1909 \textit{"in recognition of his work on catalysis and for his investigations into the fundamental principles governing chemical equilibria and rates of reaction"}.
Seven years later, in 1909, Fritz Haber designed a process for the synthesis of ammonia which was later improved by Carl Bosch.
It is today known as the Haber-Bosch process and is at the origin of the mass production of ammonia using metallic catalysts \parencite{Hosmer1917, Parsons1919}.
Fritz Haber received the Nobel prize in 1918 \textit{"for the synthesis of ammonia from its elements"} \parencite{Alexander1920} and Carl Bosch in 1931 \textit{“in recognition of his contributions to the invention and development of chemical high pressure methods”}.

\begin{figure}[!htb]
\centering
    \begin{tikzpicture}
        \node (image) [anchor=south west, inner sep=0pt] {\includegraphics[width=0.95\textwidth]{/home/david/Documents/PhD/Presentations/Slides/PhdSlides/Figures/worldindata/world-population-with-and-without-fertilizer.png}};
        \begin{scope}[x={(image.south east)}, y={(image.north west)}]
            \node [text width=4cm,align=right] at (0.18, 0.7) (OP) {\textcolor{Ostwald}{Ostwald} process\\(1902)\\ \textrightarrow Nobel prize\\(1909)};
            \draw [-latex, ultra thick, Ostwald] (0.10,0.75) to (0.10,0.12);
            \node [text width=4cm,align=right] at (0.19, 0.45) (HBP) {\textcolor{Haber}{Haber-Bosch}\\process (1908)\\ \textrightarrow Nobel prize\\(1918)};
            \draw [-latex, ultra thick, Haber] (0.13,0.4) to (0.13,0.12);
        \end{scope}
    \end{tikzpicture}
    \caption{
    Since the discovery of the Ostwald and Haber-Bosch process that allowed the mass production of nitrogen-based fertilisers, the world population increase has been relying on their production and use for agriculture.
    Today, about \qty{50}{\percent} of the worlds population relies on nitrogen-based fertiliser to produce the food necessary to their alimentation.
    Taken from \cite{WorldDataFertilizer}.
    }
    \label{fig:FertilizerWID}
\end{figure}

The oxidation of ammonia can be described by three equations that, depending on the stoechiometric ratio between \ammonia and \dioxygen, have different products.

\begin{align}
    \label{eq:AmmoniaOxidationNitrogen}
    4 \ammonia (g) + 3 \dioxygen (g) & \rightarrow 6 \water (g) + 2 \nitrogen (g) \\
    \label{eq:AmmoniaOxidationNitrousOxide}
    4 \ammonia (g) + 4 \dioxygen (g) & \rightarrow 6 \water (g) + 2 \nitrousoxide (g) \\
    \label{eq:AmmoniaOxidationNitricOxide}
    4 \ammonia (g) + 5 \dioxygen (g) & \rightarrow 6 \water (g) + 4 \nitricoxide (g)
\end{align}

The first equation (eq. \ref{eq:AmmoniaOxidationNitrogen}) yields nitrogen (\nitrogen), a naturally occurring gas that does not pollute the environment nor shows any toxic behaviour towards humans (tab. \ref{tab:NitrogenGases}).
The second equation (eq. \ref{eq:AmmoniaOxidationNitrousOxide}) yields nitrous oxide (\nitrousoxide), a powerful greenhouse effect gas (tab. \ref{tab:NitrogenGases}).
The third equation (eq. \ref{eq:AmmoniaOxidationNitricOxide}) yields nitric oxide, also called nitrogen monoxide (\nitricoxide), which is the main desired product for the subsequent production of nitrogen-based fertilisers (tab. \ref{tab:NitrogenGases}) \textit{via} the synthesis of nitric acid with the Ostwald process.

Known side reactions to the oxidation of ammonia are the recombination of nitric oxide with unreacted ammonia that leads to the production of water and nitrogen (eq. \ref{eq:SideReactions1}), and the thermal decomposition of nitric oxide (eq. \ref{eq:SideReactions2}), that both lower the total yield of the reaction when aiming at the production of nitric oxide.
It is also possible to observe the dissociation of ammonia resulting in the production of \ce{N_2} and \ce{H_2} (eq. \ref{eq:SideReactions3}).

\begin{align}
    \label{eq:SideReactions1}
    4 \ammonia (g) + 6 \nitricoxide (g) & \rightarrow 5 \nitrogen (g) + 6 \water (g)\\
    \label{eq:SideReactions2}
    2 \nitricoxide (g) & \rightarrow \nitrogen (g) + \dioxygen (g)\\
    \label{eq:SideReactions3}
    2 \ammonia (g) & \rightarrow \nitrogen (g) + 3 \ce{H_2} (g)
\end{align}

Reactions \ref{eq:AmmoniaOxidationNitrogen}, \ref{eq:AmmoniaOxidationNitrousOxide}, \ref{eq:AmmoniaOxidationNitricOxide} and \ref{eq:SideReactions1} are strongly exothermic, leading to a significant increase of the catalyst temperature during the reaction \parencite{Hatscher2008}.

The second (eq. \ref{eq:Stage2}) and third (eq. \ref{eq:Stage3}) stages of the Ostwald process stem from the production of \ce{NO} \textit{via} the first stage, \textit{i.e.} the oxidation of ammonia (eq. \ref{eq:AmmoniaOxidationNitricOxide}).

\begin{align}
    \label{eq:Stage2}
    2 \nitricoxide (g) + \dioxygen (g) & \rightarrow \nitrogendioxide (g) \\
    \label{eq:Stage3}
    3 \nitrogendioxide (g) + \water (l) & \rightarrow 2\nitricacid (aq) + \nitricoxide (g)
\end{align}

Nitrogen dioxide (\nitrogendioxide) is produced from \ce{NO} which then reacts with water to form nitric acid (\nitricacid), an important actor in multiple industrial process (tab. \ref{tab:NitrogenGases}), including fertilisers such as ammonium nitrate (eq. \ref{eq:AmmoniumNitrate}).

\begin{align}
    \label{eq:AmmoniumNitrate}
    \nitricacid + \ammonia & \rightarrow \ammoniumnitrate
\end{align}

\section{The importance of heterogeneous catalysis}\label{sec:AmoOxiHC}

\subsection{Industry conditions and catalysts}

Important selectivity towards the production of \ce{NO} must be achieved in the first stage of the Ostwald process when aiming at the production of fertilisers.
However, at low temperature or in the absence of catalyst, the production of \ce{N_2} is favoured (eq. \ref{eq:AmmoniaOxidationNitrogen}), which is at the origin of the selective, heterogeneous catalytic oxidation of ammonia towards the production of nitrogen oxide.

Since the 1930s, the presence of a platinum-rhodium (\qty{\approx 10}{\percent} Rh) knitted gauzes catalyst at favourable reaction conditions (e.g. \qty{900}{\degreeCelsius} - \qty{5}{\bar} - excess \dioxygen) allowed a \qty{98}{\percent} \ce{NO} yield to be achieved \parencite{Ostwald1908, Kaiser1909, Handforth1934, Heck1982}.

In order to understand the reaction mechanism occuring on the catalyst surface, and the role of Pt and Rh in the catalyst stability and selectivity, scientist have been studying the reaction by different methods since the beginning of last century.
A comprehensive review of the ammonia oxidation is given by Hatscher et al. \parencite*{Hatscher2008}, important literature findings from the last 100 years will be resumed below.

% roughening, oxides
The importance of high temperature for the selective production of \ce{NO} was demonstrated early by temperature dependant studies\parencite{Nutt1968, Pignet1974, Li1997}, while a restructuration of the catalyst was put into evidence \parencite{McCabe1974, FlytzaniStephanopoulos1979, McCabe1986} by the means of \textit{ex-situ} SEM imaging of industrial samples (fig. \ref{fig:Gauzes}).
This \textit{activation process} leads to an increase selectivity towards \ce{NO} while the gauzes undergo a transformation towards a roughened surface.

\begin{figure}[!htb]
    \centering
    \includegraphics[width=0.49\textwidth]{/home/david/Documents/PhD/Figures/sample/EtchedGauze.png}
    \includegraphics[width=0.49\textwidth]{/home/david/Documents/PhD/Figures/sample/ReconstructedGauze.png}
    \caption{
    SEM images of Pt-Rh reconstructed gauzes after use in industry, taken from Bergene et al. \parencite*{Bergene1996}.
    The horizontal bar is \qty{0.1}{\mm} wide.
    }
    \label{fig:Gauzes}
\end{figure}

The deactivation of Pt-Rh industrial catalysts after long exposure times have been explained by the presence of rhodium oxides \parencite{Bergene1996}.
However, a deactivation process was also reported on pure Pt calatyst \parencite{Ostermaier1974}, linked to the presence of platinum oxides \parencite{Ostermaier1976}.

The role of volatile surface oxides in the roughening and etching of the catalyst surface was theorised by Wei et al. \parencite*{Wei1996} and confirmed experimentally by Nilsen et al. \parencite*{Nilsen2001}.
The transportation of Pt and Rh was found to be permitted by oxygen \parencite{Hannevold2005a}, decreasing with increasing Rh content and decreasing oxygen pressure, only for high temperature gradient, \textit{i.e.} \qtyrange{1400}{800}{\degreeCelsius} over \qty{100}{\mm}.
\textit{Ex-situ} characterization of Pt and Pt-Rh industrial gauzes by x-ray powder diffraction and electron microscopy allowed the identification of defect sites at the origin of high temperature gradient areas during reaction.
The existence of such areas allows the restructuration process to occur nearby initial defects by the formation of \ce{PtO_2} and \ce{RhO_2} oxides, depositing metallic atoms on colder regions, and also leading to some loss of the precious metals constituting the calatyst \parencite{Hannevold2005}.

The progressive deactivation of the catalyst due to the ever increasing presence of \ce{Rh_2O_3} \parencite{McCabe1986} was refuted by Hannevold et al. \parencite*{Hannevold2005}, explaining that such oxide could only form during the cooling of the catalyst, which is a good example of the limitation of \textit{ex-situ} works.
A \ce{Pt_3O_4} catalyst used for the oxidation of ammonia was proven to be unstable at working temperature above \qty{690}{\degreeCelsius}, decomposing into a Pt phase after \qty{7}{\hour} of operation \parencite{Zakharchenko2001}.

\subsection{Reaction mechanism}

First studies performed at low working pressure but different temperatures have supported a Langmuir-Hinshelwood mechanism during which both reactants are decomposed and adsorbed on top of the catalyst surface \parencite{Nutt1969, Pignet1974, Ostermaier1974, Pignet1975, Gland1978a}, the final pathway towards the production of nitrogen or nitrous oxide depending on the ratio between \ce{O_2} and \ce{NH_3}.

The importance of adsorbed atomic oxygen and \ce{OH} in the dissociation of \ce{NH_3} on both Pt (111) \parencite{Mieher1995, Bradley1997} and Pt (100) surface was put into evidence \parencite{Bradley1995, vandenBroek1999, Kim2000}.
High oxygen coverage on Pt (100) has shown to reduce the production of \ce{N_2}, mainly produced by the dissociation of \ce{NO}, favoured at lower temperatures over \ce{NO} desorption \parencite{Bradley1995}.
Asscher et al. \parencite*{Asscher1984} have reported the existence of a mechanism involving oxygen in the gas phase reacting with adsorbed NH species on Pt (111) for the production of \ce{NO}.
At UHV, a rotated hexagonal reconstruction on top of Pt (100) was found to impinge \ce{NH_3} dissociation at low temperature (\qty{-123}{\degreeCelsius}) \parencite{Bradley1997}, the ammonia oxidation stabilising the (1x1) phase \parencite{Rafti2007} by the formation of \ce{NH_x} intermediates between \qty{125}{\degreeCelsius} and \qty{350}{\degreeCelsius}.

% steps
The importance of atomic steps in the catalytic activity was first revealed by Gland et al. \parencite*{Gland1978, Gland1980}.
The more recent study of the oxidation of ammonia on several model catalysts (Pt(533), Pt(443), Pt(865), Pt(100), Pt foil) to investigate the structure selectivity by Yingfeng \parencite*{Yingfeng2008} linked the presence of steps and kinks with higher catalytic activity in the \qtyrange{1e-9}{1e-5}{\bar} range.
Moreover, the production of \ce{N_2} was reported to be promoted by lower temperatures and a reduced \ce{O_2}/\ce{NH_3} ratio in the incoming gas flow, while higher temperatures and an elevated \ce{O_2}/\ce{NH_3} ratio tend to result in a higher selectivity towards the formation of \ce{NO} \parencite{Zeng2009}
No production of \ce{N_2O} was observed within the studied pressure range.

% tap, kinetic studies
Rebrov et al. \parencite*{Rebrov2002} detailed the reaction kinetics and mechanism with a 13 step, temperature dependent model, the parameters of which have been refined with data collected from a microreactor.
A wide range of conditions was explored, including ambient pressures of \ce{NH3} (\qtyrange{0.01}{0.12}{\bar}) and \ce{O_2} (\qtyrange{0.10}{0.88}{\bar}), in a large temperature range (\qtyrange{250}{400}{\degreeCelsius}).
A dual site mechanism was proposed, with a preference for hollow site for oxygen species, and for top or bridge sites for nitrogen species, while \ce{NO} species have been reported to exist on both top and bridge sites.

Pérez-Ramirez et al. \parencite*{PerezRamirez2004} have also attempted to study the kinetics of the reaction by directly analysing the selectivity of catalysts used in industry (\textit{i.e.} Pt-Rh and Pt gauzes) above \qty{700}{\degreeCelsius}, by the means of reactant gas pulses.
The heterogeneous catalysis reaction mechanism was found to be similar on both sample, a pre-exposition of the catalysts to oxygen facilitated the dissociation of \ce{NH_3}, without which the decomposition towards \ce{N_2} was not detected.
A high coverage of the catalyst by oxygen was linked to the formation of \ce{NO}, strongly bound oxygen-rich species favour the production of \ce{N_2}, whereas weakly bound were associated to \ce{NO} selectivity.
The importance of adsorbed oxygen species to prevent spontaneous \ce{NO} dissociation was confirmed.
In a second study \parencite{PerezRamirez2009}, they proved that increasing the \ce{O_2}/\ce{NH_3} ratio to 10 pushed \ce{NO} selectivity to almost \qty{100}{\percent}, \ce{N_2} and \ce{N2_O} production being both suppressed by favouring \ce{NO} desorption.

A similar study combining dynamical field theory (DFT) calculations and temporal analysis of products (TAP) by Baerns et al. \parencite*{Baerns2005} confirmed the importance of surface \ce{O} and \ce{OH} for the dehydrogenation of \ce{NH_3} on the calatyst surface.
\ce{N_2O} could only be detected at ambient pressures, \ce{N_2} is favoured at lower temperature, whereas \ce{NO} is favoured at higher temperatures.
Importantly for the use of model catalysts, no difference in the temperature dependant production of \ce{N_2} and \ce{NO} between a Pt(533) single crystal, Pt foil, and knitted Pt gauzes could be observed.
Surface roughnening of the catalyst at ambient pressure was linked to activation and selectivity change for a Pt foil.

The first studies performed under working conditions brought forward two recurrent problems when carrying out low pressure studies.
First, the production of \ce{N_2O} was rarely discussed because undetected at low pressures, van den Broek et al. \parencite*{vandenBroek1999} hypothetised that \ce{NO} was a precursor in the production of nitrous oxide.
\ce{N_2O} was then detected also by Pérez-Ramirez et al. \parencite*{PerezRamirez2004}, Baerns et al. \parencite*{Baerns2005} and Kondratenko et al. \parencite*{Kondratenko2007} which confirmed the precursor role of \ce{NO}, and that \ce{N_2} was produced by either dissociation of \ce{NO} or/and complete dehydrogenation of \ce{NH_3}.

Secondly, the known roughening process occuring on the catalyst at working conditions could not be reproduced without long working times.
Kinetic studies on polycrystalline Pt up to \qty{10}{\bar} but at temperatures below \qty{385}{\degreeCelsius} observed the roughening transition \parencite{Kraehnert2008}, which proved difficult to fit with develop kinetic models, attributed to local increase in temperature and surface area.

% The rate-limiting step in the ammonia decomposition has been identified to be the recombination of nitrogen atoms on Fe catalysts \parencite{Vilekar2012}

% DFT
\begin{figure}[!htb]
    \centering
    \includegraphics[width=\textwidth]{/home/david/Documents/PhD/Figures/ammonia/ReactionMechanism.png}
    \caption{
    Example of \ce{NH_3} stripping process by adsorbed \ce{O}, taken from Imbihl et al. \parencite*{Imbihl2007}.
    }
    \label{fig:ReactionMechanism}
\end{figure}

Novell-Leruth et al. \parencite*{NovellLeruth2005} confirmed the adsorption of \ce{NH_3} and \ce{NH_2} to occur respectively on top and bridge sites for both Pt (100) and Pt(111), but with a more favourable adsorption process on Pt (100).
Similar adsorption energies have been reported for \ce{NH} and \ce{N} that both adsorb on hollow sites \
Additionnal DFT studies of the reaction pathways and kinetics have confirmed a reaction mechanism following a step-by-step decomposition of ammonia on the catalyst surface, the stripping of adsorbed ammonia hydrogen atoms being facilitated by the presence of oxygen species \parencite{Offermans2006}.
A comparative study between the Pt (100), Pt (111) and Pt (211) (\textit{i.e.} stepped) surfaces did not reveal a strong structure sensitivity during the \ce{NH_3} stripping process \parencite{Offermans2007}.

Imbihl et al. \parencite*{Imbihl2007} have further improved the understanding of the production of \ce{N_2O}, which happens not only by the recombination of two adsorbed \ce{NO} species but also \textit{via} the reaction between adsorbed \ce{NO} and \ce{NH_x} species.

The first dehydrogenation step was found to be the slowest, while the desorption of \ce{NO} is the rate limiting-step \parencite{NovellLeruth2008}.
Interestingly, for the first time different oxygen species were reported for the dehydrogenation process, respectively \ce{O} for Pt(111) (fig. \ref{fig:ReactionMechanism}) and \ce{OH} for Pt(100).
A high energy barrier for \ce{NO} desorption and \ce{N_2O} formation explain the high temperature needed for their production in comparison with \ce{N_2}.

% conclude
Overall, the mechanism of the oxidation of ammonia is still unclear, different studies reporting different mechanisms that depend on the reactant partial pressures, as well as the catalyst temperature and structure.
The main drawback often common to the large amount of studies performed in the last and current century is that they have been performed at pressures and temperatures far from industrious conditions, even at UHV, sometimes \textit{ex-situ}.

It is of key importance to perform additional studies that will help build a comprehensive understanding of the reaction mechanics from low pressure and temperature to working conditions.

\section{Environmental impact}

\subsection{Greenhouse effect}

As illustrated in fig. \ref{fig:FertilizerWID}, nitrogen-based fertilisers have permitted an industrial development of agriculture.
In a remarkable review of the presence of \ce{N_2O} in the atmosphere linked to human activities, Pérer-Ramirez et al. \parencite*{PerezRamirez2003} have shown that the most important contribution to nitrous oxide in the atmosphere is from unused volumes of nutrients (fig. \ref{fig:PieGreenhouseNO2}), but also from nitric acid manufacture.
Understanding and limiting the by-products of \ce{N_2O} by controlling the process selectivity is thus capital.

For an equal amount of \ce{N_2O} and \ce{CO_2}, the amount of \ce{N_2O} will trap 298 times more heat than the amount of \ce{CO_2} over the next 100 years \parencite{MITCLIMATE}, responsible for \qty{6.2}{\percent} of the total U.S.A. greenhouse gases emissions in 2021 (fig. \ref{fig:PieGreenhouseNO2}).
\ce{N_2O} has an atmospheric lifetime of 114 years \parencite{Solomon2007}.

Davidson et al. \parencite{Davidson2009} have detailed the contribution of manure, fertilizer and nitrogen oxide to atmospheric nitrous oxide, revealing that a steady growth in its production is still expected in the years to come.
Fertilisers can be either ammonium- or nitrate-based, when plants don't fully absorb all the nutrients, a series of microbe-mediated transformations occur.
These processes lead to the release of nitrogen back into the atmosphere, primarily as nitrogen gas (\ce{N_2}) and, to a lesser extent, as \ce{N_2O}.
Moreover, nitrogen-based fertilisers and other human activities can lead to nitrogen runoff into water bodies, contributing to eutrophication (excessive growth of algae) and causing harm to aquatic ecosystems (tab. \ref{tab:NitrogenGases}).
Approximately half of the production of ammonia is lost to the environment \parencite{Erisman2007}.

\begin{figure}[!htb]
    \centering
    \includegraphics[width=\textwidth]{/home/david/Documents/PhD/Presentations/Slides/PhdSlides/Figures/ammonia/NO2pie.pdf}
    \caption{
    Pie charts underlining the importance of different gas in the total US greenhouse gas emissions in 2021 (a) and the specific contribution of nitrogen-based fertilisers to the total \ce{N_2O} emissions in 2021.
    LULUCF means Land Use, Land-Use Change, and Forestry.
    Adapted from \cite{EPAGreenhouseGases}.
    }
    \label{fig:PieGreenhouseNO2}
\end{figure}

\subsection{Pollution}

Nevertheless, it can also be interesting to be able to remove \ce{NH_3} from the atmosphere, a colourless gas with a pungent odour, that irritates the eyes, nose, throat, and respiratory tract if inhaled in small amounts due to its corrosive nature and is poisonous in large quantities.
Ammonia also pollutes and contributes to the eutrophication and acidification of terrestrial and aquatic ecosystems, and forms secondary particulate matter (PM2.5) when combined with other pollutants in the atmosphere (tab. \ref{tab:NitrogenGases}).
In that case, the reaction must be tuned towards the production of \ce{N_2} which is the only non pollutant and toxic gas.

Finally, the important effect of nitrogen oxides (\ce{NO_x}, \textit{i.e.} \ce{NO} and \ce{N_2O}) on the environment has brought forward the necessity to control their emissions, especially from the exhaust of diesel engines that are responsible for \qty{65}{\percent} of their emissions.
The selective catalytic reaction (SCR) using urea or ammonia as reductant (tab. \ref{tab:NitrogenGases}) has proven to be effective and to reach \qty{95}{\percent} efficiency \parencite{MitsubishiSCR}.
However, there can be a subsequent problem of unreacted ammonia \textit{slipping} from the reaction, which is also an important subject of study \parencite{Thermofischer}.

Today, the ammonia oxidation is an essential catalytic reaction used in the production of nitrogen-based fertilisers and in environmental applications.
In both cases, particular focus is on two products of the reaction, namely, \ce{NO} and \ce{N_2}.
Depending on the application of the ammonia oxidation, the catalytic reaction must be tuned towards a specific product, this \textit{selectivity} is controlled by the reaction temperature, pressure, the  reactant ratio, and the type of catalyst.
To be able to drive the reaction, the impact of each parameter \textit{at relevant industrial conditions} on the product pressure must be studied.

Recently, ammonia has been investigated as an energy vector for hydrogen fuel cells, which has reignited the interest in understanding the complex system drawn by the many simultaneous reactions \parencite{Afif2016, Georgina2021}.

\begin{landscape}
\begin{table}[!htb]
\centering
\resizebox{\columnwidth}{!}{%
    \begin{tabular}{@{}l|l|lll|l|l|ll@{}}
    \toprule
    Formula & \ammonia & \nitrogen & \nitrousoxide & \nitricoxide & \nitrogendioxide & \nitricacid & \ammoniumnitrate & \urea \\
    \midrule
    Name & Ammonia & Nitrogen & \begin{tabular}[c]{@{}l@{}}Nitrous oxide,\\ Laughing gas\end{tabular} & \begin{tabular}[c]{@{}l@{}}Nitrogen oxide,\\ Nitric oxide\\ Nitrogen monoxide\end{tabular} & Nitrogen dioxide & Nitric acid & \begin{tabular}[c]{@{}l@{}}Ammonium\\ nitrate\end{tabular} & Urea \\
    Origin & \begin{tabular}[c]{@{}l@{}}Haber-Bosch\\ process\end{tabular} & \begin{tabular}[c]{@{}l@{}}Naturally present\\ in the atmosphere,\\ Ammonia \\ oxidation,\\ Selective catalytic\\ reaction (SCR)\end{tabular} & \begin{tabular}[c]{@{}l@{}}Ammonia \\ oxidation,\\ Emissions from\\ nitrogen-based\\ fertilisers\end{tabular} & \begin{tabular}[c]{@{}l@{}}Ammonia oxidation,\\ Anthropogenic sources \\ (combustion process, \\ industry, agriculture, ...)\\ Naturally produced \\ from lightning or\\ volcanoes\end{tabular} & \begin{tabular}[c]{@{}l@{}}Ostwald process (step 1)\\ Anthropogenic sources \\ (combustion process, \\ industry, agriculture, ...)\\ Naturally produced \\ from lightning or \\ volcanoes\end{tabular} & \begin{tabular}[c]{@{}l@{}}Ostwald process\\ (step 2)\end{tabular} & \begin{tabular}[c]{@{}l@{}}Nitric acid\\ neutralisation \\ with ammonia\end{tabular} & Ammonia \\
    Major use & \begin{tabular}[c]{@{}l@{}}Ostwald process\\ (fertilisers),\\ Direct use in soil,\\ Fuel,\\ Hydrogen \\ carrier,\\ Cooling, SCR \end{tabular} & \begin{tabular}[c]{@{}l@{}}Ammonia \\ production\end{tabular} & \begin{tabular}[c]{@{}l@{}}Medicine,\\ Propellant\end{tabular} & Production of nitric acid & Production of nitric acid & \begin{tabular}[c]{@{}l@{}}Fertiliser production,\\ Nitration (explosives,\\ dyes, ...),\\ Propellant,\\ Etching\end{tabular} & \begin{tabular}[c]{@{}l@{}}Fertilising \\ agent,\\ Explosives\end{tabular} & \begin{tabular}[c]{@{}l@{}}Fertilising \\ agent,\\ SCR to reduce\\ NOx into N2\end{tabular} \\
    Toxicity & \begin{tabular}[c]{@{}l@{}}Dangerous for\\ the environment\\ Toxic, \\ Corrosive\end{tabular} & \begin{tabular}[c]{@{}l@{}}Asphyxiation by\\ displacing O2\end{tabular} & \begin{tabular}[c]{@{}l@{}}Anaesthetic, \\ euphoric\end{tabular} & \begin{tabular}[c]{@{}l@{}}Oxidising, Corrosive,\\ Toxic\end{tabular} & \begin{tabular}[c]{@{}l@{}}Oxidising, Corrosive,\\ Toxic, Health hazard\end{tabular} & Oxidising, Corrosive &  &  \\
    Pollution & PM formation &  & \begin{tabular}[c]{@{}l@{}}Ozone\\ depletion\end{tabular} & \begin{tabular}[c]{@{}l@{}}Smog, acid rains, \\ ozone depletion,\\ Precursor to \ce{NO_2} in\\ the atmosphere\end{tabular} & \begin{tabular}[c]{@{}l@{}}Smog, acid rains, \\ ozone formation\end{tabular} & \begin{tabular}[c]{@{}l@{}}Decomposes towards\\ \ce{NO_2}\end{tabular} & \begin{tabular}[c]{@{}l@{}}Decomposes \\ into \ce{NO_2} when\\ used as fertiliser,\\ Eutrophication\end{tabular} & \begin{tabular}[c]{@{}l@{}}Decomposes \\ into NH3 when\\ used as fertiliser,\\ Eutrophication\end{tabular} \\
    \begin{tabular}[c]{@{}l@{}}Greenhouse \\ effect\end{tabular} &  &  & \begin{tabular}[c]{@{}l@{}}Very important\\ (298 \ce{CO_2} eq.)\end{tabular} &  &  &  & \begin{tabular}[c]{@{}l@{}}Nitrogen-based\\ fertiliser release\\ N2O in the \\ atmosphere\end{tabular} & \begin{tabular}[c]{@{}l@{}}Nitrogen-based\\ fertiliser release\\ N2O in the \\ atmosphere\end{tabular} \\ \bottomrule
    \end{tabular}%
    }
    \caption{
        Nitrogen based species involved in the oxidation of ammonia, the Haber-Bosch process, the Ostwald process or nitrogen-based fertilisers.
        Information compiled from various sources: \cite{Thiemann2000, Harrison2001, Baerns2005, Imbihl2007, Hatscher2008, Davidson2009, Resta2020a, Borodin2021, Pottbacker2022}.
    }
    \label{tab:NitrogenGases}
\end{table}
\end{landscape}

\section{From industry to model catalysis}\label{sec:LiteratureAmmonia}

\epigraph{"A long standing conundrum in the catalysis community emerged at the interface between surface science and heterogeneous catalysis, better known as the pressure and materials gap."}{\textit{Nature Catalysis editorial, \cite*{NatureEditorial2018}.}}

\begin{table}[!htb]
    \centering
    \begin{tabular}{l|l|l|l}
    \toprule
                & Pressure    & Material                         &     Temperature \\
    \midrule
    Industry   & \qtyrange{1}{12}{\bar} & Knitted gauzes wires   & \qtyrange{750}{900}{\degreeCelsius} \\
               &              & (diameter \qty{\approx 80}{\um}) & \\
    \midrule
    Literature & UHV, mbar    & Single crystals                  & \qtyrange{25}{900}{\degreeCelsius} \\
    \midrule
    This study & Near ambient & Single crystals                  & \qtyrange{25}{450}{\degreeCelsius} \\
               & pressure (\qty{0.5}{\bar})  & and nanoparticles & \\
    \bottomrule
    \end{tabular}
    \caption{
        The difficulty in understanding the mechanisms at play during industrial reactions is highlighted with the example of the oxidation of ammonia using heterogeneous catalysis.
        Reproducing the same exact industrial reaction conditions \parencite{Hatscher2008} and sample can be difficult in a laboratory due to the nature of the probe, the sensitivity of the technique, and the design of reactor cells for synchrotrons.
        This is the so-called material and pressure gap in heterogeneous catalysis.
    }
    \label{tab:Gap}
\end{table}

The size of crystalline samples for x-ray studies is typically limited by the technique if not by the synthesis process.
For Bragg coherent diffraction imaging (BCDI), the coherence lengths of the instrument, the need to sample the interfringes with the detector (inversely proportional to the crystal size until not visible anymore), and the need for the beam to fully illuminate the particle are most important (sec. \ref{sec:BCDI}) and translate into an upper limit for the sample size (\qty{\approx 1}{\um}).

The advantage of the sub-micron-scale \textit{nano}-particles typically used in BCDI is that they exhibit a various amount of active sites such as clear facets, edges, corners and defects.
They are thus considered to be a good approximation of real catalysts \parencite{Somorjai2007, Molenbroek2009, Cuenya2010, Kwangjin2012, Schauermann2013}.
Moreover, the particle can be imaged in three dimensions, which allows the visualisation and study of both the surface and bulk structure.

However, samples below a certain size are not easily measured experimentally with other diffraction techniques for a simple reason which is that the outcoming photon flux is proportional to the sampled volume (sec. \ref{eq:ScatteredIntensity}).
This limitation also exists in BCDI, despite the highly focused beam, which draws a limit between experimental samples (the smallest imaged nanoparticle is \qty{60}{\nm} large - \cite{Bjorling2019, Carnis2021}) and the few \unit{\nm} large nanoparticles used in dynamical field theory (DFT) or molecular dynamics (MD), theoretical approaches only limited by the current computational power.

The single crystal samples used in this thesis for surface x-ray diffraction (SXRD) and x-ray photoelectron spectroscopy (XPS) are much larger (\qty{\approx 1}{\cm}) and tend more towards a model approach since they only exhibit a single facet on their surface \parencite{Goodman1994}.
The large sample surface area allows an improved sensitivity to the facet surface structure compared to BCDI, by increasing its scattered signal (sec. \ref{sec:SXRD}).
Nevertheless, they show some limitations as correctly highlighted by Calle-Valejo et al. \parencite*{CalleVallejo2017}, which underlines the fact that large single-crystals such at Pt(111) also exhibit a large amount of steps also contributing to the catalytic activity.
%Following the same line of thought, if working at ambient pressure is not a problem with x-rays that have a long penetration depth, the electrons that escape the probed surface in XPS are quickly absorbed by the atmosphere.
%A large sample yields increased statistics.

% Most of the catalytic activity of the single crystal can be considered to origin from its top surface, whereas in the case of BCDI, the sample is constituted of thousands of particles that all together contribute to the catalytic activity, despite only a few being imaged.

Bridging the pressure and material gap has already been the subject of several dissertations in recent years, Ackermann \parencite{Ackermann2007} has for example pushed forward the development of \textit{operando} SXRD for the study of heterogeneous catalysis, an effort recently pushed towards spectroscopy techniques by Dann \parencite*{Dann2019}.

As discussed above, using single crystals as model catalysts has been the prefered appraoch for the past 100 years.
Only recently have nanoparticles (\qty{\approx 1}{\nm} large) been used \parencite{Schaffer2013}, showing a remarkable selectivity towards \ce{NO} at high temperature and ambient pressure.
Supported palladium nanoparticles have shown high selectivity towards production of \ce{N_2} in the selective catalytic oxidation of \ce{NH_3} \parencite{Dann2019}, supported platinum nanoparticles at \qty{250}{\degreeCelsius} reached

In the frame of this thesis, both the material and pressure gap will be partly bridged by operating at temperatures after the catalyst light-off, at almost industrial pressure, by the means of \textit{operando} diffraction and spectroscopy techniques.
The difference in activity between different crystalline facets will be studied with different samples that together offer a good compromise between industrial catalysts and samples compatible with the experimental setup of synchrotron beamlines.

\begin{figure}[!htb]
    \centering
    \includegraphics[height=3.75cm]{/home/david/Documents/PhD/Presentations/Slides/PhdSlides/Figures/sample/pt_gazes.png}
    \includegraphics[trim=140 100 0 75, clip, height=3.75cm]{/home/david/Documents/PhD/Presentations/Slides/PhdSlides/Figures/bcdi_data/B7/B7_facets.png}
    \includegraphics[height=3.75cm]{/home/david/Documents/PhD/Presentations/Slides/PhdSlides/Figures/sample/sxrd_sample.png}
    \caption{
    Platinum gauzes used in industry, diameter \qty{\approx 80}{\um} (left), image taken from
    Reconstructed Pt particle measured at SixS, with its surface coloured by the displacement of surface layers from their equilibrium positions, diameter of about \qty{300}{\nm} (middle).
    The orientation of each facet on the particle surface is indicated.
    Pt $(111)$ single crystal used in SXRD and XPS experiments, diameter of about \qty{8}{\mm} (right).
    }
\end{figure}

\section{Aim and Scope}

The oxidation of ammonia is a catalytic reaction that has had an extremely high impact on the $20^{th}$ and $21^{st}$ century, being at the origin of dramatic changes in the world demography with the fertiliser industry, and being today part of numerous industrial process that not only contribute to the fast climate change with greenhouse gases, but also to the ever growing pollution of our ecosystems.

In this first chapter, the importance of the oxidation of ammonia has been underlined.
It was shown that despite being a major catalytic process in a multi-billion industry, the exact mechanisms of action are not yet understood.
The focus will be set on the sole oxidation of ammonia, the goal is also to develop experimental methods that can speak to and attract the synchrotron neophyte to the study of heterogeneous catalysis with synchrotron radiation, together with efficient and reliable computer methods that allow good data reduction and analysis.

%Indeed, if it is of prime interest for a scientist to always be able to understand the fundamentals and to be ready to answer to the question \textit{why ?}, one must also realize that the impact one has strongly depends on the application and reach of his work.
Keeping this idea in mind, this work aims first at bridging the material and pressure gap to bring forth the possibility for synchrotron users to study catalytic reactions at conditions tending to \textit{real} industry conditions.
Secondly, a focus has also been set on explaining the origin, advantages, drawbacks and workflows of each technique, especially for Bragg coherent diffraction imaging, a technique that has yet to reach its full potential through the development of $4^{th}$-generation synchrotrons and powerful computing clusters.

\section{Outline of the Thesis}

The initial chapter will provide a concise explanation of heterogeneous catalysis and of the fundamental principles governing the interaction between x-rays and matter.
This will serve to highlight the origins, benefits, and limitations of each technique employed in this study and explore how catalytic reactions can be indirectly observed using x-rays, leveraging the unique signatures they leave on the materials involved.
The SixS beamline (synchrotron SOLEIL), which served as the primary location for conducting the majority of experiments, will be presented.
Notably, the latest advancements in experimental techniques specific to this beamline will be covered.
Finally, the latest computer programs used for the analysis of the collected data and the challenges of data analysis will be presented so that the reader of this thesis has complete knowledge of every step undertaken to obtain the final results.

In a second chapter, the results obtained during the operando, near-ambient pressure study of the catalytic reaction with three different synchrotron techniques, surface x-ray diffraction, Bragg coherent diffraction imaging and x-ray photoelectron spectroscopy, will be presented.
Sample preparation, the choice of experimental conditions, the reproducibility of the results, the quality of the final data as well as the different difficulties encountered during data collection are discussed.
Moreover, to explore the correlation between surface structure and catalytic activity, mass spectrometry data was collected simultaneously to the other techniques.

Finally, in the last chapter of this thesis will be discussed quantitatively and qualitatively the relation between the result of each technique, their comparison to literature findings, as well as their reliability, representativity and validity, paving the way for any future study.
     
    \chapter{Theory and methods}
    Understanding the different mechanisms at play when photons interact with matter is of crucial importance to be able to decide how to use x-rays as a probe in material science.
Each phenomenon is at the source of different techniques (tab. \ref{tab:techniques} - fig. \ref{fig:TechniquesDataExamples}), scattering brings surface x-ray diffraction (SXRD) and Bragg coherent diffraction imaging (BCDI), two techniques used in the frame of this thesis that yield complementary information about the sample bulk (\textit{i.e.} below the surface) and surface structure.

X-ray absorption together with the photoelectric effect explained by Einstein in \cite*{Einstein1905} are at the origin of a third technique used, x-ray photoelectron spectroscopy, sensitive to the electronic structure of the adsorbates on the sample's surface.

\begin{table}[!htb]
\centering
\resizebox{\textwidth}{!}{%
    \begin{tabular}{l|l|l|l|l}
        Phenomena & Technique & Sample & Information \\
        \toprule
        Scattering & Bragg Coherent & Pt particles & Shape, 3D strain \\
        & Diffraction Imaging & (111), (110), (100), ... & and displacement \\
        & (BCDI) & & \\
        \midrule
        Scattering & Surface x-ray & Pt single crystals & Roughness, relaxation \\
        & Diffraction (SXRD) & (111), (100) & and crystallographic phases \\
        \midrule
        Absorption & X-ray Photoelectron & Pt single crystals & Species presence, \\
        & Spectroscopy (XPS) & (111), (100) & quantity, oxidation state \\
        \bottomrule
    \end{tabular}
    }
    \caption{
        X-ray techniques carried out in the frame of this thesis.
    }
    \label{tab:techniques}
\end{table}

\begin{figure}[!htb]
    \centering
    \includegraphics[height=0.25\textwidth]{/home/david/Documents/PhD/Presentations/Slides/PhdSlides/Figures/gwaihir/dp_pr.png}
    \includegraphics[height=0.25\textwidth]{/home/david/Documents/PhD/Presentations/Slides/PhdSlides/Figures/sxrd_data/Pt100/maps/Map_hkl_surf_or_2227-2283_not_patched.png}
    \includegraphics[height=0.25\textwidth]{/home/david/Documents/PhD/Presentations/Slides/PhdSlides/Figures/xps_data/transition_xps.png}
    \caption{
    Slice of 3D diffraction pattern collected at the SixS beamline at the SOLEIL synchrotron using BCDI (left).
    Reciprocal space in-plane map collected at the SixS beamline using SXRD (middle).
    Ambient pressure XPS spectra measured during a transition between two atmospheres (\argon  and \dioxygen), collected at the B07 beamline at the Diamond synchrotron using XPS (right).}
    \label{fig:TechniquesDataExamples}
\end{figure}

In this chapter, heterogeneous catalysis will be discussed, as well the main reaction mechanisms occuring during catalytic reactions.
The concept of active sites will be introduced together with the link between adsorption and surface strain, crucial for the study of catalytic reaction with x-ray diffraction techniques that are intrinsically sensitive to lattice strain in crystals.
The origin of each technique from fundamentals that not only tailor their experimental design but also their operability and the information one can gain from their use will be detailed (tab \ref{tab:techniques}).
Finally, the importance of open-source, user-friendly, comprehensive, and detailed computer programs will be adressed with \text{Gwaihir}, a \textit{Python} software developped during this thesis not only to bring new users to BCDI, but also to facilitate the data analysis workflow.
A brief description of the other sofwares used during this thesis will be given.
    \section{Heterogeneous Catalysis}\label{sec:Catalysis}

\textit{
A catalyst is a substance that speeds up a chemical reaction, or lowers the temperature or pressure needed to start one, without itself being consumed during the reaction.
Catalysis is the process of adding a catalyst to facilitate a reaction.
}

Chemical reactions involve the breaking, rearranging, and rebuilding of bonds between atoms in molecules, resulting in the formation of new molecules.
Catalysts play a crucial role in enhancing the efficiency of these reactions by providing an alternative route for the reaction with a lower \textit{activation energy} $E_a$, the minimum energy that must be attained for the reaction to take place.
This process facilitates the breaking and formation of chemical bonds, increasing the reaction rate \parencite{Schlogl2015, Hagen2016}, described by the rate equation (eq. \ref{eq:RateEquation}) as a function of the rate constant $k$, the reactant concentration $[A]$, $[B]$ and the order of reaction $a, b$ which refers to how the concentration of the reactants affects the rate of the reaction.

\begin{equation}
    \label{eq:RateEquation}
    \text{rate} = k [A]^a [B]^b
\end{equation}

The impact of changing the temperature $T$ of the reaction and the activation energy is directly set in the rate constant defined by the Arrhenius equation (eq. \ref{eq:RateConstant}), $R$ is the gas constant.

\begin{equation}
    \label{eq:RateConstant}
    k \propto e^{\frac{-E_a}{RT}}
\end{equation}

\subsection{Mechanisms}

Heterogeneous catalysis is a type of catalytic process where the catalyst exists in a different phase (solid, liquid, or gas) from the reactants.
Most commonly, the catalyst is in the solid phase, while the reactants are in either the gas or liquid phase.
This distinction sets heterogeneous catalysis apart from homogeneous catalysis, where the catalyst and reactants are in the same phase.
One of the key advantages of heterogeneous catalysis is the ease of catalyst separation and reuse.
Since the catalyst is in a different phase, it can be easily separated from the reaction mixture once the reaction is complete.
This makes the catalyst recyclable and economically attractive for industrial processes \parencite{FECHETE20122}.

There are a few key parameters to take into account when comparing the efficiency and performance between different catalysts \parencite{Boudart1995, ZHANG2019, WACHS2022}.
The \textit{turnover frequency} (TOF) refers to the number of reactant molecules catalyzed into product molecules by a single catalyst site per unit of time, expressed in units of moles of product formed per second per mole of active catalyst site, it is an indicator of the catalyst activity.
The \textit{turnover number} (TON) refers to the total number of reactant molecules that are converted into product molecules by a single catalyst site over the entire course of a reaction, and provides a measure of the catalyst's overall \textit{activity} and \textit{stability} before deactivation.
Finally, the \textit{selectivity} of the catalyst refers to its ability to favor the formation of a specific product or products over undesired products.

Developing a catalyst that simultaneously fulfills all the desired requirements of high stability, selectivity, and activity for every chemical reaction remains an ongoing scientific pursuit \parencite{Hagen2016}.

The Sabatier principle, at the origin of the Nobel price of chemistry of 1912 \parencite{Che2013}, states that when a heterogeneous catalyst is used, there exists an optimum catalyst binding energy for the reactant on the catalyst surface.
On one hand, if the catalyst binding energy is too weak, the reactant molecules do not adsorb strongly enough on the catalyst, leading to low reaction rates.
On the other hand, if the catalyst binding energy is too strong, the reactant molecules bind too tightly and do not easily react with the other reactant, resulting in low selectivity and efficiency.
Moreover, if the adsorption is too strong, the product of the reaction is not able to desorb from the active site.

\begin{figure}[!htb]
    \centering
    \includegraphics[width=0.49\textwidth]{/home/david/Documents/PhD/Figures/ammonia/ER.pdf}
    \includegraphics[width=0.49\textwidth]{/home/david/Documents/PhD/Figures/ammonia/LH.pdf}
    \caption{
        In the Elley Rideal mechanisms (a) one of the reactants interacts directly from the gas phase with an adsorbed layer of the other reactant, whereas for the Langmuir Hinshelwood mechanism (b) both of the reactants are adsorbed on the catalyst's surface.
        \textcolor{Important}{ADD ARROWS, ENERGY PATH FOR REAGENTS; ETC ...}
    }
    \label{fig:Mechanisms}
\end{figure}

In a founding study in \cite*{Langmuir1922}, Langmuir detailed the possible mechanisms of actions for heterogeneous catalysis.
When studying bimolecular reactions, the Langmuir–Hinshelwood and Elley-Rideal mechanisms are two main mechanisms of interest \parencite{catal12101134}.
In the Langmuir–Hinshelwood mechanism, both reactants are adsorbed in nearby sites at the surface of the catalyst, and react on the surface to form the product which is subsequently desorbed from the surface \parencite{Prins2018, ROSS2019}.
In the Eley–Rideal mechanism, a reactant in the gas phase comes and interacts with the adsorbed reactant without being itself adsorbed, followed by the desorption of the product \parencite{rideal_1939, Weinberg1996}  (fig. \ref{fig:Mechanisms}).

% In the third mechanism, an atom from the surface layer of the catalyst reacts with an adsorbed reactant, which creates a vacancy in his original position on the surface when the products is desorbed.
% The vacancy is then replaced by an atom that comes from from the bulk or the gas phase \parencite{MARS1954, Doornkamp2000}.
% This is called the Mars-Van Krevelen mechanism.

The Langmuir–Hinshelwood mechanism seems to be generally preferred, the bonding of the reactant molecule to the atoms weakening its bonds and preparing the molecule for reaction by reducing the activation energy required for the reaction to occur \parencite{Baxter2002}.

The complete heterogeneous catalytic reaction consists in a series of elementary steps such as reactant dissociation, adsorption, surface diffusion, surface chemical reactions, and desorption that are nowadays extensively studied \textit{via} different theoretical computing methods such as density-functional theory (DFT - \cite{Reuter2004, Molenbroek2009, Yawei2015, Gaggioli2019, Chatelier2020}) or more recently machine learning \parencite{Kitchin2018, Schlexer2019, Anstine2023} as a function of the reaction parameters.

Due to the increase of interest in the field of catalysis, synchrotron techniques have been developped to study the reaction in \textit{operando} conditions, meaning that the reaction happens while being observed \parencite{Meirer2018}.
The first step towards the use of these experimental methods is to understand what specific signatures on the working catalyst can be identified using an x-ray probe.
% (sec. \ref{sec:XRIntMatter}, tab. \ref{tab:techniques}).

\subsection{Active sites in heterogeneous catalysis}

Heterogeneous catalysis being a surface process, the tendency is to use small particles which maximize the surface area to volume ratio, increasing the number of available active sites for catalysis \parencite{Vogt2022}.
This statement is at the origin of the push towards ever smaller samples for heterogeneous catalysis, and at the origin of the increase of interest on nanoparticle catalysis \parencite{CHE1989, Molenbroek2009, vanSanten2010, Schauermann2013}.
As an example, platinum (Pt) nanoparticles (NPs) are of great interest for the petrochemical industry \parencite{Astruc2005}, fuel cell technology and for automobile exhaust gas purification \parencite{HECK2001}.
The surface of nanoparticles can be understood as constituted of facets (fig. \ref{fig:Samples}) whose crystallographic orientation depends on the nanoparticle synthesis route, its size, and on the absence (Wulff construction) or not (Winterbottom construction) of substrate \parencite{WINTERBOTTOM1967, Boukouvala2021}.
The orientation and structure of each facet on the nanoparticle surface can be described by its Miller (hkl) indices.
Between facets are the nanoparticle edges and corners, notorious for being very favorable adsorption sites due to a low local coordination geometry \parencite{Jiang2009}.
\cite{CalleVallejo2014}, \cite*{CalleVallejo2018} have used density dunctional theory (DFT) to show an approximate linear trend between the generalised coordination number $\bar{CN(i)}$ of a nanoparticle surface atom $i$ with $n_i$ neighbours (eq. \ref{eq:CN}) and its adsorption properties.
The \textit{Taylor ratio} translates this effect by defining the fraction between active sites and the total exposed surface \parencite{Taylor1925}.
\textcolor{Important}{CalleVallejo2023?}

\begin{equation}
    \bar{CN(i)} = \sum_{j=1}^{n_i} \frac{cn(j) n_j}{cn_{max}}
    \label{eq:CN}
\end{equation}
where \textit{cn} is the usual coordination number.%, \textit{i.e.} the number of neighbouring atoms.

However, strong adsorption sites are not always the most active sites, following the Sabatier principle, it is possible that if the adsorption is too strong, there is no subsequent desorption of the reaction product \parencite{Nilsson2005, Jiang2009}.

The activity of single facets can also be studied with \textit{single crystals}, much larger crystals, synthetized in a way so that their macroscopic surface consists of a single crystallographic plane (fig. \ref{fig:Samples}).

\begin{figure}[!htb]
    \centering
    \includegraphics[trim=140 75 0 100, clip, height=5cm]{/home/david/Documents/PhD/Figures/bcdi_data/B7/B7_facets.png}
    \includegraphics[height=5cm]{/home/david/Documents/PhD/Figures/sample/sxrd_sample.png}
    \caption{
        Catalytically active samples studied with x-ray \textit{operando} techniques, faceted nanoparticle (left) for which the orientation of each facet is written.
        The area in white corresponds to edges and corners.
        Round single crystal set on sample holder (right), the top of the crystal is the (111) facet that is also present on the nanoparticle.
        The area of the \{111\} facets on the nanoparticle is \qty{\approx 0.25}{\um^2} and occupies \qty{20}{\percent} of its surface, whereas the surface area of the single crystal is \qty{\approx 5e7}{\um^2}.
    }
    \label{fig:Samples}
\end{figure}

\subsection{Linking strain and reactivity}

A modern theoretical approach to the interaction between surface and adsorbates in heterogeneous catalysis has been proposed by Hammer, Norskov and Mavrikakis \parencite*{Hammer1995, Mavrikakis1998, Hammer2000} for transition metals, which is the main focus of this thesis.

In a crystalline lattice, the energy levels of electrons are so close that they form bands \parencite{Ashcroft76}, the study considers that the energy levels of the \textit{d}-bands in transition metals is responsible for bonding and adsorption in catalytic reactions, the main conclusion being that there is a direct relation between lattice strain, adsorption energy and surface reactivity, which has been confirmed by later studies \parencite{Kitchin2004, Kibler2005, Ontaneda2015}.
Moreover, d-bands narrowing effect also justifies the difference between bulk, surfaces steps and kinks in terms of reactivity \parencite{Haydock1972, DESJONQUERES1975, EGELHOFF1987, Hammer2006}.
\textcolor{Important}{Check ref CA, also d-band centre ?}
The application of this work has been demonstrated by Wang et al. in \cite*{Wang2016} and is at the origin of many scientific studies for strain engineering in catalysis, probed with techniques sensitive to the catalyst structure, \textit{i.e} to the lattice strain if the resolution is good enough, such as surface x-ray diffraction \parencite{Resta2020a}, and Bragg coherent diffraction imaging \parencite{Sneed2015, ulvestad_situ_2016, Kim2019, Bjorling2019, Passos2020, Carnis2021a, Carnis2021b, Dupraz2022} in which strain has been proven to impact the catalytic properties of the catalyst.


    \newpage
\section{X-ray interaction with matter} \label{sec:XRIntMatter}

Understanding the different mechanisms at play when photons interact with matter is of crucial importance to be able to decide how to use x-rays as a probe in material science.
Each phenomenon is at the source of different techniques, scattering brings surface x-ray diffraction (SXRD) and Bragg coherent diffraction imaging (BCDI), two techniques used in the frame of this thesis that give complementary information about the sample structure bulk and surface.

X-ray absorption together with the photoelectric effect explained by Einstein in 1905 are at the source of a third technique used, x-ray photoelectron spectroscopy, which yields information specific to the nature of the adsorbates on the sample's surface.

\begin{table}[!htb]
    \centering
    \small
    \begin{tabular}{l|l|l|l}
        Technique & Sample & Sensitivity & Information \\
        \toprule
        Bragg Coherent & Pt nanoparticles & Bragg electronic & Shape, 3D strain  \\
        Diffraction Imaging & (111), (110), (100), ... & density and phase & and displacement \\
        (BCDI) &  &  & \\
        \midrule
        Surface x-ray & Pt single crystals & Surface & Roughness, relaxation \\
        Diffraction (SXRD) & (111), (100), (311) & structure & and crystallographic phases \\
        \midrule
        x-ray Photoelectron & Pt single crystals & Surface & Species presence, \\
        Spectroscopy (XPS) & (111), (100) & species & quantity, oxidation state \\
        \bottomrule
    \end{tabular}
    \caption{Near-ambient pressure (NAP) x-ray techniques carried out at the SixS (SOLEIL) \\ or B-07 (Diamond) beamlines}
    \label{tab:techniques}
\end{table}

\begin{figure}[!htb]
    \centering
    \includegraphics[height=0.25\textwidth]{/home/david/Documents/PhD/Presentations/Slides/PhdSlides/Figures/gwaihir/dp_pr.png}
    \includegraphics[height=0.25\textwidth]{/home/david/Documents/PhD/Presentations/Slides/PhdSlides/Figures/sxrd_data/Pt100/maps/Map_hkl_surf_or_2227-2283_not_patched.png}
    \includegraphics[height=0.25\textwidth]{/home/david/Documents/PhD/Presentations/Slides/PhdSlides/Figures/xps_data/transition_xps.png}
    \caption{Slice of 3D diffraction pattern collected at the SixS beamline at the SOLEIL synchrotron (left). Reciprocal space map in HK plane collected at the SixS beamline at the SOLEIL synchrotron (middle). Ambient pressure XPS spectra collected during a transition between two atmospheres (Ar and O$_2$), collected at the B07 beamline at the Diamond synchrotron (right).}
\end{figure}


In this chapter we will discuss the origin of each technique from fundamentals that not only tailor their experimental design but also their operability and sensitivity.

\subsection{Scattering from electrons and atoms}

The duality between wave and particles was first mentionned by Max Planck and Albert Einstein in the early 20$^{th}$ century and generalized to all matter by Louis-Victor de Broglie in 1924 with the famous formula:

\begin{equation}
	\lambda = \frac{h}{p}
\end{equation}

Electromagnetic waves, \textit{i.e.} light or photons can be characterized by their energy $E$ in \si{\electronvolt} and wavelength $\lambda$ in \si{\meter}. $p$ corresponds to the momentum transfer.
The conversion between is realized thanks to Planck's constant $h = 6.626 \times 10^{-34}$ \si{\joule.\second} and the speed of light in vacuum $c = 2.9979 \times 10^{8}$ \si{\meter.\per \second} (eq. \ref{eq:EnergyLambda}).

\begin{equation}
    \label{eq:EnergyLambda}
	E = \frac{hc}{\lambda}
\end{equation}

The properties of the photon and its use in our society depends on its energy and wavelength.
If visible light is situated between $500 \, \si{\electronvolt}$ and $900 \, \si{\electronvolt}$, micro-waves used in our everyday life are situated between $10^{-6} \si{\electronvolt}$ and $10^{-4}\si{\electronvolt}$.
On the other side of the electromagnetic spectrum, we have higher energy photons such as x-rays ($\in [10^{2}, 10^{6}] \si{\electronvolt}$) and $\gamma$-rays (above $10^{6}$ $\si{\electronvolt}$).

x-rays have a wavelength of a few \si{\angstrom} ($10^{-10} \:\si{\meter}$) making them the perfect probe to study the structures of materials at the atomic scale thanks to different interactions with matter.
\textcolor{Important}{Explain why this particular wl is the best}

\subsubsection{Cross-sections}

When an electromagnectic wave interacts with matter it will be attenuated by absorption, reflection or scattering.
Each process can be quantified depending on the atoms (and thus on the electronic cloud) the beam interacts with and the energy on the incident photon, this is illustrated in figure \ref{fig:cross_sections}.
The cross-section for a particuliar process $p$ is defined as follows \parencite{Willmott}:

\begin{equation}
	\sigma_p = (\Lambda_p N_i)^{-1}
\end{equation}

$\Lambda_p$ is the attenuation length in \si{\meter}, \textit{i.e.} the length after which the beam's intensity is reduced to $1/e$, $N_i$ is the atomic number-density in atoms/unit volume.

\begin{figure}[!htb]
    \centering
    \includegraphics[width=\textwidth]{Figures/introduction/cross_sections.pdf}
    \caption{Cross-sections for platinum (Z=78) for various processes that occur when photons interact with matter. The data was taken from the NIST (National Institute of Standards and Technology) \parencite{NIST_cross_sections} website. The energy range of the SixS beamline at SOLEIL is highlighted in blue.}
    \label{fig:cross_sections}
\end{figure}

Compton scattering, also named inelastic scattering, is a process during which some of the incident electromagnetic wave energy is transfered to the atoms' electrons.
This results in a lower energy for the scattered photon (and therefore a higher wavelength) compared to the incident photon.
%This effect has a low cross-section compared to the two other processes and is therefore not taken account during the experiments.

In the frame of this thesis, the (elastic) Thomson scattering cross-sections is the most important, at the origin of x-ray diffraction.
This process is dominant for energies below $200 \si{k\electronvolt}$, together with photoelectric absorption.% for which the K, L and M edges are shown.

\subsubsection{Scattering from an electron}

We begin our discussion of x-ray scattering by considering scattering from a single free electron using classical electromagnetic theory.
During elastic scattering, the oscillating electric field of the incident x-ray wave exerts an electromagnetic force on the electron, causing it to accelerate and oscillate in the same direction as the incident electronic field.

\begin{figure}[!htb]
    \centering
    \includegraphics[trim=0 75 0 25, clip, width=\textwidth]{Figures/introduction/torus.pdf}
    \caption{
        Effect of relation between synchrotron x-ray polarization $\hat{\epsilon}$ and scattered x-ray polarization $\hat{\epsilon}'$ on the scattered field.
        Working in the vertical plane ($\vec{x}, \vec{z}$) is preferable to maximize the intensity of the scattered field, this has pratical repercussions in surface x-ray diffraction for which it is preferable to work in a vertical geometry to scan large 2D areas of the reciprocal space (a).
        The scattered field intensity is attenuated by a $\cos{(\theta)}$ factor, where $\theta$ is the angle between the plane perpendicular to the electric field and the direction of observation $\vec{R}$ (b)
    }
    \label{fig:polarization_effect}
\end{figure}

During an elastic scattering event, the oscillating electron emits a spherical electromagnetic wave with the same wavelength as the incident beam (Thomson scattering).
The scattered field $\vec{E}_{scatt}$ is then proportional to the incident electromagnetic field $\vec{E}_{in}$ as follows \parencite{NielsenMcMorrow}:

\begin{equation}
    \label{eq:scatt_field}
    \frac{\vec{E}_{scatt}(R, \, t)} {\vec{E}_{in}} = -r_0 \frac{e^{\vec{k}.\vec{R}}} {|\vec{R}|}| \hat{\epsilon}.\hat{\epsilon}'|
\end{equation}
$|\vec{R}|$ is the distance at which the scattering is detected at the time $t$, $\hat{\epsilon}$ is the synchrotron x-ray polarization vector, $\hat{\epsilon}'$ the scattered x-ray polarization vector, and $r_0$ is the Thomson scattering length or the classical radius of the electron defined as:

\begin{equation}
    \label{eq:scatt_thomson_scat_length}
    r_0 = \frac{e^2} {4\pi\epsilon_0 m_e c^2}
\end{equation}
with $m_e$ the mass of the electron and $\epsilon_0$ the permittivity of free space.

The minus sign illustrates a phase shift of $\pi$ between the incident and scattered wave, $\hat{\epsilon}$ and $\hat{\epsilon}'$ are respectively the polarization vectors of the incident and scattered electromagnectic fields.

The differential cross-section for Thomson scattering measures the efficiency of the scattering in the volume occupied by a solid angle $d\Omega$ in the direction $\vec{R}$ \parencite{NielsenMcMorrow}. It is defined as follows:
\begin{equation}
    \label{eq:dif_cross_sec_thomson1}
    \frac{d\sigma_{ts}} {d \Omega} = \frac{ |\vec{E}_{scatt}(R, t)|^2 R^2} {|\vec{E}_{in}|^2}
\end{equation}

By substituting eq. \ref{eq:scatt_field} into eq. \ref{eq:dif_cross_sec_thomson1}, it becomes clear that the scattering is proportional to the Thomson scattering length and that the intenisty is attenuated depending on the dot product between the two polarizations.
The polarization factor $P$ for scattered beams is defined as $P =  | \hat{\epsilon}.\hat{\epsilon}'|^2$ and we can write the differential cross-section as:

\begin{equation}
    \frac{d\sigma_{ts}} {d \Omega} = r_0^2 | \hat{\epsilon}.\hat{\epsilon}'|^2 = r_0^2 P
\end{equation}

The effect of the polarization of the incident beam is illustrated in figure \ref{fig:polarization_effect}.
At synchrotron sources where the incident beam is horizontally polarized, working in the vertical plane becomes more effective since the polarization factor is always equal to one.

The total cross-section for the scattering event by a single free electron can be computed by integrating the differential cross-section over all the possible scattering angles, \textit{i.e.} by averaging all possible polarization directions \parencite{Willmott}.
This yields $\sigma_{ts} = 8 \pi r_0^2 /3 = 0.665 \:\si{\barn}$, the total Thomson scattering cross-section is constant, independent of the incoming photon energy. This results holds for x-rays for which the scatterer \textit{i.e.} the electron can be considered as free \parencite{Willmott}.

\subsubsection{Scattering from a single atom}\label{sec:scattering}

As we have seen the main scatterer for Thomson scattering is the electron.
An atom can be described as a small volume $d^3\vec{r}$ in which the electrons are localized.
The electromagnetic field scattered from an atom is therefore proportional to the atomic electronic density $\rho_{atom}(\vec{r})$.
For a single atom of atomic number $Z$ we have :

\begin{equation}
    \int \rho_{atom} (\vec{r}) d^3\vec{r} = Z
\end{equation}

To describe the phase of the scattered field, which depends on the phase of each incident wave at each point in the scattering volume (fig. \ref{fig:q}), we introduce the scattering vector, $\vec{q}$ (eq. \ref{eq:Q}), to derive the amplitude of a scattering event in the direction defined by the scattering angle $2\theta$.

\begin{equation}
    \label{eq:Q}
    \vec{q}=\vec{k_s}-\vec{k_i}.
\end{equation}

\begin{figure}[!htb]
    \centering
    \includegraphics[scale=0.6]{Figures/introduction/q.pdf}
    \caption{
    Geometry of the scattering vector $\vec{q}$ in reciprocal space, $2\theta$ is the scattering angle.
    The magnitude of the scattering vector can be derived with the angle $\theta$ that draws a line cutting $\vec{q}$ at $|\vec{q}|/2$.
    }
    \label{fig:q}
\end{figure}

Therefore, the phase difference between a wave scattered at a position $O$ and a wave scattered at a distance $\vec{r}$ from $O$ is equal to $(\vec{k_i} - \vec{k_s}).\vec{r} = \vec{q}.\vec{r}$ (fig. \ref{fig:q}).

From fig. \ref{fig:q} and the relation between wavevector and wavelength ($k = 2 \pi/\lambda$), the scattering vector can be written as a function of the wavelength of the incident beam $\lambda$, and of the direction of detection defined by the scattering angle $2\theta$ (eq. \ref{eq:QSinTheta}) between the wavevector of the incident photon $\vec{k_i}$ and the wavevector of the scattered photon $\vec{k_s}$.
We assume here that the scattering event is elastic ($|\vec{k_i}|=|\vec{k_s}|$) and that the waves are plane waves parallel to each other when in the small scattering volume $d^3\vec{r}$.

\begin{equation}
    \label{eq:QSinTheta}
    |\vec{q}| = \frac{4\pi}{\lambda} \sin{\theta}
\end{equation}

Each electron in the small volume $d^3\vec{r}$ will have a contribution proportional to the Thomson scattering length $r_0$ to the scattered field with a phase $e^{i\vec{q}.\vec{r}}$.
By integrating over the volume occupied by the atom we obtain the total contribution of an atom to the scattered field in the direction $2\theta$:

\begin{equation}
    \label{eq:AtomicFormFactor}
    -r_0 \int \rho_{atom} (\vec{r}) e^{i\vec{q}.\vec{r}} d\vec{r} = -r_0 f^0(\vec{q}) = -r_0 FT [\rho_{atom} (\vec{r})]
\end{equation}

The scattering amplitude as a function of $\vec{q}$ is described by the atomic from factor $f^0(\vec{q})$, which is defined as the Fourier transform of the atomic electronic density $\rho_{atom}(\vec{r})$.
The values for the atomic from factor can be approximated by a sum of Gaussians using tabulated coefficients (eq. \ref{eq:AtomicFormFactorTab}) available online \parencite{InterTablesOfCryst}.
It decreases with $\vec{q}$ (or $\sin(\theta) / \lambda)$ as illustrated in figure \ref{fig:atomic_form_factor}.

\begin{equation}
    \label{eq:AtomicFormFactorTab}
    f^0(\vec{q}) = \sum_{i=1}^4 a_i \exp (-b_i (\frac{q} {4\pi})^2) + c
\end{equation}

The scattering intensity is equal to the square of the scattering amplitude.
For example, the scattering intensity of palladium atoms at $|\vec{q}| \approx 2.75 \AA^{-1}$ is only $\approx 31\%$ of that of platinum atoms. In the case of oxygen, the intensity falls down to $\approx 6.7\%$.
This difference in scattering intensity between elements becomes crucial when working with small objects that have a small scattering volume such as nanoparticles.

\begin{figure}[!htb]
    \centering
    \includegraphics[width=\textwidth]{Figures/introduction/atomic_form_factor.pdf}
    \caption{
    Atomic form factor calculated for Pt (Z=78) using tabulated values \parencite{InterTablesOfCryst} for equation \ref{eq:AtomicFormFactor}. The scattering intensity decreases with the scattering angle $\theta$ but increases with the incident wavelength $\lambda$. $i$ and $c$ respectively designate the Gaussian contribution and constant in eq. \ref{eq:AtomicFormFactorTab}.
    }
    \label{fig:atomic_form_factor}
\end{figure}

\subsection{Refraction, reflection and absorption effects}\label{sec:RefractionReflectionAbsorption}

For now we have considered the electron as free.
However, electrons in atoms or molecules exist in discrete bound states (characterised by a binding energy $E_b$), and their response to the incoming electromagnetic field becomes dependent on the incoming energy.
In a classical point of view, we can reconsider the bound electron as a damped harmonic oscillator with associated resonant frequencies $\omega_b$ corresponding to the electron's binding energy.
First, when approaching the binding energy, the amplitude of the electron's oscillation are reduced and so is the atomic form factor by a real energy-dependant factor $f'(\hbar\omega)$.
Secondly, when approaching the binding energy, the oscillation of the scattering electron experiences a phase lag.
This is taken into account by introducing an imaginary energy-dependent factor $f''(\hbar\omega)$.
The atomic form factor $f^0(\vec{q})$ must therefore be corrected by these two correction factors, we define the total atomic form factor $f(q, \hbar\omega)$ as \parencite{NielsenMcMorrow}:

\begin{equation}
    f(\vec{q}, \hbar\omega) = f^0(\vec{q}) + f'(\hbar\omega) + if''(\hbar\omega) = f_1(\vec{q}, \hbar\omega) + i f_2(\hbar\omega)
\end{equation}

The values of the correction factors $f'$ and $f''$, also known as the dispersion corrections to $f^0$, depend on the electronic structure of the atom and are maximum when approaching the binding energies of the electrons in the atom.
To understand the impact of the correction factors on the electromagnetic wave inside a material, we introduce the complex refraction index $n$, which for x-rays can be written as:

\begin{gather}
    n = 1-\delta+i\beta = 1 - \frac{r_0 \lambda^2}{2\pi} \sum_i N_i f_i(\vec{q}=0)\\
    \label{eq:delta}
    \delta = \frac{r_0 \lambda^2}{2\pi} \sum_i N_i f_1(\vec{q}=0)\\
    \beta = \frac{r_0 \lambda^2}{2\pi} \sum_i N_i f_2
\end{gather}
with $f_i(\vec{q}=0)$ the total atomic form factor of the atom $i$ in the forward direction ($\vec{q}=0$), and $N_i$ the number of atoms corresponding to the element $i$ per unit volume.
$\delta$ is the refractive index and $\beta$ the absorption index.

During simple experiments, when x-rays impinge on a plane surface between air ($n$=1) and another medium of refractive index $n$, x-rays are refracted, the wavelength changes to $\lambda/n$ and the direction of propagation becomes closer to the surface between both mediums.
Below a certain incident angle known as the \textit{critical angle} $\alpha_c$ between the direction of the incident x-rays and the plane interface, the x-rays will be totally reflected.
This effect is used in focussing mirrors for optics but also in techniques used to probe the surface of materials to obtain information about the surface density since the critical angle is approximately equal to $\sqrt{2\rho}$ \parencite{Willmott}.
Moreover, during surface x-ray diffraction, the incoming angle perpendicular to the surface is kept to very low values approaching the critical angle to limit the volume probed by the x-rays to increase the contrast between the amplitude of the x-rays scattered from the bulk in regard to the x-rays scattered from the latest surface layers.
The penetration depth of the x-rays in the material is detailed is sec. \ref{sec:SXRD}.

Far from the binding energies of the electron in the system, the correction factor $f'$ in the expression of $\delta$ can be ignored and the sum in eq. \ref{eq:delta} is then equal to the electronic density of the system.

Finally, to highlight the impact of the complex part of the absorption index $\beta$, we express the electric field of a plane electromagnetic wave propagating in the $\vec{z}$ direction inside a material of refractive index $n$ as follows:

\begin{equation}
    E(\vec{z}, t) = E_0 e^{i(n\vec{k}.\vec{z} - wt)} = E_0 e^{-\beta \vec{k}.\vec{z}} e^{i(\delta\vec{k}.\vec{z} - wt)}
\end{equation}

The imaginary part of the form factor translates into an exponential decay of the intensity of the incoming wave as a function of the absorption index $\beta$, the intensity of the electromagnetic field decreases; \textit{i.e.} the wave is absorbed by the material.
The absorption of the photons by the material is at the origin of many different phenomena such as the promotion of electrons to either higher energetic bound states or directly to vacuum.
The kinetic energy of the electrons released from the material in the latter scenario is equal to the incident photon energy $\hbar\omega$ minus the energy required to release the electron from its bound state $E_b$, minus the energy required for the electron to escape from the electronic cloud $W$ (also known as the work function).

\begin{equation}
    E_{kin} = \hbar\omega - W - E_b
\end{equation}

Measuring the kinetic energy of the photoelectron that escaped the material yields direct information about the different binding energies present in the material, \textit{i.e.} about its electronic structure.
The study of these kinetic energy is known as x-ray photoelectron spectroscopy (XPS),
a technique used in this thesis to study the electronic structure of the latest atomic layers of a surface, electrons released from atoms in deeper layers being ultimately absorbed(more like cancelled due to the interraction with the e cloud of) by the surrounding atoms.
The cross-section for XPS is directly proportional to the absorption cross-section $\sigma_a$ (eq. \ref{eq:AbsorptionCrossSection}, fig. \ref{fig:cross_sections}) which can be derived from the imaginary atomic form factor:

\begin{equation}
    \label{eq:AbsorptionCrossSection}
    \sigma_a(\hbar\omega) = 2 r_0 \lambda f_2(\hbar\omega)
\end{equation}

\section{Scattering from crystals}

A crystal is a solid material composed of a regular, repeating arrangement of atoms, ions, or molecules (\textit{i.e.} a pattern), exhibiting a highly ordered structure with long-range periodicity in three dimensions.

A Bravais lattice refers to an infinite array of points (nodes) that represents the basic repeating unit of a crystal lattice, defining the translational symmetry of the crystal structure.
It is characterized by a set of three basis vectors and their linear combinations, which generate the entire lattice when translated in space.
The \textit{primitive} unit cell is the smallest cell with which you can describe the crystal.

The structure of the Bravais lattice (cubic, hexagonal, ...) combined with the position of the patterns in the lattice and the symmetry relations between them defines the crystal \textit{space group}.
In total there exist 230 space group in crystallography.
In the simplest case the pattern consists of a single atom, for example pure platinum (Z=78) cristallizes at room temperature in a cubic Bravais lattice (fig. \ref{fig:fcc}).

We can define the scattering factor $F_{crystal}$ of the crystal, that describes the amplitude of the scattered waves, by the sum of the atomic scattering factor $f_j(\vec{q})$ of each atom present in the crystal, while the phases between the scattered waves depend on the positions of the atoms $\vec{R}_j$ in the direction that is perpendicular to the scattering vector $\vec{q}$ (eq. \ref{eq:FCrystal1}).
In other words, each atom can be described as a small volume $d^3\vec{r}$ in the electronic density $\rho(\vec{r})$, and the scattered field by the superposition of the contribution from the electronic cloud surrounding the atoms.
If the atomic scattering factor is equal to the Fourier transform of the electronic density of a single atom, the structure factor of a crystal can be written as the Fourier transform of the crystal electronic density.
%Thus, the total scattering amplitude at $2\theta$ will be equal to the vector sum of the scattering amplitudes in this direction from all volume elements $d^3\vec{r}$ in $\rho(\vec{r})$ taking into account the phases between them.

\begin{equation}
    \label{eq:FCrystal1}
    F_{crystal}(\vec{q}) = \sum_j^{N_{atoms}} f_j(\vec{q}) e^{i\vec{q}.\vec{R}_j} = FT[\rho(\vec{r})]
\end{equation}

$f_j(\vec{q})$ is the atomic scattering factor of the $j^{th}$ atom at position $\vec{R}_j$ in a crystal made of $N_{atoms}$ atoms.
The Thomson scattering length is left apart for simplicity.

\begin{figure}[!htb]
    \centering
    \includegraphics[width=\textwidth]{/home/david/Documents/PhDScripts/Drawing/blender/FCC.png}
    \caption{Face centered cubic (FCC) lattice of Pt (space group 225). Atoms are represented as solid balls and situated on the corners and at the middle of the faces. The lattice parameter at room temperature is $a = 3.9254 \:\AA$. Close packed direction is achieved along the diagonal of the lateral faces, the distance between the atoms then becomes is $2.71 \:\si{\angstrom}.$}
    \label{fig:fcc}
\end{figure}

The position of any atom in the crystal $\vec{R}_j$ is equal to the sum of the position of the unit cell containing the atom $\vec{R}_{uc}$, plus its position within the unit cell $\vec{r}_j$.
For a crystals made of $N_{uc}$ unit cells each composed of $N_{atoms,uc}$ we have:

\begin{equation}
    F_{crystal}(\vec{q}) = \sum_j^{N_{atoms,uc}} f_j(\vec{q}) e^{i\vec{q}.\vec{r}_j} \sum_k^{N_{uc}} e^{i\vec{q}.\vec{R_k}}
    \label{eq:Fcrystal}
\end{equation}

% As seen in Equation (\ref{eq:scatt_thomson_scat_length}), the Thomson scattering length is inversely proportional to the mass of the particle, while the scattering cross-section is proportional to the square of the Thomson scattering length.
% Hence the nucleus contributes less than one part in a million to the scattering amplitude and can be completely ignored.
% The charge distribution is therefore equal to the electronic density.

% When detecting the scattered field at a point $\vec{R}$, we assume that $\|vec{R}|$ is far greater than the electronic density, this is called the dipole approximation. NOT GOOD Each atom can be considered as a ball of certain radius (see \ref{fcc_lattice}).

% Furthermore, we assume that the source and detector are sufficiently far from the charge distribution so that the incident and scattered x-rays may be represented as plane waves, parallel to each other and perpendicular to the direction of propagation. % also more easily described
%This is the so-called far-field limit or the fraunhofer region.

% Finally, we assume that the charge distribution is small and that the scattering is weak so that the sample can be treated as a simple "perturbation" to the incident beam resulting in a linear problem \parencite{NielsenMcMorrow} and that the Born approximation is valid.

% Working under these assumption is known as kinematic diffraction for which we can easily derive the intensity of the scattered photons.
% We therefore ignore multiple-scattering of photons inside the sample and xxx.
% This is correct when working with low incident angles and penetration depths as for example in surface x-ray diffraction (see Section x).

% only valid in far-field, faunhofer region
% DWBA
% dipole approximation

The Bravais lattice is defined by three basis vectors $\vec{a},\ \vec{b},\ \vec{c}$ and three angles $\alpha$ [$\angle (\vec{b}, \vec{c})$] , $\beta$ [$\angle (\vec{a}, \vec{c})$] and $\gamma$ [$\angle (\vec{a}, \vec{b})$].
Any vector $\vec{R_k}$ describing the position of a node in the real space can be created by a linear combination of these three vectors:

\begin{equation}
    \label{eq:R_k}
    \vec{R_k}=n_1\vec{a} + n_2\vec{b} + n_3\vec{c}, \quad \ (n_1,n_2,n_3) \in \mathbb{Z}^3
\end{equation}

To understand the contribution of the second sum in eq. \ref{eq:Fcrystal} to the scattering amplitude, it is convenient to introduce the \textit{reciprocal space} which is the Fourier transform of the real space.
It is defined by three basis vectors $\vec{a*},\ \vec{b*},\ \vec{c*}$.
Similarly, the nodes of the reciprocal space can be accessed from its origin by a linear combination $\vec{G}$ of these three vectors.

\begin{equation}
    \vec{a^*}=\frac{2\pi}{V}(\vec{b}\times \vec{c}), \qquad
    \vec{b^*}=\frac{2\pi}{V}(\vec{c}\times \vec{a}), \qquad
    \vec{c^*}=\frac{2\pi}{V}(\vec{a}\times \vec{b}), \qquad
    \vec{a}_i . \vec{a}_j^* = \delta_{i,j}
\end{equation}

\begin{equation}
    \label{eq:G}
    \vec{G}=h\vec{a}^* + k\vec{b}^* + l\vec{c}^*, \quad \ (h,k,l) \in \mathbb{Z}^3
\end{equation}

By combining eq. \ref{eq:R_k} and eq. \ref{eq:G} we find that the solution to the second sum in eq. \ref{eq:Fcrystal}, \textit{i.e.} to the equation $\vec{q}.\vec{R_k} = n \times 2\pi$ $(\ n \in \mathbb{Z})$ is to write $\vec{q}$ as a linear combination of the reciprocal space vectors, \textit{i.e.} as $\vec{G}$.

\begin{equation}
    \label{eq:LaueCond}
    \vec{G} . \vec{R_k} = hn_1 + kn_2 + ln_3 = n \quad  with \ n \in \mathbb{Z}
\end{equation}

This is also known as the Laue condition, that ensures that only certain scattering vectors, corresponding to the reciprocal lattice points, fulfill the condition for constructive interference between the scattered waves.
These specific scattering vectors determine the scattering angles of the peaks observed in a diffraction pattern, also known as Bragg peaks.

\begin{equation}
    \vec{q} = \vec{G}  = h\vec{a}^* + k\vec{b}^* + l\vec{c}^*
\end{equation}

\subsection{Bragg's Law}

The intensity scattered from a crystal as a function of the scattering vector $\vec{q}$ provides information about the arrangement and spacing of crystalline planes perpendicular to $\vec{q}$.
The Miller indices are a set of integers used to represent the planes' orientation and spacing.
A plane denoted by the ($h, k, l$) indices intercepts the axes $\vec{a}$, $\vec{b}$, $\vec{c}$ on the points $|\vec{a}|/h, |\vec{b}|/k, |\vec{c}|/l$.
The direction perpendicular to the $(hkl)$ plane is written as $[hkl]$, the distance between each plane is $d_{hkl}$.

\begin{figure}[!htb]
    \centering
    \includegraphics[scale=0.6]{Figures/introduction/BraggLaw.pdf}
    \caption{The difference in the path length between plane waves scattered at an angle $\theta$ must be an integer multiple of wavelengths for constructive interference to occur. NOT FINISHED}
    \label{fig:BraggLaw}
\end{figure}

A Bragg peak results from the constructive interference between scattered waves at discrete values of the incident angle $2\theta$ on a specific set of crystalline planes.
From fig. \ref{fig:BraggLaw} we can retrieve the condition to have constructive interference, also known as Bragg law and given by eq. \eqref{eq:Bragglaw}.

\begin{equation}
    \label{eq:Bragglaw}
    n\lambda = 2d_{hkl} \sin{\theta}, \quad \ n \in \mathbb{Z}
\end{equation}

% \begin{gather}
%     \frac{(2\pi)^2}{d_{hkl}^2} = h^2 \, (\Vec{a^*}.\Vec{a^*}) + k^2 \,
%     \label{eq:dsquare}(\Vec{b^*}.\Vec{b^*}) + l^2 \, (\Vec{c^*}.\Vec{c^*}) + 2hk \, (\Vec{a^*}.\Vec{b^*}) + 2hl \, (\Vec{a^*}.\Vec{c^*}) + 2kl \, (\Vec{b^*}.\Vec{c^*})\\
%      \frac{(2\pi)^2}{d_{hkl}^2} = h^2 \, {a^*}^2 + k^2 \, {b^*}^2 + l^2 \, {c^*}^2 + 2hk \, {a^*}.{b^*}\cos{\gamma^*} + 2hl \, {a^*}.{c^*}\cos{\beta^*} + 2kl \, {b^*}.{c^*}\cos{\alpha^*}\\
%      \frac{(2\pi)^2}{d_{hkl}^2} = Ah^2 + Bk^2 + Cl^2 + Dhk + Ehl + Fkl
%      \label{eq:RecSpaceMetricTensor}
% \end{gather}{}

In the case of a cubic Bravais lattice, the distance between Miller planes can be written as:
\begin{equation}
    \label{eq:Interplanarspacing}
    d_{hkl}=\frac{2\pi}{|\vec{a^*}|\sqrt{h^2 + k^2 + l^2}}=\frac{|\vec{a}|}{\sqrt{h^2 + k^2 + l^2}}
\end{equation}

From eqs. \ref{eq:QSinTheta} and \ref{eq:Bragglaw}, we fall back on the Laue condition which is a generalization of Bragg law. The diffraction order ($n$) in the Bragg's law equation can be omitted since it is implicitly determined by the Miller indices (hkl) representing the crystalline planes involved in the diffraction.

\begin{gather}
    \label{eq:QandD}
    |\vec{Q}| = \frac{4\pi \sin(\theta)}{\lambda},\\
    \sin(\theta) / \lambda = \frac{1}{2d_{hkl}},\\
    |\vec{Q}| = \frac{2\pi}{d_{hkl}}
\end{gather}{}

When fulfilling the Laue (or Bragg) condition, the direction of the scattering vector is perpendicular to the crystalline planes represented by the ($h, k, l$) Miller indices, and its magnitude is inversely proportional to the distance between consecutive crystalline planes.

\subsection{Structure factor} \label{sec:StructureFactor}

On one hand, in real space, the crystal electronic density can be formed by repeating the unit cell at each lattice point, \textit{i.e.} it can be represented as a convolution of the lattice and unit cell functions.
This representation captures the periodicity and arrangement of the crystal.

On the other hand, in reciprocal space (Fourier space), the crystal scattering factor $F_{crystal}$ is the Fourier transform of the electronic density.
According to the convolution theorem, the Fourier transform of a convolution is equal to the product of the Fourier transforms of the individual functions.
The scattering factor can therefore be expressed as the product of the Fourier transforms of the lattice and unit cell functions.
These Fourier transforms are known as the lattice factor $F_{lat}$ and the unit cell structure factor $F_{uc}$.

\begin{equation}
    F_{crystal} = FT[\rho(\vec{r})] = F_{uc} \times F_{lat}
    \label{eq:FcrystalFlatFuc}
\end{equation}

The lattice factor $F_{lat}$ represents the periodic arrangement of the crystal lattice in Fourier space, corresponding to the reciprocal lattice.
It determines the positions and intensities of the diffraction peaks in the scattering pattern.
The unit cell structure factor $F_{uc}$ describes the distribution of electron density within the unit cell of the crystal and modulates the scattered beam amplitude depending on what atoms are in the unit cell (amplitude) and their positions (phase).
By identifying the different parts of eq. \ref{eq:Fcrystal} and eq. \ref{eq:FcrystalFlatFuc} we can first isolate the unit cell structure factor:

\begin{equation}
    \label{eq:StrucFactor}
    F_{uc} = \sum_j^{N_{atoms}} f_j(\vec{q}) e^{i\vec{q}.\vec{r}_j}
\end{equation}

Each diffraction peak in the scattering pattern corresponds to a specific Fourier component, representing a sinusoidal wave of electron density with a particular frequency and direction determined by its scattering vector $\vec{q}$, \textit{i.e.} by its position in Fourier space.
By knowing the phase relationships between these Fourier components, which can be obtained from the measured diffraction pattern, the electron density within the unit cell can be reconstructed.
The superposition of these sinusoidal waves, representing the diffraction peaks or Fourier components, recreates the original electron density distribution of the crystal.

To summarize, the diffraction peaks in the scattering pattern represent the Fourier components of the crystal's electron density distribution, and their spatial frequencies and phase relationships carry information about the crystal structure.
%The first sum is the unit cell structure factor, that modulates the scattered beam amplitude depending on what atoms are in the unit cell (amplitude) and their positions (phase).
%The second sum is the lattice factor that takes the value of $N_{uc}$ if $\vec{q}.\vec{R_k} = n \times 2\pi$ $(\ n \in \mathbb{Z})$.

The structure factor $F_{uc}$ is the summation of the contribution of each atom $j$ at the position $\vec{r}_j$ of atomic form factor $f_j(\vec{q})$ in the unit cell for a given scattering vector $\vec{q}$.
The position of each atom is defined in the unit cell by eq. \ref{eq:AtomPos}, an example with the platinum atoms in its primitive unit cell is given in tab. \ref{tab:PtAtoms}.

\begin{equation}
    \label{eq:AtomPos}
    \vec{r}_j = x_j\vec{a} + y_j\vec{b} + z_j\vec{c}
\end{equation}

\begin{table}[!htb]
    \centering
    \begin{tabular}{@{}lllllllllllllll@{}}
    \toprule
    Atom & Pt & Pt  & Pt  & Pt  & | Pt & Pt & Pt & Pt  & Pt  & Pt & Pt & Pt & Pt & Pt  \\ \midrule
    x    & 0  & 0.5 & 0.5 & 0   & | 1  & 0  & 1  & 0.5 & 1   & 0  & 1  & 0  & 1  & 0.5 \\
    y    & 0  & 0.5 & 0   & 0.5 & | 0  & 1  & 1  & 1   & 0.5 & 0  & 0  & 1  & 1  & 0.5 \\
    z    & 0  & 0   & 0.5 & 0.5 & | 0  & 0  & 0  & 0.5 & 0.5 & 1  & 1  & 1  & 1  & 1   \\ \bottomrule
    \end{tabular}
    \caption{Position of platinum atoms in the face-centered cubic (FCC) unit cell. Illustrated in fig. \ref{fig:fcc}. The first four atomic positions  lead to the other positions by symmetry.}
    \label{tab:PtAtoms}
\end{table}

Systematic extinctions ($F_{uc} = 0$) occur when certain diffraction peaks are forbidden due to the crystal's symmetry elements, such as the atomic positions.
For example in the FCC lattice of platinum, systematic extinctions occur when the value of Miller indices are not all even or all odds.
Otherwise, the intensity of a Bragg peak detected at the scattering vector $\vec{q}$ is equal to $4\times f_{Pt}(\vec{q})$ where $f_{Pt}(\vec{q})$ is the atomic scattering factor of platinum.

\subsection{Lattice factor} \label{sec:LatticeFactor}

If the scattering factor considers the composition of the unit cell, the lattice factor represents the contribution of the crystal lattice to the overall scattering amplitude.
It takes into account the arrangement of atoms or scatterers within the crystal lattice and their interaction with the incident wave.

By identifying the different parts of eq. \ref{eq:Fcrystal} and eq. \ref{eq:FcrystalFlatFuc} we can now isolate the lattice factor (eq \ref{eq:LatFactor}).
In the case of a simple cubic arrangement of unit cells in three directions perpendicular to each other, we can divide the sum over all the unit cells $N_{uc}$ in the crystal into three sums, each in one of the directions of the crystal lattice ($\vec{a}, \, \vec{b}, \, \vec{c}$) so that $N_{uc} = N_{uc, x} \times N_{uc, y} \times N_{uc, z}$.

\begin{gather}
    \label{eq:LatFactor}
    F_{lat} = \sum_k^{N_{uc}} e^{i\vec{q}.\vec{R_k}} = \sum_{n_x}^{N_{uc, x}} e^{i n_x\vec{q}.\vec{a}} \times \sum_{n_y}^{N_{uc, y}} e^{i n_y\vec{q}.\vec{b}} \times \sum_{n_y}^{N_{uc, z}} e^{i n_y\vec{q}.\vec{c}}\\
    F_{lat} = \prod_{j=\{x,y,z\}} \Bigg( \sum_{n_j}^{N_{uc, j}} e^{i n_j\vec{q}.\vec{a}_j} \Bigg)\\
    F_{lat} = \prod_{j=\{x,y,z\}} S_{N_{uc, j}}(\vec{q}.\vec{a}_j)
\end{gather}

\begin{gather}
    \label{eq:SFunction1}
    S_{N_{uc, j}}(\vec{q}.\vec{a}_j) = \sum_{n_j}^{N_{uc, j}} e^{in\vec{q}.\vec{a}_j} = \sum_{n_j}^{N_{uc, j}} e^{(i\vec{q}.\vec{a}_j)^n_j} = \frac{1-e^{(i N_{uc, j} \vec{q}.\vec{a}_j)}}{1-e^{(i\vec{q}.\vec{a}_j)}},\\
    \label{eq:SFunction2}
    |S_{N_{uc, j}}(\vec{q}.\vec{a}_j)| = \frac{\sin(N_{uc, j} \vec{q}.\vec{a}_j/2)}{\sin(\vec{q}.\vec{a}_j/2)},\\
    \label{eq:SFunction3}
    |S_{N_{uc, j}}(\vec{q}.\vec{a}_j)| \underset{N_{uc, j} \to \infty}{\longrightarrow} \approx \delta (\vec{q})
\end{gather}

Where $S_{N_{uc, j}}(\vec{q}.\vec{a})$ is a function whose modulus $|S_{N_{uc, j}}(\vec{q}.\vec{a})|$ tends towards a dirac function when $N_{uc, j} \longrightarrow \infty$.
The condition to observe peaks in eq. \ref{eq:SFunction2} is for the denominator to be equal to zero: $\frac{\vec{q}.\vec{a}_j}{2} \equiv \pi$ which brings us back to the Laue condition: $\vec{q} = \vec{G}$ where $\vec{G}$ is a reciprocal space lattice vector.
Minor peaks appear when the numerator is equal to one: $N\frac{\vec{q}.\vec{a}_j}{2} \equiv \frac{\pi}{2}$, the distance between each minor peak is equal to $\frac{2\pi}{t_j}$ where $t_j$ is the thickness of the crystal in the direction $\vec{a_j}$.

To resume, the lattice factor $F_{lat}(\vec{q})$ defines the shape of Bragg peaks.
For a large crystal (\textit{i.e.} a large number of unit cells in three directions), ideal Bragg peaks will look like Dirac functions (in reality they are widened due to different experimental factors).
For a small crystal such as nanoparticles, the fringes of the lattice factor will be further apart and their intensity more visible in contrast to the intensity of the main peak (fig. ).
In a more general sense, the lattice factor is the Fourier transform of the crystal's lattice, the effect of the Fourier transform is that fringes will be visible in the reciprocal space in directions perpendicular to the crystal's surface (fig. ).
The ability to effectively sample these fringes during experiments is at the root of the BCDI technique, to be able to perform the inverse Fourier transform of the Bragg peak to then retrieve the crystal shape (sec. \ref{sec:BCDI}).


\subsection{Coherence} \label{sec:Coherence}

Every x-ray source can be described by several parameters.
The brilliance is a measure of the distribution of flux in both space and angular range, it indicates how the flux is spread over the source area and solid angle.

\begin{equation}
    Brilliance = \frac{photons/second}{(mm^2 source \; area)(mrad)^2(0.1\% \; bandwidth)}
\end{equation}

The spectral flux, on the other hand, quantifies the number of photons passing through a defined area per unit bandwidth per second.
Synchrotron facilities aim to optimize both the photon flux and brilliance, a higher flux meaning faster experiments whereas a high brilliance is important for those that need a coherent beam.
To do so, the source size which is the convolution of the size of the photon source and the transverse size of the electron beam has been reduced by an order of magnitude in 4th generation synchrotrons compared to 3rd generation synchrotrons \parencite{Willmott}.

\begin{equation}
    \label{eq:Emittance}
    \epsilon = D \times \Delta\theta
\end{equation}
The emittance $\epsilon$ is defined as the product of the source size $D$ and the beam divergence $\Delta\theta$ (eq. \ref{eq:Emittance}), a low emittance is characteristic of a small source size and almost parallel emitted beams.

\begin{figure}[!htb]
    \centering
    \includegraphics[scale=0.9]{Figures/introduction/coherence_lengths.pdf}
    \caption{
    (a) Two waves become out of phase in the direction of the beam at a distance $L_L$ from being in phase after $N/2$ wavelengths for $\lambda$ and $(N+1)/2$ wavelengths for $\lambda - \Delta\lambda$. Since $N >> 1$, we can write $N\approx\lambda/\Delta\lambda$ and find eq. \ref{eq:LongitudinalCoL}.
    (b) A plane wave emitted from one extremity of the source is out of phase with another plane wave emitted from the other extremity with maximum angular divergence $\Delta\theta$ at a distance $L_T$ in the plane perpendicular to the direction of the beam situated at a distance $R$ from the source.
    }
    \label{fig:CoherenceLengths}
\end{figure}

The longitudinal coherence length $L_L$ describes the distance in the direction of the beam after which two waves of slighly different wavelength become out of phase due to the small extent of the beam bandwidth $\Delta\lambda$ (fig. \ref{fig:CoherenceLengths} - a).
The bandwidth is linked to the quality of the monochromator used to obtain a monochromatic beam through the relation $\lambda/\Delta\lambda$ (eq. \ref{eq:LongitudinalCoL}).

\begin{equation}
    \label{eq:LongitudinalCoL}
    L_L = \frac{\lambda^2}{2\Delta\lambda}
\end{equation}

The transverse coherence length $L_T$ describes the distance in the plane perpendicular to the beam after which two waves become out of phase due to the finite source size $D$ and to the angular divergence of the beam $\Delta\theta$ (fig. \ref{fig:CoherenceLengths} - b).
By assuming that the source has a Gaussian profile, we can relate $D$ to the standard deviation of the beam $\sigma_{x,y}$ in the $\vec{x}$ and $\vec{y}$ directions perpendicular to the beam direction (\cite{Willmott}, eq. \ref{eq:TransverseCoL}).

\begin{equation}
    \label{eq:TransverseCoL}
    L_{T,(x,y)} = \frac{\lambda}{2\Delta\theta_{(x,y)}} = \frac{\lambda R}{2 D_{(x,y)}} = \frac{\lambda R}{2\sqrt{\pi}\sigma_{x,y}}
\end{equation}

At synchrotrons, the vertical and horizontal source size are significantly different due to the shape of the electron beam, much more extended in the horizontal plane.
The horizontal coherence length is usually smaller and the limiting factor.
By multiplying the two transverse coherence lengths $L_{T,x}$, $L_{T,y}$ and the longitudinal coherence length $L_L$ we obtain the coherent volume.
An example if given in sec. \ref{sec:SIXS} for the SixS beamline of the synchrotron SOLEIL.

To perform Bragg coherent diffraction imaging experiments for which the coherent volume must be larger than the sample, the beamline must have low vertical and horizontal emittances, high monochromaticity and a sample environment situated as far as possible from the source.
Indeed, to be able to resolve the interference fringes resulting from the lattice factor, a high brilliance is crucial since the intensity of the total scattered field, proportional to the scattering volume, is weak.
A detailed explanation of how to optimize BCDI experiments is given in sec. \ref{sec:BCDI}.

\textit{
So far we have limited ourselves to the kinematical approach of diffraction in the frame of this thesis which means that the scattering is considered weak enough so that multiple scattering can be ignored.
This allows us to assume that the sample is only a perturbation to the incident field and that the structure factor can be written as the Fourier transform of the electronic density (Born approxiamtion).
We have also assumed that the incident wave is plane within the scattering volume and that the distance between the scatterer and the detector is large enough for the scattered waves to be considered planes.
}

\textcolor{Important}{
Field / amplitude / intensity consistent ?
$r_0$ use consistent ? $2\pi$ ?
hypothesis to explain: DWBA, dipole, perfect crystals
}

    \section{SXRD}


\subsection{Surface reconstructions}


\subsection{Crystal truncation rods}

Thus the diffraction intensity of the finite-sized crystal has diffuse streaks connecting all the Bragg points. The diffuse intensity far from the nodes is of order of magnitude $N^4$ compared with $N^6$ at the nodes.

Scattering that is sharp in two directions and diffuse in the third (referred to as a "rod" of scattering) must arise from a crystalline object that is localized in one dimension and extended in the other two.

We are then left with only the sixth component due to the sharply truncated surface. We will call these features "crystal truncation rods. '

We now wish to estimate the strength of the truncation rods in the Bragg geometry. We must first modify Eqs. (1) and (2) by including the x-ray coherence length, $m$ (measured in unit cells), of the experimental configuration. This broadens all the diffraction features to $\frac{1}{m}$ reciprocal units. The Bragg points then have intensity of order $N_1 N_2 N_3 m^3$, while the diffuse intensity is $N_1 N_2 m^2$

A typical penetration depth is $1 \mu m$ so $N_3 =10^3$ unit cells (perpendicular to the face). With $m \approx 100$ unit cells, this gives a relative intensity
$\frac{I(Bragg peak)}{I(truncation rod)} = N_3 m \approx 10^5$

More roughness means wider Bragg peaks and deeper valleys between the BP.

It is not clear that different detailed models of roughness could be distinguished at this level of accuracy. One central concept to all descriptions of crystal truncation rods, however, is the continuation of the crystal lattice into the roughened region (not defects, keeping symmetry).


\subsection{Reflectivety}


\subsubsection{ROD}

    \section{Bragg coherent diffraction imaging} \label{sec:BCDI}

Bragg Coherent Diffraction Imaging (BCDI) \parencite{Robinson2001, Pfeifer2006, Robinson2009} is a powerful technique for the non-destructive characterisation of material structure in three dimensions.
It reaches unparalleled spatial and strain resolution for a scattering technique, few nanometres \parencite{Labat2015, Cherukara2018a} and \num{e-4} respectively \parencite{Newton2010, Lauraux2020}.

The BCDI method is reliant on the coherence of the light available, thus only began to be exploited at third generation synchrotron sources \parencite{Miao1999, Miao2000, Robinson2001, Labat2007, Robinson2009, Vaxelaire2010, Chamard2010, Clark2012, Clark2013, Yang2013, Xiong2014}, also the subject of reviews \parencite{Nugent2010, Miao2015}.
Since then it was developed into a characterisation tool for \textit{in situ / operando} studies of materials structure, mostly applied to three fields: electrochemistry \parencite{Ulvestad2015}, heterogeneous catalysis \parencite{Ulvestad2016}, and the structure and evolution of defects \parencite{Labat2015}.
BCDI has gained from further source and beamline improvements as many synchrotrons have recently completed or are in the process of an upgrade from $3^{rd}$ to $4^{th}$ generation synchrotrons, which is predicted to allow the study of \textit{in situ} dynamical phenomena \parencite{Lo2018}.
Focusing optics that allowed the development of coherent imaging techniques by increasing the coherent flux on the probed sample will be discussed in this section.

\subsection{Scattering by strained crystals}\label{sec:StrainBCDI}

In this thesis, the intensity scattered by a single crystal of dimensions below the coherent volume of the source is considered to introduce BCDI.
The scattered waves interfere with each other, and maximum intensity is observed when the scattering vector satisfies the Laue condition.

Note that the following derivations only apply if the diffraction pattern is viewed at a distance far away from the diffracting object, in a region known as the far-field or Fraunhofer region \parencite{Willmott}, where the scattered waves can be considered as plane waves.
Moreover, the scattering volume is assumed to be in the focal plane of the focusing optics and far from the source so that the waves in the scattering volume can also be described as plane waves.
Finally, the scattering volume is assumed to be small with dimensions below the x-ray extinction length to ignore multiple scattering, and stay in the kinematical approach of diffraction.
The incident field is assumed to be perfectly coherent in the longitudinal and transverse direction, with a homogeneous amplitude in the scattering volume.
A discussion of the extent of these approximations is given in Godard et al . \parencite*{Godard2021}.
% only elastic scattering and no lattice variations

The community has recently been investigating a dynamical approach to Bragg Coherent Diffraction Imaging (BCDI) \parencite{Yan2014, Shabalin2017, Hu2018, Gao2022} using the Takagi-Taupin equations \parencite{Takagi1962, Takagi1969}.
This method tends to validate the kinematical approach when working with slightly strained sub-micron sized crystals \parencite{Karpov2019, Barringer2021}.
The phase of the incoming and scattered waves must however be corrected for refraction and absorption \parencite{Harder2007, Gao2022} which does not contradict the main hypothesis of the kinematical approach which is that in the Born approximation, the structure factor of the crystal is equal to the Fourier transform of its electronic density \parencite{Paganin}.

Another approach is to write the structure factor as the sum of the atomic scattering factor multiplied by a phase component, over each atom present in the crystal (eq. \ref{eq:FCrystal1}).
In general, the Fourier transform approach is preferred since one can use quick algorithms such as the Fast Fourier Transform \parencite{Cooley1965, Cochran1967} to compute the structure factor, whereas the atomistic approach relies on knowing the position of each atom in the crystal.

As seen in sec. \ref{sec:StructureFactor}, when considering the crystal to be made of repeating unit cells that organise themselves on the crystal lattice, the structure factor can also be written as the convolution between the Lattice factor $F_{lat}$ and the unit cell structure factor $F_{uc}$.
The shape of the scattering intensity around Bragg peaks is then given by $F_{lat}$, which is the Fourier transform of the crystal's lattice, whereas $F_{uc}$ is the Fourier transform of the crystal's unit cell electronic density.

\begin{figure}[!htb]
    \centering
    \includegraphics[width=0.66\textwidth]{/home/david/Documents/PhD/Figures/Introduction/StrainedCrystal.pdf}
    \caption{
    Simplified representation of a crystal in which each dot corresponds to a unit cell.
    The red dots represent displaced unit cells from their ideal position by a displacement vector $\vec{u}_k$, resulting in a shift in the phase of the scattered x-rays.
    }
    \label{fig:StrainedCrystals}
\end{figure}

In the case of strained crystals (fig. \ref{fig:StrainedCrystals}), eq. \ref{eq:Fcrystal} can be rewritten as follows:

\begin{align}
    \label{eq:FcrystalBCDI1}
    F_{crystal}(\vec{q}) & = \sum_j^{N_{atoms, uc}} f_j(\vec{q}) e^{i\vec{q}.(\vec{r}_{0,j} + \vec{u}_{j})} \sum_k^{N_{uc}} e^{i\vec{q}.\vec{R}_k}\\
    \label{eq:FcrystalBCDI2}
    F_{crystal}(\vec{q}) & = \sum_j^{N_{atoms, uc}} f_j(\vec{q}) e^{i\vec{q}.\vec{r}_{0,j}} \sum_k^{N_{uc}} e^{i\vec{q}.(\vec{R}_k + \vec{u}_{k})}
\end{align}

where $\vec{r}_{0,j}$ is the ideal position of the atom $j$ in the unit cell, and $\vec{u}_{j}$ the displacement from that position.
In eq. \ref{eq:FcrystalBCDI2}, all atoms in the unit cell are considered to have the \textit{same displacement} to be able to shift the displacement parameter from the atomic position to the unit cell position, $\vec{u}_{k}$ becomes the displacement of the unit cell $k$.

To reduce the complexity in the second part of eq. \ref{eq:FcrystalBCDI2}, only the intensity in a volume situated near the Bragg peak is considered so that $\vec{q} \approx \vec{G}$, $\vec{G}$ designates the reciprocal space vector that corresponds to the position of the Bragg peak in reciprocal space \parencite{Pfeifer2006, Minkevich2007, Harder2007}.

\begin{align}
    \label{eq:FcrystalBCDI3}
    F_{crystal}(\vec{q}) & = \sum_j^{N_{atoms, uc}} f_j(\vec{q}) e^{i\vec{q}.\vec{r}_{0,j}} \sum_k^{N_{uc}} e^{i\vec{q}.\vec{R}_k} e^{i\vec{q}.\vec{u}_k}\\
    \label{eq:FcrystalBCDI4}
    F_{crystal}(\vec{q}) & = \sum_k^{N_{uc}} F_{uc}(\vec{q}) e^{i\vec{q}.\vec{R}_k} e^{i(\vec{G}+\vec{q}-\vec{G}).\vec{u}_k}\\
    \label{eq:FcrystalBCDI5}
    F_{crystal}(\vec{q}) & = \sum_k^{N_{uc}} F_{uc}(\vec{q}) e^{i\vec{q}.\vec{R}_k} e^{i\vec{G}.\vec{u}_k} \times e^{i(\vec{q}-\vec{G}).\vec{u}_k}\\
    \label{eq:FcrystalBCDI6}
    F_{crystal}(\vec{q} \approx \vec{G}) & = \sum_k^{N_{uc}} F_{uc}(\vec{G}) e^{i\vec{q}.\vec{R}_k} e^{i\vec{G}.\vec{u}_k}\\
    \label{eq:FcrystalBCDI7}
    % F(\vec{q} \approx \vec{G}) & = DFT[\rho(\vec{r}) e^{i\vec{G}.\vec{u}_k}]
    F_{crystal}(\vec{q} \approx \vec{G}) & = DFT[ F_{uc}(\vec{G}) e^{i\vec{G}.\vec{u}_k}]_{N_{uc}}
\end{align}

It is assumed that the variations of the unit cell structure factor $F_{uc}(\vec{q})$ are negligible when $\vec{q} \approx \vec{G}$, so that $F_{uc}(\vec{q})$ becomes $F_{uc}(\vec{G})$.
% $F_{uc}(\vec{q})$ can be linked to its electronic density by eq. \ref{eq:AtomicFormFactor}.
% The first part of eq. \ref{eq:FcrystalBCDI6} can be considered to be the Fourier transform of the unit cell $k$ of electronic density $\rho(\vec{r})$ in which all atoms have the same displacement vector $\vec{u}_k$.
% In this equation the electronic density is allowed to vary from one unit cell to another.
One can recognise the expression of a discrete Fourier transform (DFT) of frequency $\vec{G}$, coefficient $F_{uc}(\vec{G}) e^{i\vec{G}.\vec{u}_k}$, and $N_{uc}$ samples \parencite{Cooley1965, Cochran1967, FavreNicolin2011a, Godard2021}.
The approximation $\vec{q} \approx \vec{G}$ allows the use of a Fast Fourier Transform (FFT) algorithm to compute the DFT, with a linearithmic time complexity $O(n \, log(n))$.
This approximation breaks down far away from the Laue condition, or if the strain is so important that the condition $(\vec{q}-\vec{G}).\vec{u}_k<<1$ is not anymore respected to move from eq. \ref{eq:FcrystalBCDI5} to eq. \ref{eq:FcrystalBCDI6} \parencite{Takagi1969}.

In the case of highly strained crystal, or when the intensity must be computed far away from Bragg peaks, it is still possible to resort to so-called \textit{atomistic} simulations, by computing the scattering intensity directly from the position of each atom in the crystal (eq. \ref{eq:FCrystal1}).
However, the position of each atom must be known and, for a large crystal, the sum becomes too large to handle even with high-performance clusters due to a quadratic time complexity $O(n^2)$ (n is the number of operations).
Atomistic simulations are usually used when studying crystals of dimensions below \qty{100}{\nm} \parencite{Dupraz2022}.

For a non-strained crystal, the Lattice factor is the same at each Bragg peak, equal to the Fourier transform of the crystal's lattice.
In the case of a strained crystal, the additional phase factor introduced in eq. \ref{eq:FcrystalBCDI7} breaks the centrosymmetry around each Bragg peak, due to the heterogeneous displacement field in the crystal.
The effect of the displacement can be seen in the distortion of the scattering intensity around Bragg peaks.
It is important to notice that the phase in eq. \ref{eq:FcrystalBCDI7} is equal to the dot product between the displacement vector and the reciprocal space vector.
This means that to possess the complete displacement field, one must collect the scattered intensity around three non-coplanar Bragg peaks \parencite{Newton2010} that will yield the displacement field over an orthogonal basis ($\vec{q}_x, \vec{q}_y, \vec{q}_z$).
The electronic density is also sometimes named Bragg electronic density to illustrate this sensitivity.

\subsubsection{Data collection} \label{sec:DataCollectionBCDI}

The data collection during a BCDI measurement is performed by continuously rocking the sample around the incident angle, which will result in a tilt of the associated Ewald sphere, so that the detector probes a different slice of the reciprocal space at each step.
This measurement is called a rocking curve (fig. \ref{fig:3DDP} - left), the associated 3D scattering intensity is illustrated in fig. \ref{fig:3DDP} (right).

\begin{figure}[!htb]
    \centering
    \includegraphics[width=0.39\textwidth]{/home/david/Documents/PhD/Figures/Introduction/RockingCurve.png}
    \includegraphics[width=0.59\textwidth]{/home/david/Documents/PhD/PhDScripts/SixS_2021_06_BCDI_NH3/figures/3D_DP/DP_white_1.png}
    \caption{
    Left: rocking curve for the measurement of the 3D scattered intensity around a Bragg peak. Taken from \cite{Willmott}.
    Right: 3D diffraction pattern collected at SixS (SOLEIL).
    Different iso-surfaces from high (red) to low (yellow) intensity are plotted, highlighting the direction dependent drop of the scattered intensity as a function of $q_{hkl}+\delta q$ where $\delta q$ is the distance from the centre of the Bragg peak.
    }
    \label{fig:3DDP}
\end{figure}

The measurements are commonly collected in angular space, and can be converted into $q$-space to compare different Bragg peaks in the reciprocal space.
However, this step is usually avoided since the intermediate interpolation process can be inaccurate, especially when the scattering intensity is low.
It is also possible to collect 3D measurements by tuning the energy of the incident beam \parencite{Cornelius2011}.

\subsection{Phase retrieval}\label{sec:PhaseRetrieval}

As seen in sec. \ref{sec:DataCollectionSXRD}, the diffracted intensity collected by the detector is proportional to the squared modulus of the structure factor $F_{crystal}(\vec{q})$ (eq. \ref{eq:ScatteredIntensity}), the phase of the scattered x-rays is lost during the measurement.

Without the phase information in reciprocal space, it is impossible to directly reconstruct the crystal electronic density with the help of an inverse discrete Fourier transform (IDFT, from eq. \ref{eq:FCrystal1}), and to obtain a detailed structural model of the crystal.

\begin{equation}
    \label{eq:IDFT}
    \rho(\vec{r}) = IDFT [F_{crystal}(\vec{q})]
\end{equation}

It has been demonstrated that when fulfilling the \textit{oversampling} condition, the phase lost during the measurement can be retrieved \parencite{Shannon1949, Sayre1952}.
The oversampling condition refers to the practice of sampling a signal at a higher sampling rate than the Nyquist rate \parencite{Miao2000}, which is twice the highest frequency present in the signal.
Increasing the sampling rate in reciprocal space results in a denser set of data points in real space.

In the case of low-strained faceted crystals, the oversampling $\sigma$ can be derived by computing the distance between the fringes (\textit{i.e.} the interfringe) on rods in directions parallel to the normal of crystalline facets.
It was shown that the interfringe $\delta q$ observed on the rods is proportional to the size of the crystal in that direction, $t$, following the relation $\delta q = \frac{2\pi}{t}$ (sec. \ref{sec:LatticeFactor}).
For small crystals ($t<$ few \unit{um}), the interfringe can be resolved depending on the experimental parameters.
It thus becomes possible to oversample the fringes and to recover the phase information.
The instrumental requirement for oversampling in the direction of a rod is that the detector must have at least two pixels between each fringe: $\sigma >2$.
% There is also the possibility to adapt the detector distance to get a good oversampling.
% At \qty{8}{\keV}, to get an oversampling of 3, the detector must be positioned at \qty{6.4}{\m} for a particle of \qty{6}{\um}.

In three dimensions, the oversampling condition becomes $\sigma>2^{1/3}$ \parencite{Miao1998, Miao2000, Miao2000a}.
When it comes to Fourier transforms, higher oversampling can offer several advantages such as reducing the possibility of aliasing (high-frequency components being mistakenly represented as lower frequencies due to insufficient sampling), spectral leakage (frequency components spread into neighbouring frequency bins in the Fourier transform, minimised if the frequency bins become narrower).
Moreover, with more densely spaced samples, the interpolation process can more faithfully reconstruct the original signal without introducing artefacts.
However, high oversampling rate take more measurement time, and can for example limit the number of rocking curves that can be recorded.
An oversampling of three in each direction is in general sufficient to successfully reconstruct the probed samples \parencite{Dupraz2015}.
The oversampling condition partly dictates the instrumental parameters which define the extent of the data voxels, the small volumes that constitute the 3D data array.

\begin{figure}[!htb]
   \centering
   \includegraphics[width=\textwidth]{/home/david/Documents/PhD/Figures/Introduction/PhaseRetrieval.pdf}
   \caption{
   Iterative phase retrieval algorithms converge towards a solution by the means of real and reciprocal space constraints.
   $\rho'(\vec{r})$ is the updated Bragg electronic density after application of the support constraint.
   }
   \label{fig:PRAlgo}
\end{figure}

Fienup et al. \parencite*{Fienup1978, Fienup1982, Fienup1986} have laid the ground for computer methods permitting phase retrieval, once the oversampling condition has been satisfied.
The error-reduction (ER) algorithm works by iterating between real and reciprocal space by the means of fast Fourier transforms (FFT) and inverse fast Fourier transforms (IFFT) methods.
The process starts by assigning random phases to the square root of the measured scattered intensity ($I(\vec{q})$), to create a first guess of the scattered amplitude: $A(\vec{q}) = \sqrt{I(\vec{q})}e^{i\Phi(\vec{q})}$, proportionnal to the structure factor.
The inverse Fourier transform yields a complex object $\rho(\vec{r})$ of amplitude proportionnal to the electronic density of the crystal (eq. \ref{eq:IDFT}), and with an additionnal phase factor linked to the displacement of the unit cells (eq. \ref{eq:FcrystalBCDI7}).
A real space constraint is then applied by selecting a region above a certain threshold $t$ of the modulus of the electronic density which is called the \textit{support}.
The electronic density outside the support is set to zero (eq. \ref{eq:ElectronicDensity}).
The scattered amplitude is then computed with the Fourier transform of the electronic density, the following reciprocal space constraint consisting in replacing the amplitude of the computed amplitude by the square root of the measured intensity.

\begin{equation}
    \label{eq:ElectronicDensity}
    \rho(\vec{r}) =
        \begin{cases}
            0  & \text{if} \; |\rho(\vec{r})| < t \\
            \rho(\vec{r}) & \text{if}  \; |\rho(\vec{r})| \geq t
        \end{cases}
\end{equation}

By iterating between the real and reciprocal space, the algorithm converges towards a final solution for the electronic density (fig. \ref{fig:PRAlgo}).

A classical error-metric used to follow the convergence of the algorithm is the minimisation of the residuals between the square root of the measured intensity, and the amplitude computed from the Fourier transform of the electronic density.
The main drawback of this error-metric is that each voxel is given the same weight, whereas the impact of the approximation detailed previously depend on the distance from the Bragg peak \parencite{Godard2021}.

More algorithms have since then been developed to help converge towards the final solution without being stuck in local minima such as the hybrid input-output (HIO) and relaxed averaged alternating reflections (RAAR) algorithms, by introducing a new parameter, $\beta$, which relaxes the threshold condition in real space from the ER algorithm \parencite{Marchesini2003, Luke2005, Marchesini2007}.

Finally, a major recent improvement of the phase retrieval algorithm is the possibility to take into account the partial coherence of the beam \parencite{Sinha1998, Vartanyants2001, Williams2007, Whitehead2009, Nugent2010, ChenBo2012} together with the point-spread function of the detector by using the deconvolution Richardson-Lucy algorithm \parencite{Richardson1972, Lucy1974, Fish1995, Clark2012}.
The point-spread function (PSF) must first be approximated with a 3D Gaussian, Lorentzian or pseudo-voigt shape, depending on the expected coherence of the beam.
It is refined during the last steps of the phase retrieval when the support is already well determined, to avoid any divergence from simultaneously having to refine both the support and the PSF.

Moreover, new methods such as convolution neural networks (CNN) are under study and have started to show some encouraging results, but are not yet sufficiently robust, particularly in the case of strained particles \parencite{Cherukara2018, Shen2019, Chan2021, Kim2021a, Wu2021a}.
%Shifting the use of convolutional kernels to retrieve the phase of the scattered x-rays directly in reciprocal space rather than the Bragg electronic density could also be one way phase obparticle

\subsection{Resolution}

The resolution of the reconstructed object, defined according to multiple parameters in the literature, is a prominent figure of merit for an imaging technique.
The voxel size in real space ($2\pi / \delta q_x, 2\pi / \delta q_y, 2\pi / \delta q_z$) depends on the extent of the collected volume in the reciprocal space ($\delta q_x, \delta q_y, \delta q_z$).
However, the voxel size would only be a good criteria to determine the resolution if the signal to noise ratio of the scattering intensity was important enough in all of the probed reciprocal space volume, so that all of the frequencies in the $q$-space could be effectively retrieved with the FFT algorithm.
A high-intensity dynamic range detector is therefore needed to correctly probe intensities at the center, and extremities of the Bragg peaks \parencite{Latychevskaia2018}.

This is the case in simulated data, but not with experimental data, for which the noise starts to become important far away from the Bragg peak \parencite{Bikondoa2021}, keeping in mind that the intensity decreases as a function of $q^{-4}$ as a function of the distance from the Bragg peak \parencite{Marchesini2003a}, depending also on the form factor \parencite{Croset2017}.
The community has therefore relied so far on a few methods to estimate the spatial resolution of the reconstructed electronic density.

First, the Phase Retrieval Transfer Function (PRTF) \parencite{Chapman2006} is the \textit{ratio} of the calculated amplitude $|A_k(\vec{q})|$ to the square-root of the measured intensity $\sqrt{I_k(\vec{q})}$ as a function of the resolution shell.
The resolution shell $k$ is defined as all voxels of position $q=\sqrt{q_x^2 + q_y^2 +q_z^2}$ that satisfy the condition $q \in [q_k, q_{k+1}]$.
$\delta q$ is the width of the resolution shell which divides the extent $[q_{min}, q_{max}]$ of the sampled reciprocal space in $N$ shells so that $q_{max} = q_{min} + N \times \delta q$ and $q_{k+1} - q_{k} = \delta q$.

\begin{equation}
    PRTF(q_k, q_{k+1}) = \frac{|<F_k(\vec{q})>|}{<\sqrt{I_k(\vec{q})}>}
\end{equation}

Secondly, the Fourier shell correlation (FSC) \parencite{VanHeel2005} measures the \textit{normalised cross-correlation coefficient} between the amplitude of two independent reconstructions, also as a function of the resolution shell, but from two independent datasets collected with the same instrumental parameters:

\begin{equation}
    FSC(q_k, q_{k+1}) = \frac{ |\sum F_{1,k} F_{2,k}^*| }{\sqrt{ \sum F_{1,k} F_{1,k}^* \times \sum F_{2,k} F_{2,k}^*}}
\end{equation}

Both of these methods allow us to quantify the distance from the centre of the Bragg peak after which one of the criteria diverges, which yields a corresponding real space resolution.
There are several limitations to these methods.
First, the resolution curve as a function of $|\vec{q}|$ strongly depends on the width of the bins used to create the resolution shells.
The number of voxels in each bin is dependent on the distance from the centre $q_{min}$, and the voxel size is usually not isotropic.
Finally and most importantly, it supposes that the resolution is isotropic in reciprocal space, which is false since the signal to noise ratio will be much higher around the rods perpendicular to the crystal's facets \parencite{Cherukara2018a}.
Indeed, the scattered intensity along the rods is proportional to the size (\textit{i.e.} the number of unit cells) of the crystal in the direction perpendicular to the facet.
Thus, for large facets, the rod intensity will be much more intense, increasing the resolution in that direction.

In real space, the spatial resolution can be quantified by differentiating line profiles of electron density amplitude across the object-air interface, and fitting these with a Gaussian profile.
The average 3D spatial resolution is taken as 2$\sigma$ of the fitted Gaussian \parencite{Hofmann2020}.
This method has the advantage of being straightforward, and direction sensitive.

One of the challenges encountered in Bragg coherent diffraction imaging is establishing a single, meaningful criterion for the spatial resolution.
Its anisotropic nature further complicates this task, as commonly used criteria such as the phase retrieval transfer function (PRTF) or Fourier shell correlation (FSC) are isotropic.
To address this, a new method for resolution assessment when aiming at the analysis of specific facets in real space could involve examining the extension of the corresponding crystal truncation rods in reciprocal space, before the signal becomes overwhelmed by counting noise from the lack of scattered photons.
In this approach, the resolution of the measurement would not rely on the quality of the phase retrieval or the type of algorithms utilised, but solely on the quality of the measurement itself.
Consequently, the significance of the sample's shape, as well as the measurement quality would be emphasised.
For instance, a large planar particle featuring a sizeable (111) top facet would yield a more readily resolvable signal in the (111) direction compared to a smaller, spherical particle with a diminished (111) top facet.

Concerning the resolution of the retrieved atomic displacement, Labat et al. \parencite*{Labat2015} have demonstrated a displacement field accuracy of few \unit{pm} along preferential directions, by comparing with simulated displacement fields.
Another method is to work in a similar way to the Fourier shell correlation by replacing the shells of amplitude as a function of $\vec{q}$ by shells of strain as a function of the strain magnitude, and see if there exists a strain threshold below which the correlation drops.
This method is then called the strain shell correlation \parencite{Girard2020}.

However, one must keep in mind that the upper threshold for the strain resolution exists based on the $(\vec{q}-\vec{G}).\vec{u}_k<<1$ condition, which mostly affects the quality of the phase retrieval.
Highly strained crystals are very difficult to reconstruct and represent the next milestone in Bragg coherent diffraction imaging, recent methods have been proposed for specific cases such as phase transitions \parencite{Wang2020}.
One could imagine a future in which the atomistic computation of the scattering factor has been so optimised that phase retrieval algorithms would only be needed to provide an initial guess to the electronic density, replaced in the final steps by a fitting of the computed amplitude to the square root of the measured intensity.
A reverse Monte-Carlo method used already for the simulation of powder diffractograms could also be used \parencite{RLMcGreevy2001}.

\subsection{Facet analysis} \label{sec:FacetAnalysis}

After phase retrieval, it becomes possible to exploit the electronic density of the reconstructed crystal.
Its surface corresponds to a layer of voxels in the three dimensional array.
Carnis et al. \parencite*{Carnis2019} have presented guidelines on how to correctly select the surface voxels, which are detailed below.
The histogram in fig. \ref{fig:histo} represents the number of voxels per bin of normalised electronic density, the observed peak corresponds to the bulk voxels for which the amplitude of the electronic density is close to the maximum (fig. \ref{fig:histo} - right).
An empiric criterion to select the voxels that correspond to the crystal's surface is to take the lowest foot of that peak, here around \num{0.6}, the surface being defined as the hull surrounding the bulk voxels.

\begin{figure}[!htb]
   \centering
   \includegraphics[height=4.5cm]{/home/david/Documents/PhD/PhDScripts/SixS_2021_06_BCDI_NH3/figures/B-7/hist_iso.pdf}
   \includegraphics[height=4.4cm]{/home/david/Documents/PhD/PhDScripts/SixS_2021_06_BCDI_NH3/reconstructions/vti_good_files/D6/300/AmplitudeLines.png}
   \caption{
   Histogram representing the distribution of the electronic density amplitude, normalised by its maximum (left).
   Values below 0.05 are ignored since they typically belong to voxels located far away from particle support.
   Amplitude of the electronic density in a slice of the reconstructed object (right).
   Two curves are drawn at the iso-surface values of \num{0.1} (outer curve - black ) and \num{0.6} (inner curve - red).
   }
   \label{fig:histo}
\end{figure}

Once the threshold for the iso-surface is selected, it is possible to visualise a contour of the object with specialised solutions such as \textit{Paraview} (\cite{Ahrens2001} - fig. \ref{fig:FacetsParaview}) or \textit{Gwaihir} (\cite{Simonne2022} - fig. \ref{fig:3D_object}).

The surface is usually created by using the Marching-Cubes algorithm \parencite{Lorensen1987}, which works by assigning a scalar value to each voxel of the data (in our case the amplitude of the retrieved complex electronic density).
An iso-value which acts as a threshold is selected by the user, separating the volume of the particle from the rest.
Finally, by \textit{marching} over the volume of the array, the algorithm derives a surface representation by assigning a configuration to each point in the array that depends on whether or not the 8 neighbouring cubes have values above or below the iso-value (fig. \ref{fig:MarchingCubes}).

\begin{figure}[!htb]
    \centering
    \includegraphics[width=0.66\textwidth]{/home/david/Documents/PhD/PhDScripts/SixS_2021_06_BCDI_NH3/figures/MarchingCubes.png}
    \caption{
    The Marching-cubes algorithm generates triangles connecting the vertices to form the mesh surface.
    Each selected configuration has a specific arrangement of vertices that dictate how to form triangles to create a smooth surface.
    Image taken from \url{https://fr.wikipedia.org/wiki/Marching_cubes}
    }
    \label{fig:MarchingCubes}
\end{figure}

Crystallographic facets can be identified on the surface of the reconstructed object when studying faceted objects with a highly coherent beam \parencite{Richard2018}, allowing in-depth studies of facet dependent strain and displacements.
Lattice strain and displacements are key information to retrieve during \textit{in situ} and \textit{operando} experiments in fields such as heterogeneous catalysis \parencite{Ulvestad2016, Yau2017a, Kim2018, Fernandez2019, Abuin2019, Kim2019, Kawaguchi2019, Suzana2019, Choi2020, Passos2020,  Kim2021, Carnis2021, Dupraz2022} or electrochemistry \parencite{Ulvestad2015a, Bjorling2019, Vicente2021, Kawaguchi2021, Carnis2021b, Atlan2023}.

\begin{figure}
    \centering
    \includegraphics[width=0.49\textwidth]{/home/david/Documents/PhD/PhDScripts/SixS_2021_06_BCDI_NH3/reconstructions/vti_good_files/D6/300/FacetProbability.png}
    \includegraphics[width=0.49\textwidth]{/home/david/Documents/PhD/PhDScripts/SixS_2021_06_BCDI_NH3/reconstructions/vti_good_files/D6/300/FacetIDs.png}
    \caption{
    a) Surface of a reconstructed nanoparticle ($\diameter \approx \qty{800}{\nm}$) with each surface voxel coloured by its probability to belong to a facet.
    The same surface is represented in b) with a hull in which each facet is drawn, delimited by thick white lines.
    }
    \label{fig:FacetsParaview}
\end{figure}

\begin{table}
    \begin{center}
    \resizebox{0.95\textwidth}{!}{%
    \begin{tabular}{@{}lllll@{}}
    \toprule
    Facet id & Facet normal & Relative facet size & Average displacement (\unit{\angstrom}) & Average strain (\num{e-4})\\
    \midrule
    0 & Edges and corner            & NaN         & \num{-0.224 \pm 0.259} & \num{0.20 \pm 0.86}\\
    1 & $(1 1 1)$                   & \num{0.106} & \num{0.080 \pm 0.173}  & \num{-0.28 \pm 0.51}\\
    2 & $(\bar{1} \bar{1} \bar{1})$ & \num{0.223} & \num{0.052 \pm 0.260}  & \num{-1.17 \pm 0.64}\\
    3 & $(1 \bar{1} 0)$             & \num{0.106} & \num{0.080 \pm 0.173}  & \num{-0.28 \pm 0.51}\\
    4 & $(1 0 0)$                   & \num{0.096} & \num{0.137 \pm 0.192}  & \num{0.27 \pm 0.49}\\
    \bottomrule
    \end{tabular}
    }
    \end{center}
    \caption{
    The output of \text{FacetAnalyser} is a list of values for each facet.
    The accessible features are the facet size, the average strain and displacement (along the [111] direction, for the outermost voxel), and the coordinates of the normal to the facet surface.
    The uncertainty on the average displacement and strain corresponds to the standard deviation of the displacement and strain distribution, respectively.
    }
    \label{tab:facets}
\end{table}{}

Retrieving the facets can be achieved by analysing the probability distributions of the orientations of triangle normals on a mesh representation of the object \parencite{Grothausmann2012}.
This method is used in the \textit{Paraview} plugin \textit{FacetAnalyser} \parencite{Grothausmann2015}, and yields a list of features detailed in tab. \ref{tab:facets}.
The probability of each surface voxel to belong to a facet is illustrated in fig. \ref{fig:FacetsParaview}.
The orientation of each facet is retrieved by analysing the angles between the normal of each facet, also returned by the algorithm, and a reference direction.

In the case of highly resolved measurement, another approach could be inspired from image processing and convolutional neural networks, where 3D convolutions kernels are used to identify surfaces \parencite{RaschkaMirjalili2019}.
By tuning the kernel size and values, it is possible to be sensitive to boundaries in specific directions.

In this example, the FacetAnalyser plugin is used (fig. \ref{fig:FacetsParaview}).
Note that the edges and corners of the reconstructed particle are also retrieved together as the voxels not belonging to any facets.
The contribution from atomic edges and corners to those voxels, sites of particular interest for heterogeneous catalysis \parencite{Taylor1925}, depends on the spatial resolution of the experiment since averaged in the voxel size.
The result depends on the algorithm input parameters, such as the minimum relative facet size, the angular acceptance for the facet normals, \textit{etc.}
    \section{Synchrotron radiation for the study of materials} \label{sec:SIXS}

\subsection{Synchrotron radiation}

SOLEIL (Source Optimisée d’Énergie Intermédiaire du LURE (Laboratoire pour l’Utilisation du Rayonnement Électromagnétique)) is a $3^{rd}$ generation synchrotron source facility build in 2006 and localised near Paris, France, which operates at an electron energy of \qty{2.75}{\GeV} with a storage ring diameter of \qty{354}{\m}.

\begin{figure}[!htb]
    \centering
    \includegraphics[width=0.49\textwidth]{/home/david/Documents/PhD/Figures/sixs/RingSoleil.jpg}
    \includegraphics[width=0.49\textwidth]{/home/david/Documents/PhD/Figures/sixs/BeamlineSoleil.jpg}
    \caption{
    	The storage ring of a synchrotron alternates between straight and curved sections (left), insertion devices are situated at the end of straight sections and deliver the synchrotron radiation to the beamline, separated in optical hutches, experimental end-stations and working stations (right).
    	Copyright by SOLEIL (\url{https://www.synchrotron-soleil.fr/en/research}).
    }
    \label{fig:SOLEIL}
\end{figure}

Synchrotron radiation is generated when electrons travelling at relativistic speeds are accelerated and forced to travel in curved trajectories by strong magnetic fields supplied by \textit{bending magnets}.
The radiation spans a broad spectrum, from infrared to X-rays, depending on the electrons, the acceleration, and the strength of the magnetic fields.

First, electrons are produced by an electron gun and accelerated in a linear accelerator (LINAC), then by a booster ring.
Secondly, once the electrons reach the desired energy, they are injected into a storage ring, which is an almost circular vacuum chamber (fig. \ref{fig:SOLEIL} - left) surrounded by strong magnets.
The magnetic fields in the ring act as a guiding force, bending the electrons' trajectories.

Due to their high energy, the electrons travel at speeds close to the speed of light, making them relativistic particles.
As they move through the curved paths, they are accelerated and emit electromagnetic radiation tangentially to their trajectory as a result \parencite{Willmott, NielsenMcMorrow}.

By the use of insertion devices such as \textit{wiggler} or \textit{undulators}, synchrotron radiation is exceptionally bright and brilliant, with a high flux of photons and an excellent collimation, wigglers are best used when in the need of a broad energy range, whereas undulators produce beams with extremely high brilliance and narrow bandwidth.
Each beamline at SOLEIL begins with one of those insertion devices or with a bending magnet (fig. \ref{fig:SOLEIL} - right).

\subsection{The SixS beamline}

SixS (Surfaces Interfaces X-ray Scattering) is a wide-energy range beamline dedicated to the structural characterisation of surfaces and interfaces (solid-solid or solid-liquid), as well as nano-objects in controlled environments by the means of surface-sensitive x-ray scattering techniques.

\begin{figure}[!htb]
    \centering
    \includegraphics[width=\textwidth]{/home/david/Documents/PhD/Figures/sixs/SixSBeamline.png}
    \caption{
		Schematic view of SixS, the monochromator is situated in the optical hutch, the focusing mirrors, the attenuators and the MED end-station are in the first experimental hutch, additional mirrors and the UHV end-station are in the second experimental hutch.
		Only the MED end-station was used in the frame of this thesis.
    }
    \label{fig:SixSBeamline}
\end{figure}

The beamline operates with a $U_{20}$ undulator situated before the front-end (shutter) which delivers a horizontally polarised beam characterised by its horizontal and vertical full-width at half maximum (FWHM), respectively $\sigma_H = \qty{913.7}{\um}$ and $\sigma_V = \qty{19.1}{\um}$, and by its horizontal and vertical divergences, respectively $\Delta\theta_H = \qty{14.5}{\micro\radian}$ and $\Delta\theta_V = \qty{4.61}{\micro\radian}$.

An absorber and diaphragm are placed after the undulator to cut any parasitic signal from the source, a double crystal monochromator (DCM) of silicium (111) allows the selection of a monochromatic beam of thin bandwidth $\Delta \lambda$ in the range of \qtyrange{4}{20}{\keV} after receiving the white beam.
The crystals are cooled by liquid nitrogen due to the high intensity of the white beam and produce harmonics depending on the incident angle of the beam through diffraction (Bragg's law, eq. \ref{eq:Bragglaw}), only the first harmonic is selected with the second crystal.
Primary slits are placed just before the DCM, to select a homogeneous portion of the beam to work with to avoid fluctuations in the beam intensity.
Secondary slits after the DCM are used to clean the signal after the monochromator, the vertical focus of the beam is then handled by a pair of mirrors, one of which can be slightly curved to focus anywhere between the MED and the UHV stations (fig. \ref{fig:SixSBeamline}).
Piezo-actuated and pneumatic attenuators are respectively situated before and after the mirrors, they limit the intensity of the incident beam (\textit{e.g.} when working in the direct geometry).
The first experimental end-station on the beam path is the multi-environment diffractometer, where the focused beam is approximately large by \qty{60}{\um} vertically and \qty{1500}{\um} horizontally.
Additional focusing optics used for BCDI are present in the MED end-station and are detailed below.
The distance between each element on the beamline is recapitulated in tab. \ref{tab:DistanceSixS}, a sketch of the beamline is illustrated in fig. \ref{fig:SixSBeamline}.
The second experimental hutch, mainly used for ultra-high vacuum experiments, was not used during this thesis.

\begin{table}[!htb]
	\centering
    \small{
	\begin{tabular}{@{}llllllllll@{}}
	\toprule
	$U_{20}$ & Diaphragm & Ring    & Primary     & DCM         & Secondary       & Foc.      & Plan.     & Focusing &  MED \\
	         &           & wall    & slits       & Si (111)    & slits           & mirror    & mirror    & optics        &  \\
 	\midrule
 	\qty{0}{\m}&\qty{11.7}{\m}&\qty{15}{\m}&\qty{16.7}{\m}&\qty{18.8}{\m}&\qty{\approx 19}{\m}&\qty{24.5}{\m}&\qty{25.5}{\m}&\qty{\approx 30.5}{\m}&\qty{31}{\m}
	\end{tabular}
    }
	\caption{
		Distance between each main element that allows the focusing of a monochromatic beam on the multi environment diffractometer (MED) sample stage at SixS.
	}
    \label{tab:DistanceSixS}
\end{table}

\subsection{Multi environment diffractometer}\label{sec:MED}

The multi-environment diffractometer (MED) (fig. \ref{fig:MEDDiffractometer}) at SixS can be used in either a vertical or horizontal configuration (fig. \ref{fig:Diffractometer}), and can accommodate a large variety of experimental chambers around the sample stage.

\begin{figure}[!htb]
    \centering
    \includegraphics[trim=100 225 100 275, clip, width=0.9\textwidth]{/home/david/Documents/PhD/Figures/sixs/Diffractometer.pdf}
    \caption{
        Multi-environment diffractometer (MED), here in the horizontal configuration, accommodating the XCAT reactor \parencite{VanRijn2010}, placed upon a hexapod which allows the sample to move in three orthonormal directions.\\
        Image copyright by Leiden Probe Microscopy (\url{https://leidenprobemicroscopy.com}).
    }
    \label{fig:MEDDiffractometer}
\end{figure}

Situated at a distance $R = \qty{31}{\m}$ from the source, the transverse coherence length of the beam can be computed at the sample stage when working with an energy of \qty{8.5}{\keV} (working energy of the coherence optics) following eq. \ref{eq:TransverseCoL}.
For a Gaussian distribution, the FWHM is related to the standard deviation $\sigma$ by the relation $FWHM = 2\sqrt{2 ln (2) } \sigma$.
The transverse coherent lengths are equal to $L_{T,H} = \qty{3.3}{\um}$, and $L_{T,V} = \qty{157.3}{\um}$.

The longitudinal coherence length can be computed from eq. \ref{eq:LongitudinalCoL}, the ratio $\Delta\lambda/\lambda$ approximately equal to \qty{e-4}{\keV} at the energy of \qty{8.5}{\keV} gives $L_L = \qty{729}{\nm}$.
These values corresponds to perfect coherent lengths without taking into account any imperfection in the beamline elements (such as the monochromator) that can result in increasing the virtual source size and decreasing the transverse coherence lengths \parencite{Jacques2010}.

\subsubsection{BCDI configuration}

During this thesis, the horizontal configuration was used for SXRD experiments whereas the vertical configuration was used for BCDI experiments.
Working in a vertical configuration is preferable for BCDI since it allows a larger range of incident angle ($\beta$, the incident angle in the horizontal configuration, is limited to \ang{3}).
The mirrors, that focus the beam in the horizontal configuration, are removed from the optical path, the focus of the beam is then handled by specific focusing optics.

The secondary slits, which are wide open during surface x-ray diffraction experiment, are closed to \qty{100}{\um} vertically and horizontally, acting as a secondary source during Bragg coherent diffraction imaging experiments, reducing the effective size of the source seen by the focusing optics \parencite{Jacques2010}.
So called \textit{coherence} slits are placed just before the beam stop to select a homogeneous portion of the beam that impinges on focusing refractive lenses called Fresnel zone plates (FZP), designed to work at an energy of \qty{8.5}{\keV}.
The use of Fresnel zone plates, that preserve the coherent wavefront, together with coherence slits has been proven to increase the coherent flux impinging on the sample \parencite{Schroer2008, Diaz2009, Mastropietro2011}.

An order-sorting aperture (OSA) and a beamstop block the direct beam and the higher diffraction orders from the FZP (fig. \ref{fig:OpticalSetup}).
Since the coherence slits also have the counter effect of decreasing the total photon flux on the sample, their width is subject to a compromise to increase the coherent flux on the sample.

The final beam size on the sample is of about \qty{0.8}{\um} horizontally and \qty{1.2}{\um} horizontally.

\begin{figure}[!htb]
    \centering
    \includegraphics[width=0.8\textwidth]{/home/david/Documents/PhD/Figures/sixs/OpticalSetup.png}
    \caption{
    	Focusing optics used for Bragg coherent diffraction imaging at the MED end-station, increasing the coherent flux on the sample.
        $f$ stands for focal distance and $\lambda$ for the wavelength of the incoming beam.
    }
    \label{fig:OpticalSetup}
\end{figure}

\begin{table}[!htb]
    \centering
	\begin{tabular}{l|l|l|l|l|l}
	    & Coherence slits & Beam stop & FZP & OSA & Beam \\
        \toprule
	    Diameter & \qtyproduct{60 x 20}{\um} & \qty{80}{\um} & \qty{300}{\um} & \qty{70}{\um} & \qty{\approx 1}{\um}\\
        Sample distance & \qty{-400}{\mm} & \qty{-220}{\mm} & \qty{-200}{\mm} & \qty{-50}{\mm} & \qty{0}{\mm} \\
	\end{tabular}
	\caption{
	Diameter of the coherence focusing optics used to increase the coherent flux on the sample.
	}
    \label{tab:OpticsBCDI}
\end{table}

\subsubsection{Improving Bragg coherent diffraction imaging at SixS}

One of the bottlenecks of the BCDI technique is its slow data reduction and analysis process.
On $3^{rd}$ generation synchrotrons that offer a lower coherent flux (eq. \ref{eq:CoherentFlux}) than $4^{th}$ generation synchrotrons, the measurement time can also exceed several dozens of minutes.
For example, at SixS, a rocking curve lasts between \qtyrange{20}{90}{\min} depending on the particle size, the quality of the alignment, the strain of the particle, \textit{etc.}
Once the raw data is obtained, the particle must be \textit{reconstructed}, so that the displacement and strain arrays can be retrieved.
The analysis workflow can take up to an hour, which totals to a maximum of two hours from the start of the measurement to the when the user has a good idea of the sample shape and structure.

SixS is a beamline that does not only carry out BCDI experiment, but also SXRD experiments, in the same experimental end-station, the multi-environment diffractometer (MED, sec. \ref{sec:MED}).
When aiming at performing \textit{operando} catalysis experiments, switching from one setup to another can take up to a few days.
This leaves only a limited remaining amount of time to align the sample, find a suitable nanoparticle, and carry out the experimental plan.

Li \textit{et al.} \parencite*{Li2020}, who have first shown that the SixS beamline could be used to carry out BCDI experiment, started to work on improving the BCDI measurement process.
By comparing continuous and step-by-step measurements, they have shown that continuous scanning would result in the same data quality while decreasing the measurement dead-time by \qty{30}{\percent}, thereby paving the way for quicker BCDI measurements.

During this thesis, the measurement process was further improved by taking advantage of the new possibility to perform continuous \textit{on-the-fly} scans at SixS, while moving the sample plane with the hexapod.
When using the coherence setup (fig. \ref{fig:OpticalSetup}) with the diffractometer in the vertical geometry (fig. \ref{fig:Diffractometer}), the beam focused on the sample is about \qty{1}{\um} large vertically, the horizontal footprint depending on the incident angle between the beam and the sample.

The incoming angle $\mu$ is set to a Bragg angle ($\theta$ in eq. \ref{eq:Bragglaw}) and the scattered x-rays are collected by setting the out-of-plane detector angle $\gamma$ at a position equal to $2\theta$ (similar to fig. \ref{fig:EwaldSphereSpecular}).
Finally, by simultaneously moving the sample with the hexapod and recording the Bragg scattered intensity with the detector, it is possible to map the sample surface with a sub-micron resolution (fig. \ref{fig:SampleMapping}).

An example of the mapping result is illustrated below in fig. \ref{fig:SampleMapping} with Pt nanoparticles for which the [111] crystallographic direction is parallel to the normal of the sample holder, the (111) Bragg peak was thus measured during the mapping process.

A nanoparticle with a width equal to \qty{300}{\nm} was identified with this technique (sec. \ref{sec:BCDINanoparticles}), which is a good estimate of the spatial resolution that can be attained, limited by the hexapod resolution (smallest step size equal to \qty{\approx 300}{\nm}) and the beam size.

\begin{figure}[!htb]
    \centering
    \includegraphics[height=5cm]{/home/david/Documents/PhD/Figures/sample/microscope_image.png}
    \includegraphics[height=5cm]{/home/david/Documents/PhD/Figures/sample/microscope_image_photon.png}
    \caption{
        Microscope image of the sample seen through the sapphire window of the PEEK dome (left).
        Map of the sample performed in Bragg condition (right), the high intensity (red) areas correspond to platinum nanoparticles.
        The letters, numbers and isolated nanoparticles in the centre of squares can be recognised on the sample.
    }
    \label{fig:SampleMapping}
\end{figure}

\textit{Gwaihir} (sec. \ref{sec:Gwaihir}) was developed primarily for the SixS beamline to counter the long analysis process, which allowed a significant reduction in the analysis time from around an hour to a few minutes.
The following beamtimes profit from the new software by having a more \textit{solution}-driven experimental process.
Indeed, to be measured, a nanoparticle must be isolated, not too small (weak scattered intensity, e.g. $>\qty{100}{\nm}$ at SixS), not too big (loss of coherence, fringes not visible, e.g. $<\qty{1000}{\nm}$ at SixS), and not too initially strained (difficult to obtain a good guess of the support).
These conditions are sometimes difficult to assert by simply looking at the diffraction pattern.
Therefore, quick inversion using \textit{Gwaihir} allowed a faster decision process regarding the continuation or not of the nanoparticles measurement.

Successful first measurements by \cite{Lim2021} have permitted the simultaneous use of BCDI measurements from SixS with measurements from other imaging beamlines (ID01 - ESRF, P10 - DESY), designing a robust method to identify defects in the real space with convolutional neural networks (CNN).

\subsubsection{Catalysis reactor}\label{sec:XCAT}

A catalysis reactor is mounted on the sample stage of the MED for the study of catalytic reactions.
A 3D view in the horizontal configuration with the high pressure surface diffraction reactor XCAT (X-ray CATalysis) set on the goniometer is illustrated in fig. \ref{fig:MEDDiffractometer}.

This reactor is used for the study of heterogeneous catalysis \parencite{VanRijn2010, Resta2020a}, and couples a UHV environment for classical surface science preparations (sputtering, annealing, evaporation) with an ambient pressure reactor.
In the frame of this thesis, the metallic single crystal are purchased already cut and polished for SXRD experiments, introduced in the reactor and cleaned by sputtering \parencite{Taglauer1990} before the experiment.
It is also possible to heat the sample up to \qty{900}{\degreeCelsius} for an extended period of time to anneal the sample after sputtering \parencite{Musket1982}.

\begin{figure}[!htb]
    \centering
    \includegraphics[trim=10 325 350 64, clip, width=0.8\textwidth]{/home/david/Documents/PhD/Figures/sixs/SampleHolder.pdf}
    \caption{
        Sample holder used for the XCAT catalysis reactor.
        \textcolor{Important}{TRY TO FIX WORDS, MAYBE TICKZ ?, add ceramic}
    }
    \label{fig:SampleHolder}
\end{figure}

The sample holder (fig. \ref{fig:SampleHolder}) consists of a cross-shaped graphite heater covered in boron nitride held by two conductive screws in a thick (\qty{\approx 1}{\cm}) and round ceramic slab.
The protruding ends of the screws on the opposite side of the ceramic slab serve as contact pins used for the heater power circuit, and to securely fasten the sample holder to the reactor.
A Beryllium or PEEK dome closes the reactor, the V-seal sealing the reactor atmosphere from the outside, two types of PEEK dome are available, with or without a sapphire window, allowing to observe the sample during the experiments with a microscope.

This design not only enables swift attachment and detachment of the sample holder but also ensures a robust setup in both horizontal and vertical orientations.
The sample is then placed on top of the sample holder and set in place by a wire.
All the materials used in the reactor were carefully chosen as not catalytically active for specific reactions.
For example, tantalum (Ta) parts (screws and wire) are used in oxygen rich environment since Ta is more resilient than molybdenum (Mo) to oxidation reaction, whereas Mo parts are used in hydrogen rich environment because it is more resilient to hydrides formation.
A boron nitride paint additionally covers and protects the screws and wires used to fix the sample holder and the sample from oxidation effects.
The heater power is adjusted as a function of the sample temperature read by a type C thermocouple.

\begin{figure}[!htb]
    \centering
    \includegraphics[trim=0 0 0 42, clip, width=0.49\textwidth]{/home/david/Documents/PhD/Figures/sixs/IMG_20210612_153249_240.jpg}
    \includegraphics[trim=150 0 200 0, clip, width=0.49\textwidth]{/home/david/Documents/PhD/Figures/sixs/CELL.jpg}
    \caption{
        MED end-station in vertical configuration for BCDI experiments (left), the OSA can be seen next to the reactor cell, opposite to the anti-scattering tube, the sample is visible in the absence of the reactor dome.
        Smaller version of the XCAT reactor cell with PEEK dome (right).
    }
    \label{fig:MEDV}
\end{figure}

The gas composition inside the reactor is tuned by setting the flow for each gas in standard cubic centimetres per minute (SCCM).
The inlet and outlet of the gases in the reactor volume are through holes in the ceramic slab that are set just below the sample holder.
The standard supported gases are argon (\argon), oxygen (\dioxygen), nitrogen oxide (\nitricoxide), carbon monoxide (\ce{CO}) and hydrogen (\ce{H_2}).
It is possible to change the gases used by computing the correct conversion factor for a mass flow controller, each one originally set to work with a specific gas.
For example, the bottle of \ce{H_2} was replaced by a bottle of ammonia (\ammonia) during the experiments involving the catalytic oxidation of ammonia.
Moreover, the pressure inside the volume of the reactor can be controlled by a pressure controller, from \qtyrange{100}{1200}{\milli\bar} (fig. \ref{fig:GasSupplySystem}), in a small volume (\qty{\approx 10}{\ml}), under batch or constant flow conditions.

Finally, the evolution of the product and reactant pressure during the catalytic reaction is probed by a residual gas analyser (RGA) which is connected to a leak from the reactor outlet (fig. \ref{fig:GasSupplySystem}).
The leak pressure is usually in the range of \qty{e-6}{\bar} and the assumption is made that the partial pressure of each mass detected by the RGA after the leak is equal to the partial pressure in the reactor multiplied by the same constant for all masses.

Overall, this experimental setup allows the exploration of a large, multi-dimensional parameter space during heterogeneous catalysis experiments connecting reaction kinetics and surface structure.

\begin{figure}[!htb]
    \centering
    \includegraphics[trim=75 50 100 50, clip, width=\textwidth]{/home/david/Documents/PhD/Figures/sixs/GasSupplySystem.pdf}
    \caption{
    	Gas supply system used at SixS together with the multi-environment diffractometer (MED) for the high pressure surface diffraction reactor.
    	The mixing of reacting gases is performed before the reactor (MIX), argon is used as a carrier gas (MRS).
    	The leak to the RGA is represented as a side circuit from the reactor outlet.
    	Copyright by Leiden Probe Microscopy (\url{https://leidenprobemicroscopy.com}).
    }
    \label{fig:GasSupplySystem}
\end{figure}

When performing BCDI measurements, a smaller version of the XCAT is used which is light enough to be supported vertically by the goniometer (fig \ref{fig:MEDV}).
The two reactors have the same inner volume and host the same sample holder stage (fig \ref{fig:SampleHolder}).
The main setbacks being the lack of sputtering gun and the impossibility to study catalytic reaction in batch mode (as opposite to flow mode) due to the design of the small reactor (the leak to the RGA is after the pressure controller in the small reactor).

The partial pressure of each gas detected after the leak in the main circuit (fig. \ref{fig:GasSupplySystem}) by the residual gas analyser can be converted back to partial pressures in the reactor by assuming (i) ideal gas behaviour (ii) homogeneous temperatures in the reactor and in the RGA chamber (iii) no selectivity of the gases through the leak, \textit{i.e.} the ratio between the partial pressures after the leak is the same that in the reactor (iv) all gases are detected by the RGA (v) the ratio between the incoming gases flow is the same as the ratio between the partial pressures in the reactor in the abscence of reaction.
The partial pressure of argon, which does not partake in the reaction, is then used to compute the partial pressure of each gas in the XCAT chamber.

For example, at a reactor pressure of \qty{0.5}{\bar}, a total gas flow of \qty{50}{\ml\per\min}, a partial gas flow of \ce{NH_3} equal to \qty{1}{\ml\per\min}, and a partial gas flow of argon equal to \qty{49}{\ml\per\min}, we estimate the partial pressure of argon in the reactor to be equal to \qty{490}{\milli\bar}.
The goal of the use of the mass spectrometer in this thesis is not to perform a quantitative analysis of the catalyst activity but rather to compare the evolution of the product partial pressures for different atmospheres, which is compatible with the setup.

The measurement process for the collection of the scattered intensity is detailed in sec. \ref{sec:DataCollectionSXRD} for surface x-ray diffraction and in sec. \ref{sec:DataCollectionBCDI} for Bragg coherent diffraction imaging.
Since this reactor was designed for SXRD together with a 6-circle diffractometer, the entire surface of the sample is accessible during the experiment which allows the full access to the reciprocal space to probe any particular symmetry, limited only by the extent of the Ewald sphere (\textit{i.e.} a maximum extent in $\vec{q}$).

\subsubsection{Detectors}

Different detectors can be accommodated on the diffractometer arm.
On one hand, the MERLIN \parencite{Bewley2006} combines a small pixel size (\qty{55}{\um}) with a large array (\numproduct{512 x 512} pixels), few gaps, and a large dynamic counting range, similar to the MAXIPIX detector \parencite{Ponchut2011} which was first used for coherence experiments on the beamline \parencite{Schavkan2013, Li2020}.
For example, the pixel size is of main importance when computing the oversampling requirements (sec. \ref{sec:PhaseRetrieval}).

On the other hand, the XPAD detector \parencite{Basolo2005, Dawiec2016} is used for surface x-ray diffraction experiments, its pixel size is equal to \qty{130}{\um}.
Different detectors can be used that change the array size, respectively (\numproduct{540 x 130}) for the XPAD70 and (\numproduct{540 x 260}) for the XPAD140.
The area covered by the detector is larger than with the MERLIN, but the main advantages of the XPAD are its high dynamic range and highly programmable interface \parencite{Fertey2013} which can be programmed to work with automatic attenuators \parencite{Dawiec2016}, crucial in SXRD when navigating between bulk and surface scattered signals .
The XPAD detector is usually set with its largest side along $\vec{c}$ to collect the highest range in $l$ during the measurements.
    \newpage
\section{Computer programs}

With the upgrade of synchrotrons to more brilliant sources, imaging beamlines have received increasing attention from the scientific community and yield larger and more complex data that far exceed the basic diffraction pattern (a reconstructed BCDI measurement is a complex array in 3 dimensions) (cite new beamlines that do imaging techniques like SixS, Crystal, new beamline at esrf etc).
In a more general frame, the increase in flux at synchrotrons leads to quicker experiments, which in turn leads to an increased data stream stored on the beamlines, often in different formats.

Synchrotrons have effectively stepped in the wold of "big data" \parencite{Alizada2017, Wang2018}, which refers to a vast volume of structured, semi-structured, and unstructured data that is generated at a high velocity and comes from various sources.
This data is typically too complex or large to be processed and analyzed using traditional data management tools or methods that were originally written to simulate or fit only a few datasets at a time.
In the specific case of Bragg coherent diffration imaging, new tools have been developed during the thesis (sec. \ref{sec:Gwaihir}, \cite{jerome_carnis_2021_5741935, Simonne2022}) to help create a fast and reproducible workflow, but it is a more general way of working that must be adopted in synchrotrons to adapt and profit from the transition to big data \parencite{Wang2018}.
This does not only concern scientific analysis but also the complexity of synchrotron beamlines that could benefit from the use of big data to upgrade their performance (cite projet maxime).

First, synchrotron data must be stored in a comprehensive data format \parencite{Konnecke2015} that contains not only detector images but all of the important metadata linked to the experiment, providing a more comprehensive and detailed view.
The fine selection of the beamline attributes that can be of interest for the discovery of hidden patterns, correlations, and trends that may not be apparent with smaller datasets is left to the beamline scientists.

Secondly, synchrotrons must offer high performance computing (HPC) clusters \parencite{Wang2021} compatible with the most common integrated development environment (IDE) and either have an internal strategy or trust their users to be able to develop tools that can use their power to uncover valuable insights \parencite{WangIEEE2016, Khaleghi2019}.
These insights can lead to new discoveries and enhance our understanding of complex phenomena, driving scientific progress.

% The use of big data expands research capabilities by allowing studies to be conducted on a much larger scale whereas traditional scientific experiments often rely on a limited number of observations or samples, which may not capture the full complexity of real-world phenomena.

\subsection{Data reproducibility}

On one hand, linking reproducibility and big data is crucial for ensuring the credibility and reliability of the insights derived from large and complex datasets.
Reproducibility in the context of big data refers to the ability to replicate or re-run data analysis processes to obtain consistent results.
Achieving reproducibility with big data can be challenging due to the massive volume, velocity, and variety of the data involved, in the case of Bragg coherent diffraction imaging, if measurements took a few dozens of minutes a few years ago, we are now capable of measuring a single rocking curve in less than a minut at specialized beamlines such as ID01 at the upgraded european synchrotron \parencite{Richter2019}.

On the other hand, discussions about defining a set of rules that regulate research practice \parencite{Kretser2019} and reduce the grey zone that includes scientific misconduct at all levels of academia \parencite{Kornfeld2016} are growing, raising awareness on data reproducibility in the scientific community.

Staggering numbers \parencite{Baker2016}, show that about 65\% of scientists in the field of physics and engineering struggle to reproduce others' results, and more than 50 \% fail to reproduce their own results.
These numbers can sometimes be linked to very precise environments and techniques, with experimental conditions and processes difficult to reproduce between different laboratories, and also to knowledge transfer from academia to industry \parencite{DanielSarwitz2015}.
However, according to the study  \parencite{Baker2016}, code availability, insufficient peer reviewing, and access to raw data contribute to non-reproducible research.
Moreover, it is largely accepted that the publishing industry has its own role to play by facilitating peer-reviewing to promote reproducibility \parencite{Lee2017}.

Therefore, code-availability, access to raw data combined with metadata, and reproducible workflows are goals of utmost importance for experimental science \parencite{Munafo2017}.
Reproducible data will result in a global improvement of confidence in new techniques, such as BCDI, which could subsequently result in growth of interest and community.

\subsection{The Jupyter environment}

The Jupyter Notebook environment \parencite{IPython, Kluyver2016jupyter} was chosen for the development of data analysis tools during this thesis for its versatile, user-friendly and browser-based interface.
To simplify the data analysis pipelines in fourth-generation synchrotrons, it is of critical importance to offer the possibility for external users to analyse the collected data remotely with access to computational environments
Remote access to Jupyter Notebook or JupyterLab is provided by JupyterHub, while the sharing of ressources is handled by SLURM.

Jupyter Notebook has proven to be an effective tool for the analysis of synchrotron data, in terms of graphical user interface \parencite{Martini2019a,Simonne2020}, but also in terms of supporting scientific communities looking for high-performance frameworks \parencite{jupyter_computing_4, jupyter_computing_1, jupyter_computing_3, jupyter_computing_2, 9307800}.
Moreover, computational environments and resources can be accessed from any computer \textit{via} JupyterHub, a specific interface for computing clusters.
JupyterHub offers the possibility to proceed to heavy computations without relying on specific hardware, \textit{e.g.} graphical processing units, mandatory for accelerated phase retrieval with \textit{PyNX}.

The Notebook can be used to take notes on the experiment or to create custom Python functions.
They can be shared in the \textit{.ipynb} format or as \textit{.pdf} documents.
%\textit{Gwaihir} transforms each data-set into a structured Python object that can be accessed and manipulated in the Notebook for e.g. personalized figures, parameter tests, \textit{etc.}

Moreover, researchers can create their own work-spaces from shared resources, with direct access to tailored computational environments, without having to install multiple software, keeping in mind that the use of GPUs for scripts optimisation is far from accessible.
Thus, system administrators can efficiently manage complex environments accessible to all users.
Finally, remotely accessing the data avoids data storage issues, which can quickly become problematic with current experiments.

Interactive data analysis/visualisation relies on the \textit{ipywidgets} library.
Widgets are represented in the back-end by a single object linked to a single parameter.
The front-end relies on \textit{JavaScript} code; each time a widget is displayed, a new representation of that same object is created.
The widgets' style, orientation, and layout attributes can be edited to customize the final window, \textit{e.g.}, the layout attribute exposes a number of properties that impact how widgets are laid out, such as height and width.

\subsection{Thorondor} \label{sec:Thorondor}

Due to the lack of available free data reduction and analysis software at the B07 beamline in the Diamond synchrotron, a modified version \parencite{Simonne_Thorondor_2022_Diamond} of Thorondor \parencite{Simonne2020}, a program originally written for the analysis of NEXAFS data, was used to process the XPS data.
To facilitate data reproducibility and open-source analysis workflows, an example of the data recorded during this thesis is available directly in the GitHub repository \cite{Simonne_Thorondor_2022_Diamond}.

Since the XPS measurements were performed at different total pressures, the raw data had first to be reduced in order to be be able to analyse the different spectra all together.
The classical workflow for the analysis of XPS data is to first align the recorded spectra on the Fermi edge that corresponds to the kinetic energy of the first electron that escapes the sample.
By doing so, one can be confident that any shift in the peak positions is due to chemical changes such as the oxidation state and not to charging effects of the sample.

Secondly, to be able to quantify and compare the evolution of the peak intensity, one must normalize the intensity of the detected electron beam since the electron mean free path depends on the pressure in the reaction chamber.
The range of kinetic energy just before the absorption edge of Pt 4f was chosen since it had the best signal to noise ratio and does not depend on any experimental parameter besides the pressure.

Finally, for the peaks that showed a good signal to noise ratio, the fitting of the peak shape was realised thanks to the \textit{lmfit} \parencite{Newville2016} package by the means of the Doniach-equation which is the best approximation of the asymmetric peak shape that results partly from the convolution of the analyser function and the shape linked to the photoelectron process in metals \parencite{Doniach_1970}.

\subsection{Gwaihir} \label{sec:Gwaihir}

BCDI relies on iterative algorithms to solve the phase lost during the measurement (\cite{robinson_coherent_2009} - sec. \ref{sec:BCDI}).
A 3D intensity distribution in the vicinity of a Bragg peak (stack of diffraction patterns forming a 3D reciprocal space map with the proper sampling) is collected from a sample illuminated with coherent light \parencite{robinson_coherent_2005}, and serves as input for phase retrieval.

\begin{figure}[!htb]
    \includegraphics[width=\textwidth]{/home/david/Documents/PhD/Figures/gwaihir/Packages.png}
    \caption{Flow chart illustrating the main steps in the BCDI data analysis workflow. \textit{Gwaihir} links the \textit{bcdi} and \textit{PyNX} packages \textit{via} its graphical user interface and command line scripts, resulting in a complete and easily understandable data analysis workflow. Opt. stands for optional.
    }
    \label{fig:Packages}
\end{figure}

There are three main steps to complete to image the strain (fig. \ref{fig:Packages}); it must first be pre-processed (removal of parasitic scattering intensities, \textit{etc.} \parencite{ozturk_performance_2017}).
It must then be inverted \textit{via} phase retrieval \parencite{miao_possible_2000}, the phase containing important information lost during the measurement.
Finally, it is possible to extract meaningful results from the three dimensional complex images after post-processing (removal of phase offset, interpolation in a common orthonormal frame, \textit{etc}).
Several software packages were developed to solve these steps but none bring a comprehensive pipeline from start to finish.

For example, \textit{PyNX} \parencite{pynx2020operators} focuses on the phase retrieval step, \textit{bcdi} \parencite{jerome_carnis_2021_5741935} focuses on data pre-processing and post-processing, and \textit{Cohere} \parencite{cohere_2021} focuses on pre-processing and phase retrieval.
A few graphical user interfaces (GUIs) also exist, such as \textit{Cohere} \parencite{cohere_2021}, \textit{Phasor} \parencite{dzhigaev_dzhigaevdphasor_2021}, and \textit{Bonsu} \parencite{newton_bonsu_2012}.
Providing a workflow will reduce the time spent on data analysis for newcomers, and improve results reproducibility by facilitating sharing while keeping track of analysis parameters and metadata.

Large scale facilities and institutions seek ways to provide remote-access high powered computing services to their users, which combine existing solutions in an interactive and user-friendly environment.
Jupyter (https://jupyter.org/) is particularly advantageous and was chosen by several institutions for such purposes, \textit{e.g}. Google (Google Colab), the EGI federation, or the European Synchrotron (Simple Linux Utility for Resource Management - SLURM).

\textit{Gwaihir} is a tool which brings together the functionality of the \textit{PyNX} package for phase retrieval and bcdi package for the pre- and post- processing in a graphical user interface (GUI) built for the Jupyter framework \parencite{Kluyver2016jupyter}.
It provides an interface to the bleeding edge of data analysis in BCDI, it can be used locally or remotely and offers an interactive and user-friendly interface with complex functionality satisfying both beginners and experts.

\subsubsection{Software structure}

A part of the BCDI community relies on Python, an accessible language that has gradually become one of the most popular, versatile \parencite{IPython, Newville2016} and widely-taught \parencite{Scopatz2015, McKinney2017, Boulle2019} programming languages in science.
\textit{Gwaihir} followed this initiative by regrouping data visualisation tools, workflow guidelines, and a user-friendly interface around two Python packages (\textit{PyNX} and \textit{bcdi}), that together, offer a complete data analysis suite (fig. \ref{fig:Packages}).

\textit{Gwaihir} works with Python 3.9 and is licensed under the GNU General Public License v3.0.
The source code as well as the latest developments are available on GitHub while each stable version will be released on the Python Package Index (PyPi), along with its documentation.

\textit{bcdi} \parencite{jerome_carnis_2021_5741935} tackles the pre-processing and post-processing of the BCDI data (fig. \ref{fig:Packages}).
Aside from the specific details of the experimental setup (diffractometer and setup geometry, detector type, file system), the majority of the pre-processing pipeline is actually beamline-independent.
Based on this observation, \textit{bcdi} leverages inheritance and transforms the raw data to a common data format used for phase-retrieval.
It frees the user from having to learn or remember the technical details for each beamline and sets a common strategy across beamlines.

\textit{PyNX} \parencite{pynx2011} is a toolkit with assorted Python modules and command-line scripts which can be used for the analysis of coherent X-ray imaging data, including phase retrieval (fig. \ref{fig:Packages}).
All calculations can be executed and distributed on multiple graphical processing units (GPUs) for accelerated computing using MPI.
Elementary algorithms can be easily tailored, built upon, and combined using an operator-based approach, allowing full flexibility with high-performance computing \parencite{pynx2020operators}.
Some facilities have started to make computational resources available to their users; those machines are an ideal environment for \textit{PyNX} which is nowadays available through SLURM, at the ESRF - The European Synchrotron, or GRADES at the Optimized Light Source of Intermediate Energy of LURE (SOLEIL).

\subsubsection{Workflow for Bragg coherent diffraction imaging} \label{sec:Workflow}

\begin{figure}[!htb]
    \centering
    \includegraphics[width=\textwidth]{/home/david/Documents/PhD/Figures/gwaihir/Workflow.png}
    \caption{Workflow steps taken in \textit{Gwaihir}; the circular workflow illustrates data reproducibility, a key concept, facilitated by using the \textit{cxi} architecture.}
    \label{fig:Workflow}
\end{figure}

Data pre-processing aims at improving the quality of the phase retrieval by optimising the size and content of the 3D array used as input.
The minimal processing consists in loading the raw data and stacking it together as an input for the phase retrieval.
Intermediate optional steps can be added via a YAML (YAML Ain't Markup Language) configuration file.
A more complete pre-processing is: loading the data (with flatfield correction and masking of damaged detector pixels and detector gaps), normalizing each frame by a monitor (sensor at the beamline measuring the incoming X-ray flux), binning pixels, centering the Bragg peak and cropping the dataset to a meaningful size, masking the parasitic scattering intensities coming from neighboring crystals (via the interactive user interface), interpolating missing intensities in damaged pixels, interpolating the data stack from the detector frame to an orthonormal frame (also called as geometric transformation).
The geometric transformation can be realized using either the transformation matrix \parencite{thesismark} or the existing  \textit{xrayutilities} package \parencite{kriegner_xrayutilities_2013}, depending on which reference basis is needed.

Post-processing regroups methods applied to the complex output of the phase retrieval. If the data is still in the detector frame (geometric transformation not applied during pre-processing), the data is interpolated in an orthonormal frame using the transformation matrix.
After phase unwrapping, a refraction correction is optionally applied, and the phase ramp and phase offset are removed.
At this point the displacement and the strain component are calculated from the phase.

\textit{Gwaihir} offers an interactive workflow separated in three main groups, data pre-processing, phase retrieval, and data post-processing (fig. \ref{fig:Packages}). It is meant to be reproducible, simple to share, resulting in the phase and amplitude of the probed object. To illustrate this, the results of the following procedure on a data-set collected at the P10 beamline in at PETRA III are shown in fig. \ref{fig:GUI_file} (CXI dataset ID 195).

The parameters used during the analysis are displayed in the GUI in an order following a typical workflow, to underline the evolution of the data processing.
These parameters are ultimately saved as attributes of the Python data object, to keep track of the actions applied to the raw data to obtain the final result.

\begin{figure}[!htb]
    \centering
    \includegraphics[width=0.49\textwidth]{/home/david/Documents/PhD/Figures/gwaihir/CropDiffPatternRed.png}
    %\includegraphics[width=\textwidth]{/home/david/Documents/PhD/Figures/gwaihir/CropObjectAmplutideRed2.png}
    \includegraphics[width=0.49\textwidth]{/home/david/Documents/PhD/Figures/gwaihir/crop_phase_red.png}
    \caption{
    %Each parameter used during the data analysis is stored in the same \textit{.cxi} file, along with the results, as seen in the nested architecture on the left of the Figure.
    Here a 2D slice of a 3D coherent diffraction pattern is shown (top), with the phase of the electron density (radiants) of the reconstructed object in which the facets are clearly visible (bottom). The data array are displayed via \textit{JupyterHub}, a browser-based data analysis platform available on the \textit{SLURM} computer cluster for the ESRF and the \textit{GRADES} computing cluster at SOLEIL.
    }
    \label{fig:GUI_file}
\end{figure}

\subsubsection{Pre-processing}
First, the correct beamline and instrumental parameters must be selected (\textit{e.g.} sample-detector distance, probing energy, detector pixel size, \textit{etc}...), usually constant throughout the experiment. For example, the orthogonalisation tab regroups the parameters needed to properly set up the transformation matrix between the sample frame and \textit{e.g.} the laboratory frame. By choosing to transform the final object in a common, orthogonal frame, it becomes easier to compare the object's evolution when probing different Bragg reflections \cite{crystal_sarah}.

Once the raw data was collected, different pre-processing parameters can be modified to optimize the collected 3D diffraction intensity. For example, the size and center of the array to analyze (either fixed manually as the center of the Bragg peak if known, or determined as the center of mass of the 3D array during the analysis). The array can also be cropped to reduce its size, removing the points furthest from the Bragg peak where the signal-to-noise ratio is low (Figure \ref{fig:GUI_file}).

It is important to create a detector mask prior to the experiment to correct the raw data for hypothetical hot-pixels and background.
Moreover, it is also possible to correct the raw data frame by frame for spurious data which taints the diffraction pattern.
For example, it is possible, while recording the 3D diffracted intensity of a given reflection, to have signal coming either from the substrate or from neighbouring objects that will be summed to the probed object's intensity.

Finally, it is possible to normalise the raw data by an intensity monitor, or compute the scattering vector $\vec{q}$ of the measurement, from the instrumental geometry and parameters.

\subsubsection{Phase retrieval}
The diffracted intensity collected by the detector is proportional to the squared modulus of the structure factor $F(\vec{q})$ (Equations \ref{eq:scat_int}-\ref{eq:int_detector}).
To retrieve the phase from the diffracted intensity, \textit{PyNX} uses iterative algorithms (sec \ref{sec:BCDI}).

The phase of the Bragg electronic density (Figure \ref{fig:GUI_file}) $\Phi_{hkl}$ is proportional to the scalar component of the displacement $\vec{u}(\vec{r})$ that is parallel to $ \vec{q_{hkl}}$. A phase value of $\pi$ is equivalent to a displacement from the equilibrium position of half the lattice spacing in the direction of $\vec{q_{hkl}}$.

New methods such as convolution neural networks (CNN) are under study and have started to show some encouraging results, but are not yet sufficiently robust, particularly in the case of strained particles \cite{cherukara_real-time_2018,chan_rapid_2021,Wu2021}.

The parameters necessary to phase retrieval can be modified in the GUI (Figure \ref{fig:PRT}.
For example, the support threshold
% (from 0 to 1, 1 corresponding to the maximum amplitude of the complex electronic density), under which a point is no longer considered as ``in" the object
, the number of iterations for each algorithm, the object initialisation procedure (square, sphere, auto-correlation), and so on.
Each parameter is detailed in the GUI and also discussed in the literature (see, for instance, Ref. \cite{fienup_phase_1982,fienup_reconstruction_1978,marchesini_publishers_2007,pynx2020operators}). Initial guesses are given for each parameter but must be refined by the user to optimize the results.

State-of-the-art GPUs available onsite where data collection occurs enables almost real time data reduction.
This results in quick visualisation of the probed object's amplitude and phase offering the possibility to synchrotron users to optimize their data acquisition during the experiment. Such live feedback is critical to the success of an experiment and provides broader opportunities in the type of experiments that can be carried on.


\textit{PyNX}'s recent update includes mathematical operators \cite{pynx2020operators} which represent most of the reconstruction operations traditionally used in phase retrieval \cite{gerchberg_practical_1972,fienup_reconstruction_1978,marchesini_publishers_2007} laid the foundation for quick and interactive phase retrieval in \textit{Gwaihir}.

It is possible to refine the input parameters directly in the GUI by visualizing their impact on the reconstructed object. This allows to refine the phase retrieval input parameters before submitting a batch job, that will spawn a sub-process on the computing cluster for phase retrieval. With a single click, several dozens of solutions can be computed in the matter of minutes on computing clusters.

With well-tuned parameters and high quality data-sets, phase retrieval converges towards the same solution, but with minor differences between each reconstructed object, related to the phase retrieval process.

\begin{figure}[!htb]
    \centering
    \includegraphics[width=\textwidth]{/home/david/Documents/PhD/Figures/gwaihir/phase_retrieval_tab.png}
    \caption{Phase retrieval tab in \textit{Gwaihir}. Parameters are separated into groups (files, support, point-spread function, algorithms, ...) and detailed in the README tab. The object can be reconstructed through a batch job, submitted to the computing cluster in the backend, or with operators, that will plot the evolution of the reconstruction object in the Notebook.}
    \label{fig:PRT}
\end{figure}

\textit{Gwaihir} provides a wide range of selection criteria to find the best solution.
Each reconstructed object has a list of final attributes that can be used as a criteria for selection, such as the Free log-likelihood \cite{pynx2020FLLK} or the standard deviation of the modulus of the reconstructed object.
%Another common criterion used to assess the quality of the phase retrieval is the phase retrieval transfer function (PRTF) that compares the experimental diffraction pattern with the Fourier transform of the reconstructed object
%magenta{The PRTF is above all a tool to estimate the spatial resolution of the reconstruction. I agree it is to some extent linked to the quality of the reconstruction but I think it is a bit misleading to present it as a tool to assert the quality of phase retrieval. In this paragraph I would insist more on the traditional error metrics ("Figures of merit") used to assess the quality of phase retrieval (chi-square, sharp metric, max volume... }
%\cite{chapman_high-resolution_2006,cherukara_anisotropic_2018}.
In the case of crystallographic defects, specific metrics (Chi, Sharp, Max volume, ...) were derived that perform best depending on the type of defect in the object \cite{Ulvestad2017}.

Following the selection criteria, it is possible to quickly identify the solutions of poor quality that must be ignored to create a set of best solutions.
\textit{PyNX} then offers a method that merges this set into a single solution by computing eigen-vectors for the selected solutions \cite{pynx2020FLLK}.
An alternative approach also available in \textit{Gwaihir} is to take the average of the best solutions (see for example Ref. \cite{ulvestad_nanoscale_2014}).


\subsubsection{Post-processing}

Once the solution with the best Figure of merit is selected, it is possible to use the \textit{bcdi} scripts to process the object's phase. The phase origin can first have an impact when comparing the lattice displacement between different reconstructions. In the case of weak strain, it can be sufficient to consider the center of mass of the object as the origin of phase.
However, this can become quite complex in the case of defects or defaults in the object. Special methods are defined to target this issue in \textit{bcdi}. For example, Guizar-Sicairos \textit{et al.} \cite{guizar-sicairos_phase_2011} and Hofmann \textit{et al.} \cite{hofmann_nanoscale_2020} proposed a convenient method for the numerical calculation of phase gradients in the presence of phase jumps.  More detail and guidelines can be found in \cite{Carnis2019_scientific_reports,jerome_carnis_2021_5741935} as well as in the README tab.
Different methods can be tested and their results directly compared via the visualization tab.

\begin{figure}[!htb]
   \centering
   \includegraphics[width=0.66\textwidth]{/home/david/Documents/PhD/Figures/gwaihir/norm_dist.png}
   \caption{Histogram representing the distribution of the amplitude of the electronic density.
   An empirical criterion to select the contour threshold is to take the base of the peak ($\mu - 2\sigma$), here around 0.75.
   Thus, we ensure that we are on the outer surface of the reconstructed object.
   Note that most points in the array have values close to zero, for they are not part of the reconstructed object, and are here ignored.
   }
   \label{fig:histo}
\end{figure}

%As seen in Figure \ref{fig:histo}, an empiric criterion to select the contour threshold is to take the value at the foot (left side) of the peak, here around 0.75, resulting in a clear distinction between surface and bulk.
% why ? linked to the convolution such as between the object and beam. (better here)

\subsubsection{Graphical interface}

\textit{Gwaihir} links in a unique user-friendly and interactive GUI the aforementioned packages whilst offering 2D/3D browser-based data visualisation tools.
The interface is built with the \textit{ipywidgets} library, compatible with both Jupyter Notebook and JupyterLab, web-based interactive computing platforms, browser-based data visualisation is a key concept when functioning remotely.

Further details behind these methods can be found in the packages' respective documentation.
Therefore, the main purpose of \textit{Gwaihir}, is to offer a possibility for the user to work solely in the Jupyter environment with a unique interactive GUI.
The graphical user interface is divided into 10 tabs (fig. \ref{fig:GUI_window}), separated into three groups: instrumental parameters, data processing and data analysis.
The aim of this organization is to achieve a comprehensible but fluid workflow, whilst still separating each step.
A final tab contains information about the different methods and parameters used in the workflow, as well as a tutorial on the GUI.

% \begin{figure}[!htb]
%     \centering
%     \includegraphics[width=\textwidth]{/home/david/Documents/PhD/Figures/gwaihir/GUI_tab1.png}
%     \caption{Screenshot of \textit{Gwaihir} graphical user interface displayed in Jupyter Notebook. Jupyter Notebook can be used on computing clusters via JupyterHub or on any machine for local use.}
%     \label{fig:GUI_window}
% \end{figure}

During the manipulation of the GUI, it is possible to quickly produce a final result thanks to those intuitive and interactive widgets, whose values are passed as arguments to the data analysis functions (fig \ref{fig:3D_object}).
For example, to designate the type of detector used, a value must be selected in a dropdown list that comprises the following options: Eiger2M, Maxipix, Eiger4M, Merlin, Timepix.
These correspond to the different detectors currently supported in the \textit{BCDI} package.
The final interface is created by setting the different widgets on a grid, with the output of different functions printed below.

% \begin{figure}[!htb]
%     \centering
%     \includegraphics[width=\textwidth]{/home/david/Documents/PhD/Figures/gwaihir/3D_obj.png}
%     %\includegraphics[width=\textwidth]{/home/david/Documents/PhD/Figures/gwaihir/3D_peak.png}
%     \caption{Widgets can be used to select folders and files, but also to tune the values of input parameters in functions. The selection of the data array, the colormap, and the contour of the resulting 3D object are performed through different widgets.}
%     \label{fig:3D_object}
% \end{figure}

Jupyter natively offers multiple options for interactive data plotting.
Most of the figures displayed in the GUI are based on the \textit{matplotlib} package \parencite{hunter2007matplotlib}.
For example, it is possible to select a list of 3D complex data arrays and to visualize 2D slices in each dimension of their amplitude or phase, with different colour-maps (diverging, sequential, cyclic) and scale (linear, logarithmic).

Moreover, an interactive 3D visualisation tool is provided that relies on the \textit{ipyvolume} library \parencite{ipyvolume}, itself built on top of \textit{ipywidgets}, and specifically designed to quickly render large 3D data arrays.
In the specific case of single volume rendering for the reconstructed object, the object's surface is defined by a threshold of its maximum density (fig. \ref{fig:3D_object}).
Finally, the object surface can be colour-mapped with the values of the displacement and strain retrieved during the data analysis (fig. 6).

To further interact with the figures, \textit{e.g.} zoom, set the colorbar range... , tools were implemented that rely on \textit{Bokeh} \parencite{Bokeh}, a Python library that transforms figures in interactive web-pages (fig. \ref{fig:BokehDetector}), that can also be displayed in Jupyter Notebook.

To resume, the GUI regroups data reduction, analysis and visualisation tools in the Jupyter Notebook interface via interactive methods.

\begin{figure}[!htb]
    \centering
    \includegraphics[width=0.66\textwidth]{/home/david/Documents/PhD/Figures/gwaihir/detector_slice_bokeh.png}
    \caption{Detector images can be viewed interactively with \textit{Bokeh}, to zoom on the data, and visualize, for example, the intensity collected on each pixel.}
    \label{fig:BokehDetector}
\end{figure}

\subsubsection{Command line scripts}

Nevertheless, the complete workflow for data processing can also be launched from the command-line using \textit{Python} scripts.
In \textit{Gwaihir}, the link between each package is based on \textit{BASH} scripts. This approach is both fast and versatile, designed to quickly iterate on several data-sets to test parameters values, but less intuitive compared to the interactive GUI.
cite clément

A set of default parameters stored in a configuration file is used, that can also be overwritten by directly providing keywords directly in the command line, \textit{e.g.} the scan number.
The configuration files are written in \textit{YAML} (fig. \ref{fig:YAML_file}).
In that case, the workflow is less interactive. One may still manually mask the data or manually select the best reconstruction before post-processing. This option can yield a quick result that, with fine parameters, allows a quick visualisation of the probed object.

\begin{figure}[!htb]
    \centering
    \includegraphics[width=\textwidth]{/home/david/Documents/PhD/Figures/gwaihir/yaml.png}
    \caption{Configuration file in YAML, a human-readable data-serialization language with a minimalist syntax. Each parameter used for the analysis methods are stored. A configuration file is generated for the pre-processing and post-processing scripts, as well as for the phase retrieval.}
    \label{fig:YAML_file}
\end{figure}

\subsubsection{Data storage}

Already, raw and processed data can be accessed for the BCDI technique through the \textit{.cxi} database (https://cxidb.org/) \parencite{Maia2012}, which aims at creating a single data-storing architecture/format for coherent X-ray imaging experiments.
The data input/output follows Nexus definitions \parencite{Konnecke2015} built in a CXI format hdf5 file.
Both clarity and consistency in data formatting encourages the reproducibility of the science produced (see https://www.panosc.eu/) and guarantees the workflow.

\begin{figure}[!htb]
    \centering
    \includegraphics[width=\textwidth]{/home/david/Documents/PhD/Figures/gwaihir/tree.png}
    \caption{
    Each parameter used during the data analysis is stored in the same \textit{.cxi} file, along with the results, in a  nested architecture.
    "instrument\_1" regroups parameters associated with the instrumental setup, "data\_1" is the collected diffraction pattern.
    "image\_2" regroups the parameters associated with phase retrieval; "data\_2" is the reconstructed Bragg electronic density chosen for post-processing.
    "image\_3" regroups parameters linked to data processing, as well as the processed amplitude, phase and resulting strain. "data\_3" is a link to the processed phase of the object.
    The file structure is displayed via \textit{JupyterHub}, a browser-based data analysis platform available on the \textit{SLURM} computer cluster for the ESRF and the \textit{GRADES} computing cluster at SOLEIL.
    }
    \label{fig:TREE}
\end{figure}

Since not all beamlines provide self-explaining NeXuS data-sets, it is the \textit{bcdi} package together with \textit{xrayutilities} \cite{kriegner_xrayutilities_2013} that allows the support of the coherent imaging beamlines ID01 (ESRF), P10 (PETRA III), SixS and CRISTAL (SOLEIL), NanoMAX (MAX IV) and 34-ID-C (APS).
Data pre-processing will generate two files stored as NumPy arrays \cite{NumPy}, corresponding to the diffraction intensity and mask.
These two files are then used for phase retrieval, for which the final object is saved in a \textit{cxi} file, later used for data post-processing.
In the case of simulation, the simulated diffraction intensity can be stored as a NumPy array to start the workflow from phase retrieval.

In \textit{Gwaihir}, data sharing across teams and team members is facilitated by the creation of a single output file, respecting the \textit{CXI} \cite{Maia2012}, and thus the NeXuS \cite{Konnecke2015} architectures.
As written in \cite{Konnecke2015}, authors of data-reduction and data-analysis software can use NeXus to store processed data along with metadata and a processing log. This is illustrated in fig. \ref{fig:TREE}, in which the \textit{cxi} data-set tree is displayed. It regroups parameters from phase retrieval together with parameters from pre-processing and post-processing.

The raw data, along with all the parameters associated to the analysis, and the final results are stored in a single \textit{cxi} file, allowing complete data analysis reproducibility (fig. \ref{fig:TREE}). This is possible at every step of the workflow. The aim is first to have a comprehensible architecture, and secondly to be able to reproduce anyone's result from this file. On a small scale, results are easier to share between collaborators and more understandable, while on a larger scale it will facilitate peer-review.

Key parameters for data reproducibility are the transformation matrix used for the final interpolation, the voxel size of the resulting 3D array, the probed reciprocal space range after data pre-processing $(\delta q_x , \delta q_y , \delta q_z)$, the iso-surface threshold, \textit{etc}. By sharing those parameters along with the diffraction intensity and complex Bragg electronic density, we will push towards more transparency in the methods used to obtain the final results.
Comments or metadata, such as the horizontal and vertical coherence lengths, or beam size, if determined prior to the experiment, can also be saved.
With \textit{Gwaihir} as an analysis tool, and a final \textit{cxi} file storing all the parameters and results, it becomes much more easier to review one's data and workflow.

Nevertheless, it allows the user to share the values of many parameters which can be difficult to select and which have an important influence over the final results. As mentioned earlier, the threshold used to determine the surface layer of the object has an important impact on the strain and displacement values on this surface. The center of the array, at which the phase is set to zero for phase unwrapping, becomes difficult to fix for a highly strained object.


\subsection{BINoculars} \label{sec:BINoculars}

\subsection{ROD} \label{sec:ROD}

\subsection{GenX} \label{sec:GenX}


\textcolor{Important}{
NXS, CXI, , reproducibility, repositories, PEP
}

    % \addtocontents{toc}{\protect\newpage}
    \chapter{Pt nanoparticles during ammonia oxidation}
    \section{Introduction}

% Reaction
During the catalytic oxidation of ammonia, if Pt catalysts were first used when aiming at the production of \ce{NO}, Pt-Rh catalysts have then proven to be more effective since the 1920s when Rhodium was added to the catalyst \parencite{Handforth1934, Heck1982}.
Recently, the ammonia oxidation reaction over a \ce{Pt_{25}Rh_{75}}(001) single crystal was studied with a surface science approach by operando techniques such as near-ambient pressure X-ray photoemission spectroscopy (NAP-XPS) and surface x-ray diffraction (SXRD) combined with mass spectrometry \parencite{Resta2020a}.
It was shown that \ce{NO} production coincides with significant changes of the surface structure and the formation of a \ce{RhO_2} surface oxide.
Moreover, changes in the surface relaxation related to the catalyst selectivity as well as to the presence of an hysteresis cycle in the reaction were reported.

% Single crystal ok but
Using single crystals as models to understand \textit{in-situ} and \textit{operando} the structure-activity relationship offers some perks such as a large surface area which increases the surface signal (lowering the pressure gap with some techniques such as XPS), and the possibility to isolate specific facets to progressively build up the understanding of the role of the sample surface on the catalytic activity \parencite{Hejral2016, Resta2020a}.
However, they fail when aiming at closing the material gap since they only exhibit a single crystallographic orientation on their surface, and are much larger (almost a \unit{cm} large) than the width of Pt-Rh gauzes used for the oxidation of ammonia (few 10s of \unit{\micro\meter} large, \cite{Kaiser1909}).

% use nanoparticles
In order to push for the reduction of the material gap as well as the pressure gap in heterogeneous catalysis, Pt nanoparticles have been used in this study, approaching the size of individual grains on the industrial catalysts \cite{} while exhibiting many different types of facets on their surface.
Specifically, the catalytic structure-activity relationships of Pt nanoparticles will be investigated during the oxidation of ammonia, as a stepping stone to the study of Pt-Rh nanoparticles, recently studied as well during the oxidation of carbon monoxide \parencite{Kim2021}, in which the compositional strain must also be taken into account \parencite{Kawaguchi2019}.
Moreover, studies have shown that nanoparticle surface strain can be controlled, opening up a new path to tune and optimise nanoparticle catalysts \parencite{Zhang2014, Sneed2015, Wang2016}.

% sxrd
It is first intended to measure the total signal scattered from \ce{Al_2O_3}-supported platinum particles by taking advantage of the possibility to carry out grazing-incidence diffraction measurements at the SixS beamline of synchrotron SOLEIL.
At grazing incidences the beam footprint extends across the whole sample length, the scattered beam being then proportional to the ensemble behaviour of the nanoparticles \parencite{Nolte2008, Hejral2013, Hejral2016}.

As seen in sec. \ref{sec:SXRD}, truncated surfaces such as facets give rise to crystal truncation rods (CTR) in the reciprocal space, whose intensity as a function of the scattering vector and width in reciprocal space are related to the facet size, roughness and strain.

The average nanoparticle shape and structure will be probed by studying the intensity of crystal truncation rods in directions perpendicular to the expected facets on the nanoparticle surface.
Prior experiments with similar samples \parencite{Dupraz2017, Li2020, Lim2021, Dupraz2022} have shown that the particles exhibit a Winterbottom shape \parencite{Winterbottom1967, Boukouvala2021}, typical of nanoparticles epitaxied on a substrate, with mainly \{111\}, \{110\}, and \{100\} facets, and a [111] orientation perpendicular to the substrate.

By having the incident beam illuminating all of the particles, the CTR signal in e.g. the [111] direction will be the sum of the contribution from the [111] facet of every illuminated nanoparticle (as well as their [$\bar{1}\bar{1}\bar{1}$] facets).
Therefore, phenomena inducing structural change such as particle re-faceting/reshaping at a given condition are expected to be visible by an evolution of the different CTR.

% bcdi
To obtain a more detailed picture of the phenomena at play, surface x-ray diffraction which provides ensemble-averaged properties when studying nanoparticles will be followed by the use of another technique which will bring information on the structural behaviour of single nanoparticles.

With Bragg coherent diffraction imaging (BCDI), the individual three-dimensional (3D) structural response of isolated nanoparticles to changes in the gas mixtures will be measured under \textit{in situ} and \textit{operando} conditions, enabling us to probe in detail the structure evolution of single facets during heterogeneous catalytic reactions (sec. \ref{sec:FacetAnalysis}).

A careful in situ analysis of the properties of the nanoparticles (size, shape, strain, re-faceting, support interaction, oxide formation, etc.) in 3D during the oxidation of ammonia is of essential importance to gain more understanding of the behaviour of these nanocrystals during a catalytic reaction.

However, it must first be ensured that the nanoparticles do not move or change their epitaxial relationship with the substrate during the exposition to reacting conditions at high temperatures.
Therefore, the preceding SXRD experiment will also aim at resolving the evolution of the epitaxial relationship between the Pt nanoparticles and the \ce{Al_2O_3} substrate before measuring single nanoparticles with Bragg coherent diffraction imaging, during which the beam is reduced to micrometric size, and for which a movement of the nanoparticles prevents the measurements of rocking curves due to loss of alignment.
Both BCDI and SXRD can be applied to high gas pressures due to the x-rays’ high penetration in gasses.

\subsection{$\alpha-$\ce{Al_2O_3} supported platinum nanoparticles synthesis}\label{sec:PtParticles}

The platinum nanoparticles were synthesised thanks to a collaboration with the Israel Institute of Technology (Technion, collaboration with Dr. Eugene Rabkin), their use first showcased in the work by Dupraz et al. \parencite*{Dupraz2017}.
A 30 nm thick homogeneous layer of platinum is deposited at room temperature on a [001] oriented alumina ($\alpha-$\ce{Al_2O_3}) substrate, which crystallises in a hexagonal unit cell.
The Pt nanocrystals have their c-axis oriented along the [111] direction, normal to the (0001) plane of the $\alpha-$\ce{Al_2O_3} substrate.
A mask is then applied on the sample and a lithographic process route ensures that the platinum layer transforms to nanoparticles that are between 100 and 1000 nm large, epitaxied on the substrate surface.

Two different samples are used during the oxidation of ammonia.
On one hand, in order to obtain a more important scattered intensity with SXRD, the sample surface is homogeneously covered with Pt particles.
On the other hand, there is a need to obtain isolated nanoparticles for BCDI to avoid alien signal coming from neighbouring nanoparticles.
The application of a specific patterned mask during the lithographic process then yields a patterned sample with isolated nanoparticles in the middle of \qty{100}{\um} large squares, visible in fig. \ref{fig:Mask}.

The position of the squares is designated with Arabic numbers (row), letters (column) and roman numbers (large rectangle).
All of the position indicators are constituted of platinum nanoparticle as well, which allows the mapping of the sample's surface in Bragg condition to find the nanoparticles' positions (fig. \ref{fig:SampleMapping}).
The hole diameter in the mask changes in different areas of the sample and has a direct impact on the average nanoparticle size in those areas.
After dewetting and heating at \qty{1100}{\degreeCelsius} for \qty{30}{\minute}, the platinum nanoparticles exhibit a well-faceted Winterbottom shape observed in various BCDI measurements \parencite{Dupraz2017}.

\begin{figure}[!htb]
    \centering
    \includegraphics[width=0.49\textwidth]{/home/david/Documents/PhD/Figures/sample/mask.png}
    \includegraphics[width=0.49\textwidth]{/home/david/Documents/PhD/Figures/sample/litho1_brighter.png}
    \caption{
        Mask applied during sample preparation (left) and resulting pattern on the sample surface observed with an optical microscope (right).
    }
    \label{fig:Mask}
\end{figure}

The epitaxial relationship between platinum and $\alpha$-\ce{Al_2O_3} (also known as sapphire) has been shown to consist into the superposition of the Pt [111] plane on top of the $\alpha$-\ce{Al_2O_3} [0001] plane, \textit{i.e.} Pt[111]||Al$_2$O$_3$[0001] (fig. \ref{fig:Epitaxy} - \cite{Farrow1993}).
The \{110\} platinum planes are perpendicular to the Pt (111) plane, the second-neighbour Pt atoms are arranged in a hexagonal pattern.

\begin{minipage}{0.55\linewidth}
    \centering
    \includegraphics[width=\linewidth]{/home/david/Documents/PhD/Figures/introduction/Epitaxy.pdf}
    \captionof{figure}{
    Pt[111]||Al$_2$O$_3$[0001] epitaxy relationship.
    }
    \label{fig:Epitaxy}
\end{minipage}%
\hfill% Add horizontal space between minipages
\begin{minipage}{0.44\linewidth}
    \begin{equation}
        \epsilon = \big( a_{Pt} \frac{\sqrt{3}}{\sqrt{2}} - a_{Sapphire} \big) / a_{Sapphire}
        \label{eq:MisfitStrain}
    \end{equation}
\end{minipage}%

Lattice strain in diffraction is usually defined as the difference between the reference and experimental lattice parameter values, respectively $a_{ref}$ and $a$ (eq. \ref{eq:StrainDiffraction}).

\begin{equation}
    \epsilon = \frac{a - a_{ref}}{a_{ref}}
    \label{eq:StrainDiffraction}
\end{equation}

In this case, the misfit strain at the platinum-alumina interface can be computed using eq. \ref{eq:MisfitStrain}, the reference lattice parameter is equalled to the in-plane lattice parameter of sapphire, while the lattice parameter of the larger hexagonal arrangement of the platinum atoms on the (111) surface is equal to $\sqrt{3} a_{Pt} / \sqrt{2} $.
At room temperature, the misfit strain computed from literature values ($a_{Pt} = \qty{3.9242}{\angstrom}$, $a_{Sapphire} = \qty{4.7602}{\angstrom}$) is computed to be equal to $\epsilon = \qty{0.96}{\percent}$.

\subsection{Catalysis reactor calibration for near-ambient pressure \textit{operando} studies}

The catalytic activity of the platinum nanoparticles was first studied as a function of the temperature to (i) be certain that the sample is sufficiently catalytically active for the reaction products to be detected by the mass spectrometer, (ii) identify the catalyst light-off temperature and (iii) observe an evolution of the product selectivity as a function of the temperature.
Similar scans where also performed on the empty reactor/sample holder to ensure the absence of activity towards the oxidation of ammonia.

The heater temperature was calibrated by measuring the temperature in the reactor at different pressures by the means of a type K thermocouple, as a function of the current intensity (fig. \ref{fig:TempRamps} - a), \qty{10}{\minute} separate consecutive datapoints to ensure the heater stability.
The experimental data points were fit to be able to set the reactor to any temperature from \qtyrange{25}{600}{\degreeCelsius} when working under vacuum or at ambient pressure (\qty{0.3}{\bar} or \qty{0.3}{\bar} of \ce{Ar}).
The thermal conductivity of the gases involved in the oxidation of ammonia are in the similar order of magnitude (tab. \ref{tab:ThermalConductivity}).
Moreover, the same inert gas used to set the reactor pressure - Argon - is used as a carrier gas during the experiments, constituting at least \qty{80}{\percent} of the gas flow, and allowing us to assume that the temperature in the reaction chamber is well approximated.

\begin{table}[!htb]
\centering
    \begin{tabular}{@{}llllllll@{}}
    \toprule
     & \ce{Ar} & \ce{NH_3} & \ce{O_2} & \ce{NO} & \ce{N_2O} & \ce{N_2}& \ce{H_2O} \\
    \midrule
    \qty{300}{\kelvin} & \num{17.7} & \num{25.1} & \num{26.5} & \num{25.9} & \num{17.4} & \num{26.0} & \num{18.6} \\
    \qty{400}{\kelvin} & \num{22.4} & \num{37.2} & \num{34.0} & \num{33.1} & \num{26.0} & \num{32.8} & \num{26.1} \\
    \qty{500}{\kelvin} & \num{26.5} & \num{53.1} & \num{41.0} & \num{39.6} & \num{34.1} & \num{39.0} & \num{35.6} \\
    \qty{600}{\kelvin} & \num{30.3} & \num{68.6} & \num{47.7} & \num{46.2} & \num{41.8} & \num{44.8} & \num{46.2} \\
    \bottomrule
    \end{tabular}%
\caption{Thermal conductivity in \unit{\mW \per \meter \per \kelvin} of gases \parencite{ThermalConductivityOfGases}.}
\label{tab:ThermalConductivity}
\end{table}

A temperature ramp at a reactor pressure of \qty{0.3}{\bar} (fig. \ref{fig:TempRamps} - b) was carried out to probe the evolution of the reaction products as a function of the temperature under reacting conditions, \textit{i.e.} a constant gas flow of \qty{41}{\ml\per\min} of \ce{Ar}, \qty{8}{\ml\per\min} of \ce{O_2}, and \qty{1}{\ml\per\min} of \ce{NH_3}.
In the frame of this thesis, the same ratio between incoming gas flows in assumed to be respected between the different gas partial pressures, in this case the partial pressure of oxygen is then equal to \qty{80}{\milli\bar}, the partial pressure of ammonia to \qty{10}{\milli\bar}, and the partial pressure of argon to \qty{410}{\milli\bar}.

\begin{figure}[!htb]
    \centering
    \includegraphics[width=0.495\textwidth]{/home/david/Documents/PhDScripts/SixS_2022_01_SXRD_Pt100/gas_analysis/figures/ThermocoupleCalibration.pdf}
    \includegraphics[width=0.495\textwidth]{/home/david/Documents/PhDScripts/Test_Reactor_CO2_2021_01/Figures/TempRamp2.pdf}
    \caption{
        a) Temperature inside the reactor cell measured with a type C thermocouple under vacuum and different \ce{Ar} pressures.
        b) Partial pressures evolution under a constant gas flow (\qty{41}{\ml\per\min} of \ce{Ar}, \qty{8}{\ml\per\min} of \ce{O_2}, \qty{1}{\ml\per\min} of \ce{NH_3}) at a reactor pressure of \qty{0.3}{\bar} during increasing and decreasing (low transparency) temperature ramp to \qty{650}{\degreeCelsius} with 100 steps, each lasting \qty{10}{\second}.
        The partial pressure of oxygen is omitted for simplicity.
    }
    \label{fig:TempRamps}
\end{figure}

The catalyst light off temperatures for the production of \ce{N_2}, \ce{NO} and \ce{N_2O} can respectively be identified to be around \qty{300}{\degreeCelsius}, \qty{450}{\degreeCelsius}, and \qty{550}{\degreeCelsius}.
The excess of oxygen compared to ammonia increases the production of \ce{NO} at high temperatures, with a partial pressure similar to the partial pressure of nitrogen at \qty{650}{\degreeCelsius}.
No decrease of the nitrogen partial pressure can yet be detected as a function of the increasing temperature.
The production of nitrous oxide stays at low values during the temperature ramp, almost imperceptible from the background pressure.

According to these primary results, the study of the oxidation of ammonia was originally decided to be carried out at \qtylist{300;500;600}{\degreeCelsius}, temperatures before and after the catalyst light off (tab. \ref{tab:ConditionsNanoparticles}).

\begin{table}[!htb]
\centering
\resizebox{\textwidth}{!}{%
    \begin{tabular}{@{}lll|l|lll|l@{}}
    \toprule
    \multicolumn{3}{l|}{Gas Flow (\unit{\ml\per\min})} & Total pressure (\unit{\milli\bar}) & \multicolumn{3}{l|}{Partial pressures (\unit{\milli\bar})} & Targeted information \\
    \midrule
    \midrule
    \ce{Ar} & \ce{O_2} & \ce{NH_3} &  & \ce{Ar} & \ce{O_2} & \ce{NH_3} &  \\
    \midrule
    \midrule
    50 & 0 & 0 & 300 & 300 & 0 & 0 & Catalyst state without reactants (unactive) \\
    49 & 0 & 1 & 300 & 294 & 0 & 6 & Ammonia adsorption \\
    48.5 & 0.5 & 1 & 300 & 291 & 3 & 6 &  \multirow{4}{*}{\begin{tabular}[c]{@{}l@{}}Influence of (\ce{NH_3}/\ce{O_2}) ratio on the \\ catalyst selectivity and structure \end{tabular}}\\
    48 & 1 & 1 & 300 & 288 & 6 & 6 &  \\
    47 & 2 & 1 & 300 & 282 & 12 & 6 &  \\
    41 & 8 & 1 & 300 & 276 & 18 & 6 & \\
    50 & 0 & 0 & 300 & 300 & 0 & 0 & Returning to pristine state \\
    \bottomrule
    \end{tabular}%
}
\caption{Different atmospheres used to probe the ammonia oxidation on Pt nanoparticles with BCDI.}
\label{tab:ConditionsNanoparticles}
\end{table}
    \section{Collective behaviour of Pt nanoparticles: SXRD}

\subsection{Experimental setup for SXRD experiments in the vertical geometry}\label{sec:SXRDSetupV}

The MED diffractometer was used in a vertical geometry (fig. \ref{fig:Diffractometer}).
The incident angle $\mu$ was fixed to \ang{0.3}, to illuminate a larger portion of the sample surface.
The incoming energy was set to \qty{18.44}{\keV}.
In the vertical configuration, the sample surface must be parallel to the plane defined by the incoming beam and the vertical axis, so that sample rotations do not change the orientation of the sample atomic planes.
The alignment of the sample was performed in two consecutive steps, first the out-of-plane sample position was adjusted by using the direct beam, secondly any possible tilt of the sample surface was corrected by recording the intensity of the \textit{reflected} beam as a function of the $u$ (when $\omega=\ang{0}$) or $v$ (when $\omega=\ang{90}$) angles of the hexapod (fig. \ref{fig:Diffractometer}).

\subsection{Epitaxial relationship under different atmospheres and temperature}

To ensure that the \ce{Al_2O_3} supported particles have the expected epitaxial relationship with the substrate, the scattering intensity in the plane parallel to the sample surface was measured.

\subsubsection{Room temperature study under inert atmosphere}

\begin{figure}[!htb]
    \centering
    \includegraphics[width=\textwidth]{/home/david/Documents/PhD/PhDScripts/SixS_2021_03_SXRD_NH3/figures/map/qxqyqz_38_40.pdf}
    \caption{
        In-plane reciprocal space map at room temperature.
        Starting from the centre of the map are (i) six $(2\bar{1}\bar{1}0)$-type Bragg peak from the substrate.
        (ii) thin circular powder signals for the $(111)$, $(200)$ and $(220)$ platinum Bragg peaks.
        (iii) six isolated $(220)$-type platinum Bragg peaks.
        (iv) $(3\bar{3}00)$-type Bragg peaks from the substrate next to the $(220)$-type platinum Bragg peaks.
    }
    \label{fig:QxQyMap}
\end{figure}

The unit cell used to index the different Bragg peaks of platinum in this section is the FCC unit cell presented in sec. \ref{sec:ScatCrystal}.
This first map of the reciprocal space was recorded under inert argon atmosphere at room temperature (fig. \ref{fig:QxQyMap}), and shows the six expected (220)-type in-plane peaks for [111] oriented nanoparticles (fig. \ref{fig:Epitaxy}).
The h, k, l Miller indices must have the same parity for the Bragg reflection to be allowed for face-centred cubic crystals, leading to the extinctions of (110)-type Bragg peaks.
From the average position of these peaks in $q$-space is computed the average in-plane lattice parameter of the platinum nanoparticles $a_{Pt}=\qty{3.91}{\angstrom}$, from which the scattering vectors $\vec{q}$ of the (200) and (111) Bragg peaks are computed (tab. \ref{tab:Reflections}).

Three arcs are drawn in fig. \ref{fig:QxQyMap} to underline the value of the (111), (200) and (220) scattering vectors.
If no (200) and (111) Bragg peaks are observed, powder signals are visible for each reflection, which shows that some of the nanoparticles have a random orientation.
Six (2$\bar{1}$0)-type Bragg peaks coming from the \ce{Al_2O_3} substrate are also observed at the lowest magnitude of the scattering vector ($q = \qty{2.64}{\per\angstrom}$, tab. \ref{tab:Reflections}).
From the average position of these peaks in $q$-space is also computed the in-plane lattice parameter of the substrate, $a_{Sapphire}=\qty{4.77}{\angstrom}$.
A first approximation of the in-plane misfit strain at the interface of the nanoparticles with the substrate is given following eq. \ref{eq:StrainDiffraction}: $\epsilon = \qty{0.45}{\percent}$.
This value must be considered with caution since the Pt in-plane lattice parameter is averaged over the volume of all the nanoparticles present on the sample, and is not representative of the platinum-sapphire interface.
$(3\bar{3}00)$-type Bragg peaks from the substrate are also visible, next to $(220)$-type platinum Bragg peaks.
Not all six peaks can be seen which may be due to the sample alignment, or to the sample coverage by platinum nanoparticles, important for the patterned sample.

\begin{table}[htb!]
    \begin{minipage}{.475\linewidth}
        \centering
        \resizebox{\textwidth}{!}{%
            \begin{tabular}{@{}lllll@{}}
            \toprule
            \multicolumn{5}{c}{Platinum Bragg peaks} \\
            (h k l) & $2\theta$ & q & Int & Int (\%) \\
            \midrule
            (1 1 1) & \ang{17.13} & \qty{2.78}{\per\angstrom} & \num{6362.52}  & \num{100.0} \\
            (2 0 0) & \ang{19.80} & \qty{3.21}{\per\angstrom} & \num{3251.30}  & \num{51.10} \\
            (2 2 0) & \ang{28.14} & \qty{4.55}{\per\angstrom} & \num{2399.87}  & \num{37.72} \\
            % (3 1 1) & \ang{33.13} & \qty{5.33}{\per\angstrom} & \num{2913.37}  & \num{45.79} \\
            % (2 2 2) & \ang{34.65} & \qty{5.57}{\per\angstrom} & \num{843.07}   & \num{13.25} \\
            % (4 0 0) & \ang{40.2235} & \qty{6.43}{\per\angstrom} & \num{390.01}   & \num{6.13}  \\
            % (3 3 1) & \ang{44.0121} & \qty{7.01}{\per\angstrom} & \num{1155.08}  & \num{18.15} \\
            % (4 2 0) & \ang{45.2178} & \qty{7.19}{\per\angstrom} & \num{1054.07}  & \num{16.57} \\
            % (4 2 2) & \ang{49.8120} & \qty{7.87}{\per\angstrom} & \num{755.99}   & \num{11.88} \\
            % (5 1 1) & \ang{53.0614} & \qty{8.35}{\per\angstrom} & \num{808.25}   & \num{12.70} \\
            \bottomrule
            \end{tabular}%
        }
    \end{minipage}%
    \hfill
    \begin{minipage}{.475\linewidth}
        \centering
        \resizebox{\textwidth}{!}{%
            \begin{tabular}{@{}lllll@{}}
            \toprule
            \multicolumn{5}{c}{\ce{Al_2O_3} Bragg peaks} \\
            (h k i l) & $2\theta$ & q & Int & Int (\%) \\
            \midrule
            % (1 $\bar{1}$ 2)         & \ang{11.07} & \qty{1.80}{\per\angstrom} & \num{21.50} & \num{49.52} \\
            % (1 $\bar{1}$ $\bar{4}$) & \ang{15.13} & \qty{2.46}{\per\angstrom} & \num{35.60} & \num{82.01} \\
            (2 $\bar{1}$ $\bar{1}$ 0) & \ang{16.21} & \qty{2.64}{\per\angstrom} & \num{16.81} & \num{38.73} \\
            (3 $\bar{3}$ 0 0)         & \ang{28.27} & \qty{4.57}{\per\angstrom} & \num{28.40} & \num{65.43} \\
            % (4 $\bar{2}$ $\bar{2}$ 0) & \ang{32.76} & \qty{5.27}{\per\angstrom} & \num{3.21}  & \num{7.40} \\
            % (0 0 6)                 & \ang{17.85} & \qty{2.90}{\per\angstrom} & \num{0.18}  & \num{0.42}  \\
            (2 $\bar{1}$ 1 3)         & \ang{18.52} & \qty{3.01}{\per\angstrom} & \num{38.90} & \num{89.60} \\
            % (2 $\bar{2}$ $\bar{2}$) & \ang{19.67} & \qty{3.195}{\per\angstrom} & \num{0.59}  & \num{1.35}  \\
            % (2 $\bar{2}$ 4)         & \ang{22.24} & \qty{3.607}{\per\angstrom} & \num{20.98} & \num{48.32} \\
            % (2 $\bar{1}$ $\bar{6}$) & \ang{24.21} & \qty{3.921}{\per\angstrom} & \num{43.41} & \num{100.0} \\
            % (3 $\bar{2}$ $\bar{1}$) & \ang{25.06} & \qty{4.057}{\per\angstrom} & \num{1.15}  & \num{2.65}  \\
            % (3 $\bar{1}$ $\bar{2}$) & \ang{25.60} & \qty{4.142}{\per\angstrom} & \num{1.47}  & \num{3.39}  \\
            \bottomrule
            \end{tabular}%
        }
    \end{minipage}%
    \caption{
        Scattering angle, scattering vector magnitude, and intensity of the scattered waves for selected Pt and \ce{Al_2O_3} Bragg peaks, computed for an energy of \qty{18.44}{\keV} using eq. \ref{eq:Bragglaw} and eq. \ref{eq:Fcrystal}.
    }
    \label{tab:Reflections}
\end{table}

\subsubsection{Study during ammonia oxidation before and after catalyst light off temperature}

In-plane maps were recorded on a smaller angular range (shown by the extent of black arcs in fig. \ref{fig:QxQyMap}) due to the long experimental time needed to collect large maps.
The intensity was then integrated in a \ang{0.1} wide region around the (111), (200) and (220) scattering angles.
The region around the (220) Bragg peak was carefully selected to avoid the Bragg peaks of the substrate.
These measurements are presented respectively in fig. \ref{fig:Epitaxy111}, fig. \ref{fig:Epitaxy200} and fig. \ref{fig:Epitaxy220}.

\begin{figure}[!htb]
    \centering
    \includegraphics[width=\textwidth]{/home/david/Documents/PhD/PhDScripts/SixS_2021_03_SXRD_NH3/figures/epitaxy/220.pdf}
    \caption{
        Integrated intensity in a \ang{0.1} range around the value of the (220) platinum scattering angle, as a function of the in-plane sample angle $\omega$, presented for different atmospheres.
    }
    \label{fig:Epitaxy220}
\end{figure}

Two (220)-type Bragg peaks are measured in each condition, with a stable position, shape, and intensity (fig. \ref{fig:Epitaxy220}), which shows that the crystals do not rotate around their [111] axis on the surface, and that they all share the same out-of-plane [111] orientation.
No (200) or (111)-type Bragg peaks appear during the different measurements (figures visible in the appendix \ref{fig:Epitaxy200}, and \ref{fig:Epitaxy111}).
There is no variation of the powder signal for the (200) Bragg peak, but a slight decrease of the powder signal for the (111) Bragg peak when increasing the temperature from \qtyrange{300}{500}{\degreeCelsius}.
Most importantly, these measurements confirm that the nanoparticles are still epitaxied on the substrate even at high temperature under highly aggressive environments, proving a strong epitaxial relationship of the nanoparticles with the substrate

($\bar{1}\bar{1}\bar{1}$) and (111) facets are expected to be present on all nanoparticles, respectively at the interface with the substrate and at the top of the nanoparticles.
For that reason, we followed the signal of those facets by measuring the CTR intensity as a function of $q_z$ (fig. \ref{fig:2DCTR111Particles}).

The ($\bar{1}$11) and (200) Bragg peaks can be identified respectively at $q_z \qty{\approx0.92}{\per\angstrom}$ and $q_z \qty{\approx1.85}{\per\angstrom}$.
Both peaks being visible along the same CTR shows that all nanoparticles do not share the same in plane orientation, due to different stacking of the (111) layers along the [111] axis, e.g. ABC or ACB stacking, rotated by \ang{180} \parencite{Jones2019}.
Low intensity CTRs in other directions due to other facets present on the nanoparticles can be seen around both Bragg peaks.
The peak at $q_z \qty{\approx1.46}{\per\angstrom}$ corresponds to a \{2$\bar{1}1$3\} Bragg peak from the \ce{Al_2O_3} substrate ($q = \qty{3.01}{\per\angstrom}$, tab. \ref{tab:Reflections}).
The substrate is also at the origin of its own CTR in the [111] direction, parallel to the more intense Pt CTR (fig. \ref{fig:2DCTR111Particles}).

\begin{figure}[!htb]
    \centering
    \includegraphics[height=7cm]{/home/david/Documents/PhD/PhDScripts/SixS_2021_03_SXRD_NH3/figures/ctr/CTR111_K_vmax100_edited.pdf}
    \includegraphics[height=7cm]{/home/david/Documents/PhD/PhDScripts/SixS_2021_03_SXRD_NH3/figures/ctr/CTR_300_edited.pdf}
    \caption{
        Left: Large out-of-plane reciprocal space map in which a CTR signal in the [111] direction can be identified, as well as weaker CTRs passing through the Bragg peaks.
        The ($\vec{q}_x$, $\vec{q}_y$) plane was rotated around $\vec{q}_z$ to highlight the presence of signals from facets on the particles other than the (111) facet, indicated by white arrows.
        Right: Crystal truncation rod signals in the $[111]$ direction as a function of $q_z$, presented for different atmospheres, at \qty{300}{\degreeCelsius}.
    }
    \label{fig:2DCTR111Particles}
\end{figure}

The CTR intensity as a function of $q_z$ is presented in fig. \ref{fig:2DCTR111Particles} (right).
There is no visible evolution of the CTR shape, the anti-Bragg region where the relaxation effects are the most visible near $q_z \qty{\approx1.46}{\per\angstrom}$ being masked by the Bragg peak of the substrate.
Efforts were then focused on resolving the CTR signal around the Bragg peaks coming from different facets present on the nanoparticles surfaces.

\subsection{Particle reshaping during the oxidation of ammonia}

\begin{figure}[!htb]
    \centering
    \includegraphics[width=\textwidth]{/home/david/Documents/PhD/PhDScripts/SixS_2021_03_SXRD_NH3/figures/facets/facets_together.pdf}
    \caption{
        Reciprocal space volume around the $(\bar{1}11)$ Bragg peak, collected at \qty{300}{\degreeCelsius}, under argon atmosphere.
        a) Projection perpendicular to $\vec{q}_z$.
        The three white dashed line are respectively a projection in the $(\vec{q}_x, \vec{q}_y)$ plane of the $[\bar{1}11]$, $[\bar{1}10]$ and $[010]$ directions.
        b-c) Slices in the $(\vec{q}_y, \vec{q}_z)$ plane after rotation of the $\vec{q}_x$ and $\vec{q}_y$ axis around the $\vec{q}_z$ axis by \ang{30} and \ang{60} to highlight the presence of crystal truncation rods in different directions.
        d) Slice in the $(\vec{q}_y, \vec{q}_z)$ plane to highlight the presence of a crystal truncation rod in the $[1\bar{1}1]$ direction.
        The interval delimited by the vertical white dashed lines in (b-c-d) is used for the qualitative analysis of the CTRs intensities.
    }
    \label{fig:FacetMaps}
\end{figure}

A volume of the reciprocal space was collected around the ($\bar{1}$11) Bragg peak to find the direction of the observed facet signals, and to quantify their evolution as a function of the different conditions.
Similar experiments have shown the oxygen-induced shape transformation of rhodium \parencite{Nolte2008}, or platinum nanoparticles \parencite{Hejral2013}, which was also linked to the presence of surface oxides.
The data was computed in $q$-space with different in-plane offsets, so that each rod observed in fig. \ref{fig:FacetMaps} (a) becomes parallel to the $\vec{q}_x$ direction for a certain in-plane offset, thereby facilitating the data analysis.

Beside the crystal truncation rods, there are two different parasitic signals going through the Bragg peak.
First, a curved signal almost parallel to $\vec{q}_y$ extends from $q_y = \qtyrange{-2.80}{-2.40}{\per\angstrom}$ in the ($\vec{q}_x, \vec{q}_y$) plane (in red in fig. \ref{fig:FacetMaps} - a).
This signal comes from scattered x-rays difusing towards the detector in position near the Bragg peak, and increasing the background signal.
Secondly, a very intense powder ring ($\approx \num{1e2}$) can be seen around the crystal truncation rod perpendicular to $\vec{q}_z$, especially in fig. \ref{fig:FacetMaps} (c).

\begin{figure}[!htb]
    \centering
    \includegraphics[width=\textwidth]{/home/david/Documents/PhD/Figures/Introduction/stereographic_projection_top.pdf}
    \caption{
        $[111]$ stereographic projection (for a face-centred cubic lattice).
        The circles describe the angle with the $[111]$ direction from \ang{0} (centre) to \ang{90} (outer-ring).
    }
    \label{fig:StereoTop}
\end{figure}

The identification of the crystallographic direction of each CTR present around the Bragg peak was performed thanks to the [111] stereographic projection.
This stereographic projection is presented in fig. \ref{fig:StereoTop}, which in the current experiment is parallel to $\vec{q}_z$.
Each dot represents a crystallographic direction, \textit{e.g.} [100], [110], etc.
The distance between the dots and the centre of the figure is a function of the angle with the [111] direction, whereas their angular distribution corresponds to the planar angle of the component perpendicular to the [111] direction.
Therefore, a signal that is for example at an angle of \ang{30} with the [111] direction and \ang{30} with the [$\bar{1}$01] direction can be identified to be in the [113] direction thanks to fig. \ref{fig:StereoTop}.

A CTR parallel to $\vec{q}_y$ is visible in fig. \ref{fig:FacetMaps} - (a), with an angle of \ang{72} with the [111] direction (fig. \ref{fig:FacetMaps} - (d)), thus corresponding to a signal in the [1$\bar{1}$1] direction.
After having first identified the [111] and [$1\bar{1}1$] directions, the identification of the direction of the other CTRs becomes straightforward.

Two CTRs in the [$1\bar{1}3$] and [$0\bar{1}1$] directions are identified, both have an angle of \ang{30} with the [$1\bar{1}1$] direction (visible in  fig. \ref{fig:FacetMaps} - a), and an angle of respectively \ang{30} and \ang{90} with the [$\bar{1}\bar{1}\bar{1}$] direction in the ($\vec{q}_x, \vec{q}_{y, 30}$) plane visible in fig. \ref{fig:FacetMaps} - (b).

A [001] oriented CTR visible at \ang{60} from the [$1\bar{1}1$] direction (fig. \ref{fig:FacetMaps} - a), has an angle of \ang{60} with the [$\bar{1}\bar{1}\bar{1}$] direction in the ($\vec{q}_x, \vec{q}_{y, 60}$) plane visible in fig. \ref{fig:FacetMaps} - (c).

Two [$11\bar{1}$] and [$\bar{1}11$]-oriented CTRs are expected at \ang{120} in the ($\vec{q}_x, \vec{q}_y$) plane with the [$1\bar{1}1$]-oriented CTR.
Each [$1\bar{1}1$]-oriented CTR is also linked to a [$\bar{1}1\bar{1}$]-oriented CTR in the opposite direction, which is why there is a [$1\bar{1}1$]-oriented rod at \ang{60} from the [$1\bar{1}1$] direction (fig. \ref{fig:FacetMaps} - a), with an angle of \ang{108} with the [$\bar{1}\bar{1}\bar{1}$] direction in the ($\vec{q}_x, \vec{q}_{y, 60}$) plane visible in fig. \ref{fig:FacetMaps} - (c).

\begin{figure}[!htb]
    \centering
    \includegraphics[width=\textwidth]{/home/david/Documents/PhD/PhDScripts/SixS_2021_03_SXRD_NH3/figures/facets/facet_signal_evolution_together.pdf}
    \caption{
    Evolution of the scattered intensity taken along a square area perpendicular to the [111] direction at three different positions in the ($\vec{q}_x, \vec{q}_z$) plane to probe the intensity evolution of crystal truncation rods.
    }
    \label{fig:FacetSignal}
\end{figure}

Despite the high intensity parasitic signal, it was possible to determine the orientation of 5 crystal truncation rods.
From these results, the average shape of the particles on the sample is expected to be highly faceted, with not only small (\{100\}, \{110\}, \{111\}) but also high indices facets (\{113\}) present on their surface.

The intensity of the crystal truncation rods listed above was studied by measuring the scattered intensity as a function of $q_z$ in a square area in the ($\vec{q}_x, \vec{q}_y$) plane.
The same width was used for each study, the region of integration is detailed in fig. \ref{fig:FacetMaps} - (b-c-d) with white dashed lines.
The evolution of the integrated scattered intensity as a function of $q_z$ is presented if fig. \ref{fig:FacetSignal} for different atmospheres and temperatures.

The scattered intensity does not evolve at \qty{300}{\degreeCelsius} and \qty{500}{\degreeCelsius} as a function of the atmosphere in the reactor cell.
However, there is a transition between \qty{300}{\degreeCelsius} and \qty{500}{\degreeCelsius} when looking at the signal from the ($0\bar{1}1$) facets, that evolves from a double to a single peak when increasing the reactor temperature.
This can be due to defects present in the nanoparticles at lower temperature despite the high temperature annealing, that are removed when heating the sample (sec. \ref{sec:TempRampBCDI}).
%For example, \textcolor{Important}{Ask yves to detail again}

At \qty{600}{\degreeCelsius}, a progressive increase of the intensity is observed for the [$1\bar{1}1$], [$\bar{1}\bar{1}1$] and [$001$]-oriented rods, whereas a progressive decrease of the intensity is observed for the [$0\bar{1}1$] and [$1\bar{1}3$]-oriented rods, starting after the introduction of oxygen in the reactor.
There is an increase of scattered intensity in fig. \ref{fig:FacetSignal} (i) at $q_z \qty{\approx 0.95}{\per\angstrom}$ which could be from (110) facets at \ang{35} with the [111] direction (fig. \ref{fig:StereoTop}).
The low signal-to-background ratio in that region due to the powder ring makes it challenging to be certain of the existence of this CTR.

Moreover, the introduction of oxygen induces a global shift of the facet signals towards higher $q_z$ values which can be linked to a homogeneous compressive lattice strain in the [111] direction, shifting the position of the Bragg peak.
The position of the [$0\bar{1}1$]-oriented CTR, whose direction is perpendicular to the [111] direction, follows that shift in $q_z$, but its intensity vanishes after having more than \qty{1}{\ml\per\min} of oxygen in the reactor cell.
A qualitative evolution of the strain in the [111] direction can be obtained by looking at the remaining CTR position in $q_z$ (fig. \ref{fig:FacetSignal} - g\&i).
The gradual increase of the oxygen partial pressure in the cell is accompanied by two additional changes of strain (i) from compressive to tensile strain when increasing the oxygen partial pressure from \qty{5}{\milli\bar} to \qty{10}{\milli\bar} in the reactor and (ii) from tensile to compressive strain when increasing the oxygen partial pressure from \qty{20}{\milli\bar} to \qty{80}{\milli\bar}, the reference position being taken under argon and ammonia atmosphere.

These result indicate a progressive reshaping of the nano-catalysts towards particles exhibiting mostly \{111\} and \{100\} facets, together with the increase of the \ce{O_2}/\ce{NH_3} ratio.
A roughening of the particles could explain the loss of intensity from the [$0\bar{1}1$]-oriented CTR but not the increased intensity of the [$1\bar{1}1$], [$\bar{1}\bar{1}1$], and [001]-oriented CTRs.

The catalytic activity of the particles was recorded \textit{via} the use of a mass spectrometer and shows the production of nitrogen, nitrous oxide and nitrogen oxide (fig. \ref{fig:RGASXRDNanoparticlesComparison}), proving the activity of the catalyst.
At \qty{600}{\degreeCelsius}, the production of \ce{N_2} is favoured at lower oxygen partial pressures in comparison with \ce{N_2O} and \ce{NO}.
The opposite effects takes place at high oxygen partial pressures, with a clear transition between the main products from \ce{N_2} towards \ce{NO} when the oxygen to ammonia ratio is equal to 8.
The original mass spectrometer signal as a function of time can be seen in the appendix \ref{fig:RGA300SXRDNanoparticles}, \ref{fig:RGA500SXRDNanoparticles}, and \ref{fig:RGA600SXRDNanoparticles}.

\begin{figure}[!htb]
    \centering
    \includegraphics[width=\textwidth]{/home/david/Documents/PhD/PhDScripts/SixS_2021_03_SXRD_NH3/figures/rga/product_comparison_no_norm.pdf}
    \caption{
        Evolution of reaction product partial pressures during the SXRD experiment on the non-patterned sample containing Pt nanoparticles, at \qty{300}{\degreeCelsius}, \qty{500}{\degreeCelsius}, and \qty{600}{\degreeCelsius}.
        Mean partial pressures during \qty{1}{\minute} at the end of each condition, recorded from a leak in the reactor output by a residual gas analyser (RGA).
        The background pressure in the absence of reaction has been subtracted.
    }
    \label{fig:RGASXRDNanoparticlesComparison}
\end{figure}

The same measurements at \qty{600}{\degreeCelsius} were not repeated at a fixed atmosphere which prevents us from knowing whether or not the observed phenomena would have occurred as a function of time under a fixed atmosphere.
Moreover, the reversibility of this effect and the formation of possible platinum oxides on the nanoparticles (\textit{via} large in-plane maps) could not be investigated from the lack of available experimental time.

Similar experiments on Pt nanoparticles (average size about \qty{50}{\nm}) have been carried at \qty{6.5e-6}{\bar} and \qty{0.5}{\bar} of oxygen by \cite{Hejral2013} which show the formation of \ce{Pt_3O_4} and $\alpha-$\ce{PtO_2} bulk oxides, the formation of high indices facets, a decrease of \{111\} facet signals and an increase of \{100\} facet signals.
Different effects have been observed in this study with an increase in intensity of \{111\} and \{100\} facets, and the decrease in intensity of \{110\} facets during the oxidation of ammonia, highlighting the impact of the reaction on the particles' surface.

A study of the structure/activity relationship of each facet must be undertaken during the ammonia oxidation to better understand the role of the surface structure, and of potential surface oxides in the catalytic reaction, which is the subject of chapter 4.
    \section{Temperature ramp}\label{sec:TempRampBCDI}

% strain field energy: kinda useless
% need to put scan table in appendix
% specific rocking curves

To be able to de-corellate the effect of temperature from the effect of the catalytic reaction on the nanoparticle structure, the alumina-supported nanoparticles were gradually heated from \qtyrange{25}{600}{\degreeCelsius} under a constant \argon-based gas flow (\qty{50}{\ml\per\min}) and at a pressure of \qty{0.3}{\bar}.

The same nanoparticle (so-forth called nanoparticle Amaterasu) was tracked during the heating process, rocking-curves around the (111) Bragg peak were measured at different temperatures to probe its structural evolution, with \qty{201}{steps}, each counting for \qty{5}{\second}, resulting in an angular step of \ang{0.005}.
The measurement of a rocking curve took approximately \qty{22}{\minute}, with \qty{25}{\percent} of dead time.
Two rocking curves were measured at each temperature during the ramp to have improved statistics about the catalyst structure.
The alignement of the sample was performed with the procedure as detailed in sec. \ref{sec:BCDISetup}.

\subsection{Experimental setup for BCDI experiments in the vertical geometry}\label{sec:BCDISetup}

The BCDI experiment was performed at the SixS (Surface Interface X-ray Scattering) beamline of synchrotron SOLEIL, France (sec. \ref{sec:SIXS}).
As detailed in sec. \ref{sec:Gwaihir}, one of the bottlenecks of the BCDI technique is its slow data reduction and analysis process.
Moreover, on $3^{rd}$ generation synchrotrons that offer a lower coherent flux (eq. \ref{eq:CoherentFlux}) than $4^{th}$ generation synchrotrons, the measurement time can also be very long.
At SixS, a rocking curve lasts between \qtyrange{20}{90}{\min} depending on the particle size, the quality of the alignment, the strain of the particle, \textit{etc.}
Once the raw data is obtained, the particle must be \textit{reconstructed} (sec. \ref{sec:PhaseRetrieval}), and - if of interest - the displacement and strain arrays can be retrieved.
The analysis workflow can take up to an hour, which totals to a maximum of two hours from the start of the measurement to the moment when the user has a good idea of the sample shape and structure.

SixS is a beamline that does not only carry out BCDI experiment, but also SXRD experiments, in the same experimental end-station, the multi-environment diffractomer (MED, sec. \ref{sec:MED}).
When aiming at performing \textit{operando} catalysis experiments, switching from one setup to another can take up to a few days.
This leaves only a limited remaining amount of time to align the sample, find a suitable nanoparticle, and carry out the experimental plan.

Li \textit{et al.} \parencite*{Li2020}, who have first shown that the SixS beamline could be used to carry out BCDI experiment, also started to work on improving the BCDI measurement process.
By comparing continuous and step-by-step measurements, they have shown that continous scanning would result in the same data quality while decreasing the measurement dead-time by \qty{30}{\percent}, thereby paving the way for quicker BCDI measurements.

During this thesis, the measurement process was further improved by taking advantage of the new possibility to perform continuous \textit{on-the-fly} scans at SixS, while moving the hexapod holding the sample.
When using the coherence setup (fig. \ref{fig:OpticalSetup}) with the diffractomer in the vertical geometry (fig. \ref{fig:Diffractometer}), the beam focused on the sample is about \qty{1}{\um} large vertically, the horizontal footprint depending on the incident angle between the beam and the sample.
The incoming angle $\mu$ is set to the Bragg angle ($\theta$ in eq. \ref{eq:Bragglaw}) and the scattered x-rays are collected by setting the out-of-plane detector angle $\gamma$ at a position equal to $2\theta$ (similar to fig. \ref{fig:EwaldSphereSpecular}).
Finally, by simultaneously moving the sample with the hexapod and recording the Bragg scattered intensity with the detector, it is possible to map the sample surface with a sub-micron resolution (fig. \ref{fig:SampleMapping}).
A nanoparticle with a width equal to \qty{300}{\nm} was identified with this technique, which is a good estimate of the spatial resolution that can be attained, limited by the hexapod resolution (\qty{\approx 500}{\nm}) and the beam size.
A new generation of motots with a nanometric resolution will further improve this technique in the future.

During this experiment, the simplest possible measurement geometry is to measure the (111) Bragg peak, perpendicular to the sample surface since the nanoparticles have their c-axis oriented along the [111] direction, parallel to the normal of the sample holder.
% ($\vec{z}$ direction in the laboratory frame).
Both in-plane angles are thereby kept to zero during the mapping process.

\begin{figure}[!htb]
    \centering
    \includegraphics[height=5cm]{/home/david/Documents/PhD/Figures/sample/microscope_image.png}
    \includegraphics[height=5cm]{/home/david/Documents/PhD/Figures/sample/microscope_image_photon.png}
    \caption{
        Microscope image of the sample seen through the sapphire window of the PEEK dome (left).
        Map of the sample performed in Bragg condition (right), the high intensity (red) areas correspond to platinum nanoparticles.
        The letters, numbers and isolated nanoparticles in the centre of squares can be recognized on the sample.
        \textcolor{Important}{colorbar}
    }
    \label{fig:SampleMapping}
\end{figure}

On the other hand, \textit{Gwaihir} (sec. \ref{sec:Gwaihir}) was developped primarily for the SixS beamline to counter the long analysis process, which allowed a significant reduction in the analysis time from around an hour to a few minutes.
The following beamtimes profit from the new software by having a more \textit{solution}-driven experimental process.
Indeed, to be measured, a nanoparticle must be isolated, not too small (weak scattered intensity), not too big (loss of coherence, fringes not visible), and not too initially strained (difficult to obtain a good guess of the support).
These conditions are sometimes difficult to assert by simply looking at the diffraction pattern.
Therefore, quick inversion using \textit{Gwaihir} allowed a faster decision process regarding the continuation or not of the nanoparticles measurement.

Successfull measurements by \cite{Lim2021} have permitted the simulatenous use of BCDI measurements from SixS with measurements from other imaging beamlines (ID01 - ESRF, P10 - DESY), designing a robust method to identify defects in the real space with convolutional neural networks (CNN).

In the frame of this thesis, the required beam size was obtained with a Fresnel zone-plate (focal distance of \qty{20}{\cm}), which focused the beam down to \qty{1}{\um} (horizontally) $\times$ \qty{2}{\um} (vertically).
A coherent portion of the beam was selected with high precision slits by matching their horizontal and vertical gaps with the transverse coherence lengths of the beamline: \qty{20}{\um} (horizontally) and \qty{100}{\um} (vertically).
\textcolor{Important}{check consistency}
A circular beam-stop, and a circular order-sorting aperture, were used to block the transmitted beam, and higher diffraction orders, respectively (fig. \ref{fig:OpticalSetup}).
The sample was mounted in the dedicated XCAT reactor with the substrate surface oriented in the vertical plane on a hexapod mounted vertically on the MED diffractometer (fig. \ref{fig:MEDV}).

The BCDI experiment was performed in a vertical specular geometry at a beam energy of \qty{8.5}{\keV} (wavelength of \qty{1.46}{\angstrom}).
Three-dimensional (3D) diffraction data were collected with rocking curves of the rotation angle around the normal of the sample, the diffracted beam was recorded with a 2D MAXIPIX photon-counting detector (\numproduct{515 x 515} square pixels, \qtyproduct{55 x 55}{\um} wide) positioned on the detector arm at a distance of \qty{1.22}{\meter}.
The in-plane ($\omega$) and out-of-plane ($\mu$) angles of the sample at respectively equal to \ang{0} and \ang{18.6} when measuring at room temperature.
The in-plane ($\delta$) and out-of-plane ($\gamma$) angles of the detector are respectively equal to \ang{0} and \ang{37.2}.
The value of the scattering angle $\mu$ ($\theta$ in eq. \ref{eq:Bragglaw}) and the out-of-plane detector angle ($2\theta$) are expected to vary as a function of the temperature during the experiment due to the thermal expansion of the sample.
The alignement of the sample was performed with the same procedure as detailed in sec. \ref{sec:SXRDSetupV}.

\textit{Gwaihir} was used to process the data using the workflow detailed in sec. \ref{sec:Gwaihir}, each dataset was saved as a \textit{cxi} file containing all of the data and metadata related to the data reduction process.
The use of iterative algorithm for phase retrieval is detailed in tab. \ref{tab:ReconstructionProcess}.

\begin{table}[!htb]
\centering
\resizebox{\textwidth}{!}{%
    \begin{tabular}{@{}llll@{}}
    \toprule
    Iteration & Algorithm & PSF & Description \\
    \midrule
    0-199 & HIO & False & Work on the gross identification of the support (typical support voxel \% = 20) \\
    200-599 & RAAR & False & Refining the support (typical support voxel \% = 20) \\
    600 & RAAR & True & \begin{tabular}[c]{@{}l@{}}Activate the use of a point-spread function (PSF) to take into account the partial coherence \\ and the response of the detector.\end{tabular} \\
    601-999 & RAAR & True & Refine the PSF shape, the support must be already well defined to avoid diverging. \\
    1000-1200 & ER & True & Further refine the support by reducing the algorithm flexibility. \\
    \bottomrule
    \end{tabular}%
}
\caption{Example of algorithm chain used in BCDI for the phase retrieval.}
\label{tab:ReconstructionProcess}
\end{table}

\subsection{Shape evolution}

The surface of the Pt nanoparticle coloured by the surface voxel strain values is presented in fig. \ref{fig:AmaterasuA} and \ref{fig:AmaterasuB}, from \qtyrange{150}{600}{\degreeCelsius}, to show the evolution of the particle shape under an inert gas flow as a function of the temperature.
A dislocation identified at the interface with the substrate at \qty{150}{\degreeCelsius} from its strain signature and missing pipe of electronic density \parencite{Dupraz2015} was removed by an instrumental mistake which resulted in a flash heating procedure, the heater going from \qtyrange{150}{800}{\degreeCelsius} for a few minutes before coming back to \qty{150}{\degreeCelsius}.
The sample was realigned with the direct beam after the error of manipulation.

The displacement field could not be perfectly unwrapped at \qty{150}{\degreeCelsius}, \qty{350}{\degreeCelsius} and \qty{450}{\degreeCelsius}, resulting in saturated strain regions on the [$\bar{1}\bar{1}\bar{1}$] facet visible in the bottom view of the nanoparticle.

\begin{figure}[!htb]
    \centering
    \includegraphics[width=0.49\textwidth]{/home/david/Documents/PhDScripts/SixS\_2021\_01/paraview/1414top.png}
    \includegraphics[width=0.49\textwidth]{/home/david/Documents/PhDScripts/SixS\_2021\_01/paraview/1414bottom.png}
    \includegraphics[width=0.49\textwidth]{/home/david/Documents/PhDScripts/SixS\_2021\_01/paraview/1483top.png}
    \includegraphics[width=0.49\textwidth]{/home/david/Documents/PhDScripts/SixS\_2021\_01/paraview/1483bottom.png}
    \includegraphics[width=0.49\textwidth]{/home/david/Documents/PhDScripts/SixS\_2021\_01/paraview/1534top.png}
    \includegraphics[width=0.49\textwidth]{/home/david/Documents/PhDScripts/SixS\_2021\_01/paraview/1534bottom.png}
    \includegraphics[width=0.49\textwidth]{/home/david/Documents/PhDScripts/SixS\_2021\_01/paraview/1563top.png}
    \includegraphics[width=0.49\textwidth]{/home/david/Documents/PhDScripts/SixS\_2021\_01/paraview/1563bottom.png}
    \caption{
        Surface of the reconstructed Amaterasu Pt nanoparticle at \qty{150}{\degreeCelsius} (before and after the flash annealing), \qty{250}{\degreeCelsius}, and \qty{350}{\degreeCelsius}.
        The surface is coloured by the values of the heterogeneous strain as described in eq. \ref{eq:StrainTensorzz} at each surface voxel, the limit of each facet is delimited by thick white lines.
    }
    \label{fig:AmaterasuA}
\end{figure}

\begin{figure}[!htb]
    \centering
    \includegraphics[width=0.49\textwidth]{/home/david/Documents/PhDScripts/SixS\_2021\_01/paraview/1588top.png}
    \includegraphics[width=0.49\textwidth]{/home/david/Documents/PhDScripts/SixS\_2021\_01/paraview/1588bottom.png}
    \includegraphics[width=0.49\textwidth]{/home/david/Documents/PhDScripts/SixS\_2021\_01/paraview/1631top.png}
    \includegraphics[width=0.49\textwidth]{/home/david/Documents/PhDScripts/SixS\_2021\_01/paraview/1631bottom.png}
    \includegraphics[width=0.49\textwidth]{/home/david/Documents/PhDScripts/SixS\_2021\_01/paraview/1675top.png}
    \includegraphics[width=0.49\textwidth]{/home/david/Documents/PhDScripts/SixS\_2021\_01/paraview/1675bottom.png}
    \caption{
        Surface of the reconstructed Amaterasu Pt nanoparticle at \qty{450}{\degreeCelsius} and \qty{600}{\degreeCelsius} under both inert argon atmosphere and \qty{60}{\minute} after the introduction of \ammonia in the reactor.
        The surface is coloured by the values of the heterogeneous strain as described in eq. \ref{eq:StrainTensorzz} at each surface voxel, the limit of each facet is delimited by thick white lines.
    }
    \label{fig:AmaterasuB}
\end{figure}

The nanoparticle shape stayed roughly the same during the beginning of the temperature ramp and after the introduction of ammonia, the triangular [111] facet at the top of the particle can be recognized, with the same agency of the surrounding facets (fig. \ref{fig:AmaterasuA} - [111] facet surrounded by six facets, [100], [010], [001], [1$\bar{1}$1], [11$\bar{1}$], [$\bar{1}$11]).
The bottom facet, in contact with the substrate, has a [$\bar{1}\bar{1}\bar{1}$] orientation.

The low spatial resolution of the reconstructed nanoparticle due to the quality of the measurement (low coherent flux, \textit{i.e.} difficulty to resolve the scattered intensity far from the Bragg peak) resulted in difficulties when trying to identify the smallest facets on the nanoparticles, as well as their orientation.

\begin{figure}[!htb]
    \centering
    \includegraphics[width=\textwidth]{/home/david/Documents/PhD/Figures/introduction/stereographic_projection_bottom.pdf}
    \caption{
        Stereographic projection perpendicular to $[\bar{1}\bar{1}\bar{1}]$ crystallographic orientation.
        The circles describe the angle with the $[\bar{1}\bar{1}\bar{1}]$ direction from \ang{0} (center) to \ang{90} (outer-ring).
    }
    \label{fig:StereoBottom}
\end{figure}

Indeed, the crystal truncation rods perpendicular to small facets are less intense in comparison to large facets, which is at the origin of the direction-dependent spatial resolution in BCDI \parencite{cherukara_anisotropic_2018}.
Moreover, when identifying the facets in real space by the means of a surface recognition algorithm such as \textit{FacetAnalyser} (sec. \ref{sec:FacetAnalysis}), the normal to the smaller facets is more sensitive to deviation from its true value due to the low amount of voxels on the particle surface.

In fig. \ref{fig:StereoBottom} is presented the stereographic projection perpendicular to the $[\bar{1}\bar{1}\bar{1}]$ crystallographic orientation which can be used to identify the type of facets present on the nanoparticles in real or reciprocal space \parencite{Richard2018}.

In the frame of this thesis, the problem of facet identification was often encountered when facets were either too small to be resolved by the recognition algorithm, or when the uncertainty in the direction of the normal to the smallest facets is too important to be certain of the facet type.
For example, the area in orange in fig. \ref{fig:StereoBottom} represents a region in which facets with different structures have a similar orientation and for which it can be complicated to make the difference between $[\bar{1}\bar{1}1]$, $[\bar{2}\bar{2}1]$ and $[\bar{3}\bar{3}1]$ facets.

In this study, facets around the nanoparticle interface with the substrate have been identified as the smallest on the nanoparticle surface and the most subject to structural evolutions.
The orientation of all the facets present of the nanoparticle was performed in real space by the use of \textit{FacetAnalyser} taking into account the expected symmetry of the nanoparticle.

% 150 -> 150
After the flash annealing, the nanoparticle shape near the interface has become less round and more faceted from the disappearance of the defect, the surface strain between the interface and the particle has practically disappeared, a [$10\bar{1}$] facet appeared on the side of the nanoparticle near a [2$\bar{1}\bar{2}$] facet that replaced a [1$\bar{1}\bar{1}$] facet (fig. \ref{fig:AmaterasuDefect}).
A [$\bar{2}\bar{1}2$] facet is still present on the other side of the particle.
The total surface area occupied by the [$\bar{1}\bar{1}\bar{1}$] facet, ($1\bar{1}1$)-type and (100)-type facets has increased (fig. \ref{fig:AmaterasuFacetsEvolution} - a), the removal of the defect directly increasing the amount of voxels in the neighbouring facets (fig. \ref{fig:AmaterasuDefect}), \textit{i.e.} almost \qty{10}{\percent} more of the sample surface was recognized as part of the different facets by the algorithm.

\begin{figure}[!htb]
    \centering
    \includegraphics[width=0.49\textwidth]{/home/david/Documents/PhDScripts/SixS_2021_01/paraview/WireframeStrain1414.png}
    \includegraphics[width=0.49\textwidth]{/home/david/Documents/PhDScripts/SixS_2021_01/paraview/WireframeStrain1483.png}
    \caption{
        The dislocation results in a large volume of the particle that is not visible, due to the large strains \textcolor{Important}{Better explain this}.
        Removing the dislocation increase the area of the particle that is then recognized as facets, especially for the $[\bar{1}\bar{1}\bar{1}]$ facet, at the interface with the substrate.
    }
    \label{fig:AmaterasuDefect}
\end{figure}

%150->450, loose all 212, get -101 that then change side
When heating from \qty{150}{\degreeCelsius} to \qty{250}{\degreeCelsius} , the [$\bar{2}\bar{1}2$] facet present from before the defect removal is transformed towards a [$\bar{1}\bar{1}1$] facet.
When heating to \qty{350}{\degreeCelsius}, the [$10\bar{1}$] facet disappeared and a [$\bar{1}01$] facet appeared on the opposite side of the particle, this transformation is visible in fig. \ref{fig:AmaterasuA}.
If one expects a certain degree of symmetry from the equilibrium Winterbottom shape of a particle on a substrate \parencite{WINTERBOTTOM1967, Boukouvala2021}, this symmetry seems to only be reached at \qty{450}{\degreeCelsius} with three [$\bar{1}1\bar{1}$], [1$\bar{1}\bar{1}$] and [$\bar{1}\bar{1}$1] facets around the substrate for a total of 11 facets, which is also the only reconstruction without any ($\bar{2}\bar{1}2$), ($\bar{1}01$) and ($2\bar{1}\bar{1}$)-type facets (fig. \ref{fig:AmaterasuFacetsEvolution}).

Overall, the surface occupied by the facets around the top of the particle increase with heating, while the surface area that is not recognized as faceted by the algorithm decreases, the respective number of each \{$1\bar{1}1$\}, \{100\}-type facet does not change during the experiment.
% was multipied by \qty{50}{\percent} between \qty{25}{\degreeCelsius} and \qty{450}{\degreeCelsius}, which is also the temperature at which most of the particle surface was recognized as faceted (fig. \ref{fig:AmaterasuFacetsEvolution}).
The facets that are around the bottom of the particle are the most subject to change during this temperature ramp while the facets at the top of particle are more stable (fig. \ref{fig:AmaterasuFacetsEvolution}).
It is possible that due to the low imaging resolution of the experiment, smaller facets that exist on the particle surface cannot be distinguished.
Indeed, fig. \ref{fig:AmaterasuFacetsEvolution} shows that the total surface area not recognized as facets from the algorithm is always near \qty{40}{\percent}.

\begin{figure}[!htb]
    \centering
    \includegraphics[width=\textwidth]{/home/david/Documents/PhDScripts/SixS_2021_01/FacetAnalyser/FacetSizeEvolution.pdf}
    \caption{
        a) Evolution of the number of facet near the bottom (interface with the substrate) of the particle.
        b) Evolution of the particle surface area occupied by specific facets (indicated with []) or facet families (indicated with \{\}).
        The surface not recognized as part of a facet by the algorithm (e.g. rough areas of the particle surface, edges and corners) is taken into account and shown in black.
    }
    \label{fig:AmaterasuFacetsEvolution}
\end{figure}

% 450 -> 600
Interestingly, after the reshaping of the particle at \qty{450}{\degreeCelsius} that shows the disappearance of the higher Miller indices facets (fig. \ref{fig:AmaterasuFacetsEvolution}), the following rocking curves where very difficult to reconstruct until \qty{600}{\degreeCelsius}.
The increase of temperature could have been the trigger of thermally-induced mobility of the Pt atoms on the particle surface, making it impossible for the algorithm to determine a fixed support for the reconstruction.
\textcolor{Important}{find citations here}

A reconstruction at \qty{600}{\degreeCelsius} shows the appearance of a defect at the interface with the substrate (fig. \ref{fig:AmaterasuB}), and the appearance of a [$\bar{2}\bar{1}2$] facet that was present on the nanoparticle at \qty{150}{\degreeCelsius} before the removal of the defect (fig. \ref{fig:Amaterasu110}).
After the introduction of \ammonia at \qty{600}{\degreeCelsius}, the reconstruction of the particle shows the appearance of a [2$\bar{1}\bar{1}$] facet for the first time, together with the appearance of a [$\bar{1}01$] facet and the disappearance of the [$\bar{2}\bar{1}2$] facet, similarily to the changes that occured at \qty{150}{\degreeCelsius} after the removal of the defect (fig. \ref{fig:Amaterasu110}).
The change is also visible in fig. \ref{fig:AmaterasuFacetsEvolution}, for example the relative surface occupied by the \{100\}-type facets decreases, occupied by the [$\bar{1}01$] facet.

\begin{figure}[!htb]
    \centering
    \includegraphics[width=0.32\textwidth]{/home/david/Documents/PhDScripts/SixS\_2021\_01/paraview/110facet1588.png}
    \includegraphics[width=0.32\textwidth]{/home/david/Documents/PhDScripts/SixS\_2021\_01/paraview/110facet1631.png}
    \includegraphics[width=0.32\textwidth]{/home/david/Documents/PhDScripts/SixS\_2021\_01/paraview/110facet1675.png}
    \caption{
        Surface of the reconstructed Amaterasu Pt nanoparticle at \qty{450}{\degreeCelsius} under Argon and at \qty{600}{\degreeCelsius}, \qty{60}{\minute} after the introduction of \ammonia, highlighting the appearance of a [110] facet.
        The surface is coloured by the values of $\epsilon_{zz}$ at each surface voxel, the limit of each facet is delimited by white tubes.
    }
    \label{fig:Amaterasu110}
\end{figure}

When introducing \dioxygen at \qty{600}{\degreeCelsius} to study the oxidation of ammonia, the nanoparticle was definitely lost during the measurement, not to be found again, which underlines the difficulty to study a highly exothermic reaction with BCDI.
Indeed, according to \cite{}, the reaction can heat the catalyst to very high temperatures which in our case could have been the trigger for the loss of the particle during the measurement, especially since the particle had so far resisted to the beam during the temperature ramp to \qty{600}{\degreeCelsius}.
From this first set of results, the measurement of Pt nanoparticles with BCDI during the oxidation of ammonia was decided to be carried out at lower temperatures, e.g. \qty{300}{\degreeCelsius} and \qty{450}{\degreeCelsius}.

\subsection{Strain evolution} \label{sec:StrainTempRamp}

\subsubsection{Determination of strain}

Lattice strain in diffraction is usually defined as the difference between the  reference and experimental lattice parameter values, respectively $a_{ref}$ and $a$ (eq. \ref{eq:StrainDiffraction}).

\begin{equation}
    \epsilon = \frac{a - a_{ref}}{a_{ref}}
    \label{eq:StrainDiffraction}
\end{equation}

The values of the lattice parameter can be extracted from the position $\vec{G}$ of the Bragg peaks in reciprocal space \textit{via} eq. \ref{eq:QandD3} and eq. \ref{eq:Interplanarspacing}.
However, the information extracted from the retrieved phase in BCDI is more complex since, by measuring three non-coplanar Bragg peaks, one may retrieve the full strain tensor (eq. \ref{eq:StrainTensor}) from the reconstructed displacement field $\vec{u}_{\hat{q_x}, \hat{q_y}, \hat{q_z}}$ (eq. \ref{eq:DisplacementField}), $\hat{q_x}, \hat{q_y}, \hat{q_z}$ being three orthogonal basis vector in the reciprocal space \parencite{Karpov2019}.

\begin{equation}
    \vec{u}_{\hat{q_x}, \hat{q_y}, \hat{q_z}} =
     \begin{pmatrix}
        \vec{u}_{\hat{q_x}} \\
        \vec{u}_{\hat{q_y}} \\
        \vec{u}_{\hat{q_z}} \\
     \end{pmatrix}
     \label{eq:DisplacementField}
\end{equation}

\begin{multicols}{3}
    \begin{equation}
        \epsilon =
        \begin{bmatrix}
            \epsilon_{xx} & \epsilon_{yx} & \epsilon_{zx}\\
            \epsilon_{xy} & \epsilon_{yy} & \epsilon_{zy}\\
            \epsilon_{xz} & \epsilon_{yz} & \epsilon_{zz}
        \end{bmatrix}
        \label{eq:StrainTensor}
    \end{equation}
    \break
    \begin{equation}
      \epsilon_{ij} = \frac{1}{2}
        \Bigg(
        \frac{\partial \vec{u}_{\hat{q_i}}}{\partial \hat{q_j}}
        +
        \frac{\partial \vec{u}_{\hat{q_j}}}{\partial \hat{q_i}}
        \Bigg)
        \label{eq:StrainTensorIJ}
    \end{equation}
    \break
    \begin{equation}
      \epsilon_{zz} =
        \Bigg(
        \frac{\partial \vec{u}_{\hat{q_z}}}{\partial \hat{q_z}}
        \Bigg)
        \label{eq:StrainTensorzz}
    \end{equation}
\end{multicols}

The idea behind using the strain tensor is to identify shear components in the displacement field, \textit{i.e.} see if components in one direction depend on other directions, which can also help to identify defects present in nanoparticles \parencite{Lauraux2021}.
Using a single Bragg peak, we can only retrieve one component of the displacement field, obtained from the division of the retrieved phase $\Phi$ of the scattered x-rays by the value of the scattering vector at the position of the Bragg peak (\textit{i.e.} $\vec{q} = \vec{G}$), following the assumptions for phase retrieval detailed in sec. \ref{sec:StrainBCDI}, eq. \ref{eq:FcrystalBCDI3} - \ref{eq:FcrystalBCDI7}.

In our case, the direction of the [111] scattering vector is perpendicular to the sample, along the $\vec{z}$ axis of the laboratory frame.
Therefore, the [111] scattering vector, $\vec{G}_{111}$, is so forth described as $\vec{q}_z$, of magnitude  $|\vec{q}_z|$, with the direction described by the unit vector  $\hat{q_z}$, to be consistent with the equations detailed above.
Our approach to the strain is considerably simplified since we can only correctly derive one component of the strain tensor, $\epsilon_{zz}$ (eq. \ref{eq:StrainTensorzz}) as detailed below in eq. \ref{eq:StrainFromPhase1} - \ref{eq:StrainFromPhase3}.

\begin{align}
    \label{eq:StrainFromPhase1}
    & \Phi =  \vec{q}_z.\vec{u} \\
    \label{eq:StrainFromPhase2}
    & \frac{\Phi}{|\vec{q}_z|} = \frac{\vec{q}_z}{|\vec{q}_z|}.\vec{u} = \hat{q_z}.\vec{u} = \vec{u}_{\hat{q_z}} \\
    \label{eq:StrainFromPhase3}
    & \vec{\nabla} \vec{u}_{\hat{q_z}} = \frac{\partial u_{\hat{q_z}}}{\partial z} = \epsilon_{zz}
\end{align}

\subsubsection{Homogeneous strain}

The deviation of the interplanar spacing from the room temperature value due to the thermal expansion of the crystal is expected to be homogeneous within the particle.
This isotropic \textit{homogeneous} strain ($\epsilon_{hmg}$) in the particle is removed by centering the Bragg peak before phase retrieval, otherwise resulting in a linear phase ramp \parencite{}.
It is of utmost important to study the evolution of the material first under an inert atmosphere to see if there is an evolution of the homogeneous strain not only due to the thermal expansion of the crystal, but also to global relaxation phenomena induced e.g. by the adsorption of molecules involved in the catalytic reaction.

% \begin{figure}[!htb]
%     \centering
%     \includegraphics[width=\textwidth]{/home/david/Documents/PhDScripts/SixS\_2021\_01/Max.png}
%     \caption{
%         The maximum value of the 3D Bragg peak is used to center the intensity before phase retrieval and to compute the average interplanar spacing $d_{111}$.
%     }
%     \label{fig:MaxPeak}
% \end{figure}

\begin{figure}[!htb]
    \centering
    \includegraphics[width=\textwidth]{/home/david/Documents/PhDScripts/SixS\_2021\_01/RockingCurves.pdf}
    \includegraphics[width=\textwidth]{/home/david/Documents/PhDScripts/SixS\_2021\_01/HomoStrain.pdf}
    \caption{
        a) Evolution of the integrated scattered intensity in the detector as a function of the scattering angle during the rocking scans, fitted (b) with a lorentzian profile to retrieve the peak positions and full width at half maxima (FWHM).
        Evolution of the $d_{111}$ interplanar spacing (c), and associated homogenous strain values as a function of the temperature.
        The reference for the computation of homogeneous strain is taken at \qty{25}{\degreeCelsius}.
        Evolution of the Bragg peak FWHM as a function of the temperature (d).
        \textcolor{Important}{Errorbars}
        %(111) Bragg peak angular position (c),
    }
    \label{fig:HomoStrain}
\end{figure}

The average interplanar spacing between (111) crystallographic planes in the particle was computed from the angular position of the Bragg peak \textit{via} eq. \ref{eq:QandD3} (fig. \ref{fig:HomoStrain}).
The interplanar spacing follows approximately two linear increases from \qtyrange{25}{150}{\degreeCelsius} before the defect removal and from \qtyrange{150}{600}{\degreeCelsius} after the defect removal, which had the effect of relaxing the particle, increasing the average interplanar spacing and decreasing the peak FWHM, signature of the decrease of heterogeneous strain inside the particle.
The particle was realigned after the removal of the defect to be sure of the absence of any angular offset, the flash heating process could have moved the sample positions due thermal expansion.
The measurement at \qty{550}{\degreeCelsius} shows a very large FWHM and could be the start of the appearance of a new defect at the interface with the substrate, imaged at \qty{600}{\degreeCelsius} with BCDI.
The measurement at \qty{550}{\degreeCelsius} could not be reconstructed.
At \qty{600}{\degreeCelsius}, the two values of the average interplanar spacing under inert atmosphere are only slightly higher or lower than the value at \qty{550}{\degreeCelsius}, and decrease after the introduction of ammonia.
The FWHM decreases approximately one hour after having reached \qty{600}{\degreeCelsius} (each data point is separated by \qty{\approx30}{\min}) from its value at \qty{550}{\degreeCelsius}, returning to the previous values observed between \qty{150}{\degreeCelsius} and \qty{450}{\degreeCelsius}.
The introduction of ammonia has no visible effect on the FWHM.
The decrease of the interplanar spacing between \qty{550}{\degreeCelsius} and \qty{600}{\degreeCelsius} occurs together with the decrease of the FWHM, as if an equilibrium state has been reached.\textcolor{Important}{Why}

\subsubsection{Heterogeneous strain}

The remaining strain after phase retrieval is called the \textit{heterogeneous} strain ($\epsilon_{htg}$) \parencite{GREDIAC1996, FAVIER2007, Atlan2023}, the total strain observed during this experiment being equal to $\epsilon_{tot} = \epsilon_{hmg} + \epsilon_{zz, htg}$ when considering the changes from a reference (usually room temperature data).

In BCDI, it is the displacement of small unit blocks making up the crystal lattice that is observed.
Depending on the instrumental parameters, these small unit blocks, called \textit{voxels}, have a more or less large size.
In this study, they are approximately \qtyproduct{10x10x10}{\nm} large.
The strain is thus not directly related to the deviation between the interreticular planes, but to the gradient of the displacement field of these unit blocks from their equilibrium position.
The outermost layers of the crystal only constitute a low percentage of the surface voxels (\qty{\approx 11}{\percent} if we consider 5 atomic layers separated by the value of the interplanar spacing $d_{111}$), which lowers the contribution of the surface strain to the total strain contained in the surface voxels, reducing the ability to properly resolve e.g. surface relaxation effects.
The use of padding during the reconstruction algorithm decreases the voxel size, but without relying on the sampling of high-frequency components of the scattering amplitude and therefore does not increase the strain resolution.

% \begin{figure}[!htb]
%     \centering
%     \includegraphics[width=0.6\textwidth]{/home/david/Documents/PhDScripts/SixS\_2021\_01/FacetAnalyser/Angles.pdf}
%     \caption{
%         Angles between the $[111]$ direction (normal to the top facet and opposite to the bottom facet of the particle) and other crystallographic directions that describe the normals to the other facets of the particle.
%     }
%     \label{fig:Angles}
% \end{figure}

If equivalent orientation translates into equal surface atomic structures (e.g. [111] and [$\bar{1}\bar{1}\bar{1}$]), the environment of equivalent facets is not always the same.
For example, it is important to differ between [1$\bar{1}$1] and [$\bar{1}$1$\bar{1}$] facets, the higher the value of the interplanar angle, the closer the facet is to the interface with the substrate, and thus the more its influence in prominent.
Moreover, if the [111] facet is at the top of the crystal and the furthest away from the substrate, the [$\bar{1}\bar{1}\bar{1}$] facet is expected to be fully in contact with the substrate.

Therefore, the mean value and standard deviation of the heterogeneous strain distribution on each facet is presented on fig. \ref{fig:AmaterasuStrain}, as a function of the angle between the facets normals and the [111] direction, to highlight this difference.
The crystallographic orientation of a facet can be determined by computing the angle between the [111] direction and its normal, which is then related to a direction in real space (fig. \ref{fig:StereoBottom}).
Facets with the same orientation (e.g. [100], [010], [001]) are grouped around the same value of the interplanar angle.
To be more precise, the angular value is computed as the average of the angle between the facet normal and the [111] direction plus the angle between the facet normal and the [$\bar{1}\bar{1}\bar{1}$] direction minus \ang{180}, \textit{i.e.} $(\angle (111, \vec{n}) + (\angle (\bar{1}\bar{1}\bar{1}, \vec{n}) -180))/2$.

\begin{figure}[!htb]
    \centering
    \includegraphics[width=\textwidth]{/home/david/Documents/PhDScripts/SixS\_2021\_01/FacetAnalyser/FacetStrainEvolution.pdf}
    \caption{
        Mean value and standard deviation of the heterogeneous strain ($\epsilon_{zz, htg}$) distribution as a function of the angle between the normal of each facet on the particle surface and the $[111]$ direction.
        Upwards and downwards arrow are represented for respectively positive and negative strain.
    }
    \label{fig:AmaterasuStrain}
\end{figure}

% introduce in plane out of plane
It is important to realize that the displacement observed is \textit{only} in the [111] direction, sensitive to deviations of the crystal structure perpendicular to the [111] and [$\bar{1}\bar{1}\bar{1}$] facets, but parallel to [$\bar{1}$10]-type facets, themselves perpendicular to the [111] direction (fig. \ref{fig:StereoBottom}).
On one hand, if there existed an out-of-plane displacement of the atoms on [$\bar{1}$10]-type facets, \textit{i.e.} in the direction perpendicular to the [$\bar{1}$10] planes, its contribution to the displacement field observed in this experiment would not be directly visible.
On the other hand, the in-plane displacement of the atoms on [$\bar{1}$10]-type facets is visible, whereas the in-plane displacement on the [111] and [$\bar{1}\bar{1}\bar{1}$] facets is not visible.

The $\epsilon_{zz}$ component of the strain tensor on the [111] and [$\bar{1}\bar{1}\bar{1}$] facets is easily assimilated to variation of the interplanar spacing $d_{111}$ (positive strain is tensile strain, negative strain is compressive).
However, the physical meaning of the strain becomes more complex when observing facets that are neither parallel nor perpendicular to the [111] direction, such as [$1\bar{1}1$]-type or \{100\} facets.
For this reason, the following analysis of the heterogeneous strain is meant to be qualitative, by comapring the strain evolution of the same facets at different conditions but not between different facets, a quantitative analysis could only be performed with the full strain tensor.
% The relation between in-plane and out-of-plane strain can be rationalized by the Poisson effect \parencite{Atlan2023}, which relates an in-plane tensile (compressive) strain to an out-of-plane compressive (tensile) strain (depending on the origin of the displacement field).

The values of the strain are in general quite low, always below \qty{0.1}{\percent}, while the highest value of the displacement field on the particle facets is of about \qty{1}{\angstrom} (app. \ref{fig:AmaterasuDisplacement}).
It is difficult to be certain of the maximum value of the displacement due to some difficulties when unwrapping the phase for some measurements.
Let's take the example of the particle at \qty{600}{\degreeCelsius} under Argon flow, the magnitude of the scattering vector at the position of the Bragg peak is $|\vec{G}| \qty{\approx 2.772}{\angstrom}$.
The magnitude of the scattering vector probed furthest from the Bragg peak during the measurement is along the [111] crystal truncation rod, $|\vec{q}| \qty{\approx 2.742}{\angstrom}$.
We have $\delta q = \qty{0.03}{\angstrom}$ which gives $(\vec{q}-\vec{G}).\vec{u} = 0.03 <<1$, which puts us far away from the BCDI limit discussed in sec. \ref{sec:StrainBCDI}.
\textcolor{Important}{Add a figure with orthogonalized data}

% Decribe removal of defect strain evolution 150°C -> 150 °C
The very large standard deviations seen in fig. \ref{fig:AmaterasuStrain} at \qty{150}{\degreeCelsius} for the [$\bar{1}\bar{1}\bar{1}$], [$\bar{1}1\bar{1}$], [$1\bar{1}\bar{1}$] and [$11\bar{1}$] facets can be explained by the presence of the defect, around which the phase was not well unwrapped as seen in fig. \ref{fig:AmaterasuDefect}.
Defects are also expected to be at the origin of regions with large strain deviations from their impact on the displacement field, as observed for example in the middle of the [$\bar{1}\bar{1}\bar{1}$] facet in fig. \ref{fig:AmaterasuB} at \qty{600}{\degreeCelsius}.
At \qty{150}{\degreeCelsius}, after removal of the defect, all of the facets show low strain values as well as low standard deviation.

% 250, % 350, % 450
The strain values are relatively low and stable at \qty{250}{\degreeCelsius}, \qty{350}{\degreeCelsius}, and \qty{450}{\degreeCelsius}.

% but when comparing to the measurement at \qty{150}{\degreeCelsius}, we can see in fig. \ref{fig:AmaterasuB} that the extremities of the [$\bar{1}\bar{1}\bar{1}$] facet are in tension, whereas the center is in compression.
% The opposite effect seems to take place around the [111] facet, visible both in fig \ref{fig:Amaterasu} and \ref{fig:AmaterasuStrain}.
% \textcolor{Important}{Why ?}

% 600°C
% Decribe strain change between 450°C under Argon and 600°C under ammonia
The most interesting changes in the facet strain occur at \qty{600}{\degreeCelsius}.
The appearance of a new defect at the interface with the substrate has the effect of \textit{splitting} the particle in two states, with negative strain for the facets that are the closest to the interface and positive strain for the facets that are on the top of the particle.
This effect is translated into respectively tensile and compressive strain on the [111] and [$\bar{1}\bar{1}\bar{1}$] facets, and is reversed by the introduction of ammonia at \qty{600}{\degreeCelsius}, visible also in fig. \ref{fig:AmaterasuB} and \ref{fig:Amaterasu110}.
It is also during the presence of ammonia in the reactor that a [2$\bar{1}\bar{1}$] facet appears for the first time, perpendicular to the [111] direction.
Both facets perpendicular to the [111] direction, [2$\bar{1}\bar{1}$] and [$\bar{1}01$] have low strain values, showing no displacement of the atoms parallel to their surfaces.

\begin{figure}[!htb]
    \centering
    \includegraphics[width=0.32\textwidth]{/home/david/Documents/PhDScripts/SixS\_2021\_01/paraview/1675_strain_and_facets1.png}
    \includegraphics[width=0.32\textwidth]{/home/david/Documents/PhDScripts/SixS\_2021\_01/paraview/1675_strain_and_facets2.png}
    \includegraphics[width=0.32\textwidth]{/home/david/Documents/PhDScripts/SixS\_2021\_01/paraview/1675_strain_and_facets3.png}
    \caption{
        View of the Amaterasu particle at \qty{600}{\degreeCelsius} after the introduction of ammonia.
        The three \{100\} facets are surrounded by either 4 \{111\} facets and one \{110\} facets (left), by 3 \{111\} facets and one \{110\} facets (middle) or by 4 \{111\} facets (right).
    }
    \label{fig:AmaterasuStrain1675}
\end{figure}

The alternating highly positive and negative strain regions in the middle of the [$\bar{1}\bar{1}\bar{1}$] facet (fig. \ref{fig:AmaterasuB}) is the signature of a dislocation network forming in the interface with the substrate \parencite{Dupraz2015}, also responsible for the high strain standard deviation of this facet (fig. \ref{fig:AmaterasuStrain}).
A higher resolution in a 3D displacement field together with simulations of the impact of defect at the interface are needed to fully characterized this network, which is outside the scope of this thesis.

On fig. \ref{fig:AmaterasuStrain1675} we can observe that the three \{100\}-type facets have the same strain values, none of which have the same environment (fig. \ref{fig:AmaterasuStrain1675}), while the edges and corners around the [111] and \{100\} facet show a slightly positive strain.

At first sight, all facets that have a similar orientation show a similar strain (e.g. [1$\bar{1}$1]-type, \{100\}-type).
However, the two \{111\} facets that are close to the substrate do not show the same strain than the three \{111\} facets closer to the top of the particle.
The strain being in the [111] direction with the origin of the displacement field in the center of the particle, if all of these facet had experienced the same compressive (tensile) strain perpendicular to their surface \textit{via} the adsorption of ammonia, they would all show a negative (positive) strain (of lower amplitude due to the angle between these facets and the [111] direction).

This behaviour could be linked to a strong effect of the support which has the dual effect of first preventing the [$\bar{1}\bar{1}\bar{1}$] facet from being exposed and secondly forcing them to accomodate the strain linked to the substrate.
There are three possibilities, first ammonia is effectively adsorbed on these facets but the influence of the substrate hides a possible common signature in the strain.
Secondly ammonia is not adsorbed and the strain difference is due to the influence of the substrate on the particle structure.
Thirdly, ammonia is adsorbed but only on the three top facets or only on the 2 bottom facets due to the influence of the substrate that either limitates or facilitates the adsorption process.

% Globally, the top of the particle seems to be in compression along $\vec{z}$ (fig. \ref{fig:Amaterasu}, \ref{fig:AmaterasuStrainSlices}), whereas the bottom of the particle is in tension, with the exception of the one [$1\bar{1}3$] facet (fig. \ref{fig:AmaterasuStrainSlices}) which is in compression.

% \begin{figure}[!htb]
%     \centering
%     \includegraphics[width=0.49\textwidth]{/home/david/Documents/PhDScripts/SixS\_2021\_01/paraview/1675_clipx.png}
%     \includegraphics[width=0.49\textwidth]{/home/david/Documents/PhDScripts/SixS\_2021\_01/paraview/1675_clipy.png}
%     \caption{
%         Particle slices perpendicular to the $\vec{x}$ and $\vec{y}$ directions at center of mass, a white line delimitates the particle surface.
%     }
%     \label{fig:AmaterasuStrainSlices}
% \end{figure}

Overall, it is difficult to conclude on any potential effect of the absorption of ammonia on the particle due to the presence of the dislocation network at the interface for which it is unclear whether or not it's origin if due to the adsorption of ammonia or from interfacial strain.
The fact that we only possess one component of the strain tensor, makes it difficult to quantify any relaxation effect on the facets that are not perpendicular to the [111] direction.
Bragg peaks from the substrate could have been measured to have a better idea of the interfacial structure.
Additionnal simulations of the particle equilibrium shape and the impact if interfacial defects are needed to have a better understanding of the structural dynamics at work.

% The shape of the particle in the ($\vec{x}, \vec{y}$) plane at \qty{600}{\degreeCelsius} can be roughly approximated as a hexagon with 3 sides occupied by the [$\bar{1}11$], [$1\bar{1}1$] and [$11\bar{1}$] facets.
% The three other sides are occupied by the [100], [010], [001] facets near the top of the particle, under which are situated either 1) one large [$\bar{1}10$] facet, 2) one small [$1\bar{1}0$] facet under which in turn are two smalls [113] and [$\bar{1}1\bar{1}$] facets, 3) a large [$1\bar{1}\bar{1}$] facet.

%%
% The negative strain of the [$1\bar{1}1$] facets and the positive strain on the [$\bar{1}1\bar{1}$] facets (fig. \ref{fig:AmaterasuFacetsEvolution}) follows the respective evolution of the [111] and [$\bar{1}\bar{1}\bar{1}$] facets with a lower amplitude.
% An increase of the negative strain can also be seen on all three [100] facets but, based on the angle between those facets and the [111] direction (\ang{54.7}), the interpretation of this value is difficult to understand.

    \subsection{Particle structure evolution during ammonia oxidation at \qty{300}{\degreeCelsius} and \qty{400}{\degreeCelsius}}\label{sec:BCDIAmmoniaOxidation}

Particle \textit{A} being unstable on the substrate during the oxidation of ammonia at \qty{600}{\degreeCelsius}, the reaction was measured at lower temperatures before and after the catalyst light off, respectively \qty{300}{\degreeCelsius}, and \qty{400}{\degreeCelsius}.
Two particles, namely \textit{B} and \textit{C}, were identified on the substrate, first measured at room temperature under inert argon atmosphere.

The support of particle \textit{B} (fig. \ref{fig:B7Facets}) was well determined at room temperature, exhibiting a round shape with a total of 30 facets.
Approximately \qty{76}{\percent} of the particle surface voxels were recognised to belong to facets.

\begin{figure}[!htb]
    \centering
    \includegraphics[width=0.32\textwidth]{/home/david/Documents/PhDScripts/SixS_2021_06_BCDI_NH3/reconstructions/vti_good_files/B7/RT/RT_100_B7.png}
    \includegraphics[width=0.32\textwidth]{/home/david/Documents/PhDScripts/SixS_2021_06_BCDI_NH3/reconstructions/vti_good_files/B7/RT/RT_010_B7.png}
    \includegraphics[width=0.32\textwidth]{/home/david/Documents/PhDScripts/SixS_2021_06_BCDI_NH3/reconstructions/vti_good_files/B7/RT/RT_001_B7.png}
    \caption{
        Surface of the reconstructed particle B at \qty{25}{\degreeCelsius} under inert argon atmosphere.
        The surface is coloured by the values of the heterogeneous out-of-plane strain as described in eq. \ref{eq:StrainTensorzz} at each surface voxel, the limit of each facet is delimited by thick white lines.
    }
    \label{fig:B7Facets}
\end{figure}

About \qty{10.5}{\percent} of the particle surface is covered by \{100\} facets, \qty{23.1}{\percent} by \{110\} facets, \qty{35.2}{\percent} by \{111\} facets, and \qty{7.3}{\percent} by \{113\} facets.
The particle has a round shape, with a diameter of \qty{300}{\nm}.

From the 3D representation of the particle in fig. \ref{fig:B7Facets}, the existence of the ($3\bar{1}1$) and ($\bar{1}13$) facets seem to be highly probable when considering the symmetry and shape of the particle.
Those facets could not be detected by the recognition algorithm, probably due to a lack of spatial resolution, necessary to detect the smallest facets on the particle.

In contrast with particle \textit{B}, the support of particle \textit{C} could not be determined at room temperature, parts of the particle systematically missing from the reconstructions.
No dislocation could be identified at the interface with the substrate but rather a lack of convergence of the support.
Heating the sample to \qty{300}{\degreeCelsius} under Argon did not improve the quality of the reconstruction.
Interestingly, the introduction of ammonia in the reactor at the beginning of the oxidation cycle had the effect of lowering the average strain inside the particle (discussed further), and allowed the convergence of the support.
The particle \textit{C}, shown in fig. \ref{fig:D6Facets}, was thus first reconstructed at \qty{300}{\degreeCelsius} under \qty{294}{\milli\bar} of argon and \qty{6}{\milli\bar} of ammonia in the reactor.

\begin{figure}[!htb]
    \centering
    \includegraphics[width=0.32\textwidth]{/home/david/Documents/PhDScripts/SixS_2021_06_BCDI_NH3/reconstructions/vti_good_files/D6/300/300_ammonia_100.png}
    \includegraphics[width=0.32\textwidth]{/home/david/Documents/PhDScripts/SixS_2021_06_BCDI_NH3/reconstructions/vti_good_files/D6/300/300_ammonia_010.png}
    \includegraphics[width=0.32\textwidth]{/home/david/Documents/PhDScripts/SixS_2021_06_BCDI_NH3/reconstructions/vti_good_files/D6/300/300_ammonia_001.png}
    \caption{
        Surface of the reconstructed particle C at \qty{300}{\degreeCelsius} under a partial pressure of ammonia equal to \qty{10}{\milli\bar}.
        The surface is coloured by the values of the heterogeneous out-of-plane strain as described in eq. \ref{eq:StrainTensorzz} at each surface voxel, the limit of each facet is delimited by thick white lines.
    }
    \label{fig:D6Facets}
\end{figure}

Particle \textit{C} exhibits a more rectangular shape without any \{113\} facets, \qty{10.8}{\percent} of its surface is covered by six \{100\} facets, \qty{6.4}{\percent} by \{110\} facets, and \qty{36.3}{\percent} by \{111\} facets.
The particle is at the largest \qty{800}{\nm} wide.
Only \qty{53.4}{\percent} of the surface voxels were determined to belong to facets, there is for example a large area between the (111) and ($1\bar{1}1$) facets (fig \ref{fig:D6Facets} - left) that could correspond to the missing (101) facet.
Compared to particle \textit{B}, the area between the (100)-type facets and ($\bar{1}1\bar{1}$)-type facets is clearly determined, and not expected to contains higher index facets.

The two nanoparticles have a very different shape and facet coverage, resumed in tab. \ref{tab:FacetCoverage}.
The percentage of the total surface area covered by \{100\} and \{111\} facets is similar (respectively about \qty{10}{\percent} and \qty{36}{\percent}).
The largest difference in the \{110\} facet coverage can be partially explained by the non-detection of the (101) facet.
Most importantly, the large particle does not exhibit high Miller index facets.

\begin{table}[!htb]
    \centering
    \begin{tabular}{@{}llllll@{}}
    \toprule
        & \{100\} & \{110\} & \{111\} & \{113\} & Undefined \\
    \midrule
    Particle B (\qty{300}{\nm} wide) & \qty{10.5}{\percent} (6) & \qty{23.1}{\percent} (9) & \qty{35.2}{\percent} (8) & \qty{7.3}{\percent} (7) & \qty{23.9}{\percent} \\
    Particle C (\qty{800}{\nm} wide) & \qty{10.8}{\percent} (6) &  \qty{6.4}{\percent} (8) & \qty{36.3}{\percent} (8) & \qty{0}{\percent}       & \qty{46.5}{\percent} \\
    \bottomrule
    \end{tabular}%
    \caption{
        Particle surface area covered by facet type, the number in parenthesis designates the amount of each facet identified on the surface.
    }
    \label{tab:FacetCoverage}
\end{table}

% particule a et rajouter conditions

Together with particle \textit{A}, this results shows that the sample used during this experiment is covered with many different nanoparticles that all together contribute to the catalytic activity, and probably show an even greater variance of shape and size.
The current analysis using BCDI aim to reveal the structure variation during a catalytic reaction of a few selected nanoparticles.
Statistically representing the population on the sample is impossible, the measurement process being currently too time consuming to probe more than a few nanoparticles.

The ammonia oxidation cycle was started at \qty{300}{\degreeCelsius}, with the introduction of ammonia.
The different atmospheres used to probe the relation between particle structure and selectivity have been introduced in tab. \ref{tab:ConditionsNanoparticles}.
A maximum of three scans were recorded under each condition (tab. \ref{tab:ConditionsNanoparticles}) at \qty{300}{\degreeCelsius} and \qty{400}{\degreeCelsius} for each particle to probe for any evolution of the particles surface structure during a fixed atmosphere as a function of time, and to otherwise demonstrate the reproducibility of the measurements.

A 3D view of the reconstructed nanoparticles at \qty{300}{\degreeCelsius} and \qty{400}{\degreeCelsius} during the ammonia oxidation cycle is presented in appendix \ref{app:AppendixB7}.
The surface is coloured by the values of the heterogeneous out-of-plane strain as described previously in eq. \ref{eq:StrainTensorzz} at each surface voxel.
The view perpendicular to the three axes of the laboratory frame is represented, the particles are tilted due to the incoming angle $\theta$ for the measurement of the (111) Bragg peak.
Both the retrieved displacement and strain fields are represented with their respective colour-bars, the ammonia to oxygen ratio is represented on the right part of the figure.

\subsubsection{Particle B}

A different behaviour during the oxidation of ammonia was observed on both nanoparticles.
Particle \textit{B}, which is smaller and thus scatters less than particle \textit{C}, was difficult to reconstruct once the temperature increased to \qty{300}{\degreeCelsius}.

\begin{figure}[!htb]
    \centering
    \includegraphics[width=\textwidth]{/home/david/Documents/PhD/Figures/bcdi_data/B7/hist_B7.pdf}
    \caption{
        Histograms of normalised Bragg electronic density amplitude, for particle \textit{B}.
        Values below 0.05 are ignored since they typically belong to voxels located far away from particle support.
        The difference between bulk, surface, and voxels outside the support is smeared out in b).
    }
    \label{fig:B7Histo}
\end{figure}

Normalised histograms of the particle reconstructed Bragg electronic density showed no more clear distinctions between regions inside and outside the support (fig. \ref{fig:B7Histo}), and choosing a threshold value to draw the particle iso-surface no longer became evident.

Choosing a low threshold for the particle iso-surface could still be done, but resulted in a round object without clear facets.
The small volume of the particle also probably impacts the quality of the measurements, increasing the spatial resolution and leading to this \textit{rounding} of the particle.
Some reconstructions showed a good quality at \qty{300}{\degreeCelsius}, but in a larger sense recovering the facets on the particle surface in a reproducible way became impossible.

To be able to quantify the strain inside the particle, the 3D diffraction patterns were orthogonalised in $q$-space (fig. \ref{fig:B7Ortho}).
Increasing the temperature to \qty{400}{\degreeCelsius} under inert atmosphere was at the origin of the appearance of a defect in the particle, visible directly in reciprocal space by the shape of the 3D diffraction patterns, split in two parts (fig. \ref{fig:B7Ortho}, transition from a to b), which did not improve the phase retrieval quality.

\begin{figure}[!htb]
    \centering
    % \includegraphics[width=\textwidth]{/home/david/Documents/PhD/Figures/bcdi_data/B7/B7_orthogonalized_3479.pdf}
    % \includegraphics[width=\textwidth]{/home/david/Documents/PhD/Figures/bcdi_data/B7/B7_orthogonalized_3544.pdf}
    \includegraphics[width=\textwidth]{/home/david/Documents/PhD/Figures/bcdi_data/B7/B7_orthogonalized_3804_edited.pdf}
    \includegraphics[width=\textwidth]{/home/david/Documents/PhD/Figures/bcdi_data/B7/B7_orthogonalized_3863_edited.pdf}
    \includegraphics[width=\textwidth]{/home/david/Documents/PhD/Figures/bcdi_data/B7/B7_orthogonalized_4098_edited.pdf}
    % \includegraphics[width=\textwidth]{/home/david/Documents/PhD/Figures/bcdi_data/B7/B7_orthogonalized_4110.pdf}
    \caption{
        Sum of the intensity scattered from particle \textit{B} around the (111) platinum Bragg peak in the $\vec{q}_x$, $\vec{q}_z$ and $\vec{q}_y$ directions after orthogonalisation of the diffraction pattern.
        The [111]-oriented crystal truncation rod is visible parallel to the $\vec{q}_z$ axis.
        Detector gaps lead to reactangular areas of missing intensity.
    }
    \label{fig:B7Ortho}
\end{figure}

A signal parallel to the crystal truncation rod in the [111] direction is observed slightly shifted in $\vec{q}_y$, interfringes in the [111] direction are also visible with a slightly lower interval than on the original CTR.
The distance between the [111] oriented rod and the second signal increases in $q_y$ during the ammonia oxidation cycle, at the end of which a third signal parallel to the two others can be observed (fig. \ref{fig:B7Ortho} - c).
The maximum of both signals seems to be at the same value in $q_z$, which can be related to the average out-of-plane lattice parameter inside the particle.

The defects are also visible perpendicularly to the particle substrate and through the particle in 3D in appendix \ref{app:AppendixB7}, ultimately resulting in holes in the Bragg electronic density \parencite{Clark2015, Dupraz2015}.

The intensity scattered from particle \textit{B} around the (111) platinum Bragg peak was summed along the ($\vec{q}_x$, $\vec{q}_y$), ($\vec{q}_y$, $\vec{q}_z$) and ($\vec{q}_z$, $\vec{q}_x$) planes, and respectively fitted in the $\vec{q}_z$, $\vec{q}_x$ and $\vec{q}_y$ directions.
($\vec{x}$, $\vec{y}$, $\vec{z}$) designates the sample frame, where the $\vec{z}$ direction is parallel to the [111] direction from the epitaxy relationship of the particle on the substrate.
The sum in the planes perpendicular to the $\vec{q}_x$ and $\vec{q}_y$ directions carry information about the in-plane particle shape and strain, whereas the sum perpendicular to the $\vec{q}_z$ direction carries information about the out-of-plane particle shape and strain.

The summed intensities were then fitted in each direction by a Lorentzian shape function to extract the Bragg peak full width at half maximum (FWHM) in each direction: FWHM$_x$, FWHM$_y$, FWHM$_z$.
The average interplanar spacing was also computed from the peak position in $q$-space.
Once the second signal appeared in $\vec{q}_y$, two Lorentzian profiles were used to fit both peaks, FWHM$_y$ in fig. \ref{fig:B7FWHM} is equal to the sum of FWHM$_y$ from each peak.

A typical fit is shown in appendix \ref{fig:FitB73D}, the retrieved FWHM are shown in fig. \ref{fig:B7FWHM}, the average interplanar spacing and homogeneous strain are shown in fig. \ref{fig:B7Latpara}.

\begin{figure}[!htb]
    \centering
    \includegraphics[width=\textwidth]{/home/david/Documents/PhD/Figures/bcdi_data/B7/fwhm_B7_two_peaks.pdf}
    \caption{
        Evolution of the full width at half maximum (FWHM) for particle B in the $\vec{q}_x$, $\vec{q}_y$ and $\vec{q}_z$ directions at \qty{25}{\degreeCelsius}, \qty{300}{\degreeCelsius} and \qty{400}{\degreeCelsius} as a function of the ammonia to oxygen ratio.
        The reproduction of measurements at fixed conditions yields multiple data points, lines at \qty{300}{\degreeCelsius} and \qty{400}{\degreeCelsius} link the data points in the order of the measurements.
    }
    \label{fig:B7FWHM}
\end{figure}

A decrease of FWHM$_z$ can be observed between successive measurements under argon after heating the sample to \qty{300}{\degreeCelsius}, whereas the values of FWHM$_x$ and FWHM$_y$ are stable, which could be due to the stabilisation of the particle after the temperature increase.
Introducing ammonia in the reactor has the effect of increasing FWHM$_x$ and FWHM$_y$ without increasing FWHM$_z$, thus probably related to an increase of the in-plane strain in the particle.
A decrease of the FWHM can be seen in all three directions when oxygen is introduced in the reactor, which may be associated to a global decrease of the particle strain.
The range of values in FWHM$_z$ is very low and if a slight increase is observed as a function of the ammonia to oxygen ratio, it is only for very small values.
No clear evolution can be observed in FWHM$_x$ or FWHM$_y$.

On one hand, the values of FWHM$_z$ at \qty{400}{\degreeCelsius} are stable until the ammonia to oxygen ratio becomes equal to one, which can be linked to an increase of the out-of-plane strain in the particle.
Higher amounts of oxygen further continue this increase.

A large jump is observed for FWHM$_y$ after the increase of temperature under inert atmosphere, being multiplied by \qty{50}{\percent} and \qty{100}{\percent} in the first and second measurements.
No jump is observed for FWHM$_x$ in the first measurement, an increase is observed in the second measurement, but not as important as for FWHM$_y$.
This coincides with the appearance of a new signal observed in fig. \ref{fig:B7Ortho} perpendicular to $\vec{q}_z$, \textit{i.e.} in the ($\vec{q}_x, \vec{q}_y$) plane, mostly shifted in $q_y$ in comparison to the centre of the Bragg peak.
This signal is linked to the appearance of a defect in the particle.

The introduction of ammonia does not have a visible effect in $\vec{q}_z$, but decreases the FWHM in both $\vec{q}_x$ and $\vec{q}_y$ in comparison to the second measurement under argon.
Both FWHM$_x$ and FWHM$_y$ increase together with the ammonia to oxygen ratio, FWHM$_y$ plateaus when the ratio is equal to 1, whereas FWHM$_x$ largely increases when the ratio becomes equal to 8.

The FWHM under argon after and before the oxidation cycle are similar at \qty{300}{\degreeCelsius} but not at \qty{400}{\degreeCelsius}, a temperature at which the oxidation cycle completely changed the shape and strain state of the particle, visible in app. \ref{app:AppendixB7}.
The particle shape is not reversible at \qty{400}{\degreeCelsius} after the oxidation cycle.

\begin{SCfigure}
    \centering
    \includegraphics[width=0.6\textwidth]{/home/david/Documents/PhD/Figures/bcdi_data/B7/lattice_parameter_B7.pdf}
    \caption{
        Evolution of the interplanar spacing $d_{111}$ and homogeneous strain $\epsilon_{111, homo}$ for particle B as a function of the ammonia to oxygen ratio.
        The reference for the computation of $\epsilon_{111, homo}$ was taken as the mean value at \qty{25}{\degreeCelsius}.
        The reproduction of measurements at fixed conditions yields multiple data points, lines at \qty{300}{\degreeCelsius} and \qty{400}{\degreeCelsius} link the data points in the order of the measurements.
    }
    \label{fig:B7Latpara}
\end{SCfigure}

Two main phenomena of interest must be kept in mind from the evolution of the interplanar spacing in fig. \ref{fig:B7Latpara}, besides the expected increase with temperature, verified in the previous section.
First there is a \qty{0.08}{\percent} increase of the homogeneous strain following the introduction of ammonia at \qty{300}{\degreeCelsius}.
The interplanar spacing stays approximately equal to the same value afterwards during the ammonia oxidation cycle.
Secondly, there is an increase of about \qty{0.02}{\percent} of the homogeneous strain at \qty{400}{\degreeCelsius} when the ammonia to oxygen ratio reaches the value of one.
Both of these phenomena are observed at conditions for which changes in the FWHM are also seen.

Finally, to be able to give a first estimation of the amount of strain inside the particle, the strain field energy defined as the sum of the product of stress and strain integrated over the particle \parencite{Cahn1959} was computed following the work performed by Ulvestad et al \parencite*{Ulvestad2015a} and by Kim et al. \parencite*{Kim2019}.

Assuming cubic symmetry and isotropic shear-free conditions within the unit cell, the strain field energy is simplified to eq. \ref{eq:SFE}.
Here, $G$ and $L$ represent the material's Lame constants, which can be estimated from Young's modulus $E$ and Poisson ratio $\nu$, as described in eq. \ref{eq:GandNu} \parencite{Mavko2020}.

\begin{equation}
    E_s = \frac{2G + 3L}{2} \int \Big( \frac{\partial u_{111}}{\partial x_{111}}\Big)^2 dV
    \label{eq:SFE}
\end{equation}

\begin{align}
    G = \frac{E}{2(1+\nu)} \qquad & \qquad L = \frac{E \nu}{(1+\nu)(1-2\nu)}
    \label{eq:GandNu}
\end{align}

The evolution of the strain field energy as a function of each scan to be able to compare with the 3D reconstructions of the particle shown in app. \ref{app:AppendixB7} is presented in fig. \ref{fig:B7SFE}.
For a qualitative approximation of the strain field energy, the Young's modulus was taken equal to \qty{153.3}{\giga\pascal} and the Poisson ratio to \num{0.401} based on literature values at \qty{400}{\degreeCelsius} \parencite{Matthey2022}.

\begin{figure}[!htb]
    \centering
    \includegraphics[width=\textwidth]{/home/david/Documents/PhD/Figures/bcdi_data/B7/strain_energy_B7.pdf}
    \caption{
        Evolution of the strain field energy for particle B as a function of the scan number.
        The evolution of the \ce{NH_3} and \ce{O_2} gas flows in the reactor are also plotted to highlight the effect of the simultaneous presence of ammonia and oxygen and of the increasing amount of oxygen on the values of the strain field energy.
    }
    \label{fig:B7SFE}
\end{figure}

The strain field energy shows very little evolution at \qty{300}{\degreeCelsius}, decreasing a little under the presence of ammonia and further under reacting conditions.
However, the values of the strain field energy are observed to increase under reacting conditions at \qty{400}{\degreeCelsius}, multiplied approximately by \num{3} when the \ce{O_2}/\ce{NH_3} ratio is equal to 1, and by \num{9} when equal to 8.
This result is consistent with the evolution of the different FWHM in fig. \ref{fig:B7FWHM} and with the 3D shape of the particle.

To summarise, the introduction of ammonia at \qty{300}{\degreeCelsius} has resulted in (i) a slight increase of the in-plane strain, observed from the evolution of the peak FWHM in the $\vec{q}_x$ and $\vec{q}_y$ directions in fig. \ref{fig:B7FWHM}; in (ii) a large increase of the homogeneous strain, observed from the value of the $d_{111}$ interplanar spacing in fig. \ref{fig:B7Latpara}, which was not reversed after the removal of ammonia following the end of the oxidation cycle.
These changes can be linked to the first reduction of the particle, previously only exposed to air and argon, since such changes are not observed at \qty{400}{\degreeCelsius}.
No changes in the particle strain and structure as a function of the ammonia to oxygen ratio could be detected at \qty{300}{\degreeCelsius}.

The most important changes in the particle structure occur at \qty{400}{\degreeCelsius}.
Heating the sample under inert atmosphere has resulted in the appearance of a defect visible directly in reciprocal space (fig. \ref{fig:B7Ortho}), and by the evolution of FWHM$_y$ (fig. \ref{fig:B7FWHM}).
The increase of strain inside the particle is also visible in 3D by the progressive importance of holes in the reconstructed volume (app. \ref{app:AppendixB7}), and by the evolution of the strain field energy (fig. \ref{fig:B7SFE}).
Interestingly, a transition is observed during the simultaneous presence of both ammonia and oxygen in the reactor, when \ce{O_2}/\ce{NH_3} = 1 (scans 3951, 3958 and 4002).
This transition is also visible in FWHM$_z$ (fig. \ref{fig:B7FWHM}), which can be related to an increase of the out-of-plane strain inside the particle.

\subsubsection{Particle C}

As mentioned earlier, particle \textit{C} could only be reconstructed under the presence of ammonia in the reactor.
Interestingly, no changes in the particle shape or surface strain values have been observed during the exposition to the reactants (appendix \ref{app:AppendixB7}).

\begin{figure}[!htb]
    \centering
    \includegraphics[width=\textwidth]{/home/david/Documents/PhD/Figures/bcdi_data/D6/D6_orthogonalized_3517_edited.pdf}
    \includegraphics[width=\textwidth]{/home/david/Documents/PhD/Figures/bcdi_data/D6/D6_orthogonalized_3572_edited.pdf}
    \includegraphics[width=\textwidth]{/home/david/Documents/PhD/Figures/bcdi_data/D6/D6_orthogonalized_3823_edited.pdf}
    \caption{
        Sum of the intensity scattered from particle $C$ around the (111) platinum Bragg peak in the $\vec{q}_x$, $\vec{q}_z$ and $\vec{q}_y$ directions after orthogonalisation of the diffraction pattern.
        The [111]-oriented crystal truncation rod is visible parallel to the $\vec{q}_z$ axis.
        Detector gaps lead to rectangular areas of missing intensity.
    }
    \label{fig:D6Ortho}
\end{figure}

When removing ammonia and oxygen at \qty{300}{\degreeCelsius}, the particle could not be reconstructed anymore, reproducing the same behaviour as before the introduction of ammonia.
Increasing the temperature to \qty{400}{\degreeCelsius} did not improve the reconstruction quality as the support did not converge during phase retrieval.
Using the support of the particle reconstructed at \qty{300}{\degreeCelsius} under ammonia as a starting guess for the algorithms did not suffice to reach a good final solution.
The reconstructions under inert atmosphere before introducing ammonia or after removing both reactants from the cell presented in appendix \ref{app:AppendixB7} have been performed by limiting the support update to every 200 iteration (for a total of 1100 algorithm iteration during phase retrieval, meaning that the support is only changed 5 times), and still present a highly heterogeneous Bragg electronic density.

To understand why it was impossible to successfully carry out the phase retrieval process, the diffraction pattern around the (111) Bragg peak under the presence of both reacting gases and after their removal was orthogonalised in $q$-space, presented in fig. \ref{fig:D6Ortho}.

The diffraction pattern before the introduction of ammonia (fig. \ref{fig:D6Ortho} - a) exhibits a slightly distorted shape along $\vec{q}_z$, which can be related to the presence of heterogeneous strain in the particle, typically at the interface with the substrate.
The introduction of ammonia in the reactor has removed this distortion (fig. \ref{fig:D6Ortho} - b), which generally improves the quality of phase retrieval.

\begin{figure}[!hbt]
    \centering
    \includegraphics[width=\textwidth]{/home/david/Documents/PhD/Figures/bcdi_data/D6/fwhm_D6.pdf}
    \caption{
        Evolution of the full width at half maximum (FWHM) for particle C in the $\vec{q}_x$, $\vec{q}_y$ and $\vec{q}_z$ directions at \qty{25}{\degreeCelsius}, \qty{300}{\degreeCelsius} and \qty{400}{\degreeCelsius} as a function of the ammonia to oxygen ratio.
        The reproduction of measurements at fixed conditions yields multiple data points, lines at \qty{300}{\degreeCelsius} and \qty{400}{\degreeCelsius} link the data points in the order of the measurements.
        No scans were performed at \qty{400}{\degreeCelsius} under Argon after the oxidation cycle.
    }
    \label{fig:D6FWHM}
\end{figure}

Interestingly, the diffraction pattern after the removal of both gases in the reactor (fig. \ref{fig:D6Ortho} - c) is not similar to the measurement before the oxidation cycle (fig. \ref{fig:D6Ortho} - b), both diffraction patterns are collected under argon.
In general, the fringes are less visible in all direction, especially in the ($\vec{q}_x$, $\vec{q}_y$) plane, areas of higher intensity also appeared around the Bragg peak.
Most importantly, the coherence fringes situated in the [111] direction, \textit{i.e.} along $\vec{q}_z$ show a completely different pattern, as if two CTRs were on top of each other, one with larger rectangular fringes, while the other one has smaller round fringes.
The particle clearly suffered a change of shape after the removal of the reacting gases, meaning that different phenomena before and after the oxidation cycle impinge on the reconstruction quality.
It is possible that the first exposition to ammonia has the effect of reducing the particle that was beforehand only exposed to air and argon.
Thus, removing ammonia and oxygen at the end of the cycle can yield a different state than before the oxidation cycle.

The same kind of diffraction patterns are observed at \qty{400}{\degreeCelsius} after heating under argon as in fig. \ref{fig:D6Ortho} (c), and after the introduction of ammonia as in fig. \ref{fig:D6Ortho} (b), meaning that the presence of \ce{NH_3} in the reactor lowers the strain in the particle, which can then easily be reconstructed.

The particle FWHM in all three directions follows the same evolution as a function of the different atmospheres, meaning that the same phenomena occurs in all direction of the particle.
The FWHM decreases following the introduction of ammonia both at \qty{300}{\degreeCelsius} and \qty{400}{\degreeCelsius}, and only slightly increases during the exposition to oxygen.
Returning to argon atmosphere, \textit{i.e.} removing ammonia and oxygen yields FWHM values similar to those measured after heating to \qty{300}{\degreeCelsius} under argon.

\begin{SCfigure}
    \centering
    \includegraphics[width=0.6\textwidth]{/home/david/Documents/PhD/Figures/bcdi_data/D6/lattice_parameter_D6.pdf}
    \caption{
        Evolution of the interplanar spacing $d_{111}$ and homogeneous strain $\epsilon_{111}$ for particle C as a function of the ammonia to oxygen ratio.
        The reference for the computation of $\epsilon_{111}$ was taken at \qty{25}{\degreeCelsius}.
        The reproduction of measurements at fixed conditions yields multiple data points, lines at \qty{300}{\degreeCelsius} and \qty{400}{\degreeCelsius} link the data points in the order of the measurements.
        No scans were performed at \qty{400}{\degreeCelsius} under Argon after the oxidation cycle.
    }
    \label{fig:D6Latpara}
\end{SCfigure}

The evolution of the $d_{111}$ interplanar spacing (fig. \ref{fig:D6Latpara}) is not the same as a function of different temperatures.
The interplanar spacing of the particle under the sole presence of ammonia at \qty{300}{\degreeCelsius} is remarkably reproducible with all three measurements yielding approximately the same value as under argon atmosphere.
A large increase of the interplanar spacing is observed when oxygen is introduced in the reactor, at a ratio of \num{0.5} with ammonia, which gives a homogeneous strain evolution of approximately \qty{0.07}{\percent} in reference to the value measured at room temperature.
The strain at the interface with the substrate also becomes more homogeneous at this condition, slightly negative (appendix \ref{app:AppendixD6}).
The interplanar spacing then stays approximately the same at \qty{300}{\degreeCelsius} during the oxidation cycle.
This transition is not observed at \qty{400}{\degreeCelsius}, a temperature for which an increase of the interplanar spacing can be observed under argon, without any large deviations during the oxidation cycle.

\begin{figure}[!htb]
    \centering
    \includegraphics[width=\textwidth]{/home/david/Documents/PhD/Figures/bcdi_data/D6/strain_energy_D6.pdf}
    \caption{
        Evolution of the strain field energy for particle C as a function of the scan number.
        The evolution of the \ce{NH_3} and \ce{O_2} gas flows in the reactor are also plotted to highlight the effect of the presence of ammonia on the values of the strain field energy.
    }
    \label{fig:D6SFE}
\end{figure}

The strain field energy was also computed for particle \textit{C}, transitions between large and low magnitudes can be linked to the presence of ammonia in the reactor (fig. \ref{fig:D6SFE}).
No measurement was performed at \qty{400}{\degreeCelsius} at the end of the oxidation cycle which prevents us from being certain that the decrease in the strain field energy observed at \qty{300}{\degreeCelsius} could be reproduced, similar as for the evolution of the FWHM in fig. \ref{fig:D6FWHM}.

To summarise, particle \textit{C} seems to have suffered a transition at \qty{300}{\degreeCelsius} following the introduction of ammonia in the reactor, which helped improved the quality of phase retrieval, probably from a decrease of strain.
The diffraction pattern was shown to not be exactly reproducible when returning to inert atmosphere, but still showing a decrease of sharpness regarding the position and intensity of the interfringes.
% It is difficult to be certain of the impact of the adsorption of ammonia since there is a lack of good quality reconstruction for comparative studies before the introduction of ammonia in the reactor, preventing an analysis of the facet strain.

The transition can be observed in both the strain field energy (fig. \ref{fig:D6SFE}), and in the FWHM of the three-dimensional Bragg peak (fig. \ref{fig:D6FWHM}).
The interplanar spacing computed from the position of the Bragg peak in $q$-space shows an additional transition when first introducing oxygen in the cell, not visible in the strain field energy or in the FWHMs, which can tentatively be linked to a change of strain at the interface with the substrate.
% The influence of the introduction of oxygen on the interfacial strain is unclear since no clear and reproducible evolution of the facet strain is visible on the particle surface (app. \ref{app:AppendixD6}).
% The values of the heterogeneous surface strain $\epsilon_{zz}$ being very low (almost always below \qty{0.1}{\percent}, fig. \ref{fig:D6Facets}), it is also complicated to quantify the strain evolution.

\subsubsection{Mass spectrometry results}

The catalytic activity of the particles was recorded \textit{via} the use of a mass spectrometer and shows the production of nitrogen, nitrous oxide and nitrogen oxide (fig. \ref{fig:RGANanoparticlesBCDIComparison}), proving the activity of the catalyst.
The original mass spectrometer signal as a function of time is available in the appendix \ref{fig:RGA300BCDINanoparticles}, and \ref{fig:RGA400BCDINanoparticles}.

\begin{figure}[!htb]
    \centering
    \includegraphics[width=\textwidth]{/home/david/Documents/PhDScripts/SixS_2021_06_BCDI_NH3/gas_analysis/figures/product_comparison_no_norm.pdf}
    \caption{
        Evolution of reaction products partial pressures during the BCDI experiment on the patterned sample containing Pt nanoparticles, at \qty{300}{\degreeCelsius} and \qty{400}{\degreeCelsius}.
        Mean partial pressures during \qty{1}{\minute} at the end of each condition, recorded from a leak in the reactor output.
        The pressure under \qty{49}{\ml\per\min} of argon and \qty{1}{\ml\per\min} of ammonia has been subtracted.
    }
    \label{fig:RGANanoparticlesBCDIComparison}
\end{figure}

Interestingly, a transition was recorded at both \qty{300}{\degreeCelsius}, and \qty{400}{\degreeCelsius} approximately \qty{4}{\hour} after the increase of the oxygen partial pressure from \qty{10}{\milli\bar} to \qty{20}{\milli\bar}, leading to a \ce{NH_3}/\ce{O_2} ratio equal to 2.
The partial pressure of \ce{N_2} is approximately divided by two while the partial pressure of \ce{NO} is inversely multiplied by two.

In a general sense, the production of \ce{N_2} is favoured at lower oxygen partial pressures in comparison with \ce{N_2O} and \ce{NO}.
The opposite effects takes place at high oxygen partial pressures, with a clear transition between the main products from \ce{N_2} towards \ce{NO} when the oxygen to ammonia ratio is equal to 8, this transition happens at a lower \ce{NH_3}/\ce{O_2} ratio at \qty{300}{\degreeCelsius}.

    \section{Conclusion and perspectives}

This experiment has both demonstrated the value and limits of Bragg coherent diffraction imaging when studying heterogeneous catalytic reactions.
The instrumental procedure is very time-consuming, especially at the SixS beamline for which two weeks of beamtime are in general needed to perform an experiment, the first week being used to install the BCDI setup and to align the sample in the focal plane of the beam.
The recent improvement of the instrumental setup (rocking curves performed in Fly mode instead of step-by-step mode, sample scanning in Bragg conditions to find the nanoparticles, sapphire window in the dome) tend in the right direction since, as was demonstrated in this study, BCDI is of little use if one does not have the time to measure multiple reflections during an experiment.
The planned upgrade of SOLEIL to a $4^{th}$ generation synchrotron will also play a key role in the hierarchy of SixS in the very competitive list of coherent imaging beamlines, the current resolution of the experiment being too low to properly resolve the smallest facets present of the particles, to distinguish between facets with similar orientations,to observe the growth of surface oxides or the strain in the topmost atomic layers.

Nevertheless, the evolution of Pt nanoparticles during a temperature ramp from \qty{25}{\degreeCelsius} to \qty{600}{\degreeCelsius} was successfully measured in which the reshaping of the particles as a function of temperature was put into evidence, with the appearance of dislocation loops and facets, together with the importance of the presence of a substrate on the displacement field.
The importance of probing different nanoparticles before drawing a conclusion to their global behaviour was also put into perspective by the study of the oxidation of ammonia at \qty{300}{\degreeCelsius} and \qty{400}{\degreeCelsius}, as a function of the ratio between \ce{Ar} and \ammonia.
The role of different facets in the stability of the particle was shown.

If the oxidation of ammonia was never before studied using BCDI, the exposition of platinum nanoparticles to oxygen has been probed first by Fernandez et al. \cite*{Fernandez2019} which showed that it coincided with an increase of the

An upgraded beamline could allow future BCDI studies of the oxidation of ammonia, revealing the 3D structural dynamics of the hysteresis cycle identified on single crystal by \cite{Resta2020a}, or give additional detail regarding the presence of surface oxides on specific facets, as seen in the next chapter.

    \chapter{Pt single crystals during ammonia oxidation}
    \section{Experimental setup}

Past ammonia oxidation experiments on \ce{Pt_{25}Rh_{75}} (100) \parencite{Resta2020} single-crystals gave significant new insights on the relaxation dynamics when the system crosses the catalyst activation barrier.
The experiment focused on measuring before and after the catalyst light off temperature, it was found that the transition from a nitrogen-rich to oxygen-rich input mixture coincides with a relaxation change in the topmost metallic layer: moving from outward to inward, but only for temperatures above the light off.
Thanks to NAP-XPS experiments, nitrogen rich species (such as \nitrogen) were associated to positive relaxation of the topmost layer, expansion; while oxidative species (such as NO and oxygen) were linked to negative relaxation, contraction.
This change in relaxation has an impact on the band structure and therefore on the catalytic properties as proposed in the d-band model as detailed in sec. \ref{sec:Catalysis}.
It is therefore our intent to study if this modification in the relaxation is present only in the \ce{PtRh} alloy or if it has a more general value also on the constituent elements of the alloy to fully understand the impact of the Pt surface orientation on the catalyst selectivity.
Moreover, different relaxation state and transient structures are expected to exist as a function of the ammonia to oxygen ratio.

Moreover, we have recently probed the relaxation of single multi-faceted nanoparticles at SixS under the same conditions using Bragg Coherent Diffraction Imaging.
The 3D displacement and strain evolution of two different facets (\{100\} and \{111\}) will be compared to our results obtained on Pt 100 and 111 crystal surfaces.

\section{Crystal structures}

Platinum crystallizes in a face-centered cubic structure with a lattice parameter $|\vec{a}|$ equal to \qty{3.9254}{\angstrom} at room temperature.
The structure of the [001] and [111] planes differ greatly as shown in fig. \ref{fig:Cubic100Hex111}.
The shortest distance between the atoms on the [001] and [111] surfaces is $a/\sqrt{2} = \qty{2.78}{\angstrom}$.
However, the atoms on the [001] plane still follow a cubic structure whereas the atoms on the [111] plane follow a hexagonal structure.

\begin{figure}[!htb]
    \centering
    \includegraphics[width=0.49\textwidth]{/home/david/Documents/PhD/Figures/introduction/100.pdf}
    \includegraphics[width=0.49\textwidth]{/home/david/Documents/PhD/Figures/introduction/111.pdf}
    \caption{
        Face-centered cubic unit cell of Pt with two [001] crystallographic planes drawn in green at z=0 and z=1 (left).
        The $\vec{a}$, $\vec{b}$, and $\vec{c}$ vectors are used classically to decribe the cubic structure, the $\vec{a}_s$ and $\vec{b}_s$ in-plane vectors separated by \ang{90} and of magnitude \qty{2.78}{\angstrom} can also be used specifically to describe the structure of the [001] planes.
        Face-centered cubic unit cell of Pt with $[111]$ crystallographic plane drawn in green (right).
        The arrangement of the Pt atoms on the $[111]$ crystal planes is hexagonal, which leads to a new definition of the unit cell with the $\vec{a}$ and $\vec{b}$ in-plane vectors separated by \ang{60} and of magnitude \qty{2.78}{\angstrom} and $\vec{c}$ perpendicular to the $[111]$ plane of magnitude \qty{6.81}{\angstrom}.
        There are three $[111]$ planes spanned by the magnitude of $\vec{c}$ (blue, red and green on the figure).
    }
    \label{fig:Cubic100Hex111}
\end{figure}

\subsection{XPS results} \label{sec:XPS111}

On the Pt-Rh alloy an important body of work has been
already carried out, revealing two oxide structures present in the temperature and
pressure ranges we intend to study [6].
It is our goal to retrieve the surface moieties on PtRh nanoparticles as function of the reactive atmosphere’s composition.
The particles are lithography patterned on the sapphire surface, allowing to retrieve each particle position across different setups.
Through Bragg coherent diffraction imaging (BCDI) particles will be studied under reactive condition, then we plan to measure those very same particles at B07. The BCDI measurements give the possibility to follow the particle internal stresses as a function of the gas composition and the behavior of each facet at the same time as difference to the single crystal approach that can probe only one facet at the time.
The B07 setup is sensitive to the surface moieties on the particle across the different for each gas condition.
The colored insert in In Figure 1 is one particle from lithographed sample in operando condition: SixS beamline data (SOLEIL synchrotron).
The particles are sub-micron, like the grains in the metallic gauzes industrially used for this reaction. We plan to expose the particles to five ammonia to oxygen ratios, spanning from oxygen rich to oxygen poor conditions.

The goal of the experiment is to provide new insights on the PtRh catalysts during the NH3 oxidation.
Core level photoemission spectroscopy provides different signatures for many moieties adsorbed on the surface during reactions and therefore the possibilityto strengthen the links between structure, surface moieties and reaction products.
Our recently published results from both surface x-ray diffraction experiments (SXRD) and NAP-XPS [7] link surface relaxation to changes in surface moieties.
The intent of the experiment is to extent such concept to a particle and link internal stresses/surface-relaxations obtained from BCDI with the chemistry on its surfaces.

The B07 NAP-PXS end station can control the xyz positioning of the sample with a precision in the order of 10um, allowing to position the 50um beam[5] in the center of
any of the 100x100um squares of the sample.
At the center of each square is located one isolated particle.
This approach allows to study single particle behavior as function of the NAP-XPS atmosphere but also to retrieve the very same particle studied with BCDI on the SixS beamline.
As for the experiment in Ref. [7], we intent to apply the same temperatures and pressures for BCDI and NAP-XPS.
Temperatures and gas pressures will range between 450, and 650 K and from 0.5 to few tens of mBar respectively.
For each combination of temperature and ammonia to oxygen ratio, we will measure: O1s, N1s, Pt4d, Pt4f and Rh3d core levels that should carry information of the surface state.
Essential for the experiment is also to collect the gas phase composition for each condition to compare particles reactivity.

We expect to establish a connection between selectivity and surface moieties but also between surface moieties and particle internal stress. We also hope to be able to detect variations in the rhodium surface concentration related to the imposed gas conditions

B07 and its recent ambient pressure setup [5], that has now benefit from years of experience, push the high-pressure limits well into the mBar regime.
This new limit and the wide range of achievable temperatures/energies/pressures make it a unique tool for modern surface science and catalysis.
Timewise we expect to be able to characterize the nanoparticle(s) of interest as soon as they entered the analysis chamber from atmosphere and take reference scans after few oxidation/reduction cycles.
Reference scans will be taken on “as inserted” samples and on freshly oxidized/reduced particles with and without adsorbates reactants/products.
Similar scan will also be collected for the bare sapphire substrate. It is estimated that the first 2/3 days will suffice to collect references.
The operando experiments will be performed during the remaining 3/4 days of the beamtime.

\subsubsection{Experimental setup}

Since the XPS measurements were performed at different total pressures, the raw data had first to be reduced in order to analyse the different spectra in a systematic way.
The adopted workflow for the analysis of XPS data is to first align the recorded spectra on the Fermi edge that corresponds to the kinetic energy of the first electron that escapes the sample.
By doing so, one can be confident that any shift in the peak positions is due to chemical changes, such as the oxidation state of the sample, and not to charging effects of the sample (cite).

Secondly, to be able to quantify and compare the evolution of the peak intensity, one must normalize the intensity of the detected electron beam since the electron mean free path depends on the pressure in the reaction chamber \parencite{Willmott}.
The range of kinetic energy just before the absorption edge of Pt 4f was chosen since it had the best signal to noise ratio and does not depend on any experimental parameter besides the pressure.

Finally, for the peaks that showed a good signal to noise ratio, the fitting of the peak shape was realised thanks to the \textit{lmfit} \parencite{Newville2016} package by the means of the Doniach-equation which is the best approximation of the asymmetric peak shape resulting from the convolution of the analyser function and the photoelectron process in metals \parencite{Doniach_1970}.

    \section{Surface x-ray diffraction on a Pt (111) single crystal} \label{sec:SXRD111}

% \textcolor{red}{This chapter covers three areas: analysis of the data; discussion of the results of the analysis; and how your findings relate to the literature. The analysis of the data can be discussed here but the details of any analysis, such as statistical calculations, should be shown in the appendices. You should present any discussion clearly and logically and it should be relevant to your research questions/hypotheses or aims and objectives. Insert any tables or figures that you decide are important in a relevant part of the text not in the appendices, and discuss them fully. Make sure that you relate the findings of your primary research to your literature review. You can do this by comparison: discussing similarities and particularly differences. If you think your findings have confirmed some literature findings, say so and say why. If you think your findings are at variance with the literature, say so and say why.}

% idk what to say why no peak at l = 0 and why intensity is constant for l scans
% need to fit ctr again with good structure of pt 111, see why beta increases under only ammonia, maybe fit only one ctr
% q maps ?
Platinum crystallizes in a face-centered cubic structure with a lattice parameter $a_{Pt}$ equal to \qty{3.94}{\angstrom} at room temperature.
Its structure was first presented in sec. \ref{sec:ScatCrystal} to introduce the notions of crystals.

The arrangement of the Pt atoms on the (111) surface is hexagonal, which leads to the definition of the surface unit cell shown in fig. \ref{fig:SurfaceUnitCellPt111} to be able to better represent the surface arrangement of the Pt atoms.
The in-plane vectors $\vec{a}_{(111)}$ and $\vec{b}_{(111)}$ are of equal magnitude ($a_{Pt} / \sqrt{2} = \qty{2.78}{\angstrom})$, separated by \ang{120}.
The out-of-plane vector $\vec{c}_{(111)}$ is perpendicular to the (111) plane, and of magnitude $3 a_{Pt} / \sqrt{3} = \qty{6.81}{\angstrom}$.

\begin{SCfigure}
    \centering
    \includegraphics[trim=0 2cm 0 2cm, clip, width=0.70\textwidth]{/home/david/Documents/PhD/Figures/introduction/111.pdf}
    \caption{
        Face-centered cubic unit cell of Pt with (111) crystallographic plane drawn in green.
        $\vec{a}_{(111)}$, $\vec{b}_{(111)}$ and $\vec{c}_{(111)}$ are the (111) surface unit cell vectors.
        There are three (111) planes spanned by the magnitude of $\vec{c}_{(111)}$ (blue, red and green on the figure).
    }
    \label{fig:SurfaceUnitCellPt111}
\end{SCfigure}

% \begin{minipage}{.35\textwidth}
%     \centering
%     \includegraphics[width=\textwidth]{/home/david/Documents/PhD/Figures/introduction/111.pdf}
%     \captionof{figure}{
%         Face-centered cubic unit cell of Pt with (111) crystallographic plane drawn in green (b).
%         There are three (111) planes spanned by the magnitude of $\vec{c}_{(111)}$ (blue, red and green on the figure).
%     }
%     \label{fig:SurfaceUnitCellPt111}
% \end{minipage}
%     \hfill

% \begin{minipage}{.35\textwidth}
%     \centering
%     \begin{tabular}{@{}llllll@{}}
%     \toprule
%      & $\alpha$ & $\beta$ & $\gamma$ & $|\vec{a}| = |\vec{b}|$ & $|\vec{c}|$ \\
%     \midrule
%     $[111]$ & \ang{90} & \ang{90} & \ang{120} & \qty{2.7748}{\angstrom} & \qty{6.7969}{\angstrom} \\
%     \bottomrule
%     \end{tabular}%
%     %Lattice parameters and angles used for the Pt (111) surface unit cell.
%     \label{tab:Pt111UnitCell}
% \end{minipage}

\subsection{Ammonia oxidation cycle}

To be able to identify the presence of surface reconstructions and/or surface oxides, in-plane reciprocal space maps were collected by rotating the in-plane sample and detector angles ($\omega$ and $\gamma$) from \ang{0} to \ang{120} to collect a third of the reciprocal space in the ($\vec{q}_x$, $\vec{q}_y$) plane, considering a hexagonal symmetry in the position of the Bragg peaks.

The reciprocal space map were computed in both $q$-space (to obtain the interplanar spacing related to the observed signals) and ($hkl$)-space (fig. \ref{fig:MapsPt111A} and \ref{fig:MapsPt111B}) to visualize the arrangement of surface structures or surface relaxations in comparison with the hexagonal structure of the Pt atoms on the (111) surface, the $h$ and $k$ values being computed using the hexagonal surface unit cell described above.

% The alignment of the sample was often lost during reciprocal space maps which prevented long continuous measurements, even under inert atmosphere (e.g. fig. \ref{fig:MapsPt111B} - \qty{500}{\milli\bar} of Argon).
% The sample proved to be more stable after having lowered the pressure of oxygen in the reactor (fig. \ref{fig:MapsPt111B}).

% Despite an extra layer of Boron Nitride that was applied around the screws used to fix the sample to the sample holder (fig. \ref{fig:SampleHolder}), the high temperature and highly oxidating atmosphere managed twice to corrode the conducting screws which resulted in a contact loss with the heater, thus a change of heater and realignment of the sample surface.
% However, most of the experimental plan was still carried out, lacking only in-plane measurements at \qty{600}{\degreeCelsius} under \qty{8}{\ml\per\min} of oxygen and \qty{1}{\ml\per\min} of ammonia (tab. \ref{tab:Conditions}), due to a lack of experimental time.

\begin{figure}[!htb]
    \centering
    \includegraphics[width=0.495\textwidth]{/home/david/Documents/PhDScripts/SixS_2023_04_SXRD_Pt111/figures/map_hkl_76-115.pdf}
    \includegraphics[width=0.495\textwidth]{/home/david/Documents/PhDScripts/SixS_2023_04_SXRD_Pt111/figures/map_hkl_285-320.pdf}
    \includegraphics[width=0.495\textwidth]{/home/david/Documents/PhDScripts/SixS_2023_04_SXRD_Pt111/figures/map_hkl_481-516_patched.pdf}
    \includegraphics[width=0.495\textwidth]{/home/david/Documents/PhDScripts/SixS_2023_04_SXRD_Pt111/figures/map_hkl_681-689_patched.pdf}
    \caption{
        Reciprocal space maps collected under different atmospheres measured at \qty{25}{\degreeCelsius} under UHV and at \qty{450}{\degreeCelsius} under different atmospheres, with a total pressure equal to \qty{500}{\milli\bar}, computed using the hexagonal lattice of Pt (111).
    }
    \label{fig:MapsPt111A}
\end{figure}

The first map was collected at low pressure (\qty{<1e-5}{\milli\bar}) after the cleaning of the sample by sputtering and annealing.
The ($\bar{1}$$\bar{1}$0) and ($\bar{2}$10) Bragg peaks can be observed, together with the extremity of [111]-oriented crystal truncation rods going through the [0, $\bar{2}$, 0], [0, $\bar{1}$, 0], [$\bar{1}$, 0, 0], [$\bar{1}$, 1, 0], and [$\bar{2}$, 0, 0] positions in reciprocal space.

New peaks can first be observed under \qty{80}{\milli\bar} of oxygen, identified by different colour circles in fig. \ref{fig:MapsPt111A}, the corresponding interplanar spacings are given in tab. \ref{tab:InterplanarSpacingsPt111Oxygen}.
The angle between the peaks circled in grey and white and the peaks circled in red and purple is equal to \ang{60}, which could be the signature of the existence of two hexagonal surface structures, each rotated by $\pm \ang{6}$ with respect to the Pt (111) surface unit cell and with a larger in-plane lattice parameter.
No apparent second order peak can be linked to those rotated hexagonal structures.
Summing the two vectors going from the center of the reciprocal space to the white and purple circled peaks yields the position of the black circled peak, the angle between both vectors is then equal to \ang{42.65}.

The measurement of the entire reciprocal space map took about \qty{3}{\hour}\qty{35}{\minute}, the entire map being a concatenation of multiple $\omega$ scans during which the plane of the sample is rotated, the first scan starting near the center of the reciprocal space.
Therefore, there is approximately \qty{1}{\hour} between the measurement of the peaks circled in red, grey, purple and white, and the peak circled in black.

A second map was measured in a smaller region of the reciprocal space \qty{9}{\hour}\qty{30}{\minute} after the introduction of oxygen in the cell (fig. \ref{fig:MapsPt111A}) to see if the intensity and position of the newly found peaks changed.
The measurement of this map took about \qty{1}{\hour}\qty{15}{\minute}.
Three additional peaks (circled in black) appeared on a (8X8) organisation with comparison to the hexagonal lattice of Pt (111), describing a hexagonal lattice (drawn in black dotted lines) of lattice parameter equal to $|\vec{a_{hex}}| = \qty{3.118}{\angstrom}$.

The grey circled peak seems to have disappeared from the reciprocal map, but an additionnal peak circled in red is revealed at \ang{120} from the other peak circled in red, at the same distance from the center.
From these first measurement, a first hypothesis to explain the position of the observed peaks is that (i) there first appears two rotated hexagonal structures with a larger in-plane lattice parameter in comparison to Pt (111), (ii) there secondly appears a (8X8) hexagonal structure while measuring the large reciprocal space map, also with a larger in-plane lattice parameter, but in the same direction as the in-plane unit cell vectors of Pt (111), explaining the position of the black circled peak in the first map.
It is possible that the disappearance of the grey circle peak is linked to some alignment problems, visible from the large low intensity region in fig. \ref{fig:MapsPt111A}.

Out-of-plane measurements were performed perpendicular to peaks belonging to the (8X8) hexagonal structure ([H, K] = [-0.89, 0], fig. \ref{fig:LScans80} - a) and rotated hexagonal structures ([H, K] = [-0.15, -0.74], fig. \ref{fig:LScans80} - b) to probe their 3D structures.
A second out-of-plane measurement was performed at ([H, K] = [-1.78, 0]) to probe the intensity of a potential second-order peak for the (8X8) hexagonal structure (fig. \ref{fig:LScans80} - c), \qty{10}{\hour} after the measurement perpendicular to ([H, K] = [-0.89, 0]).
The background-subtracted scattered intensity was integrated as a function of $L$ using the \textit{fitaid} module of \textit{binoculars} (fig. \ref{fig:LScans80}).

\begin{figure}[!htb]
    \centering
    \includegraphics[width=\textwidth]{/home/david/Documents/PhDScripts/SixS_2023_04_SXRD_Pt111/figures/l_scans_high_oxygen.pdf}
    \caption{
        In-plane map at $L=0$ (a-b-c) and out-of plane measurements for three different positions under \qty{80}{\milli\bar} of \dioxygen.
    }
    \label{fig:LScans80}
\end{figure}

The integrated intensity perpendicular to ([H, K] = [-0.15, -0.74]) shows a small peak at $L\approx 0.2$ followed by a constant intensity, characteristic of non-periodic out-of-plane structures such as monolayers \parencite{Robinson1991}.

Four different bulk platinum oxides have been reported to exist in litterature, $\alpha$-\ce{PtO_2}, $\beta$-\ce{PtO_2}, \ce{Pt_3O_4} and \ce{PtO}, the experimental structures are describes in tab. \ref{tab:PtOxides} together with the expected equilibrium structure found by \cite{Seriani2006, Seriani2008} from first-principles atomistic thermodynamics calculations and molecular dynamics simulations based on density functional theory.

\begin{table}[!htb]
\centering
% \resizebox{\textwidth}{!}{%
    \begin{tabular}{@{}lllll@{}}
    \toprule
     & \ce{PtO} & \ce{Pt_3O_4} & $\beta$-\ce{PtO_2} & $\alpha$-\ce{PtO_2} \\ \midrule
    Bravais lattice & Tetragonal & Cubic & Orthorhombic & Hexagonal \\
    a (\unit{\angstrom}) & 3.10 (3.08) & 5.65 (5.59) & 4.49 (4.48) & 3.14 (3.10) \\
    b (\unit{\angstrom}) & 3.10 (3.08) & 5.65 (5.59) & 4.71 (4.52) & 3.14 (3.10) \\
    c (\unit{\angstrom}) & 5.41 (5.34) & 5.65 (5.59) & 3.14 (3.14) & 5.81 (4.28-4.40) \\
    $\alpha$ & \ang{90} & \ang{90} & \ang{90} & \ang{90} \\
    $\beta$ & \ang{90} & \ang{90} & \ang{90} & \ang{90} \\
    $\gamma$ & \ang{90} & \ang{90} & \ang{90} & \ang{120}\\
    \bottomrule
    \end{tabular}%
    %}
    \caption{
    Experimental values taken from \cite{McBride1991} for \ce{PtO}, from \cite{MULLER1968} for \ce{Pt_3O_4} and $\alpha$-\ce{PtO_2}, and from \cite{Weber1990} for $\beta$-\ce{PtO_2}.
    Theoretical values and table adapted from \cite{Seriani2006}.
    }
\label{tab:PtOxides}
\end{table}

The lack of peaks at $L=0$ in fig. \ref{fig:LScans80} for the (8X8) hexagonal structure shows that we are not in the presence of a bulk structure.
However, Seriani \textit{et al.} \cite*{Seriani2006} have also found that the most stable surface oxide on the Pt (111) surface is expected to be hexagonal $\alpha$-\ce{PtO_2}.
A high uncertainty regarding the lattice parameter in the $\vec{c}$ direction is reported in experimental and theoretically studies due to a poor crystallization of the bulk oxide \parencite{MULLER1968}, and to underestimation of interlayer forces in theoretical studies.
The formation of hexagonal on hexagonal surface $\alpha$-\ce{PtO_2} on Pt (111) is expected to induce a high amount of in-plane compressive strain in the oxide layer from the large difference between the Pt-Pt distance on the (111) surface (\qty{2.82}{\angstrom}) and the in-plane lattice parameter of $\alpha$-\ce{PtO_2} (\qty{3.14}{\angstrom}).
A rotation of \ang{30} is expected to reduce the in-plane strain, resulting in a (2X2) arrangement.
The presence of $\alpha$-\ce{PtO_2} on the Pt (111) surface was also shown to provide favourable special sites that could contribute to the catalytic activity \parencite{LI2005}.
The 3D structure of bulk $\alpha$-PtO$_2$ is presented in fig. \ref{fig:AlphaPtO2}.

\begin{SCfigure}
    \centering
    \includegraphics[trim=0 2.5cm 0 2.5cm, clip, width=0.35\textwidth]{/home/david/Documents/PhD/Figures/introduction/AlphaPtO2.pdf}
    \caption{
        $\alpha$-\ce{PtO_2} bulk unit cell.
        Platinum atoms are situated on the unit cell corners while the two oxygen atoms are at the positions (1/3, 2/3, 1/4) and (2/3, 1/3, 3/4).
    }
    \label{fig:AlphaPtO2}
\end{SCfigure}

Ellinger \textit{et al.} \parencite*{Ellinger2008} reported the existence of surface $\alpha$-\ce{PtO_2} on Pt (111) in an experimental study using SXRD at similar conditions (\qty{500}{\milli\bar} of oxygen pressure of, temperature between \qty{245}{\degreeCelsius} and \qty{635}{\degreeCelsius}), with an in-plane lattice parameter equal to \qty{3.15}{\angstrom}, close to the reported bulk value of $\alpha$-PtO$_2$ (\qty{3.14}{\angstrom}).
Similar results were found by Ackermann \parencite*{Ackermann2007}.
Ellinger \textit{et al.} \parencite*{Ellinger2008} proposed a surface model for surface $\alpha$-\ce{PtO_2} consisting of a full unit cell, with an out-of-plane lattice parameter reduced by \qty{15}{\percent} in comparison with the bulk structure (\qty{3.62}{\angstrom}), and with a lateral displacement of \qty{\approx33}{\percent} in the [110] direction of the Platinum atoms on the top layers with respect to the Pt atoms of the bottom layer, in contact with the Pt (111) surface.
The question of the epitaxial relation between both structures remains, since the proposed structure does come with a large misfit strain (\qty{13.5}{\percent}) when in a hexagonal on hexagonal arrangement.

Similarly to the current experiment, an (8X8) coincidence is reported between the Pt (111) hexagonal surface and surface $\alpha$-\ce{PtO_2}.
The out-of-plane measurements in \cite{Ellinger2008} also performed perpendicular to ([H, K] = [-0.89, 0]) presents a peak at $L=0.65$ ($L$ computed using the distorted $\alpha$-\ce{PtO_2} unit cell).

The two $L$-scans performed in the current study were fitted using two Gaussian peaks of same full width at half maxima (FWHM) and a constant background (fig. \ref{fig:LScans80Fit}), the second peak near $L=1$ in the $L$-scan perpendicular to ([H, K] = [-0.89, 0]) comes from a Pt powder signal.

\begin{figure}[!htb]
    \centering
    \includegraphics[width=\textwidth]{/home/david/Documents/PhDScripts/SixS_2023_04_SXRD_Pt111/figures/ctr_reconstructions_fitting_result.pdf}
    \caption{
        Fit of out-of-plane measurements perpendicular to peaks observed in in-plane maps under \qty{80}{\milli\bar} of \dioxygen.
    }
    \label{fig:LScans80Fit}
\end{figure}

We report a peak at [-0.89, 0, 0.65] but using the Pt (111) surface unit cell, with an out-of-plane lattice parameter equal to \qty{6.81}{\angstrom}, incompatible with the results presented in \cite{Ellinger2008}.
It seems that a different out-of-plane structure is present for the $\alpha$-\ce{PtO_2} surface oxide in this study.
Additional simulations are needed to fully understand the structures present in this study, with different surface models including the \textit{spoked-wheel} superstructure identified \textit{via} scanning tunneling microscopy (STM) by \cite{VanSpronsen2017, Boden2022}, and the different surface oxides identified in \cite{Miller2011,Miller2014} and \cite{Fantauzzi2017, FARKAS2017} at similar temperature and oxygen partial pressures, respectively \qty{6.66}{\milli\bar} of \dioxygen at \qty{450}{\degreeCelsius} and \qty{1}{\milli\bar} of \dioxygen between \qty{160}{\degreeCelsius} and \qty{410}{\degreeCelsius}.

\begin{figure}[!htb]
    \centering
    \includegraphics[width=\textwidth]{/home/david/Documents/PhDScripts/SixS_2023_04_SXRD_Pt111/figures/reflecto_80.pdf}
    \caption{
    	X-ray reflectivity curves measured in a specular geometry (full lines) under \qty{80}{\milli\bar} of oxygen.
    	Curves fitted using \text{GenX} are shown an order of magnitude below as dotted lines.
    }
    \label{fig:Reflecto80}
\end{figure}

Reflectivity curves measured in a specular geometry and fitted using \textit{GenX} are presented in fig. \ref{fig:Reflecto80}.
A simple model was used consisting of a slab of platinum on top of which is present a homogeneous layer.
The density of the layer was set free during the minimization process between zero and \qty{0.021172}{FU\per\cubic\angstrom}, the value for $\alpha$-PtO$_2$.
The fitted density values is divided by two between the first (\qty{0.00727}{FU\per\cubic\angstrom}) and second (\qty{0.00358}{FU\per\cubic\angstrom}) measurement, and is always an order of magnitude below the expected value for a homogeneous layer of $\alpha$-PtO$_2$.
The layer was found to be \qty{13.89}{\angstrom} and \qty{14.185}{\angstrom} thick in the first and second reflectivity curves, with a very low roughness (\qty{<1e-7}{\angstrom}).

By combining the information obtained from in-plane, out-of-plane and reflectivity measurements, it is probable that the platinum surface is first covered by two types of domains corresponding to the two rotated hexagonal structures, and that there secondly appears a $\alpha$-\ce{PtO_2} surface oxide, with an average thickness of \qty{\approx 14}{\angstrom}, \textit{i.e.}, a few monolayers thick, confirmed by the presence of peaks on the L-scans, in contrast with the rotated structures.

% \begin{figure}[!htb]
%     \centering
%     \includegraphics[trim=0 1cm 0 1cm, clip, width=\textwidth]{/home/david/Documents/PhD/Figures/introduction/EpitaxyPtO2Pt111.pdf}
%     \caption{
%     	Epitaxy relationship between Pt(111) surface unit cell (in red) and $\alpha$-\ce{PtO_2} basal unit cell.
%     }
%     \label{fig:PtO2OnPt111}
% \end{figure}

Ammonia was subsequently introduced in the reactor cell to investigate the evolution of the platinum oxide structures, as well as the presence of additionnal surface structures or surface reconstructions during the oxidation of ammonia on the platinum surface (fig. \ref{fig:MapsPt111B}).

\begin{figure}[!htb]
    \centering
    \includegraphics[width=0.495\textwidth]{/home/david/Documents/PhDScripts/SixS_2023_04_SXRD_Pt111/figures/map_hkl_830-865.pdf}
    \includegraphics[width=0.495\textwidth]{/home/david/Documents/PhDScripts/SixS_2023_04_SXRD_Pt111/figures/map_hkl_1041-1076.pdf}
    \includegraphics[width=0.495\textwidth]{/home/david/Documents/PhDScripts/SixS_2023_04_SXRD_Pt111/figures/map_hkl_1413-1448.pdf}
    \includegraphics[width=0.495\textwidth]{/home/david/Documents/PhDScripts/SixS_2023_04_SXRD_Pt111/figures/map_hkl_1458-1493.pdf}
    \caption{
        Reciprocal space maps collected under different atmospheres measured at \qty{450}{\degreeCelsius}, computed using the hexagonal lattice of Pt (111).
    }
    \label{fig:MapsPt111B}
\end{figure}

None of the previously observed peaks could be discerned anew, and no additional peaks emerged.
Even when the oxygen pressure within the cell was reduced to \qty{5}{\milli\bar} (favoring alternate products) or eliminated at \qty{0}{\milli\bar}, there was no induction of surface reconstructions or surface structures.
All reciprocal space maps following the introduction of ammonia to a final atmosphere of \qty{500}{\milli\bar} of Argon showed a pristine surface structure, proving the oxidation cycle is reproducible and that ammonia effectively removed the different surface oxides that grew under \qty{80}{\milli\bar} of oxygen.

\begin{figure}[!htb]
    \centering
    \includegraphics[width=\textwidth]{/home/david/Documents/PhDScripts/SixS_2023_04_SXRD_Pt111/figures/reflecto_cycle.pdf}
    \caption{
    	X-ray reflectivity curves measured in a specular geometry (full lines) under different atmospheres during the oxidation of ammonia.
    	Curves fitted using \text{GenX} are shown an order of magnitude below as dotted lines.
    }
    \label{fig:ReflectoCycle}
\end{figure}

Reflectivity curves were also measured and fitted with \textit{Genx} under pure Argon atmosphere (before and after the oxidation cycle), under different reacting conditions, and with only ammonia and argon in the reactor, presented in fig. \ref{fig:ReflectoCycle}.
In accordance with the results of the in-plane reciprocal space maps in fig. \ref{fig:MapsPt111B}, no oxide layer could be detected to exist on top of the platinum surface.
The roughness of the [111]-oriented platinum crystal decreases following the introduction of ammonia after having increased to \qty{\approx 2.61}{\angstrom} when being exposed to \qty{80}{\milli\bar} of oxygen for more than \qty{20}{hours}, almost reaching its value before the start of the oxidation cycle.

\subsection{Oxide growth monitored under \qty{5}{\milli\bar} of \dioxygen}

Following the complete ammonia oxidation cycle, the sample was put under \qty{5}{\milli\bar} of oxygen to see if any of the detected surface structures would come to grow again, and to decorrelate the effect of the simultaneous presence of ammonia and oxygen in the cell on the surface when favouring the production of nitrogen at lower oxygen pressures.
The sample was cleaned with two sputtering and annealing cycles before the introduction of oxygen.
Large and small reciprocal space maps were measured to have the best compromise between a higher temporal resolution of the oxide growth, and the possibility to detect other peaks to corresponding surface unit cells.
The larger reciprocal space maps are presented in fig. \ref{fig:LargeMapsPt111LowOxygen}.

\begin{figure}[!htb]
    \centering
    \includegraphics[width=\textwidth]{/home/david/Documents/PhDScripts/SixS_2023_04_SXRD_Pt111/figures/large_maps_05O2.pdf}
    \caption{
        Reciprocal space maps collected under \qty{495}{\milli\bar} of argon and \qty{5}{\milli\bar} of oxygen at \qty{450}{\degreeCelsius}, for different exposure times.
    }
    \label{fig:LargeMapsPt111LowOxygen}
\end{figure}

In the first map, that started directly after the stabilization of the oxygen pressure at \qty{5}{\milli\bar} in the reactor cell, three low intensity peaks can be detected that correspond to the same peaks detected previously shortly after having a pressure of \qty{80}{\milli\bar} in the cell (fig. \ref{fig:LargeMapsPt111LowOxygen}).
Only the grey circled peak cannot be detected until \qty{4}{hours} of measurement, which is also when some peaks start to split in a similar pattern, exhibiting three distinct peaks around a more diffuse region.
This splitting of the peaks could be the signature of a Pt (111) surface covered by different domains exhibiting similar hexagonal structures, slightly different from each other, while being themselved rotated by $\pm \ang{6}$ with respect to the Pt (111) surface unit cell with a higher magnitude of the in-plane lattice parameter, similarly to what has been observed under an oxygen pressure of \qty{80}{\milli\bar}.
A peak can be observed at (-0.89, -0.89) in the only large reciprocal space map collected between \qty{41}{\minute} and \qty{4}{\hour}\qty{3}{\minute} after the start of the condition, at the same position as in the first large reciprocal space map measured under \qty{80}{\milli\bar} of oxygen.
The peak position corresponds to the (-1, -1, 0)$_{\alpha-PtO_2}$ Bragg peak of bulk $\alpha-PtO_2$, but no peaks at (-0.89, 0, 0) and (0, -0.89, 0) positions can be detected (respectively (-1, 0, 0)$_{\alpha-PtO_2}$ and (0, -1, 0)$_{\alpha-PtO_2}$).

The area sampled during the measurement was extended in the two last maps which allows us to observe more signals around the ($1\bar{1}0$) region.
For each map, two peaks circled by the same colour draw two vectors of the same magnitude going through the center of the reciprocal space and separated by \ang{120}, in the last two maps there are three doublet of peaks separated by \ang{120}, in red, grey and purple ($q$-map available in appendix \ref{fig:2064QSpace}).
Moreover, the purple and red circled peaks are separated by \ang{60}, likewise for the grey and white circled peaks.
During this set of measurement, the grey circled peak did not disappear on the contrary to the measurements carried out under \qty{80}{\milli\bar} of oxygen, which furthermore supports the existence of at least two rotated hexagonal structures.

The ($\bar{1}0$0)$_{\alpha-PtO_2}$, ($0\bar{1}$0)$_{\alpha-PtO_2}$ and ($\bar{1}\bar{1}$0)$_{\alpha-PtO_2}$ peaks corresponding to surface ${\alpha-PtO_2}$ are detected in the large maps after \qty{24}{hours} of measurements, circled in black, smaller maps taken with a shorter time interval are presented in fig. \ref{fig:SmallMapsPt111LowOxygen} in which the ($\bar{1}0$0)$_{\alpha-PtO_2}$ peak is also first detected after \qty{24}{hours} of measurements.

\begin{figure}[!htb]
    \centering
    \includegraphics[width=\textwidth]{/home/david/Documents/PhDScripts/SixS_2023_04_SXRD_Pt111/figures/small_maps_05O2.pdf}
    \caption{
        Reciprocal space maps collected under \qty{495}{\milli\bar} of argon and \qty{5}{\milli\bar} of oxygen at \qty{450}{\degreeCelsius}, for different exposure times.
    }
    \label{fig:SmallMapsPt111LowOxygen}
\end{figure}

A more quantitative analysis of the different structures appearing during the exposition to oxygen was performed by integrating the scattered intensity around the white, red and black circled peak, present in the ($\bar{1}$10) region.
The average background was subtracted to each reciprocal space voxel before integration.
%, while the intensity of the crystal truncation rod ($\bar{1}$10) Bragg peak was used for normalisation to correct possible miss-alignement effects on the integrated intensity.
The starting time of each reciprocal space map is used as estimate for the time since the introduction of oxygen, the evolution of each peak is presented in fig. \ref{fig:HexBraggPeaks}.

\begin{figure}[!htb]
    \centering
    \includegraphics[width=\textwidth]{/home/david/Documents/PhDScripts/SixS_2023_04_SXRD_Pt111/figures/intensity_comparison_hex_reconstructions.pdf}
    \caption{
        Intensity evolution for peaks corresponding to the hexagonal surface structure and to the surface oxide.
    }
    \label{fig:HexBraggPeaks}
\end{figure}

The ${\alpha-PtO_2}$ surface oxide only appears after \qty{\approx 23}{\hour} of exposition to \qty{5}{\milli\bar} of oxygen (fig. \ref{fig:HexBraggPeaks} - a), when it was already detected at the very least after \qty{10}{hour} under \qty{80}{\milli\bar} of oxygen.
The nature of the peak situated at (-0.89, -0.89) is not yet clear since it appears before the (-0.89, 0) and (0, -0.89) peaks, all three positions corresponding to the (8X8) hexagonal structure.
As mentionned before, it is possible to draw a surface unit cell including this peak together with the white and purple circled peak, but with an in-plane angle $\gamma^*$ equal in this case to \ang{42.13}.
The growth seems to be exponential after \qty{23}{\hour} of exposition.
The white peak is visible from the start of the exposition to oxygen (fig. \ref{fig:SmallMapsPt111LowOxygen}), the red peak is only observed \qty{7}{hour} later with a lower intensity.
The intensity of both peaks (related to the other rotated hexagonal structures) quickly increases in the first hours and then plateaus until \qty{\approx 23}{\hour}, at which it increases together with the appearance of the $\alpha$-\ce{PtO_2} surface oxide.
Finally, the intensity starts to decreases after \qty{\approx 25}{\hour} of exposition, which coincides with the duration after which the non-rotated hexagonal structure peak intensity increases exponentially.

Several out-of-plane measurement have been carried out perpendicular to the white and red circled peaks, presented in appendix \ref{fig:LScans05}.
The peak intensity is constant as a function of $L$, similarly to the out-of-plane measurements presented in \ref{fig:LScans80}, which shows that we are not either in the presence of bulk oxides for the rotated hexagonal structures.
No peak was measured perpendicularly to the (8X8) hexagonal structure.

Lowering the oxygen partial pressure in the reactor cell from \qtyrange{80}{5}{\milli\bar} has demonstrated that similar in-plane structures appear, but with a slower growth rate.
Additionnal studies are needed understand if wether or not the rotated hexagonal structures act only as precursors to the appearance of surface $\alpha$-\ce{PtO_2}, or if those peaks are still present after an extended exposure to oxygen.

% Finally, reflectivity curves were measured before the introduction of oxygen, as well as \qty{14}{\hour}\qty{30}{\minute} and \qty{23}{\hour}\qty{30}{\minute} after the introduction of oxygen, shown in fig. \ref{fig:reflecto}.
% The curves were fitted with \textit{Genx} following the method described earlier.
% The hexagonal structure was not yet detected when measuring the second reflectivity curve, a very low density oxide layer can be
% There is a significant drop of intensity in the second reflectivity curve,

% \begin{figure}[!htb]
%     \centering
%     \includegraphics[width=\textwidth]{/home/david/Documents/PhDScripts/SixS_2023_04_SXRD_Pt111/figures/reflecto_05.pdf}
%     \caption{
%     	X-ray reflectivity curves measured in a specular geometry (full lines) under different atmospheres during the oxidation of ammonia.
%     	Curves fitted using \text{GenX} are shown an order of magnitude below as dotted lines.
%     }
%     \label{fig:reflecto_5}
% \end{figure}

\subsection{Surface roughness and surface relaxation effects}

\begin{figure}[!htb]
    \centering
    \includegraphics[width=\textwidth]{/home/david/Documents/PhDScripts/SixS_2023_04_SXRD_Pt111/figures/ctr_a.pdf}
    \includegraphics[width=\textwidth]{/home/david/Documents/PhDScripts/SixS_2023_04_SXRD_Pt111/figures/ctr_b.pdf}
    \caption{
        Evolution of [111] crystal truncation rods measured perpendicular to three different Bragg peaks under different atmospheres.
    }
    \label{fig:CTRPt111}
\end{figure}

Crystal truncation rods have been measured perpendicularly to the ($\bar{2}10$), ($\bar{1}00$) and ($\bar{1}\bar{1}0$) positions \qty{6}{\hour} after the start of each condition, each measurement lasted for \qty{2}{\hour}.
The background-subtracted intensity of the CTR was integrated using the \textit{fitaid} module of \textit{binoculars} as a function of $L$ with the same integration range.% ($\Delta H, \Delta K$).
The CTR intensity is presented in fig. \ref{fig:CTRPt111}, no additionnal peak could be detected at any condition, besides the Pt Bragg peaks at higher $L$ values.
At first sight, the rugosity seems to evolve during the exposition to different atmospheres, visible from the evolution of the minimum of intensity near $L=1.5$ (or $L=2.5$ for the CTR recorded perpendicular to the ($\bar{1}$00) Bragg peak).
The intensity of the CTR under \qty{500}{\milli\bar} of argon before the oxidation cycle was progressively lost during the measurements due to problems with the sample heater, also visible in the reciprocal space map under the same condition in fig. \ref{fig:MapsPt111A}, and is therefore not shown in the data besides for the ($\bar{1}\bar{1}$0) Bragg peak.

\begin{figure}[!htb]
    \centering
    \includegraphics[trim=0 5cm 0 4.5cm, clip, width=0.45\textwidth]{/home/david/Documents/PhD/Figures/introduction/Pt111HexA.pdf}
    \includegraphics[trim=0 5cm 0 4.5cm, clip, width=0.45\textwidth]{/home/david/Documents/PhD/Figures/introduction/Pt111HexB.pdf}
    \caption{
        View from above (a) and from the side (b) of the Pt (111) surface with the atoms belonging to the A, B and C layers respectively coloured in green, red and blue.
        The size of the Pt atoms has been tuned from (a) to (b) to be able to visualize the arrangement of the A, B and C layers.
    }
    \label{fig:Pt111StructureSideAndTop}
\end{figure}

To investigate potential surface relaxation effects, the three CTR were fitted together to increase the number of data points at each condition using \textit{ROD}, with a simple model consisting of three ABC layers of [111] oriented platinum, C being the topmost layer (fig. \ref{fig:Pt111StructureSideAndTop}).
These three layers are on top of the rest of the crystal, so forth denominated as the \textit{bulk}, separating it from the \textit{surface} of the crystal in which relaxation effects can be detected.

Different models have been tested to fit the CTR intensity as a function of $L$.
Atoms on the same layer always shared the same out-of-plane position, an atomic displacement parameter per layer was introduced to see if surface relaxations effect could be detected at different atmospheres.
The different models had for difference that either one, two, or three layers were allowed to have a out-of-plane displacement.
Another model was tested in which the C and B layer shared the same displacement.
During the fitting process, in-plane atomic displacement were excluded as one of the parameters to be adjusted because the presence of too many parameters posed challenges in achieving a successful convergence of the fitting routine.
The roughness parameter $\beta$ was set free between 0 and 0.5.% (unitless parameter, see $\beta$ roughness model detailed in sec. \ref{sec:CTR}).

\begin{figure}[!htb]
    \centering
    \includegraphics[width=\textwidth]{/home/david/Documents/PhDScripts/SixS_2023_04_SXRD_Pt111/figures/fit_comparison_1_last_layer_free.pdf}
    \caption{
        Evolution of [111] crystal truncation rods measured perpendicular to three different Bragg peaks under different atmospheres.
    }
    \label{fig:CTRFit111}
\end{figure}

The best fit was found for a model in which only the C layer was allowed to have a common out-of-plane atomic displacement $\delta_z$ between \qty{-0.05}{\angstrom} and \qty{0.05}{\angstrom}, the position of the Pt atoms on the A and B layers were fixed following the position of the atoms in the bulk.

The strain with respect to the bulk was computed following eq. \ref{eq:StrainDiffraction}, the reference was set to the magnitude of the Pt (111) out-of-plane vector, \textit{i.e.} $|\vec{c}_{(111)}| = \qty{6.81}{\angstrom}$.
The evolution of the CTR roughness and of the strain of each layer is shown in fig. \ref{fig:CTRFit111}.

The roughness of the Pt (111) surface (fig. \ref{fig:CTRFit111} - a) increases with the introduction of Argon in the cell at \qty{450}{\degreeCelsius}, which is probably due to the presence of impurities in the gas flow.
The introduction of oxygen in the cell further increases the surface roughness, as expected from the formation of the different surface oxides visible in the in-plane reciprocal space maps.
Adding ammonia in the reactor cell, which was seen to remove the different surface oxides, has also the effect of decreasing the surface roughness.
The surface roughness increases slightly again when oxygen is removed, but reaches a value of 0 when both gases are removed from the reactor, falling back to an inert atmosphere, with a lower surface roughness than at the beginning of the measurement (visible also in fig. \ref{fig:CTRPt111}).
From those results, it seems that the ammonia oxidation cycle has effectively \textit{cleaned} the surface from the presence of impuritied or surface oxides.
Finally, the re-introduction of \qty{5}{\milli\bar} of oxygen increases the surface roughness again, in accordance with the formation of surface oxides.

Overall, a very low amount of strain is detected on the surface, almost imperceptible when observing position of the CTR minimum in fig. \ref{fig:CTRPt111}.
From the fitting results, the topmost layer is already under tensile strain at UHV, further increased by the presence of Argon and possible impurities from the opening of the gas valves.
The largest evolution in the strain values comes from the high oxygen atmosphere, which has the effect of decreasing the surface strain, with an out-of-plane lattice parameter almost equal to the bulk value.
The formation of surface oxides observed under this atmosphere does not seem to have a very important effect on the surface relaxation state.
The introduction of ammonia increases again the surface strain, higher ammonia to oxygen ratio coincide with higher tensile strain.
Since the reacting conditions and the sole presence of ammonia have had the effect of removing the different surface oxides present on the platinum surface, as well as decresing the roughness quantified \textit{via} reflectivity measurements and the evolution of the $\beta$ parameter in fig. \ref{fig:CTRFit111} (a), it is possible that the tensile strain in the last layer under Argon after the oxidation cycle corresponds to the equilibrium state of a clean Pt (111) surface.
Finally, the re-introduction of \qty{5}{\milli\bar} of oxygen decreases slightly the tensile strain on the topmost layer, a weaker but similar effect to the presence of \qty{80}{\milli\bar} of oxygen.
To conclude, no changes from tensile to compressive strain are observed during the reaction, the presence of oxygen alone in the reactor cell during which the growth of surface oxides has been monitored has the effect of lowering the surface strain, while the presence of ammonia has the opposite effect.
Different reacting conditions are related to the same direction of displacement but with a lower magnitude in a when lowering the partial pressure of oxygen.

\subsection{Surface species presence}

\begin{table}[!htb]
\centering
\resizebox{\textwidth}{!}{%
    \begin{tabular}{@{}lllllllll@{}}
    \toprule
    \multirow{3}{*}{Partial pressures (mbar)} & Ar & 1 & 0 & 0 & 0 & 0 & 1 & 0 \\
     & NH3 & 0 & 0 & 1.1 & 1.1 & 1.1 & 0 & 0 \\
     & O2 & 0 & 8.8 & 8.8 & 0.55 & 0 & 0 & 0.55 \\
    \midrule
    % \multicolumn{2}{l}{\begin{tabular}[c]{@{}l@{}}Peak position in N1s spectra (eV)\\ Eph = 550 eV\end{tabular}} & No peak & No peak & No peak & \begin{tabular}[c]{@{}l@{}}404.09, 399.81,\\ 397.54\end{tabular} & 405.32, 400.79 & No clear peak & No peak \\
    % \multicolumn{2}{l}{Corresponding surface moieties} &  &  & N, NH, NH2, NO ? & N, NH, NH2, NO ? & N, NH, NH2 ? &  &  \\
    \multicolumn{2}{l}{\begin{tabular}[c]{@{}l@{}}Peak position in N1s spectra (eV)\\ Eph = 700 eV\end{tabular}} & No data & No peak & No peak & \begin{tabular}[c]{@{}l@{}}404.12, 399.73,\\ 397.54\end{tabular} & 405.32, 400.83 & 400.43, 398.38 & No peak \\
    \multicolumn{2}{l}{Corresponding surface moieties} &  &  &  & N, NH, NH2, NO ? & N, NH, NH2, NO ? & N, NH, NH2 ? &  \\
    \multicolumn{2}{l}{\begin{tabular}[c]{@{}l@{}}Peak position in O1s spectra (eV)\\ Eph = 700 eV\end{tabular}} & 532.33 &  & 534.05  & 534.05 & 532.40 & 532.40 & 531.42, 529.71 \\
    \multicolumn{2}{l}{Corresponding surface moieties} &  &  &  &  &  &  &  \\
    \multicolumn{2}{l}{\begin{tabular}[c]{@{}l@{}}Peak position in Pt4f spectra (eV)\\ Eph = 550 eV\end{tabular}} &  &  &  &  &  &  &  \\
    \multicolumn{2}{l}{Corresponding surface moieties} &  &  &  &  &  &  &  \\ \bottomrule
    \end{tabular}%
    }
\caption{}
\label{tab:XPSPt111}
\end{table}

\begin{figure}[!htb]
    \centering
    \includegraphics[width=\textwidth]{/home/david/Documents/PhDScripts/B07_2022_04_XPS/Figures/Pt111/O1sN1s_700.pdf}
    \caption{
        Spectra collected around the N1s edge (tabulated binding energy equal to \qty{409.9}{\eV}) under different atmospheres at \qty{450}{\degreeCelsius} with an incoming photon energy of \qty{700}{\eV}.
        The spectra are normalized and shifted in intensity to highlight the presence of different peaks.
    }
    \label{fig:O1sN1sPt111}
\end{figure}

\begin{figure}[!htb]
    \centering
    \includegraphics[width=\textwidth]{/home/david/Documents/PhDScripts/B07_2022_04_XPS/Figures/Pt111/Pt4f_550_no_fit.pdf}
    \caption{
        Spectra collected around the Pt4f edge doublet (tabulated binding energy equal to \qty{74.5}{\eV} and \qty{71.2}{\eV}) under different atmospheres at \qty{450}{\degreeCelsius} with an incoming photon energy of \qty{550}{\eV}.
        A Shirley-type background has been subtracted from all XP spectra, that have then been normalised by the pre-edge intensity.
    }
    \label{fig:Pt4fPt111}
\end{figure}



\cite{Miller2011} have measured both O1s and Pt4f edges at under \qty{6.66}{\milli\bar} of \dioxygen at \qty{450}{\degreeCelsius}, and under \qty{0.66}{\milli\bar} of \dioxygen at \qty{350}{\degreeCelsius}.
    \section{Pt(100) single crystal studied at \qty{450}{\degreeCelsius}} \label{sec:SXRD100}

A similar experiment was carried out on a Pt(100) sample.
The arrangement of the Pt atoms on this surface is square, the distance between in-plane neighbouring Pt atoms is smaller than between out-of-plane atoms.
A surface unit cell must be derived, shown in fig. \ref{fig:SurfaceUnitCellPt100}, to represent the surface arrangement of the Pt atoms with the smallest unit vectors possible.
The in-plane vectors $\vec{a}_{(100)}$ and $\vec{b}_{(100)}$ are of equal magnitude ($a_{Pt} / \sqrt{2} = \qty{2.77}{\angstrom})$, separated by \ang{90}.
The out-of-plane vector $\vec{c}_{(100)}$ is perpendicular to the (100) plane, and of magnitude $a_{Pt} = \qty{3.92}{\angstrom}$.

\begin{SCfigure}
    \centering
    \includegraphics[trim=0 2cm 0 2cm, clip, width=0.70\textwidth]{/home/david/Documents/PhD/Figures/Introduction/100.pdf}
    \caption{
        Face-entered cubic unit cell of Pt with $(100)$ crystallographic plane drawn in green.
        $\vec{a}_{(100)}$, $\vec{b}_{(100)}$ and $\vec{c}_{(100)}$ are the $(100)$ surface unit cell vectors.
    }
    \label{fig:SurfaceUnitCellPt100}
\end{SCfigure}

Reciprocal space in-plane maps were collected to probe the structural evolution of the sample during ammonia oxidation, using the same experimental setup and atmospheres as for Pt(111) (tab. \ref{tab:ConditionsSXRD}).
Considering the square symmetry in the position of the Bragg peaks, the in-plane reciprocal space maps were collected by rotating the in-plane sample and detector angles ($\mu$ and $\gamma$) from \ang{0} to \ang{90} to collect a quarter of the reciprocal space in the ($\vec{q}_x$, $\vec{q}_y$) plane

The reciprocal space in-plane maps were computed in ($hkl$)-space (fig. \ref{fig:MapsPt100A}) to visualise the arrangement of surface structures or surface relaxations in comparison with the hexagonal surface unit cell described in fig. \ref{fig:SurfaceUnitCellPt100}.

\subsection{Oxide growth under 80 mbar of oxygen}

\begin{figure}[!htb]
    \centering
    \includegraphics[width=0.495\textwidth]{/home/david/Documents/PhD/PhDScripts/SixS_2022_01_SXRD_Pt100/figures/Map_hkl_surf_or_1335-1375.pdf}
    \includegraphics[width=0.495\textwidth]{/home/david/Documents/PhD/PhDScripts/SixS_2022_01_SXRD_Pt100/figures/Map_hkl_surf_or_1596-1635_patched.pdf}
    \caption{
        Reciprocal space in-plane maps collected under different atmospheres at \qty{450}{\degreeCelsius}, computed using the surface lattice of Pt(100).
    }
    \label{fig:MapsPt100A}
\end{figure}

The first map was collected under \qty{500}{\milli\bar} of Ar, after the cleaning of the sample by sputtering and annealing (fig. \ref{fig:MapsPt100A} - a).
The (200), (1$\bar{1}$0) and (0$\bar{2}$0) Bragg peaks can be observed, together with the intersection of the [1, 0, L] and [0, $\bar{1}$, L] CTRs with the [H, K, 0] plane (circled in white).

The sample was then first exposed to \qty{80}{\milli\bar} of oxygen, while keeping the total pressure to \qty{500}{\milli\bar} by compensation with argon (fig. \ref{fig:MapsPt100A} - b).
The signals generated by the rich oxygen condition can be divided into two groups, identified by red and green circles.
The red peaks are situated approximately at [H, K] = [1, -0.5], [H, K] = [1, -1.5], [H, K] = [0.5, $\bar{1}$] and [H, K] = [0.5, -1.5].
These signal follow a Pt(100)-($2\times2$) reconstruction, keeping in mind that it is related to surface unit cell of Pt(100).
A ($2\times2$) reconstruction of the Pt(100) surface has been reported for nanoparticles measured under \qty{1e-3}{\milli\bar} of oxygen, by the use of environmental TEM \parencite{Li2016}.

The green peaks are slightly shifted from the red peaks by the same amount $\delta = 0.07$ in either H or K.
No simple surface unit cell could be derived from their in-plane positions.

Two super structure rods (SSRs) were performed on each group of signal, up to $L=3.5$, to probe the related out-of-plane structure.
The background-subtracted intensity was integrated as a function of $L$ using the \textit{fitaid} module of \textit{BINoculars} (fig. \ref{fig:LScansHighOxygenPt100}), following the same kind of procedure performed on Pt(111).

\begin{figure}[!htb]
    \centering
    \includegraphics[width=\textwidth]{/home/david/Documents/PhD/PhDScripts/SixS_2022_01_SXRD_Pt100/figures/l_scans_high_oxygen_no_map.pdf}
    \caption{
        SSRs for four different positions in [H, K] under \qty{80}{\milli\bar} of \ce{O_2}, and \qty{420}{\milli\bar} of \ce{Ar}.
    }
    \label{fig:LScansHighOxygenPt100}
\end{figure}

The intensity distribution as a function of $L$ for the shifted peaks quickly decreases down to zero, which shows they correspond to flat monolayers (simulation example in fig. \ref{fig:SimROD}, Robinson et al. \cite*{Robinson1991}).
However, four peaks at different $L$ values are visible on both SSRs measured on the Pt(100)-($2\times2$) reconstruction in fig. \ref{fig:LScansHighOxygenPt100}, including at $L=0$, which is characteristic of a bulk structure, \textit{i.e.} more than a few unit cells thick (simulation example also in fig. \ref{fig:SimROD}, Robinson et al. \cite*{Robinson1991}).
The peaks were fitted using a model with four Gaussian peaks and a constant background, the same full width at half maxima was used for all four peaks since they are assumed to be linked to the same structure (fig. \ref{fig:FitPt100LScans}).
A $r^2$ value of \num{0.86} and \num{0.98} was reached for the SSRs at [H, K] = [0.5, $\bar{1}$], and [H, K] = [0.5, -1.5], respectively.
In the second fit, the maximum value of each peak is better represented which explains the difference in the $r^2$ values.

Corresponding interplanar spacings were computed from the position of the peaks, corresponding to a slightly distorted cubic unit cell with an in-plane lattice parameter equal to \qty{5.60}{\angstrom}, and an out-of-plane lattice parameter equal to \qty{5.64}{\angstrom}.
This unit cell was used to index each Bragg peak in fig. \ref{fig:FitPt100LScans}.

\begin{figure}[!htb]
    \centering
    \includegraphics[width=\textwidth]{/home/david/Documents/PhD/PhDScripts/SixS_2022_01_SXRD_Pt100/figures/ctr_reconstructions_fitting_result}
    \caption{
        Fit results for the two SSRs performed at $[H, K] = [0.5, -1]$ and $[H, K] = [1, -1.5]$.
        $\mu$, and $\sigma$ are the Bragg peak positions and standard deviations in $L$.
        The different peaks were indexed using a distorted cubic structure, corresponding to bulk \ce{Pt_3O_4}, $d$ are the corresponding interplanar spacings.
    }
    \label{fig:FitPt100LScans}
\end{figure}

\ce{Pt_3O_4} has a cubic structure (space group $P_{m\bar{3}n}$) with a lattice parameter equal to \qty{5.65}{\angstrom} \parencite{Galloni1941, Galloni1952, Muller1968}, presented in fig. \ref{fig:Pt3O4}.
Recent theoretical studies have proposed a lattice parameter equal to \qty{5.59}{\angstrom}, (tab. \ref{tab:PtOxides}, Seriani et al. \cite*{Seriani2006}), which is also close to the values reported in this study.
Moreover, the (212) and (232) reflections are not allowed for the proposed due to the symmetry of the \ce{Pt_3O_4} structure.
Those reflections are indeed absent from fig. \ref{fig:FitPt100LScans}.

\begin{SCfigure}
    \centering
    \includegraphics[trim=0 1cm 0 1cm, clip, width=0.35\textwidth]{/home/david/Documents/PhD/Figures/Introduction/Pt3O4.pdf}
    \caption{
        \ce{Pt_3O_4} bulk unit cell.
        Platinum atoms are situated on the faces on the cubic unit cell (e.g. $(0, 1/2, 1/4)$, $(0, 1/2, 3/4)$), while the eight oxygen atoms are inside the unit cell at the positions $(1/4, 1/4, z)$, $(1/4, 2/4, z)$, $(2/4, 1/4, z)$, $(2/4, 2/4, z)$ for $z=1/4$ and $z=3/4$.
    }
    \label{fig:Pt3O4}
\end{SCfigure}

The lattice parameter of \ce{Pt_3O_4} is approximately twice that of the Pt(100) surface lattice parameter, which allows the formation of a cubic on cubic coherent interface.
The nature of the atoms at the interface is addressed later in this chapter.
To match perfectly with the Pt(100) surface, the in-plane lattice parameter of \ce{Pt_3O_4} would need to reach the value of $2\times a_{100} =\qty{5.55}{\angstrom}$, which is lower than the reported literature values for the cubic unit cell.
This interfacial strain can be the reason for the distorted structure observed in this study.
The possibility to perform out-of-plane measurements with SXRD allows us to link the Pt(100)-($2\times2$) reconstruction with certainty to bulk \ce{Pt_3O_4}, which exhibits a slightly distorted lattice.

A first estimate of the oxide layer thickness can be obtained with the full-width at half maxima of the Bragg peaks in $q_z$: $FWHM_{q_z}$, and using the equation $2\pi/FWHM_{q_z}$ \parencite{Patterson1939, Warren1990}.
This yields a thickness of \qty{\approx 16.3}{\angstrom}, which for bulk \ce{Pt_3O_4} would correspond to approximately \num{3} unit cells.

\subsection{Near ambient pressure ammonia oxidation cycle}

\begin{figure}[!htb]
    \centering
    \includegraphics[width=0.495\textwidth]{/home/david/Documents/PhD/PhDScripts/SixS_2022_01_SXRD_Pt100/figures/Map_hkl_surf_or_1880-1902_patched.pdf}
    \includegraphics[width=0.495\textwidth]{/home/david/Documents/PhD/PhDScripts/SixS_2022_01_SXRD_Pt100/figures/Map_hkl_surf_or_1930-1936_patched.pdf}
    \caption{
        Reciprocal space in-plane maps collected under different atmospheres at \qty{450}{\degreeCelsius}, computed using the surface lattice of Pt(100).
        The heater has failed during the reacting conditions (a), the sample was then cleaned and put again under oxygen atmosphere (b).
    }
    \label{fig:MapsPt100B}
\end{figure}

The combination of high temperature, oxygen and ammonia define a very harsh environment, that can result in the oxidation of the screws responsible for the contact on the sample holder.
Fig \ref{fig:MapsPt100B} (a) shows a map in which the heating system failed after the introduction of ammonia in the reactor, reaching reacting conditions.
Nevertheless, half of the large reciprocal space in-plane maps could be measured, the red circled peaks corresponding to \ce{Pt_3O_4} are still visible, showing that \ce{Pt_3O_4} is not immediately removed during reacting condition.
The two other peaks at [0.93, -0.5] and [0.5, -0.93] have disappeared, the related structure is not stable during the reaction.

After repair of the sample heater, the sample was cleaned by sputtering and annealing cycles, and replaced in the reactor.
In order to probe for the reproducibility of the \ce{Pt_3O_4} growth on the surface, the partial pressure of oxygen was set to \qty{80}{\milli\bar}, while measuring a small area of the reciprocal space to detect the same peaks (fig. \ref{fig:MapsPt100A} - b).
Such peaks could indeed be detected, but with a shift between the green and red groups of peaks equal to $0.09$, more important than the previously measured value of $0.07$.
The second structure is possibly linked to the exposure to oxygen atmosphere, the smaller in-plane map was collected faster, and directly after the introduction of oxygen in the reactor.

Ammonia was subsequently introduced in the cell to probe the relation between surface structure and selectivity during ammonia oxidation.

\begin{figure}[!htb]
    \centering
    \includegraphics[width=0.53\textwidth]{/home/david/Documents/PhD/PhDScripts/SixS_2022_01_SXRD_Pt100/figures/Map_hkl_surf_or_1953-1981_patched.pdf}
    \includegraphics[width=0.46\textwidth]{/home/david/Documents/PhD/PhDScripts/SixS_2022_01_SXRD_Pt100/figures/l_scans_high_oxygen_ammonia_no_map.pdf}
    \caption{
        (a) Reciprocal space in-plane maps collected under different atmospheres at \qty{450}{\degreeCelsius}, computed using the surface lattice of Pt(100).
        (b) SSR for four different positions in [H, K] under \qty{80}{\milli\bar} of \ce{O_2}, \qty{10}{\milli\bar} of \ce{NH_3}, and \qty{410}{\milli\bar} of \ce{Ar}.
    }
    \label{fig:MapsAndLScansPt100HighOxAmmonia}
\end{figure}

Signals separated by $0.1$ in H or in K were observed around the platinum Bragg peaks in a square arrangement, compatible with a Pt(100)-($10\times10$) surface reconstruction.
Most of the peaks are observed at $H={0,1,2}$ or $K={0, -1}$ (fig. \ref{fig:MapsAndLScansPt100HighOxAmmonia} - a).
The other features are observed in rows or columns \textit{shifted} by $0.09$ with respect to the (1, $\bar{1}$, 0) Bragg peak, following the shift of the green circled peaks observed while only oxygen is present in the cell (fig. \ref{fig:MapsPt100B} - b).

SSRs at [H, K] = [1.9, 0], [H, K] = [0.5,  $\bar{1}$], and [H, K] = [0, -1.2] are performed on the Pt(100)-($10\times10$) reconstructions (fig. \ref{fig:MapsAndLScansPt100HighOxAmmonia} - b).
The peaks measured earlier at [H, K] = [0.5, $\bar{1}$] at $L=0.7$, $L=1.4$, and $L=2.1$ corresponding to bulk \ce{Pt_3O_4} are no longer visible.
For [H, K] = [1.9, 0], two peaks are visible near $L=0.15$ and $L=2.15$.
The reconstructed signal effectively disappears near $L=1.9$, but reappears with a strong intensity at $L=2$, following the same decrease as a function of $L$ as from $L=0$.
For [H, K] = [0, -1.2], large oscillations are visible with a minimum at $L=1.2$, which could be the signature of a multi-layer structure (e.g. simulated \ce{Pt_3O4} in fig. \ref{fig:SimROD} - a).
Large oscillations are also visible for [H, K] = [0.5, -1], with a minimum between $L=1$ and $L=1.7$, and local maxima at $L=0.4$, and $L=2.5$.

The SSR at [H, K] = [0.5, -0.91] (fig. \ref{fig:MapsAndLScansPt100HighOxAmmonia} - b) is performed on the shifted Pt(100)-($10\times10$) reconstruction.
The intensity decreases gradually with increasing $L$, characteristic of a monolayer, no signal can be detected above $L=2$.
Note that such an intensity decrease was also observed for the shifted peak at [H, K] = [0.5, -0.93] in the previous condition (fig. \ref{fig:FitPt100LScans}).

It is not certain that the same thickness of \ce{Pt_3O_4} was present on the catalyst surface at the beginning of the reacting conditions.
Indeed, the sample was exposed to \qty{80}{\milli\bar} of oxygen for \qty{6}{\hour} before SSRs could be measured.
The second exposition only lasted for \qty{1}{\hour}, during which the same in-plane peaks could effectively be measured, but without out-of-plane information.
The thickness of \ce{Pt_3O_4} is very likely to be related to the time spent under \qty{80}{\milli\bar} of \ce{O_2}.
However, it is clear that i) new signals appear during reacting conditions, related to structures that themselves differ depending on the thickness of \ce{Pt_3O_4}, and ii) the bulk oxide cannot grow under reacting conditions.
Therefore, different types of structures coexist on the Pt(100) surface when the \ce{O_2}/\ce{NH_3} ratio is equal to \num{8}, possibly linked to different adsorbed species participating in the catalytic reaction.

\begin{figure}[!htb]
    \centering
    \includegraphics[width=0.53\textwidth]{/home/david/Documents/PhD/PhDScripts/SixS_2022_01_SXRD_Pt100/figures/Map_hkl_surf_or_2227-2283_patched.pdf}
    \includegraphics[width=0.46\textwidth]{/home/david/Documents/PhD/PhDScripts/SixS_2022_01_SXRD_Pt100/figures/l_scans_low_oxygen_ammonia.pdf}
    \caption{
        (a) Reciprocal space in-plane maps collected under different atmospheres at \qty{450}{\degreeCelsius}, computed using the surface lattice of Pt(100).
        (b) SSRs under \qty{5}{\milli\bar} of \ce{O_2}, \qty{10}{\milli\bar} of \ce{NH_3}, and \qty{485}{\milli\bar} of \ce{Ar}.
    }
    \label{fig:MapsAndLScansPt100LowOxAmmonia}
\end{figure}

Lowering the partial pressure of oxygen from \qty{80}{\milli\bar} to \qty{5}{\milli\bar} in the reactor has completely removed the Pt(100)-($10\times10$) reconstruction (fig. \ref{fig:MapsAndLScansPt100LowOxAmmonia} - a).

Co-existing domains with the same hexagonal lattice on top on the Pt(100) surface are revealed, with an in-plane lattice parameter equal to \qty{2.68}{\angstrom}, \qty{\approx 3.36}{\percent} lower than the distance between neighbouring Pt atoms on the Pt(100) surface (\qty{2.77}{\angstrom}).
Second order peaks can also be seen at the borders of the reciprocal space in-plane map, each domain has one axis parallel to either $\vec{a}_{(100)}$ or $\vec{b}_{(100)}$.
%, \textit{i.e.} respectively in the [011] and [01$\bar{1}$] directions.
The hexagonal structures can be described as follows: Pt(100)-p$\begin{pmatrix} 0.96 & 0.00\\ -0.48 & 0.83 \end{pmatrix}$.
% , the closest superstructure is Pt(100)-(33x25).
For simplicity, this structure will be named Pt(100)-Hex in the rest of this thesis.
% A sketch describing the real space unit cell can be found in appendix \ref{fig:Pt100UnitCellsReconstruction}.
% precursor to bulk is chemisorbed \parencite{BradleyShumbera2007}

Different hexagonal surface structures on the Pt(100) surface have been reported in literature, in the same directions, and under ultra-high vacuum.
This has been summarised by Hammer et al. \parencite*{Hammer2016}, based on an important body of work \parencite{Heilmann1979, Vanhove1981, Heinz1982, Mase1992, Kuhnke1992, Borg1994, VanBeurden2004, Havu2010}.
Hexagonal structures have been measured to evolve to rotated hexagonal reconstructions with angles between \ang{0.77} and \ang{0.94} depending on the sample temperature, and the previous temperature treatment.
The unit cell describing those reconstructions with respect to the Pt(100) surface varies, if first contained to (5XN) where N = 20–30, the latest study reports a c(26x118) superstructure.
Exposition of the rotated hexagonal structures to oxygen at \qty{450}{\degreeCelsius} studied by low energy electron diffraction (LEED) removes the hexagonal structures, and precipitate the growth of surface oxides \textit{via} different phases \parencite{BradleyShumbera2007, BradleyShumbera2007a}.
A similar conclusion was reached by \cite{Deskins2005} by DFT studies.
Exposition to \ce{NO} has been reported to stabilise the clean Pt(100) phase \parencite{Heinz1982}, while exposition to \ce{CO} removes the hexagonal reconstruction.
An oscillatory behaviour between a clean Pt(100)-(1x1) surface and the rotated hexagonal reconstruction was reported by Cox et al. \parencite*{Cox1983} during CO oxidation.

Subsurface oxygen has also been predicted to exist on Pt(100) \parencite{Gu2007}, reported at a pressure of \qty{0.133}{\milli\bar} \parencite{McMillan2005}, and participating in the catalytic oxidation of \ce{CO}.
It was identified as a precursor to a stable surface oxide during the catalytic oxidation of \ce{CO} by Dicke et al. \parencite*{Dicke2000}, at an oxygen pressure of \qty{0.09}{\milli\bar}.
The appearance of subsurface oxygen was linked to the lifting of surface reconstructions on the clean Pt(100) surface from the adsorption of \ce{CO}, which then allowed oxygen atoms to penetrate under the topmost layer of platinum \parencite{Rotermund1993, Lauterbach1994}.

Overall, few works at ambient pressure have been found to exist, even less during the catalytic oxidation of ammonia.
SSRs were collected at the position of the two peaks (fig. \ref{fig:MapsAndLScansPt100LowOxAmmonia} - b), no bulk structure could be detected.
A different behaviour is measured when lowering the \ce{O_2}/\ce{NH_3} ratio from \num{8} to \num{0.5}, which also has an impact on the product selectivity.
The partial pressure of \ce{NO} and \ce{N_2O} decreases quickly after the change of condition (fig. \ref{fig:RGA450Pt100Cycle}), whereas the partial pressure of \ce{N_2} is more stable.

\begin{figure}[!htb]
    \centering
    \includegraphics[width=0.495\textwidth]{/home/david/Documents/PhD/PhDScripts/SixS_2022_01_SXRD_Pt100/figures/Map_hkl_surf_or_2520-2570_patched.pdf}
    \includegraphics[width=0.495\textwidth]{/home/david/Documents/PhD/PhDScripts/SixS_2022_01_SXRD_Pt100/figures/Map_hkl_surf_or_2719-2767.pdf}
    \caption{
        Reciprocal space in-plane maps collected under different atmospheres at \qty{450}{\degreeCelsius}, computed using the surface lattice of Pt(100).
    }
    \label{fig:MapsPt100C}
\end{figure}

After removing oxygen from the reactor, only the first order reflections corresponding to the hexagonal structures could be detected during the measurement of the reciprocal space in-plane map (fig. \ref{fig:MapsPt100C} - a), the following measurement under Argon showed a clean Pt(100) surface (fig. \ref{fig:MapsPt100C} - b).

\subsection{Oxide growth under 5 mbar of oxygen}

\begin{figure}[!htb]
    \centering
    \includegraphics[width=0.495\textwidth]{/home/david/Documents/PhD/PhDScripts/SixS_2022_01_SXRD_Pt100/figures/Map_hkl_surf_or_2905-2953_patched.pdf}
    \includegraphics[width=0.495\textwidth]{/home/david/Documents/PhD/PhDScripts/SixS_2022_01_SXRD_Pt100/figures/Map_hkl_surf_or_3154-3169.pdf}
    \caption{
        Reciprocal space in-plane maps collected under different atmospheres at \qty{450}{\degreeCelsius}, computed using the surface lattice of Pt(100).
        The area covered in (b) is drawn by a black rectangle in (a).
    }
    \label{fig:MapsPt100D}
\end{figure}

To make sure that the hexagonal structures are related to the reaction conditions, and not only to a lower pressure of oxygen in the reactor, a partial pressure of \qty{5}{\milli\bar} of oxygen was set.
The following measurement showed the presence of many peaks that appeared in the first hour of the measurement (fig. \ref{fig:MapsPt100D} - a).
A second map was measured directly after the end of the first map, in which none of the newly detected peaks could be observed (fig. \ref{fig:MapsPt100D} - b).

These two measurements showed first that the hexagonal structures observed under reacting conditions when the \ce{O_2}/\ce{NH_3} ratio is equal to \num{0.5} are linked to the \textit{simultaneous} presence of ammonia and oxygen in the reactor.
Secondly, the dynamics of some of the observed phenomena under \qty{5}{\milli\bar} of oxygen are too short to be effectively measured with the current time-resolution of in-plane reciprocal space maps.
Weak signals in the H and K directions can be seen in the second map, starting from the (1, $\bar{1}$, 0) Bragg peak (fig. \ref{fig:MapsPt100D} - b).

\subsection{Surface roughness and surface relaxation effects}

\begin{figure}[!htb]
    \centering
    \includegraphics[width=\textwidth]{/home/david/Documents/PhD/PhDScripts/SixS_2022_01_SXRD_Pt100/figures/ctr_a.pdf}
    \includegraphics[width=\textwidth]{/home/david/Documents/PhD/PhDScripts/SixS_2022_01_SXRD_Pt100/figures/ctr_b.pdf}
    \caption{
        Evolution of crystal truncation rods under different atmospheres.
    }
    \label{fig:CTRPt100}
\end{figure}

The (20L), ($1\bar{1}L$) and (0$\bar{1}L$) crystal truncation rods have been measured \qty{6}{\hour} after the start of each condition, each measurement lasted for \qty{2}{\hour}.
The background-subtracted intensity of the CTR was integrated using the \textit{fitaid} module of \textit{BINoculars} as a function of $L$ with the same integration range.
The CTR intensity is displayed in fig. \ref{fig:CTRPt100}, additional peaks are detected under exposition to \qty{80}{\milli\bar} of \ce{O_2}, confirming that bulk \ce{Pt_3O_4} is epitaxied on the Pt(100) surface following a Pt(100)-($2\times2$) arrangement.

Those peaks prevent us from resolving the minimal position of the CTR intensity, highly sensitive to surface relaxation effects.
The absence of oscillations in the evolution of the CTR intensity as a function of $L$ also shows that \ce{Pt_3O_4} is not homogeneously covering the Pt(100) substrate (such signal is presented in fig. \ref{fig:SimROD} - b).
This phenomena is here explained by \ce{Pt_3O_4} islands of different thickness covering the substrate, the fringes linked to islands of same thickness being smeared out from the contribution of others.

The catalyst surface can be divided in three different components as a function of $\vec{z}$.
First, bulk Pt(100).
Secondly, the interface between bulk platinum and the different islands of \ce{Pt_3O_4}, in which the platinum atoms are expected to be displaced from their equilibrium positions due to the interaction with the oxide layer.
% , and where there could also be the presence of some subsurface oxygen atoms.
Thirdly, the \ce{Pt_3O_4} islands, on top of the previous layers, with different thickness.
A first approach to approximate the catalyst surface is discussed below.

\begin{figure}[!htb]
    \centering
    \includegraphics[width=\textwidth]{/home/david/Documents/PhD/PhDScripts/SixS_2022_01_SXRD_Pt100/figures/fit_8o2.pdf}
    \caption{
        Fitting results for crystal truncation rods collected under a \qty{80}{\milli\bar} of \ce{O_2}.
        \ce{Pt_3O_4} Bragg peaks are indicated on the figure.
    }
    \label{fig:CTRFitHighOxygen}
\end{figure}

The structure factors $F_i$ resulting from the presence of one to nine unit cells of \ce{Pt_3O_4} on the Pt(100) surface was simulated with \textit{ROD}, the same interface with the substrate is used in each simulation.
The fitting routine consists in minimising the square root difference between the CTR structure factors $F_{obs}$, and the square root of the squared simulated structure factors sum by adjusting the weight $W_i$ of each signal in the total signal $F_{calc}$ (eq. \ref{eq:Fcalc}).

\begin{equation}
    F_{calc} = \sqrt{\sum_{i=1}^{9} W_i F_i^2}
    \label{eq:Fcalc}
\end{equation}

% Structure factors resulting from different domains are not coherently interfere since the beam in grazing incidence, the beam the distance between each domain may be larger than the beam coherence lengths.

In this first hypothesis, the topmost Pt(100) layers were assumed to be strain free.
Different $\beta$ roughness parameters (explained in fig. \ref{fig:CTRSimulation}) were used, the best fit is with $\beta = \num{0.6}$.
The $\beta$ roughness parameter only lowers the intensity of the bulk signal and not that of the surface layers.
Moreover, it is assumed that \ce{Pt_3O_4} islands have the same out-of-plane and in-plane lattice parameter in each layer, independently of the island thickness.
For example, the 7$^{th}$ unit cell in a \num{8} unit cell thick island has the same lattice parameters than the second unit cell in a \num{4} unit cell thick island.

The relation between each weight was adjusted by following a Gaussian distribution, expecting the different islands to not differ too much in thickness.
The result of the fitting routine is shown in fig. \ref{fig:CTRFitHighOxygen}, with a Gaussian distribution centred around $0$ unit cells, and with a standard deviation $\sigma$ equal to $2$ unit cells.
The presence and width of the \ce{Pt_3O_4} peaks is well adjusted, see for example the presence of the (405) \ce{Pt_3O_4} Bragg peak in the [2, 0, L] CTR.
The intensity shape near the platinum Bragg peak is also quite accurate, but tends to fall too quickly.
This can be adjusted by lowering the $\beta$ roughness, but low intensity regions then become impossible to represent.
Indeed, the simulation struggles to accurately reproduce the intensity when below a certain threshold.
The position of some peaks is very well reproduced, but some peaks are a little shifted, such as the (402) and (2$\bar{2}2$) \ce{Pt_3O_4} Bragg peaks in the [2, 0, L], and [1, $\bar{1}$, L] CTRs.
The ratio of intensity between the \ce{Pt_3O_4} peaks is also not perfect, and could possibly be further adjusted by changing the type of atoms present at the interface, as well as the distance between the oxide and surface.
The positions of the Pt atoms at the interface are explained by a sketch in fig. \ref{fig:Pt3O4onPt100}.

The distance between the Pt atoms in the Pt(100) bulk and \ce{Pt_3O_4} was chosen equal to \num{0.5} unit cell lattice length, following the spacing between Pt atoms in \{100\} planes.
The fitting result may be further improved by adjusting this parameter, as well as the interfacial strain.
Overall, the existence of bulk \ce{Pt_3O_4} islands is confirmed to remove the oscillations visible when a homogeneous layer is present of the sample surface (fig. \ref{fig:SimROD}).
Moreover, the average thickness retrieved by fitting the SSRs in fig. \ref{fig:LScansHighOxygenPt100}, \qty{\approx16.3}{\angstrom}, \textit{i.e.} 3 unit cells, is in good agreement with the mean value and standard deviation of the Gaussian distribution returned here.

\begin{figure}[!htb]
    \centering
    \includegraphics[width=\textwidth]{/home/david/Documents/PhD/Figures/Introduction/Pt3O4onPt100.pdf}
    \caption{
    Epitaxy relationship used when simulating structure factors for \ce{Pt_3O_4} on Pt(100) in a Pt(100)-($2\times2$) arrangement.
    }
    \label{fig:Pt3O4onPt100}
\end{figure}

When observing the CTR intensity under other atmospheres, a clear evolution in the position of the minimum intensity between both reacting conditions (in red and green in fig. \ref{fig:CTRPt100}) is visible.
The CTR recorded after exposition to reacting conditions under ammonia or under argon have similar intensities and are almost indistinguishable.

\begin{figure}[!htb]
    \centering
    \includegraphics[width=\textwidth]{/home/david/Documents/PhD/PhDScripts/SixS_2022_01_SXRD_Pt100/figures/fit_comparison.pdf}
    \caption{
        Fitting results for roughness parameter $\beta$ (a) and out-of-plane strain $\sigma_z$ (b) as a function of the experimental conditions.
        The innermost layer is named $A$, which has the same lattice parameter as in the bulk.
        The second and first topmost layers of the Pt(100) single crystal are respectively named $B$ and $C$.
    }
    \label{fig:CTRFit100}
\end{figure}

All three CTRs were fitted together using \textit{ROD}, at each condition besides the oxygen rich atmosphere.
Different models were tested, adding the possibility of in-plane lattice displacement did not show any improvement of the fit quality, and a simple surface model was kept consisting of two Pt(100) layers, each sharing an out-of-plane lattice displacement parameter, on top of bulk Pt(100).
The evolution of the surface roughness and out-of-plane strain is shown in fig. \ref{fig:CTRFit100}, the $C$ layer is the topmost layer, the $A$ layer is in the bulk.

The surface roughness evaluated \textit{via} the $\beta$ parameter in fig. \ref{fig:CTRFit100} (a) is consistent with the evolution of the CTR intensity in fig. \ref{fig:CTRPt100}.
The initial value for the surface roughness is equal to \num{0.44} when exposed to a total pressure of argon equal to \qty{500}{\milli\bar}.
The average roughness after introducing a partial pressure of \qty{80}{\milli\bar} of oxygen was estimated to increase to \num{0.6}, coherent with the creation of a bulk oxide layer on the sample surface.
The minimal intensity can be seen to decrease after introduction of oxygen, see for example at $L=1$ for the [$\bar{1}$, $\bar{1}$, L] CTR in fig. \ref{fig:CTRPt100}.

The following crystal truncation rods having been measured after the second introduction of ammonia in the cell, \textit{i.e.} after sample cleaning and short exposition to high oxygen atmosphere.
The roughness under reacting conditions, when observing the Pt(100)-($10\times10$) reconstructions (fig. \ref{fig:MapsAndLScansPt100HighOxAmmonia}), decreases only slightly in comparison to having only oxygen in the cell, and is still more important than under inert atmosphere.
Lowering the amount of oxygen in the cell decreases the surface roughness, coming back to the initial value.
In this condition, the hexagonal structures are observed (fig. \ref{fig:MapsAndLScansPt100LowOxAmmonia}).

The lowest roughness value is reached when only ammonia is present in the reactor, without oxygen.
The hexagonal structures slowly disappear (fig. \ref{fig:MapsPt100C}), their presence linked to increased roughness on the sample surface.
Removing ammonia does not change the surface roughness, which shows that oxygen plays a key role in the increased surface roughness.

Finally, the presence of \qty{5}{\milli\bar} of oxygen after the ammonia oxidation cycle increases the surface roughness to higher values, almost equal to the maximum value reached under the presence of \qty{80}{\milli\bar} of oxygen.
Bulk \ce{Pt_3O_4} was not observed in this condition, small peaks were visible linked to transient structures (fig. \ref{fig:MapsPt100D}).

% strain
The strain of the second topmost layer, ($B$ in fig. \ref{fig:CTRFit100} - b) is always very close to \qty{0}{\percent}.
The strain of the topmost layer ($C$ in fig. \ref{fig:CTRFit100} - b) is always compressive with respect to bulk parameters, and seen to increase during reacting conditions in comparison with the initial values under argon atmosphere.
A slight increase is reported when changing the \ce{O_2}/\ce{NH_3} ratio from \num{8} to \num{0.5}.
Keeping only ammonia in the reactor after the oxidation reaction brings the lowest strain observed, both layers almost exhibit the bulk lattice parameters.

The presence of ammonia can be associated with reduced roughness and strain, whereas the presence of oxygen is associated with increased roughness and strain.
Different reacting conditions, which can be linked to a different selectivity in the product, show a larger difference in the surface roughness than in the surface strain.
For example, when observing the [2, 0, L] CTR in fig. \ref{fig:CTRFit100}, the local minima positions do not show a significant change.
However, the change in intensity indicates different surface roughness.
The shape of the [0, $\bar{1}$, L] CTR when the \ce{O_2}/\ce{NH_3} ratio is equal to \num{0.5} is interesting (fig. \ref{fig:CTRPt100}).
The intensity minimum at $L=2.1$ is very low, whereas the decrease of intensity around the Bragg peaks is not very strong, in comparison with the same CTR but when the \ce{O_2}/\ce{NH_3} ratio is equal to \num{8}.
It is possible that low intensity peaks are present near $L=1.8$ or $L=2.5$ that would explain this shape, linked to the hexagonal structures observed in fig. \ref{fig:MapsAndLScansPt100LowOxAmmonia}.

\subsection{Surface species}

As for the Pt(111), the Pt(100) crystal was measured with NAP-XPS at B07 beamline.
The conditions have been resumed in tab. \ref{tab:ConditionsXPS}.
The RGA signals are presented in fig. \ref{fig:XPS100RGA}.

\begin{figure}[!htb]
    \centering
    \includegraphics[width=\textwidth]{/home/david/Documents/PhD/PhDScripts/B07_2022_04_XPS/Figures/pt100_time.pdf}
    \caption{
        Evolution of reaction product partial pressures as a function of time during the XPS experiment on the Pt(100) single crystals at \qty{450}{\degreeCelsius}.
        Transition between conditions are indicated with dashed vertical lines.
    }
    \label{fig:XPS100RGA}
\end{figure}

% Describe rga
The high \ce{O_2}/\ce{NH_3} ratio favours the production of \ce{NO} as expected, accompanied by water.
With the gas probe near the sample surface, it is possible to detect a different selectivity between the two surfaces: for the high oxygen condition Pt(100) produces mostly NO, whereas \ce{N_2} was still detected for Pt(111) (fig. \ref{fig:XPS111RGA}).
% No \ce{N_2} or \ce{N_2O} can be measured.

Lowering the amount of oxygen by reducing the \ce{O_2}/\ce{NH_3} ratio to \num{0.5} has the remarkable effect of shifting the reaction selectivity entirely towards \ce{N_2}, just like for the Pt(111) surface.
Oxygen being undetected by the mass spectrometer, all of the available oxygen at the surface dissociates on the catalyst surface and participates in the production of \ce{N_2} and \ce{H_2O} \textit{via} the oxidation reaction.
\ce{H_2} coming from the simultaneous dissociation of \ce{NH_3} can also be measured, not observed under a higher pressure of oxygen.
Thus, the more oxygen is present in the reactor, the less the dissociation of \ce{NH_3} is favoured.
Ammonia can be considered to be in excess, and partly decomposing towards \ce{N_2}.

Under pure ammonia there is still production of \ce{N_2}, with an increase in \ce{H_2} production and essentially no \ce{H_2O} yield.

\subsubsection{N 1s and O 1s levels}

The evolution of the N 1s and O 1s XPS spectra for different atmospheres is presented in fig. \ref{fig:O1sN1sPt100}.
Binding energy are given with reference to the Fermi level, all the reported peaks and corresponding species are detailed in tab. \ref{tab:XPSPt100}.

During the study of the Pt(100) oxidation by Derry et al. \parencite*{Derry1984}, O 1s peaks between \qty{530.4}{\eV} and \qty{530.5}{\eV} are reported as a function of the oxygen surface coverage.
Sugai et al \parencite*{Sugai1993} report a low intensity peak at \qty{531.5}{\eV} during the dissociation of \ce{NO} at \qty{400}{\degreeCelsius}, which they cannot specifically assign to either \ce{NO_a} or \ce{O_a}.
Different types of adsorbed oxygen species are observed during the dissociation of \ce{NO} by Rienks et al. \parencite*{Rienks2003} at \qty{527.7}{\eV} and \qty{529.7}{\eV}, the latter only shortly distant from \ce{NO_a} reported at \qty{529.8}{\eV}.
Kondratenko et al. \parencite*{Kondratenko2006} report \ce{O_a} and \ce{OH_a} at \qty{529.7}{\eV} and \qty{532}{\eV} respectively during the decomposition of \ce{N_2O} after exposition to \ce{H_2}, based on previous works by Wild et al. \parencite*{Wild2000}.

\begin{figure}[!htb]
    \centering
    \includegraphics[width=\textwidth]{/home/david/Documents/PhD/PhDScripts/B07_2022_04_XPS/Figures/Pt100/O1sN1s_700.pdf}
    \caption{
        XPS spectra collected at the O 1s (a) and N1 s (b) levels under different atmospheres at \qty{450}{\degreeCelsius} with an incoming photon energy of \qty{700}{\eV}.
        The spectra are normalised by the pre-edge intensity and shifted to highlight the presence of different peaks.
    }
    \label{fig:O1sN1sPt100}
\end{figure}
\begin{table}[!htb]
\centering
\resizebox{\textwidth}{!}{%
    \begin{tabular}{@{}ll|lllllll@{}}
    \toprule
    \multirow{3}{*}{Partial pressures (\unit{\milli\bar})} & \ce{Ar}   & 1 & 0   & 0   & 0    & 0   & 0    & 1    \\
                                              & \ce{NH_3} & 0 & 0   & 1.1 & 1.1  & 1.1 & 0    & 0    \\
                                              & \ce{O_2}  & 0 & 8.8 & 8.8 & 0.55 & 0   & 0.55 & 0    \\
    \midrule
    Gas presence (decreasing & & Ar & \ce{O_2} & \ce{O_2}, \ce{H_2O}, \ce{NO}   & \ce{H_2O}, \ce{NH_3} & \ce{H_2}, \ce{NH_3} & \ce{O_2} & Ar \\
    pressure order)          & &    &          & \ce{NH_3}, \ce{N_2}, \ce{N_2O} & \ce{N_2}, \ce{H_2}   & \ce{N_2}            &          &    \\
    \midrule
    \multicolumn{2}{l|}{N 1s: peak positions}
        & \qty{403.6}{\eV} & No peak          & \qty{404.3}{\eV} & \qty{404.3}{\eV} & \qty{404.9}{\eV} & No peak          & No peak        \\
     &  &                  &                  &                  & \qty{400.0}{\eV} & \qty{400.6}{\eV} &                  &                \\
    \multicolumn{2}{l|}{Attributed surface species}
        & Not assigned     &                  & \ce{N_{2,g}} \& \ce{NO_g} & \ce{N_{2,g}}     & \ce{N_{2,g}}     &                  &       \\
     &  &                  &                  &                           & \ce{N_a}         & \ce{NH_{x,a}}    &                  &       \\
    \midrule
    \multicolumn{2}{l|}{O 1s: peak positions}
        & \qty{531.6}{\eV} & \qty{538.3}{\eV} & \qty{538.4}{\eV} & \qty{534.3}{\eV} & \qty{532.4}{\eV} & \qty{538.3}{\eV} & \qty{532}{\eV} \\
     &  &                  & \qty{537.2}{\eV} & \qty{537.3}{\eV} & \qty{532.1}{\eV} &                  & \qty{537.2}{\eV} &                \\
     &  &                  & \qty{529.7}{\eV} & \qty{533.7}{\eV} &                  &                  & \qty{531.4}{\eV} &                \\
     &  &                  &                  & \qty{529.7}{\eV} &                  &                  & \qty{530.6}{\eV} &                \\
     &  &                  &                  &                  &                  &                  & \qty{529.7}{\eV} &                \\
    \multicolumn{2}{l|}{Attributed surface species}
        & \ce{H_2O_a}      & \ce{O_{2,g}}     & \ce{O_{2,g}}     & \ce{H_2O_g}      & \ce{H_2O_a}      & \ce{O_{2,g}}     & \ce{H_2O_a}    \\
     &  &                  & \ce{O_{2,g}}     & \ce{O_{2,g}}     & \ce{H_2O_a}      &                  & \ce{O_{2,g}}     &                \\
     &  &                  & \ce{O_a}         & \ce{H_2O_g}      &                  &                  & \ce{O_a}         &                \\
     &  &                  &                  & \ce{O_a}         &                  &                  & \ce{O_a}         &                \\
     &  &                  &                  &                  &                  &                  & \ce{O_a}         &                \\
    \bottomrule
    \end{tabular}%
    }
    \caption{Indexing of peaks measured during ammonia oxidation of the Pt(100) surface.}
\label{tab:XPSPt100}
\end{table}

% argon

% high ox
The Pt(100) surface was exposed to \qty{8.8}{\milli\bar} of oxygen after argon.
The peaks at \qty{538.3}{\eV} and \qty{537.2}{\eV} can be attributed to \ce{O_{2,g}}.
As observed for Pt(111), the positions are shifted in energy with respect to literature \parencite{Avval2022}.
Oxide signals were measured after \qty{1}{\hour} under \qty{80}{\milli\bar} of oxygen with SXRD (fig. \ref{fig:MapsPt100B}), but no stable surface of bulk oxide was observed at an oxygen pressure of \qty{5}{\milli\bar}.
Some transient peaks were measured but disappeared after a few hours of measurements (fig. \ref{fig:MapsPt100D}) .
By comparing with the previously reported literature values, the large peak at \qty{529.7}{\eV} is attributed to adsorbed oxygen.
It is possible that the oxygen atoms linked to the transient surface structure yield some intensity around \qty{530.5}{\eV}, where a small region of higher intensity is observed.
No peak near \qty{527.7}{\eV} was observed throughout the experiment.

% ratio 8
Adding ammonia in the cell brings us to reacting conditions with \ce{O_2}/\ce{NH_3} = \num{8}.
Interestingly, this does not remove the \ce{O_a} peak in the O 1s level, but introduces a gas phase water (\ce{H_2O_g}) peak at \qty{533.7}{\eV} \parencite{Weststrate2006, Linford2019}, probably generated by the oxidation reaction.
The \ce{O_a} peak has a lower intensity, which can be linked to the oxidation reaction utilising oxygen to produce water and nitric oxide.
A \qty{1}{\eV} wide peak at \qty{404.3}{\eV} is reported in the N 1s level, linked to \ce{NO_{g}} and possibly to \ce{N_{2, g}}.
\ce{NO_{g}} is expected between \qty{404.5}{\eV} and \qty{406.7}{\eV}, while \ce{N_{2,g}} between \qty{403.9}{\eV} and \qty{404.8}{\eV} \parencite{Ivashenko2021}, which could explain the broad peak shape.
% The fact that the partial pressure of \ce{NO} is higher than for Pt(111) tends to support this hypothesis.
Since gas phase peaks are shifting when changing the partial pressure of oxygen, it is difficult to be certain of the peak nature.
Most importantly, the presence of the peak at \qty{529.7}{\eV} shows that oxygen is adsorbed on Pt(100), even during reacting conditions, which was not as clear for Pt(111) (fig. \ref{fig:O1sN1sPt111}).
Novell-Leruth et al. \parencite*{NovellLeruth2008} have reported that the de-hydrogenation process of ammonia is favoured to occur \textit{via} \ce{O_a} rather than \ce{OH_a} on Pt(100) in comparison with Pt(111).
No peak could be linked to \ce{OH_a}, reported at \qty{532}{\eV} on Pt(100).
It is possible that a second peak exist near \qty{530.5}{\eV}, linked also to adsorbed oxygen as reported by Derry et al. \parencite*{Derry1984}, possibly on different sites.
This differs from the Pt(111) surface for which both \ce{O_a} and \ce{OH_a} could be clearly identified, and with similar intensity.
This difference in chemisorbed oxygen presence could explain the Pt(100) surface higher selectivity towards the production of \ce{NO}.
All the \ce{NH_x} species are probably quickly oxidised towards atomic nitrogen, which then reacts with a neighbouring oxygen to form \ce{NO_a}.
Adsorbed \ce{NO} at \qty{401.3}{\eV} was not observed during the dissociation of \ce{NO} on Pt(100) above room temperature by XPS \parencite{Rienks2003}.
\ce{NO} desorbs too quickly from the catalyst surface to be here detected, as reported by Ivashenko et al. \parencite*{Ivashenko2021}.

% ratio 05
Reducing the partial pressure of oxygen to obtain \ce{O_2}/\ce{NH_3} = 0.5 induces a remarkable change in the O 1s spectrum, removing the \ce{O_a} peak, introducing a \ce{H_2O_a} peak at \qty{532.1}{\eV}, and shifting the \ce{H_2O_g} peak by \qty{0.6}{\eV}.
It is difficult to differ between \ce{OH_a} and \ce{H_2O_a}.
\ce{OH_a} is reported at \qty{532.0}{\eV}, but such a peak also exists under argon atmosphere, characteristic of water contamination in the reactor.
A binding energy of \qty{400.0}{\eV} in the N 1s spectrum is linked to adsorbed atomic nitrogen on Pt(100) \parencite{Sugai1993, vandenBroek1999}.
Since the only product of the ammonia oxidation is now nitrogen, the peak at \qty{404.3}{\eV} can safely be assigned to \ce{N_{2,g}}, similar to reported values \parencite{Ivashenko2021}.
Ammonia being in excess in the surface vicinity, all of the oxygen is used during ammonia oxidation yielding \ce{H_2O}, since \ce{NO} or \ce{N_2O} can not be detected.
This can be explained by the lack of available atomic oxygen adsorbed on the surface to form \ce{NO_a}, which is also the first step in the production of \ce{N2O}.
%, only observed without oxygen in the reactor.
% All of the oxygen is used in the first steps of the ammonia oxidation to remove hydrogen atoms from \ce{NH_x} species.
%, adsorbed hydrogen atoms then recombine and desorb from the surface.
%When \ce{OH_a} groups recombine to produce water, one oxygen atom is left on the catalyst surface.
\ce{H_2} is also produced from the dissociation of \ce{NH_3} on the catalyst.
It seems that in this condition, the oxygen atoms are preferably used for the de-hydrogenation process rather than the formation of \ce{NO} by reacting with adsorbed nitrogen.
The recombination of two nitrogen atoms on the catalyst surface is probably energetically preferred to the reaction of adsorbed nitrogen with adsorbed oxygen at low \ce{O_2}/\ce{NH_3} ratios.
% At this condition, most of the surface sites are hypothesised to be occupied by adsorbed water or nitrogen.
\ce{NH_{3,g}} reported between \qty{400.4}{\eV} and \qty{400.7}{\eV} \parencite{Ivashenko2021} was not clearly detected.
Nevertheless, it is possible that it also gives a contribution to the peak at \qty{400.0}{\eV}.

% ammonia
The next condition is pure ammonia.
Removing oxygen from the reactor removes the gas phase water peak from the O 1s level, adsorbed water is still present but shifted in energy, with the same peak shape, maybe from a lower presence of \ce{OH_a} groups since no more oxygen is present.
The dissociation of ammonia towards \ce{N_2} and \ce{H_2} is measured with the mass spectrometer.
\ce{N_{2,g}} is measured, shifted by \qty{0.6}{\eV}.
The peak at \qty{400.6}{\eV} can tentatively be linked to a superposition of gas phase ammonia and adsorbed \ce{NH_{x,a}} species.
No \ce{N_a} peak can be observed which supports a slow dissociation process of ammonia of the Pt(100) surface.

% low o2
Removing ammonia and adding \qty{0.55}{\milli\bar} of oxygen in the reactor shows at least five peaks in the O 1s level.
Both peaks at \qty{538.3}{\eV} and \qty{537.2}{\eV} are from gas phase oxygen.
At least three different oxygen species are present on the catalyst.
The peak at \qty{529.7}{\eV} exhibits a high intensity and low width, attributed to the presence of chemisorbed oxygen on Pt(100).
At least two other peaks can be detected at higher energy, around \qty{530.6}{\eV} and \qty{531.4}{\eV}.
If only one peak was present at \qty{531.4}{\eV} the intensity would drop between at \qty{530.6}{\eV} which is not the case here.
\ce{OH_a} reported at \qty{532}{\eV} \parencite{Kondratenko2006} is too far from the energies indexed here to be linked to the peak at \qty{531.4}{\eV}.
Thus, various oxygen species, adsorbed on different sites (e.g. hollow, on-top, bridge) may be linked to the peaks at \qty{530.6}{\eV}, and \qty{531.4}{\eV}.
The peak at \qty{530.6}{\eV} may also correspond to the transient structure measured with SXRD under \qty{5}{\milli\bar} of oxygen, also measured near \qty{530.4}{\eV} with XPS under \qty{8.8}{\milli\bar} of oxygen partial pressure.
The oxygen specie at \qty{529.7}{\eV} is predominant compared to the other species.
% The ammonia oxidation may prevent the other species from existing due to the high mobility of adsorbed atoms on the catalyst surface.

\subsubsection{Pt 4f level}

The Pt 4f level was also measured, presented in fig. \ref{fig:Pt4fPt100}.
No bulk oxide could be detected during reacting conditions, or under the presence of \qty{5}{\milli\bar} of oxygen with SXRD.
However, many different surface structures were identified under reacting conditions, changing as a function of the \ce{O_2}/\ce{NH_3} ratio (fig. \ref{fig:MapsAndLScansPt100HighOxAmmonia} - \ref{fig:MapsAndLScansPt100LowOxAmmonia}).

\begin{figure}[!htb]
    \centering
    \includegraphics[width=\textwidth]{/home/david/Documents/PhD/PhDScripts/B07_2022_04_XPS/Figures/Pt100/Pt4f_550_no_fit_merged.pdf}
    \caption{
        XPS spectra collected at the Pt 4f level under different atmospheres at \qty{450}{\degreeCelsius} with an incoming photon energy of \qty{550}{\eV}.
        A Shirley-type background has been subtracted from all XPS spectra.
        Normalisation performed first by the background intensity and secondly by the maximum intensity to allow a qualitative comparison between different total pressures.
        Spectra before normalisation by maximum intensity are shown on the top left.
    }
    \label{fig:Pt4fPt100}
\end{figure}

The introduction of \qty{8.8}{\milli\bar} of oxygen in the reactor induces a shift of \qty{0.1}{\eV} in the main peak position compared to pure argon atmosphere.
Two new features can be detected, at \qty{71.8}{\eV} and \qty{71.4}{\eV}, linked to the presence of oxygen atoms on top of the catalyst surface.
One new peak was clearly identified to appear in the O 1s level during this condition at \qty{529.7}{\eV} (fig. \ref{fig:O1sN1sPt100}), attributed to chemisorbed oxygen based on reported literature values \parencite{Rienks2003, Kondratenko2006}.
One of the two peaks at \qty{71.8}{\eV} and \qty{71.4}{\eV} is probably from the same specie.
The other may be linked to the moieties yielding peaks near \qty{530.5}{\eV} in the O 1s level.
% The presence of \ce{OH_a}, involved in the reaction mechanism but not measured in the O 1s level, is thus also rules out here.

The next condition is \ce{O_2}/\ce{NH_3} = \num{8}.
The main peak is shifted towards lower binding energy, possibly from the appearance of a new peak linked to adsorbed water as discussed below.
Both peaks identified under oxygen do not disappear but have a decreased intensity, similarly to what has been observed in the O 1s level, as a result of the participation of oxygen species in the catalytic reaction.

Reducing the partial pressure of oxygen to \qty{0.55}{\milli\bar} so that \ce{O_2}/\ce{NH_3} = \num{0.5} has the effect of removing the oxygen induced features at high binding energy.
The peak region near \qty{71.5}{\eV} closely resembles that observed under argon atmosphere.
The \ce{O_a} peaks in the O 1s level also disappeared during this condition.
Interestingly, the peak shape is broader and extended towards lower binding energy.
Adsorbed water is detected in the O 1s level at this condition.
\ce{H_2O_a} could not be detected in the O 1s level when \ce{O_2}/\ce{NH_3} = 8.
However, it is not possible to rule out its presence due to the low intensity of the XPS spectrum.
Thus, the shift towards lower energy in the Pt 4f level may be due to adsorbed water for both reacting conditions, more important when \ce{O_2}/\ce{NH_3} = 0.5.
The same reasoning can be applied to adsorbed nitrogen.

The following condition is pure ammonia atmosphere, the broadening effect at low binding energy is reverted.
Removing oxygen from the reactor removed the \ce{N_a} peak from the N 1s level, adsorbed water does not disappear.
However, the related energy position in the O 1s level is slightly shifted, from \qty{532.1}{\eV} to \qty{532.4}{\eV}, which underlines different chemical environments.
A new signal near \qty{71.6}{\eV} can safely be attributed to the presence of adsorbed \ce{NH_{x,a}} species, also identified in the N 1s level.
The peak intensity with respect to the maximum intensity is the highest measured, translating a large coverage of the platinum surface by \ce{NH_{x,a}}.

Changing the reactor atmosphere to \qty{0.55}{\milli\bar} of oxygen removes this \ce{NH_{x,a}} signal at \qty{71.6}{\eV}.
Overall the peak shape is similar to that under \qty{8.8}{\milli\bar} of oxygen.
The weaker intensity at higher binding energy is probably linked to a lower presence of the species measured under \qty{8.8}{\milli\bar} of oxygen.

Finally, reverting to argon atmosphere removed the peaks linked to \ce{O_a}.
However, the peak has not completely reverted to its shape before the oxidation cycle, and is very similar to the shape under a \ce{O_2}/\ce{NH_3} ratio equal to \num{0.5}.
The adsorbed nitrogen is absent from the Pt surface, only adsorbed water is present at both conditions, which was also present before the oxidation cycle but in a lower amount.
Therefore, adsorbed water is linked to the peak at lower binding energy, rather than \ce{N_a}.
    \section{Discussion}

%%%%%%%%%%%%%%%% high oxygen%%%%%%%%%%%%%%%%
% oxides Pt(111)
In this chapter was presented the structural and chemical evolution of Pt(111) and Pt(100) single crystals during ammonia oxidation.
The pre-oxidation of the platinum surfaces under \qty{80}{\milli\bar} of oxygen has allowed the identification of different surface structures and reconstructions.
For Pt(111), a Pt(111)-($8\times8$) superstructure was measured (fig. \ref{fig:MapsPt111A} - d), preceded by two Pt(111)-($6\times6$)-R\ang{\pm 8.8} superstructures (no visible second order peak, fig. \ref{fig:MapsPt111A} - c).
% Some of the peaks belonging to the rotated structures can also be described with the following matrix notation: Pt(111)-p$\begin{pmatrix} 1.08 & -0.21 \\ -0.21 & 1.08 \end{pmatrix}$, effectively describing a unit cell with a second order peak (fig. \ref{fig:MapsPt111A} - c).
Out-of-plane measurements have revealed a multi-layer thick structure (fig. \ref{fig:LScans80}), possibly corresponding to a $\alpha$-\ce{PtO_2} surface oxide.
The Pt(111)-($6\times6$)-R\ang{\pm 8.8} structures have been linked to monolayers (fig. \ref{fig:LScans80} - app. \ref{fig:LScans05}).
The complete understanding of the out-of-plane structures will be the subject of additional work.
A precursor relation was hypothesised between the Pt(111)-($6\times6$)-R\ang{\pm 8.8} and Pt(111)-($8\times8$) structures by time-resolved diffraction studies, under a lowered oxygen atmosphere (fig. \ref{fig:HexBraggPeaks}).
The importance of the partial pressure of oxygen in the growth kinetics has also been highlighted.

% oxides Pt(100)
On Pt(100), a bulk \ce{Pt_3O_4} oxide was identified during exposure to \qty{80}{\milli\bar} of oxygen (fig. \ref{fig:MapsPt100A} - b, fig. \ref{fig:FitPt100LScans}), but not under \qty{5}{\milli\bar} (fig. \ref{fig:MapsPt100D}), which shows the importance of the total pressure on the growth of \ce{Pt_3O_4}.
\ce{Pt_3O_4} follows a Pt(100)-($2\times2$) epitaxial relationship with the Pt(100) surface (fig. \ref{fig:Pt3O4onPt100}).
A second family of peaks was measured under high oxygen pressure (fig. \ref{fig:MapsPt100A} - b), slightly shifted in H or K with respect to the \ce{Pt_3O_4} peaks measured at $L=0$.
The shift was determined to depend on the elapsed time under oxygen atmosphere, two values were measured, $\delta_{H,K}=0.07$ and $\delta_{H,K}=0.09$, respectively after \qty{3}{\hour} (fig. \ref{fig:MapsPt100A} - b) and \qty{1}{\hour} (fig. \ref{fig:MapsPt100B} - b) of measurement.
No clear unit cell could be associated to the shifted peaks, the related out-of-plane signal is compatible with the existence of a monolayer on the Pt(100) surface (fig. \ref{fig:LScansHighOxygenPt100}).
Transient structures were observed at low oxygen pressure, which disappeared in the second measurement after \qty{5}{\hour} of elapsed time (fig. \ref{fig:MapsPt100D}).

% CTR
The measurement of crystal truncation rods revealed more important intensity changes on the Pt(100) surface (fig. \ref{fig:CTRPt100}) in comparison with the Pt(111) surface during the oxidation cycle (fig. \ref{fig:CTRPt111}).
For both surfaces (fig. \ref{fig:CTRFit111} - a, \ref{fig:CTRFit100} - a), the maximal roughness value is reached during exposition to high oxygen atmosphere, consistent with the formation of oxides under oxygen pressure.
Out-of-plane surface strain on Pt(111) was estimated to be contained on the topmost atomic layer, the presence of surface oxide structures was associated with compressive strain with respect to the initial values under argon atmosphere (fig. \ref{fig:CTRFit111} - b).
The influence of bulk \ce{Pt_3O_4} grown on Pt(100) under high oxygen atmosphere on the surface strain was not yet estimated, a lower oxygen atmosphere was related to compressive out-of-plane strain.
\ce{Pt_3O_4} was shown to not grow in a homogeneous layer, but rather in terms of islands with different thickness (fig. \ref{fig:CTRFitHighOxygen}).

% XPS
The XPS experiment was conducted at similar oxygen to ammonia partial pressure ratio, but with a decreased total pressure.
The highest oxygen pressure reached during the XPS experiment is \qty{8.8}{\milli\bar}, in the same order of magnitude as the low oxygen pressure condition for SXRD (\qty{5}{\milli\bar}).
Adsorbed oxygen was identified by XPS on both surfaces during this condition (fig. \ref{fig:O1sN1sPt111}, \ref{fig:O1sN1sPt100}), as well as under reduced oxygen pressure (\qty{0.55}{\milli\bar}).
The normalised intensity of the related peaks is approximately three times higher on the Pt(100) surface compared to the Pt(111) surface, which shows that the Pt(100) surface is more easily oxidised than the Pt(111) surface.
The duration of the high oxygen condition of Pt(111) was not long enough to enable the growth of the Pt(111)-($8\times8$) structure if dynamics at \qty{8.8}{\milli\bar} can be translated to \qty{5}{\milli\bar} (\qty{23}{\hour} vs. \qty{5}{\hour}).
Therefore, only the Pt(111)-($6\times6$)-R\ang{\pm 8.8} structure is expected to yield an additional peak in the O 1s and Pt 4f levels.
A signal was effectively identified in the Pt 4f level at \qty{71.6}{\eV}, but with low intensity (fig. \ref{fig:Pt4fPt111}).
Interestingly, two more peaks are identified in the O 1s level on Pt(100) at low oxygen pressure (\qty{0.55}{\milli\bar}) in addition to the \ce{O_a} signal, which could be linked to the transient structures measured with SXRD under \qty{5}{\milli\bar} of oxygen (fig. \ref{fig:MapsPt100D}).
% Two peaks was also measured in the Pt 4f level at \qty{71.8}{\eV} (fig. \ref{fig:Pt4fPt100}).

%%%%%%%%%%%%%%%%%%%% o2/nh3=8 %%%%%%%%%%%%%%%%%%%%
% Pt(111)
None of the aforementioned structures were measured under reacting conditions on Pt(111), for both \ce{O_2}/\ce{NH_3} ratios (\num{8} and \num{0.5}, fig. \ref{fig:MapsPt111B} - a \& b), which underlines that platinum oxides are not stable during the catalytic oxidation of ammonia on Pt(111) under those conditions.
Adsorbed nitrogen species (\ce{N_a}, \ce{NH_3}) were observed in the Pt4f (fig. \ref{fig:Pt4fPt111}) and N 1s (fig. \ref{fig:O1sN1sPt111}) levels, but without related surface reconstructions observable by SXRD on Pt(111) (fig. \ref{fig:MapsPt111B} - c).
Only adsorbed water was measured during the oxidation of ammonia in the O 1s level.

% Pt(100)
In contrast with Pt(111), reacting conditions were found to be linked with the appearance of surface reconstructions on Pt(100), depending on the oxygen exposure time.
After oxygen exposure for approximately \qty{17}{\hour}, the introduction of ammonia was linked to the removal of the shifted peaks near \ce{Pt_3O_4} signals (fig. \ref{fig:MapsPt100B} - a).
Both \ce{Pt_3O_4} and shifted signals could again be measured during a second exposure of the cleaned sample to oxygen for approximately \qty{1}{\hour} (fig. \ref{fig:MapsPt100B} - b).
The following introduction of ammonia was seen to induce signals compatible with Pt(100)-($10\times10$) surface reconstructions (fig. \ref{fig:MapsAndLScansPt100HighOxAmmonia} - a).
Some periodicity can be observed in the related out-of-plane signals at [H, K] = [0, -1.2] and [H, K] = [0.5, -1], high intensity peaks are observed at $L=1$ and $L=2$ for the measurement at [H, K] = [1.9, 0] (fig. \ref{fig:MapsAndLScansPt100HighOxAmmonia} - b), which can link the Pt(100)-($10\times10$) surface reconstructions with the presence of multilayers.
Additional in-plane signals are observed in fig. \ref{fig:MapsAndLScansPt100HighOxAmmonia} (a), shifted by the same amount $\delta_{H, K}=0.09$ as the signals present under high oxygen atmosphere (fig. \ref{fig:MapsPt100B} - b).
The intensity of this signal as a function of $L$ was measured at [H, K] = [0.5, -0.91] (fig. \ref{fig:MapsAndLScansPt100HighOxAmmonia} - b), decreasing with $L$, and thus showing no out-of-plane periodicity.

%%%%%%%%%%%%%%%%%%%% o2/nh3=0.5 %%%%%%%%%%%%%%%%%%%%

Lowering the \ce{O_2}/\ce{NH_3} ratio to \num{0.5} has removed the Pt(100)-($10\times10$) surface reconstructions, and induced the appearance of a hexagonal Pt(100)-Hex superstructure, with an out-of-plane signal consistent with that of monolayers (fig. \ref{fig:MapsAndLScansPt100LowOxAmmonia}).
Removing oxygen from the reactor resulted in the progressive removal of the hexagonal structures (fig. \ref{fig:MapsPt100C}), clearly linked to the simultaneous presence of both reactants.
Another key difference with Pt(111) is reported in the XPS spectra for Pt(100), with the visible presence of oxygen species in the O 1s level during reacting condition when the \ce{O_2}/\ce{NH_3} ratio is equal to \num{8} (fig. \ref{fig:O1sN1sPt100}), but not when this ratio is reduced to \num{0.5} by lowering the oxygen pressure.
No adsorbed nitrogen specie was identified when the \ce{O_2}/\ce{NH_3} ratio is equal to \num{8} in the N 1s level  (fig. \ref{fig:O1sN1sPt100}).
Therefore, the Pt(100)-($10\times10$) surface reconstruction can be linked to the presence of a surface oxide on the Pt(100) surface, which is only permitted by the simultaneous presence of ammonia, \textit{i.e.} linked to the reaction mechanism.
Additionally, the selectivity towards \ce{NO} in comparison with \ce{N_2} is three times more important than for Pt(111) (fig. \ref{fig:RGA450Pt111AndPt100}), which highlights the importance of oxygen adsorbed on the catalyst surface to facilitate the production of \ce{NO}.

\begin{figure}[!htb]
    \centering
    \includegraphics[width=\textwidth]{/home/david/Documents/PhDScripts/SixS_2023_04_SXRD_Pt111/figures/product_comparison_single_crystals.pdf}
    \includegraphics[width=\textwidth]{/home/david/Documents/PhDScripts/B07_2022_04_XPS/Figures/product_comparison_single_crystals.pdf}
    \caption{
        Evolution of reaction product partial pressures during the SXRD (top) and XPS (bottom) experiment on the Pt(111) and Pt(100) single crystals at \qty{450}{\degreeCelsius}.
        Mean partial pressures during \qty{1}{\minute} at the end of each condition, recorded from a leak in the reactor output by a residual gas analyser (RGA).
        The partial pressures have been normalised by the partial pressure of nitrogen.
    }
    \label{fig:RGA450Pt111AndPt100}
\end{figure}

This difference in selectivity is not observed during the SXRD experiments, this can be related to the position of the gas outlet far away from the sample (fig. \ref{fig:SampleHolder}), in contrast with the XPS experiment where the analyser nose is situated a few millimetres away from the sample surface (fig. \ref{fig:SampleSXRD} - c).
The contribution of the sample edges (fig. \ref{fig:SampleSXRD} - b), as well as the presence of steps on the single crystal surface may lower the specific contribution of the sample facet (\textit{i.e} (111) or (100)).

% roughness
The persisting presence of oxygen on the Pt(100) surface can also be linked to different evolution of the catalysts surface roughness.
The roughness on the Pt(111) surface decreases during ammonia oxidation (fig. \ref{fig:CTRFit111} - a), consistent with the removal of surface oxides following the introduction of ammonia.
Indeed, low roughness is already achieved when the \ce{O_2}/\ce{NH_3} ratio is equal to \num{0.5} (fig. \ref{fig:CTRFit111} - a).
For Pt(100), reacting conditions can also be associated to a decreasing roughness following pure oxygen atmosphere (fig. \ref{fig:CTRFit100} - a).
However, the minimal roughness is only reached once oxygen is completely removed from the reactor, \textit{i.e.} after the oxidation cycle.
The  Pt(100)-($10\times10$) surface reconstructions and Pt(100)-Hex hexagonal structures present on the Pt(100) surface during reacting conditions (fig. \ref{fig:MapsAndLScansPt100HighOxAmmonia} - a, \ref{fig:MapsAndLScansPt100LowOxAmmonia} - a), can be linked to the remaining high roughness, whereas Pt(111) exhibits a bulk terminated surface (fig. \ref{fig:MapsPt111B} - a \& b).

% surface strain
A similar behaviour is observed for the out-of-plane lattice strain.
The introduction of ammonia on Pt(111) was accompanied by a progressive reversal of the strain state, returning to the initial values measured under inert atmosphere when the \ce{O_2}/\ce{NH_3} ratio is equal to \num{0.5} (fig. \ref{fig:CTRFit111} - b).
For Pt(100), both reacting conditions show a strain state similar to the initial state under inert atmosphere (fig. \ref{fig:CTRFit100} - b).
The difference in out-of-plane strain compared to \ce{O2}/\ce{NH3} = 8 is in the same order of magnitude for both surfaces, but of different nature, \qty{\approx 0.06}{\percent} tensile / compressive strain on Pt(100) /  Pt(111).
Overall, the strain measured in this experiment is more important on Pt(100) than Pt(111).

%%%%%%%%%%%%%%%%%%%% Pure nh3 %%%%%%%%%%%%%%%%%%%%

Stopping the oxidation reaction by removing oxygen has finally resulted in the removal of surface structure on the Pt(100) surface (fig. \ref{fig:MapsPt100C} - a).
If the oxygen pressure is seen to decrease directly after the change of condition in the RGA signal (fig. \ref{fig:RGA450Pt100Cycle}), the hexagonal structures are still measured during the first hours of measurement.
It is possible that this progressive removal of the hexagonal surface structure is due to a slow transition between adsorbed nitrogen (\ce{N_a}) measured during the reaction, and adsorbed de-hydrogenated species observed only in the absence of oxygen (fig. \ref{fig:O1sN1sPt100} \& \ref{fig:Pt4fPt100}).
The largest transition in terms of surface strain is also during this change of condition for Pt(100), as the topmost layer goes from important compressive strain to almost bulk atomic positions for the Pt atoms (fig. \ref{fig:CTRFit100} - b, difference of \qty{\approx 0.33}{\percent}).
Atomic nitrogen is also observed in the N 1s level of Pt(111), but together with other nitrogen species (fig. \ref{fig:O1sN1sPt111} - \ref{fig:Pt4fPt111}), possibly impinging on its long-range ordering.
No important change is measured in strain while removing oxygen on Pt(111) (fig. \ref{fig:CTRFit111} - b), supporting the link between the Pt(111)-($8\times8$) surface structure, and changes in the out-of-plane strain.

%%%%%%%%% Argon %%%%%%%

The return to inert atmosphere after the oxidation cycle has shown that the final catalyst surface is more smooth from the roughness evolution (fig. \ref{fig:CTRFit111} - b, \ref{fig:CTRFit100} - b), without any adsorbed nitrogen or oxygen species (fig. \ref{fig:O1sN1sPt111}, \ref{fig:O1sN1sPt100}), and overall weaker strain values than during the reacting conditions.
For both surfaces, the surface roughness and strain are lower than measured at first under argon, before the reaction cycle.

% Discuss role of bulk oxides
\begin{table}[!htb]
\centering
\resizebox{\textwidth}{!}{%
\begin{tabular}{@{}llll@{}}
   \toprule
   \ce{O2} (\unit{\milli\bar}) & \ce{NH3} (\unit{\milli\bar}) & Pt(111) & Pt(100) \\
   \midrule
   0  & 0 & No oxide layer / no reconstruction                            & No oxide layer / no reconstruction\\
   80 & 0 & Two hexagonal rotated Pt(111)-($6\times6$)-R\ang{\pm8.8}            & Epitaxial Pt(100)-($2\times2$) \\
    &    & superstructures (monolayer)                                    & structure (bulk \ce{Pt3O4}),      \\
    &    & After \qty{9}{\hour}\qty{30}{\minute}: Pt(111)-($8\times8$) & and signals shifted in H or K     \\
    &    & superstructure (multilayer)                                    & (monolayer)                       \\
   \midrule
   80 & 10 & No oxide layer / no reconstruction                           & Pt(100)-($10\times10$) reconstruction\\
    &    & Weak signals in O 1s                                           & (multilayer)                      \\
    &    &                                                                & and signals shifted in H or K     \\
    &    &                                                                & (monolayer)                       \\
    &    &                                                                & Higher \ce{NO} selectivity, and   \\
    &    &                                                                & High signal in O 1s in            \\
    &    &                                                                & comparison with Pt(111)           \\
   10 & 10 & No oxide layer / no reconstruction                           & Pt(100)-Hex structure (monolayer) \\
    &    & No signal in O 1s level                                        & No signal in O 1s level           \\
    &    & \ce{N_a} and \ce{NH_{3,a}} in N 1s                             & Only \ce{N_a} in N 1s             \\
   \midrule
   0  & 10 & No oxide layer / no reconstruction                           & Progressive removal of Pt(001)-Hex structure \\
   0  & 0  & No oxide layer / no reconstruction                           & No oxide layer / no reconstruction\\
   \midrule
   5  &  0 & Back to the same structures as at \qty{80}{\milli\bar}       & Transient structures, different from signals \\
      &    & of \ce{O2}, but with decreased kinetics                      & observed at \qty{80}{\milli\bar} of \ce{O2} \\
   \bottomrule
\end{tabular}%
}
\caption{Brief summary of surface structures identified with SXRD, and relevant changes in XPS or RGA signals.}
\label{tab:RecapSXRD}
\end{table}

The different structures observed during this experiment are resumed in tab. \ref{tab:RecapSXRD}.
There is no apparent role of surface and bulk oxides on the Pt(111) surface, as the structures growing under high oxygen pressure have been removed directly during the reaction.
Pt(111)-($8\times8$) and Pt(111)-($6\times6$)-R\ang{\pm 8.8} cannot be linked to an increase or decrease of the catalytic activity, and thus do not seem to take part in the reaction mechanism.

\ce{Pt_3O_4}, observed here in the Pt(100)-($2\times2$) arrangement, has proven to be a source of oxygen atoms sustaining the catalytic oxygenation of \ce{CO} \textit{via} a Mars Van Krevelen mechanism \parencite{Seriani2006, Seriani2008}.
In the current experiment, a second structure possibly related to \ce{Pt_3O_4} was measured.
When under reacting conditions favouring \ce{NO}, which is the desired product for the Ostwald process, the in-plane peaks related to both structures were still measured, but being part of a Pt(100)-($10\times10$) reconstruction.
The presence of bulk \ce{Pt_3O_4} is clearly ruled out.
Nevertheless, it is possible that the reconstruction is due to the interaction of surface \ce{Pt_3O_4} with ammonia, resulting in an increased selectivity towards \ce{NO}.
Reducing the oxygen to ammonia ratio is linked to the reconstruction removal, the stopping of \ce{NO} production, as well as the removal of the adsorbed oxygen peak in the O 1s level.

The formation of \ce{RhO2} was observed on a Pt-Rh(100) single crystal by Resta et al. \parencite*{Resta2020a} at similar reactant partial pressures (total pressure of \qty{300}{\milli\bar}, \qty{3.5}{\milli\bar} of ammonia, and from \qtyrange{0}{20}{\milli\bar} of oxygen).
The rhodium surface oxide signal is only clearly measured when \ce{O2}/\ce{NH3} $>2$, and escalates when the temperature is increased from \qtyrange{175}{375}{\degreeCelsius}.
Both change of conditions are also linked to higher selectivity towards \ce{NO}.
The presence of this rhodium surface oxide is possibly impinging on the formation of \ce{Pt3O4} on Pt(100), observed in the current experiment, but at an oxygen partial pressure \num{4} times higher.
This would explain the role of rhodium in stabilising the catalyst surface \parencite{Fierro1990, Fierro1992, Bergene1996}, since platinum oxides have been reported to be more volatile at higher temperature \parencite{Alcock1960}.
\ce{Rh2O3} has been observed at elevated oxygen pressure on the same Pt-Rh(100) model crystal \parencite{Westerstrom2008}.
It is of key importance to increase the reactant partial pressure in future studies to be able to understand the role of platinum and rhodium oxides in the ammonia oxidation, facilitated now that in-plane signals have been detected.
This can help in the comprehension of the mechanisms driving the selectivity towards \ce{NO}.

    \chapter{Conclusion}
    \section{Research aim and results}

This thesis was aimed at understanding, \textit{operando}, the evolution of catalysts surfaces during ammonia oxidation.
To do so, mainly three techniques have been used, Bragg coherent diffraction imaging, surface x-ray diffraction, and x-ray photoelectron spectroscopy, combined with mass spectrometry measurements.
These techniques are compatible with near-ambient pressure conditions, allowing to reduce the pressure gap in heterogeneous catalysis.
To bridge the material gap, large Pt nanoparticles and single crystals were used.
Throughout this work, low \ce{O_2}/\ce{NH_3} partial pressure ratios and temperatures were linked with selectivity towards \ce{N_2}, whereas high \ce{O_2}/\ce{NH_3} ratios and temperatures were linked to an increased selectivity towards \ce{NO} and \ce{N_2O}.
This is consistent with the reaction selectivity reported for industrial catalysts, which supports the use of nanoparticles and single crystals as model catalysts to improve the system understanding.

Pt nanoparticles epitaxied on sapphire were first measured at \qtylist{300;500;600}{\degreeCelsius} with SXRD during ammonia oxidation, starting with the introduction of ammonia in the reactor, and followed by a progressive increase of the oxygen pressure (\ce{O_2}/\ce{NH_3} = \{0, 0.5, 1, 2, 8\}).
The average shape of the Pt nanoparticles is mainly constituted of \{111\}, \{110\}, \{100\} and \{113\} facets.
A progressive reshaping of the nanoparticles was revealed at \qty{600}{\degreeCelsius} during reacting conditions.
Once \ce{O_2}/\ce{NH_3} = 1, and for increasing oxygen to ammonia ratios, the \{110\} and \{113\} facets surface coverage decreases, replaced by \{111\}, and \{100\} facets.

A single nanoparticle was then imaged with BCDI under inert atmosphere, at temperatures between \qtylist{25;600}{\degreeCelsius}.
A dislocation present at the interface from room temperature to \qty{125}{\degreeCelsius} was successfully removed by flash annealing above \qty{800}{\degreeCelsius}.
This dislocation was shown to have an impact of the thermal relaxation of the particle.
The \{111\} facets close to the interface evolve between \{211\}, \{221\} and \{110\} facets as a function of the sample temperature, and interfacial strain.
The introduction of ammonia in the reactor at \qty{600}{\degreeCelsius} resulted in a reversal of the facet strain, likely at the origin of the creation of an interfacial dislocation network, highlighting the link between facet and interfacial strain in nanoparticles.

Two additional nanoparticles were measured at \qtylist{300;400}{\degreeCelsius} during ammonia oxidation.
Different behaviours were revealed on the two nanoparticles, that exhibit a different size, shape, facet coverage, and initial strain state.
The surface of the larger nanoparticle (particle \textit{C}, \qty{800}{\nm} wide) is constituted by  \{111\}, \{110\}, and \{100\} facets.
Particle \textit{C} showed a reversible decrease/increase of homogeneous strain when ammonia was introduced/removed from the reactor, and could not be reconstructed without it.
A non-reversible increase in heterogeneous strain was measured at \qty{300}{\degreeCelsius} when first exposing the sample to reacting condition (\ce{O_2}/\ce{NH_3} = 0.5), no such increase was reproduced at the same conditions at \qty{400}{\degreeCelsius}.

The increase of homogeneous strain linked to the presence/absence of ammonia was not reproduced on the second nanoparticle (particle \textit{B}, \qty{300}{\nm} wide), and with \{113\} facets also present on its surface.
A non-reversible increase of heterogeneous strain was also measured at \qty{300}{\degreeCelsius}, but induced by the presence of ammonia, without oxygen.
The possibility of having an oxidised sample can explain this evolution under a reducing atmosphere, but is not observed for particle \textit{C}.
The appearance of a defect at \qty{400}{\degreeCelsius} was linked to a non-reversible and important increase of homogeneous strain during the oxidation of ammonia, which further increased as a function of the \ce{O_2}/\ce{NH_3} ratio.
This structural evolution is clearly visible in 3D, with volumes of missing Bragg electronic density.
The presence of defects could have a prominent role in the catalyst strain field and thus catalytic properties.
A global activation of the catalyst at \qtylist{300;400}{\degreeCelsius} was reported while the \ce{O_2}/\ce{NH_3} ratio was kept equal to 2, with an increased production of \ce{NO}, and in a lesser extent \ce{N2O}, at the detriment of \ce{N2}.

To better understand the role of each facet in the behaviour of the Pt nanoparticles, SXRD and XPS experiments were carried out at \qty{450}{\degreeCelsius} on Pt(111) and Pt(100) single crystals.
The samples were first exposed to high oxygen atmosphere (\qty{80}{\milli\bar}) to oxidise the Pt surfaces (total pressure always kept to \qty{500}{\milli\bar} by the use of argon for SXRD).
Two different oxygen to ammonia ratios are used after surface oxidation, by first introducing \qty{10}{\milli\bar} of ammonia (\ce{O_2}/\ce{NH_3} = \num{8}), and then reducing the oxygen pressure to \qty{5}{\milli\bar} (\ce{O_2}/\ce{NH_3} = \num{0.5}).
The XPS experiment was performed with the same ammonia to oxygen ratios, but at lower partial pressures (approximately \qty{10}{\percent} regarding the SXRD experiment).

The oxidation of both surfaces is first discussed.
Bulk \ce{Pt_3O_4} was identified on Pt(100), in a Pt(100)-($2\times2$) arrangement of mean thickness equal to \qty{16}{\angstrom} (\textit{i.e.} 3 unit cells thick).
Signals shifted in reciprocal space in comparison to the \ce{Pt_3O_4} peaks are also measured.
Both the \ce{Pt_3O_4} and shifted peaks were measured \qty{1}{\hour} after the introduction of oxygen, but could not be detected under a reduced oxygen atmosphere (\qty{5}{\milli\bar}), even after several hours.
Transient signals are measured instead, underlying the importance of the oxygen pressure in the surface oxidation.
Bulk \ce{Pt3O4} was not observed on Pt(111).
Nevertheless, a Pt(111)-($8\times8$) commensurate superstructure was clearly identified after \qty{9}{\hour}\qty{30}{\minute} of elapsed time under high oxygen atmosphere, and after \qty{23}{\hour}\qty{30}{\minute} under reduced oxygen atmosphere.
From the intensity distribution of the related out-of-plane signals, this structure was determined to be a few layers thick.
Additional in-plane signals are also detected as soon as oxygen is introduced in the reactor, at both pressures, linked to a Pt(111)-($6\times6$)-R\ang{\pm 8.8} structure.
The intensity of those signals decreases once the Pt(111)-($8\times8$) structure is present, which supports a precursor link between both structures.
% Moreover, when describing the signals with the Pt(111)-p$\begin{pmatrix} 1.08 & -0.21 \\ -0.21 & 1.08 \end{pmatrix}$ matrix, which results in a real space angle of \ang{137.4}, a second-order signal is shared.
Signals measured in the XPS O 1s level at reduced oxygen atmosphere and attributed to surface oxygen species are more important for Pt(100) than Pt(111), showing that Pt(100) surface is more readily oxidised than Pt(111).
Oxidation of Pt(111) and Pt(100) was linked to increased surface roughness, and to compressive out-of-plane strain with respect to inert atmosphere for Pt(111).

Different behaviours were measured on both surfaces during reacting conditions.
Under high \ce{O_2}/\ce{NH_3} ratio, which favours the production of \ce{NO}, surface oxides are directly removed from Pt(111), but reconstruct with a (10x10) arrangement on Pt(100).
Oxygen species signals in the O 1s level are weak for Pt(111), difficult to dissociate from the background, whereas peaks are clearly detected for Pt(100).
The difference in surface oxygen presence is already observed during previous surface oxidation.
The reconstruction of the Pt(100) surface is linked to the persisting presence of surface oxygen species during ammonia oxidation, thereby differing from the Pt(111) surface.
Moreover, the selectivity towards \ce{NO} is increased for Pt(100), also linked to the more importance presence of surface oxygen, which is crucial in the production of \ce{NO} during the reaction mechanism \parencite{NovellLeruth2005, Offermans2006, Offermans2007, Imbihl2007, NovellLeruth2008}.
Increased roughness is also observed during this condition for both surfaces.

Pt(111) and Pt(100) show a similar selectivity towards \ce{N_2} when lowering the \ce{O_2}/\ce{NH_3} ratio to \num{0.5}.
The magnitude of the out-of-plane strain (surface relaxation), and the surface roughness, are already reduced in this condition for Pt(111), but stay at similar values for Pt(100).
On both surfaces, oxygen species are absent from the O 1s level, and adsorbed atomic nitrogen is measured, linked to the Pt(100)-Hex surface superstructure for Pt(100).
Adsorbed ammonia is also observed on Pt(111).

The difference in out-of-plane strain (surface relaxation) when lowering the \ce{O2}/\ce{NH3} ratio from \num{8} to \num{0.5} is in the same order of magnitude, but of different nature.
\qty{\approx 0.06}{\percent} tensile / compressive strain on Pt(100) /  Pt(111).
Most importantly, the strain is contained on the topmost layers of the platinum catalysts.
If one considers a Pt(100) surface voxel \qty{10}{\nm} thick, the corresponding averaged voxel strain would be equal to \qty{0.002}{\percent} with respect to the bulk lattice.
Therefore, the change of strain between reacting conditions becomes difficult to resolve, which can explain why no differences are observed on particle \textit{C} during ammonia oxidation.
Another factor to take into account is the increased surface roughness linked to high oxygen pressure, which has the effect of decreasing the intensity of the scattered photons far from the Bragg peak, and will thus reduce the experimental resolution.
Additional high-resolution studies could contribute to understanding the effect of adsorbates on surface relaxation.

The thesis demonstrates the importance of combining structural measurements at both the individual particle and particle assembly levels.
With Bragg coherent diffraction imaging (BCDI), we have demonstrated that depending on the morphology (shape, size, facet type and coverage, \textit{etc.}) and initial strain state of the Pt particles, different structural evolution are observed (in terms of strain, morphology, and defects) during the ammonia oxidation.
A better representation of the different behaviours followed by Pt particles during reaction, by measuring multiple individual particles, appears mandatory for a comprehensive understanding of the structure-activity relationship.
It is noteworthy that the strain at the particle/support interface seems to exert a pronounced influence on the particle behaviour.
Nanoparticles serve as a platform to explore the simultaneous structural evolution of diverse crystallographic facets at once.
Single crystal studies allow to better isolate the behaviours of single facets on nanoparticles.

\section{Perspectives}

This thesis has underlined the importance of bridging the material gap, by combining model catalysts with samples approaching the shape of industrial samples.
The structure evolution of Pt nanoparticles was shown to strongly depend on their size, shape, facet coverage and initial strain state.
The ammonia oxidation, which posses at least four different products, and a selectivity dependant on the reactant ratio, working temperature, and total pressure, draws a multi-dimensional parameter space only correctly addressed with information of different nature.

BCDI has proven to be powerful, but limited by the reconstruction process, as well as the time needed to find and image isolated nanoparticles.
Significant progress has already been made at the European synchrotron at the ID01 beamline, where rocking curves can be acquired under a minute.
\textit{Operando} measurements have proven to be challenging, Pt nanoparticles moving on the substrate at \qty{600}{\degreeCelsius} under reacting conditions, and sample environments failing during SXRD measurements.
Closing the pressure gap can be achieved for SXRD and BCDI by designing reactors resistant to highly oxidising atmosphere, and in this case, compatible with the ammonia/nitrogen oxide gas family, often banned due to their corrosive and toxic nature.
Future experiments must move towards industrial condition to be relevant, the importance of total pressure was demonstrated in this thesis with e.g. oxide formation on Pt(100) at \qty{80}{\milli\bar} and not at \qty{5}{\milli\bar}.
Simpler reaction cycles (e.g. by reducing the amount of \ce{O2}/\ce{NH3} ratios) will also help reduce the complexity of the data analysis.
A future reaction cycle is proposed.
First exposing the sample to increased oxygen pressure will reduce the time needed to detect potentially relevant oxides, as showcased here with the detection time of the Pt(111)-($8\times8$) structure, divided by \num{2.5} under \qty{80}{\milli\bar} of oxygen compared to \qty{5}{\milli\bar}.
A significant thickness of those platinum oxides could also be sufficient to be detected during BCDI experiments, especially now that the corresponding interplanar spacings are clearly determined.
Their impact of the facet strain in BCDI could then be determined by reconstructing the 3D strain tensor.
Higher oxygen pressure may reveal a role of platinum oxides during ammonia oxidation, invisible at lower pressure.
This will also remove the unknown regarding initial oxidised/contaminated states when introducing ammonia, and set a more systematic starting point for chemical cycles.
Ending the reaction cycle by having only ammonia in the reactor is also at the origin of large strain variations measured with SXRD, that could be more easily detected with BCDI.
In the future, surface x-ray experiments on Pt(110) and Pt(113) samples will add to our knowledge, and help better understand the structural behaviour of Pt nanoparticles.

Nevertheless, experimentally resolving the role of interfacial strain \textit{operando} could be performed by e.g. (i) exposing a nanoparticle that possesses an interfacial defect to reacting atmosphere, (ii) subsequently removing this defect by annealing, and (iii) repeating the exposition to reacting condition to observe a potential difference in the facet strain.
However, a certain control of the initial particle strain is needed.
The development of efficient simulation workflows to understand and de-correlate the effect of interfacial strain from adsorbant strain is also part of the solution.
The development of x-ray techniques that can bring higher statistics in the study of individual nanoparticles must be pushed forward.
In this perspective, x-ray Bragg ptychography is a technique that should receive increased interest since the measurement process can provide a unique solution.

The upgrade of synchrotrons is of key importance to efficiently navigate between different temperatures/pressures/reactant ratios.
For SXRD and BCDI, a reduced counting time will help catch transient structures, while increased photon count will increase the possibility to detect weak second and third order signals.
Moreover, BCDI gains from the increased coherent flux, while NAP-XPS also benefits from the increased photon flux.
Reaction cycles could also be repeated to increase the results significance, while changing the steps order.
Furthermore, a software upgrade is also needed, the surface x-ray diffraction community lacks an efficient and complete software combining data reduction and analysis.
For BCDI, the development of \textit{Gwaihir} was undertaken to respond to such needs from the community.
Additional work may address the data reduction automatisation by developing error-metrics for phase retrieval, methods to find the best solution, automatic facet detection in Python, \textit{etc}.
This would largely reduce the data reduction time, facilitating the analysis of multiple particles at various conditions.
A detailed study of the impact of trying to improve the image resolution by increasing the distance from the centre of the Bragg peak by combining simulations and experimental reconstructions would be of great interest to the community.
    
    \phantomsection
    \printbibliography[
        heading=bibintoc,
        title={Bibliography}
    ]

    \begin{appendices}
    \chapter{Nanoparticles study}
    \begin{figure}[!htb]
    \centering
    \includegraphics[width=\textwidth]{/home/david/Documents/PhDScripts/SixS\_2021\_01/FacetAnalyser/FacetDispEvolution.pdf}
    \caption{
        Mean value and standard deviation of the displacement ($\vec{u}_{\hat{q_z}}$) distribution as a function of the angle between the normal of each facet on the particle surface and the [111] direction.
        Upwards and downwards arrow are represented for respectively positive and negative displacement.
    }
    \label{fig:AmaterasuDisplacement}
\end{figure}

\begin{figure}[!htb]
    \centering
    \includegraphics[width=0.8\textwidth]{/home/david/Documents/PhDScripts/SixS_2021_03_SXRD_NH3/figures/epitaxy/200.pdf}
    \caption{
        Integrated intensity in a \ang{1} range around the value of the (200) scattering angle, as a function of the in-plane sample angle $\omega$.
    }
    \label{fig:Epitaxy200}
\end{figure}

\begin{figure}[!htb]
    \centering
    \includegraphics[width=0.8\textwidth]{/home/david/Documents/PhDScripts/SixS_2021_03_SXRD_NH3/figures/epitaxy/111.pdf}
    \caption{
        Integrated intensity in a \ang{1} range around the value of the (111) scattering angle, as a function of the in-plane sample angle $\omega$.
    }
    \label{fig:Epitaxy111}
\end{figure}

% \begin{figure}[!htb]
%     \centering
%     \includegraphics[width=\textwidth]{/home/david/Documents/PhDScripts/SixS_2021_03_SXRD_NH3/figures/epitaxy/powder_300.pdf}
%     \includegraphics[width=\textwidth]{/home/david/Documents/PhDScripts/SixS_2021_03_SXRD_NH3/figures/epitaxy/powder_500.pdf}
%     \includegraphics[width=\textwidth]{/home/david/Documents/PhDScripts/SixS_2021_03_SXRD_NH3/figures/epitaxy/powder_600.pdf}
%     \caption{
%         Scattered intensity as a function of $\delta$ showing the Bragg peak originating from the substrate near \ang{16}, the
%         (220) powder ring near \ang{27.9} and the (220) Bragg peak from the [111] oriented nanoparticles near \ang{28.1}.
%     }
%     \label{fig:PowderNanoparticles}
% \end{figure}

\begin{figure}[!htb]
    \centering
    \includegraphics[width=\textwidth]{/home/david/Documents/PhD/Figures/introduction/stereographic_projection_bottom.pdf}
    \caption{
        Stereographic projection perpendicular to $[\bar{1}\bar{1}\bar{1}]$ crystallographic orientation.
        The circles describe the angle with the $[\bar{1}\bar{1}\bar{1}]$ direction from \ang{0} (centre) to \ang{90} (outer-ring).
    }
    \label{fig:StereoBottom}
\end{figure}

\begin{figure}[!htb]
    \centering
    \includegraphics[width=\textwidth]{/home/david/Documents/PhD/Figures/bcdi_data/B7/Fit3D.png}
    \caption{
        Sum of Bragg peak intensity perpendicular to $\vec{q}_x$ and $\vec{q}_y$ and $\vec{q}_z$, fitted by Lorentzian profiles.
    }
    \label{fig:FitB73D}
\end{figure}


\begin{figure}[!htb]
    \centering
    \includegraphics[width=\textwidth]{/home/david/Documents/PhDScripts/SixS_2021_03_SXRD_NH3/figures/rga/rga_300.pdf}
    \caption{
        Time dependent partial pressures recorded from a leak in the reactor output by a residual gas analyser (RGA) during the SXRD experiment on the non-patterned sample containing Pt[111]||Al$_2$O$_3$[0001] nanoparticles at \qty{300}{\degreeCelsius}.
        Vertical dotted lines indicate transitions between two conditions for which the \ce{NH_3} and \ce{O_2} flow is indicated in the legend.
    }
    \label{fig:RGA300SXRDNanoparticles}
\end{figure}

\begin{figure}[!htb]
    \centering
    \includegraphics[width=\textwidth]{/home/david/Documents/PhDScripts/SixS_2021_03_SXRD_NH3/figures/rga/rga_500.pdf}
    \caption{
        Time dependent partial pressures recorded from a leak in the reactor output by a residual gas analyser (RGA) during the SXRD experiment on the non-patterned sample containing Pt[111]||Al$_2$O$_3$[0001] nanoparticles at \qty{500}{\degreeCelsius}.
        Vertical dotted lines indicate transitions between two conditions for which the \ce{NH_3} and \ce{O_2} flow is indicated in the legend.
    }
    \label{fig:RGA500SXRDNanoparticles}
\end{figure}

\begin{figure}[!htb]
    \centering
    \includegraphics[width=\textwidth]{/home/david/Documents/PhDScripts/SixS_2021_03_SXRD_NH3/figures/rga/rga_600.pdf}
    \caption{
        Time dependent partial pressures recorded from a leak in the reactor output by a residual gas analyser (RGA) during the SXRD experiment on the non-patterned sample containing Pt[111]||Al$_2$O$_3$[0001] nanoparticles at \qty{600}{\degreeCelsius}.
        Vertical dotted lines indicate transitions between two conditions for which the \ce{NH_3} and \ce{O_2} flow is indicated in the legend.
    }
    \label{fig:RGA600SXRDNanoparticles}
\end{figure}

\begin{figure}[!htb]
    \centering
    \includegraphics[width=\textwidth]{/home/david/Documents/PhDScripts/SixS_2021_06_BCDI_NH3/gas_analysis/figures/rga_300.pdf}
    \caption{
        Time dependent partial pressures recorded from a leak in the reactor output by a residual gas analyser (RGA) during the BCDI experiment on the patterned sample containing Pt[111]||Al$_2$O$_3$[0001] nanoparticles at \qty{300}{\degreeCelsius}.
        Vertical dotted lines indicate transitions between two conditions for which the \ce{NH_3} and \ce{O_2} flow is indicated in the legend.
    }
    \label{fig:RGA300BCDINanoparticles}
\end{figure}

\begin{figure}[!htb]
    \centering
    \includegraphics[width=\textwidth]{/home/david/Documents/PhDScripts/SixS_2021_06_BCDI_NH3/gas_analysis/figures/rga_400.pdf}
    \caption{
        Time dependent partial pressures recorded from a leak in the reactor output by a residual gas analyser (RGA) during the BCDI experiment on the patterned sample containing Pt[111]||Al$_2$O$_3$[0001] nanoparticles at \qty{400}{\degreeCelsius}.
        Vertical dotted lines indicate transitions between two conditions for which the \ce{NH_3} and \ce{O_2} flow is indicated in the legend.
    }
    \label{fig:RGA400BCDINanoparticles}
\end{figure}

    \clearpage
\includepdf[pages=-]{/home/david/Documents/PhD/PhDScripts/SixS_2021_06_BCDI_NH3/figures/Poster3D/3D_B7.pdf}\label{app:AppendixB7}
\includepdf[pages=-]{/home/david/Documents/PhD/PhDScripts/SixS_2021_06_BCDI_NH3/figures/Poster3D/3D_D6_edited.pdf}\label{app:AppendixD6}

    \chapter{Single crystals study}
    \begin{table}[!htb]
    \centering
    \resizebox{\textwidth}{!}{%
    \begin{tabular}{@{}|l|l|lllllllllll|@{}}
        \toprule
        Structure & Interplanar & \multicolumn{11}{c|}{Oxygen pressure} \\
        \midrule
          & spacing (\unit{\angstrom}) & \multicolumn{2}{l|}{80 mbar} & \multicolumn{9}{l|}{5 mbar} \\
        \midrule
         & & \multicolumn{11}{c|}{Time since gas introduction (end of measurement)} \\
        \midrule
                                     &                  & 03h23 & 09h53 & \multicolumn{1}{|l}{00h34} & 04h03 & 08h00 & 15h57 & 22h56 & 24h08 & 25h43 & 26h36 & 27h28 \\
        \midrule % first struc split low O2
        Pt(111)-($6\times6$)-R\ang{\pm 8.8} & $3.01 \pm 0.01$ & \yes  & \yes  & \multicolumn{1}{|l}{\yes}  & \yes  & \yes  & \yes  & \yes  & \yes  &                              \yes & \yes & \yes \\
        \midrule % first struc not split high O2
        Pt(111)-($6\times6$)-R\ang{\pm 8.8} & $2.92 \pm 0.03$ & \yes  & \yes  & \multicolumn{1}{|l}{\no}   & \no   & \no   & \no   & \no   & \no   &                              \yes & \no & \no \\
        \midrule % first struc not split low O2
        Pt(111)-($6\times6$)-R\ang{\pm 8.8} & $2.87 \pm 0.02$ & \no   & \no   & \multicolumn{1}{|l}{\yes}  & \yes  & \yes  & \yes  & \yes  & \yes  &                              \yes & \yes & \yes \\
        \midrule % first struc split low O2
        Pt(111)-($6\times6$)-R\ang{\pm 8.8} & $2.79 \pm 0.07$ & \no   & \no   & \multicolumn{1}{|l}{\no}   & \no   & \yes  & \yes  & \yes  & \yes  &                              \yes & \yes & \yes \\
        \midrule % hex struct
        Pt(111)-($8\times8$)       & $2.69 \pm 0.02$ & \no   & \yes  & \multicolumn{1}{|l}{\no}   & \no   & \no   & \no   & \no   & \yes  &                              \yes & \yes & \yes \\
        % \midrule % surf oxide
        % Pt(111)-($6\times6$)-R\ang{\pm 8.8} & $2.31 \pm 0.04$ & \no & \yes & \multicolumn{1}{|l}{\no} & \no & \no & \yes & \yes & \yes & \yes & \yes & \yes \\
        \midrule % first struc not split high O2
        % NRHS & $1.550 \pm 0.000$ & \yes & nv & \multicolumn{1}{|l}{nv} & \no & nv & nv & nv & nv & nv & nv & nv \\
        % \midrule % first struc not split low O2
        Pt(111)-($8\times8$)       & $1.53 \pm 0.00$ & \yes & nv     & \multicolumn{1}{|l}{nv}    & \yes  & nv    & nv    & nv & nv & nv & nv & nv \\
        % NRHS & $1.35 \pm 0.000$ & \no & nv & \multicolumn{1}{|l}{nv} & \yes & nv & nv & nv & nv & nv & nv & nv \\
        \bottomrule
    \end{tabular}
    }
    \caption{
        Interplanar spacings computed from signals observed during in-plane reciprocal maps, for different oxygen pressure and exposition times for Pt(111).
        The three markers (\yes, \no, \text{nv}) correspond to observed, non-observed, and non visible signals (\textit{i.e.} not in the area spanned by the map).
        The errors on the interplanar spacings are computed by considering the positions of similar peaks in q-space.
    }
    \label{tab:InterplanarSpacingsPt111Oxygen}
\end{table}

\begin{figure}[!htb]
    \centering
    \includegraphics[width=0.49\textwidth]{/home/david/Documents/PhDScripts/SixS_2023_04_SXRD_Pt111/figures/rotated_reconstruction_cells.pdf}
    % \includegraphics[width=0.49\textwidth]{/home/david/Documents/PhDScripts/SixS_2023_04_SXRD_Pt111/figures/reconstruction_42punto65.pdf}
    \caption{
    Sketch of Pt(111) surface and Pt(111)-($6\times6$) structure unit cell vectors, in real and reciprocal space (a).
    % Pt(111)-p$\begin{pmatrix} 1.08 & -0.21 \\ -0.21 & 1.08 \end{pmatrix}$ structure (b).
    }
    \label{fig:Pt11180O2Structures}
\end{figure}

\begin{figure}[!htb]
    \centering
    \includegraphics[width=0.49\textwidth]{/home/david/Documents/PhDScripts/SixS_2023_04_SXRD_Pt111/figures/alphapto2_reconstruction_cell.pdf}
    \caption{
    Sketch of Pt(111) surface and Pt(111)-($8\times8$) structure unit cell vectors, in real and reciprocal space (a).
    }
    \label{fig:Pt111AlphaPtO2}
\end{figure}

\begin{figure}[!htb]
    \centering
    \includegraphics[width=\textwidth]{/home/david/Documents/PhDScripts/SixS_2023_04_SXRD_Pt111/figures/l_scans_low_oxygen.pdf}
    \caption{
        SSR observed under \qty{5}{\milli\bar} of \dioxygen.
    }
    \label{fig:LScans05}
\end{figure}

\begin{figure}[!htb]
    \centering
    \includegraphics[width=\textwidth]{/home/david/Documents/PhDScripts/SixS_2023_04_SXRD_Pt111/figures/T450_1.pdf}
    \caption{
        Time dependent partial pressures recorded from a leak in the reactor output by a residual gas analyser (RGA) during the SXRD experiment on the Pt(111) single crystal at \qty{450}{\degreeCelsius}.
        Vertical dotted lines indicate transitions between two conditions for which the \ce{NH_3} and \ce{O_2} flow is indicated in the legend.
        The RGA electron multiplier is off for all the masses.
    }
    \label{fig:RGA450Pt111_1}
\end{figure}

\begin{figure}[!htb]
    \centering
    \includegraphics[width=\textwidth]{/home/david/Documents/PhDScripts/SixS_2023_04_SXRD_Pt111/figures/T450_2.pdf}
    \caption{
        Time dependent partial pressures recorded from a leak in the reactor output by a residual gas analyser (RGA) during the SXRD experiment on the Pt(111) single crystal at \qty{450}{\degreeCelsius}.
        Vertical dotted lines indicate transitions between two conditions for which the \ce{NH_3} and \ce{O_2} flow is indicated in the legend.
        The RGA electron multiplier is on for all the masses besides oxygen and argon.
    }
    \label{fig:RGA450Pt111_2}
\end{figure}

\begin{figure}[!htb]
    \centering
    \includegraphics[width=\textwidth]{/home/david/Documents/PhDScripts/SixS_2023_04_SXRD_Pt111/figures/T450_3.pdf}
    \caption{
        Time dependent partial pressures recorded from a leak in the reactor output by a residual gas analyser (RGA) during the SXRD experiment on the Pt(111) single crystal at \qty{450}{\degreeCelsius}.
        Vertical dotted lines indicate transitions between two conditions for which the \ce{NH_3} and \ce{O_2} flow is indicated in the legend.
        The RGA electron multiplier is on for all the masses besides argon.
    }
    \label{fig:RGA450Pt111_3}
\end{figure}

\begin{figure}[!htb]
    \centering
    \includegraphics[width=\textwidth]{/home/david/Documents/PhDScripts/SixS_2023_04_SXRD_Pt111/figures/T450_4.pdf}
    \caption{
        Time dependent partial pressures recorded from a leak in the reactor output by a residual gas analyser (RGA) during the SXRD experiment on the Pt(111) single crystal at \qty{450}{\degreeCelsius}.
        Vertical dotted lines indicate transitions between two conditions for which the \ce{NH_3} and \ce{O_2} flow is indicated in the legend.
        The RGA electron multiplier is on for all the masses besides argon.
    }
    \label{fig:RGA450Pt111_4}
\end{figure}
    
\begin{figure}[!htb]
    \centering
    \includegraphics[width=0.49\textwidth]{/home/david/Documents/PhDScripts/SixS_2022_01_SXRD_Pt100/figures/hex_reconstruction.pdf}
    \caption{
    Sketch of Pt(100) surface and Pt(100)-p(33x25)-R\ang{0} structure unit cell vectors, in real and reciprocal space (a).
    }
    \label{fig:Pt100UnitCellsReconstruction}
\end{figure}

\begin{figure}[!htb]
    \centering
    \includegraphics[width=\textwidth]{/home/david/Documents/PhDScripts/SixS_2022_01_SXRD_Pt100/gas_analysis/figures/exp_before_heater_change_no_norm.pdf}
    \caption{
        Time dependent partial pressures recorded from a leak in the reactor output by a residual gas analyser (RGA) during the SXRD experiment on the Pt(100) single crystal at \qty{450}{\degreeCelsius}, before the change of heater.
        Vertical dotted lines indicate transitions between two conditions for which the \ce{NH_3} and \ce{O_2} flow is indicated in the legend.
    }
    \label{fig:RGA450Pt100BeforeHeaterChange}
\end{figure}

\begin{figure}[!htb]
    \centering
    \includegraphics[width=\textwidth]{/home/david/Documents/PhDScripts/SixS_2022_01_SXRD_Pt100/gas_analysis/figures/exp_cycle_no_norm.pdf}
    \caption{
        Time dependent partial pressures recorded from a leak in the reactor output by a residual gas analyser (RGA) during the SXRD experiment on the Pt(100) single crystal at \qty{450}{\degreeCelsius}, after the change of heater.
        Vertical dotted lines indicate transitions between two conditions for which the \ce{NH_3} and \ce{O_2} flow is indicated in the legend.
    }
    \label{fig:RGA450Pt100Cycle}
\end{figure}

\begin{figure}[!htb]
    \centering
    \includegraphics[width=\textwidth]{/home/david/Documents/PhDScripts/SixS_2022_01_SXRD_Pt100/gas_analysis/figures/exp_cond_g_no_norm.pdf}
    \caption{
        Time dependent partial pressures recorded from a leak in the reactor output by a residual gas analyser (RGA) during the SXRD experiment on the Pt(100) single crystal at \qty{450}{\degreeCelsius}, after ammonia oxidation cycle.
        Vertical dotted lines indicate transitions between two conditions for which the \ce{NH_3} and \ce{O_2} flow is indicated in the legend.
    }
    \label{fig:RGA450Pt100CondG}
\end{figure}
    \begin{figure}[!htb]
    \centering
    \includegraphics[width=\textwidth]{/home/david/Documents/PhD/PhDScripts/SixS_2023_04_SXRD_Pt111/figures/product_partial_pressure_comparison_vs_time.pdf}
    \caption{
        Evolution of reaction product partial pressures as a function of time during the SXRD experiment on the Pt(111) and Pt(100) single crystals at \qty{450}{\degreeCelsius}.
        Mean partial pressures during \qty{1}{\minute} at the end of each condition, recorded from a leak in the reactor output by a residual gas analyser (RGA) as detailed in sec. \ref{sec:XCAT}.
    }
    \label{fig:RGA450SXRDPt111AndPt100Time}
\end{figure}

\begin{figure}[!htb]
    \centering
    \includegraphics[width=\textwidth]{/home/david/Documents/PhD/PhDScripts/B07_2022_04_XPS/Figures/product_partial_pressure_comparison_vs_time.pdf}
    \caption{
        Evolution of reaction product partial pressures as a function of time during the XPS experiment on the Pt(111) and Pt(100) single crystals at \qty{450}{\degreeCelsius}.
        Mean partial pressures during \qty{1}{\minute} at the end of each condition, recorded from a leak in the reactor output by a residual gas analyser (RGA) as detailed in sec. \ref{sec:XCAT}.
    }
    \label{fig:RGA450XPSPt111AndPt100Time}
\end{figure}



\end{appendices}

\end{document}

%%%%%%%%%%%%%%%%%%%%%%%%%%%%%%%%%%%%%%%%%%%%%%%%%%%%%%%%%%%%%%